\documentclass{situationist-preprint}
\usepackage{longtable,array,multicol}

\title{A Modal Ontological Argument Without the S5 Axiom}
\msid{REL-01-b}
\msfamily{REL-01}
\msdiscipline{Philosophy of Religion}
\mstype{Companion manuscript}
\mskeywords{ontological argument; modal logic; S5; Brouwerian system; frame correspondence; ancestral accessibility; possibility premise}

% generated by gen/generate.py -- do not edit
\newcommand{\stTotal}{3{,}160}
\newcommand{\stNOne}{2}
\newcommand{\stNTwo}{10}
\newcommand{\stNThree}{104}
\newcommand{\stNFour}{3{,}044}
\newcommand{\stNumB}{118}
\newcommand{\stNumC}{1{,}296}
\newcommand{\stNumU}{127}
\newcommand{\stNumFive}{50}
\newcommand{\stNumFour}{291}
\newcommand{\stNumSFive}{11}
\newcommand{\stNumCnotB}{1{,}178}
\newcommand{\stNumUnotFive}{77}
\newcommand{\stNumBnotU}{80}
\newcommand{\stNumBnotFour}{93}
\newcommand{\stNumSOneNotSThree}{80}
\newcommand{\stNumSThreeNotSOne}{89}
\newcommand{\stNumArgStar}{1{,}296}
\newcommand{\stNumArgOrig}{38}
\newcommand{\stNumSSeven}{609}
\newcommand{\stNumSEight}{1{,}264}
\newcommand{\stNumBnotSSeven}{10}
\newcommand{\stNumBnotSEight}{0}
\newcommand{\stNumCnotSEight}{32}
\newcommand{\stChecks}{50{,}560}
\newcommand{\stValuations}{49{,}580}
\newcommand{\stQValuations}{786{,}088}

\newcommand{\MG}{\mathrm{MG}}
\newcommand{\MGs}{\mathrm{MG}^{*}}
\newcommand{\Rp}{R^{+}}
\newcommand{\nplus}{\nec^{+}}

\begin{document}
\maketitle

\begin{abstract}
The main argument of \emph{On the Existence of God} (REL-01-a) derives the necessary existence of a maximally great being from the possibility that one exists, using the S5 conditions on accessibility. We determine exactly which frame conditions each step of that argument needs, and we give a version that does not use the S5 axiom. With the original possibility premise, the actual existence of a maximally excellent being requires exactly symmetry, and necessary maximal greatness requires exactly a condition we call uniformity: worlds accessible from a common world see the same worlds. With a stronger possibility premise, that possibly some being exists and is maximally excellent at every world reachable by any finite chain of accessibility, the whole conclusion follows on every frame in which each step of accessibility can be retraced, and in particular on every symmetric frame, with or without reflexivity and transitivity. All correspondence claims are proved in general and checked by exhaustive enumeration of the \stTotal{} frames with at most four worlds, listed in the appendix.
\end{abstract}

\section{Introduction}

\subsection{The S5 premise and its critics}

REL-01-a states the modal ontological argument in the possible-worlds framework of \citet{Plantinga1974}. Its formal premises are the S5 conditions on accessibility (Assumption~2.1 of REL-01-a: reflexivity, symmetry and transitivity) and the possibility premise (Assumption~3.1), that at the actual world it is possible that some being is maximally great. A being is maximally great at a world if it exists and is maximally excellent at every world accessible from that world (Definition~2.4). The argument passes through two lemmas: whatever is maximally great at a world accessible from $w$ is maximally great at every world accessible from $w$ (Lemma~3.2), and whatever is possibly necessary is necessary (Lemma~3.3). It concludes that necessarily some being is maximally great (Theorem~3.4), and hence that some being actually exists and is maximally excellent (Corollary~3.5).

The S5 conditions have been doubted for metaphysical modality. \citet{Chandler1976} and \citet{Salmon1989} argued, from cases of tolerance in the origins of artefacts and organisms, that what is possible may itself be contingently possible, so that accessibility among metaphysically possible worlds is not transitive and the S4 principle $\nec p \to \nec\nec p$ fails. Others have accepted S5 as the logic of metaphysical necessity on the ground that what is metaphysically possible does not vary from world to world \citep{Williamson2013}. REL-01-a takes the second view and notes that, since its proof uses Lemma~3.3, the proof does not survive the move to a weaker logic; it assigns to the present manuscript a version of the argument that does not use the S5 axiom, at the cost of a stronger possibility premise. REL-01-g observes in passing that the principle $\pos\nec p \to p$, which needs only symmetry, would yield actual existence directly. The question that these remarks leave open is precise: which conditions on accessibility does each conclusion of the argument actually need?

\subsection{The modal premise in the history of the argument}

The dependence of the modal ontological argument on a strong modal logic was visible from its revival. \citet{Malcolm1960} argued that Anselm's second argument shows the existence of God to be either necessary or impossible, so that it is necessary unless it is impossible. \citet{Hartshorne1962} gave the argument a formal statement in which the decisive step is a principle he called Becker's postulate, that modal status is itself necessary: if something is possible or necessary, it is necessarily so. That postulate is characteristic of S5, and it plays in Hartshorne's proof the role that Lemma~3.3 plays in REL-01-a. \citet{Plantinga1974} restated the argument in terms of possible worlds and took the logic of broadly logical necessity to be S5, so that what is possible does not vary from world to world; his maximal greatness, like Definition~2.4 of REL-01-a, builds necessary existence into the property whose possibility is premised.

Critics have pressed on the modal premise from two directions. The first questions S5 itself. \citet{Chandler1976} and \citet{Salmon1989} argued that accessibility among metaphysically possible worlds is not transitive. On their view an artefact could have originated from slightly different material, and could have had, had it so originated, a further slight difference, without its being possible for it to have originated from material that differs by the sum of the two: a world possible relative to a possible world need not be possible relative to the actual world. If they are right, S4 fails, and with it the S5 principle. The second direction accepts S5 and turns it against the argument: in S5 the possibility of a maximally great being and the possibility of its absence are incompatible, and the argument from the second is as valid as the argument from the first \citep{Oppy1995}. REL-01-a records this as Proposition~4.1 and assigns the choice between the premises to Volume~III.

The present manuscript addresses the first line of criticism only, and it does so without taking sides on whether metaphysical accessibility is transitive. A parallel question arises for Gödel's version of the argument. \citet{Benzmuller2014} reported, from their computer-assisted analysis of Scott's version of Gödel's proof, that the modal logic KB suffices for its main theorem. REL-01-c shows by a checked derivation that the same holds for the emended version it studies, whose only relational step is symmetry. Our results for the Plantinga-style argument are of the same kind but sharper: for each step, they identify the exact condition on frames, not only a sufficient system.

\subsection{Thesis}

We answer the question completely for the frames of propositional and constant-domain quantified modal logic. Our results are of two kinds.

\emph{Correspondence results.} Keeping the possibility premise of REL-01-a, Corollary~3.5 (actual existence and maximal excellence) follows on a frame if and only if the frame is symmetric, and Theorem~3.4 (necessary maximal greatness) follows if and only if the frame is \emph{uniform}: any two worlds accessible from a common world have the same successors. Uniformity is strictly weaker than the Euclidean condition that corresponds to the S5 axiom, so the original argument is valid on exactly the symmetric uniform frames, a class that properly includes the equivalence relations and contains frames that are neither reflexive nor transitive (Corollary~\ref{c:orig}).

\emph{The argument without the S5 axiom.} Replace maximal greatness by \emph{ancestral} maximal greatness: existence and maximal excellence at every world reachable by a finite chain of accessibility steps. From the possibility of ancestral maximal greatness, the conclusions of Theorem~3.4 and Corollary~3.5 both follow on every frame in which every accessibility step can be retraced by some chain of steps (condition C), and in particular on every symmetric frame, whether or not it is reflexive or transitive (Theorem~\ref{t:main}). Condition C is also necessary (Proposition~\ref{p:exact}). On equivalence relations the strengthened premise coincides with the possibility premise of REL-01-a (Proposition~\ref{p:coincide}), so the version agrees with the main argument where the main argument applies.

The general results are proved below. They are also checked by a program that enumerates every frame with at most four worlds up to isomorphism, \stTotal{} frames in all, and decides by brute force over all valuations whether each of eight argument schemas is valid on each frame. The listing in Appendix~\ref{app:listing} records the verdicts frame by frame.

\subsection{Plan}

Section~\ref{s:framework} sets out the framework and the argument schemas. Section~\ref{s:corr} proves the correspondence results. Section~\ref{s:main} states and proves the argument without the S5 axiom. Section~\ref{s:computation} summarizes the computation and discusses six frames from the listing. Section~\ref{s:objections} answers five objections, and Section~\ref{s:not} records what the results do not establish.

\section{Framework}\label{s:framework}

\subsection{Frames and conditions}

A \emph{frame} is a non-empty set $W$ of worlds with a relation $R$ on $W$. We write $R(w)$ for the set of worlds accessible from $w$. A proposition is a set of worlds; $\nec p$ holds at $w$ if $p$ holds at every world in $R(w)$, and $\pos p$ if $p$ holds at some world in $R(w)$. A schema is \emph{valid on a frame} if every instance holds at every world under every valuation \citep{Kripke1963, HughesCresswell1996}. For the purposes of the argument the actual world is any world; validity on a frame therefore requires the argument to succeed whichever world is actual.

\begin{definition}\label{d:std}
The five standard conditions are: T (reflexivity: $wRw$); B (symmetry: $wRv$ implies $vRw$); 4 (transitivity: $wRv$ and $vRu$ imply $wRu$); 5 (the Euclidean condition: $wRv$ and $wRu$ imply $vRu$); D (seriality: $R(w)$ is non-empty). The S5 conditions of Assumption~2.1 of REL-01-a are T, B and 4 together, which hold exactly when $R$ is an equivalence relation.
\end{definition}

\begin{definition}\label{d:ancestral}
$\Rp$ is the transitive closure of $R$: $w\Rp v$ if there is a chain $w = w_0 R w_1 R \cdots R w_k = v$ with $k \geq 1$. We write $\nplus p$ for the proposition that holds at $w$ if $p$ holds at every world in $\Rp(w)$.
\end{definition}

\begin{definition}\label{d:conds}
A frame is \emph{uniform} (condition U) if, for all $w$ and all $u, v \in R(w)$, $R(u) = R(v)$. It is \emph{ancestrally uniform} (condition U$^+$) if, for all $w$ and all $u, v \in R(w)$, $\Rp(u) = \Rp(v)$. It is \emph{retraceable} (condition C) if $wRv$ implies $v \Rp w$: every step of accessibility can be undone by a chain of steps.
\end{definition}

Retraceability is equivalent to the symmetry of $\Rp$, and on a finite frame to the condition that every arrow lies on a cycle. Symmetry is the special case in which every step can be undone in one step.

\subsection{Maximal greatness and its ancestral}

We keep the constant domain of beings, the existence predicate and the definitions of maximal excellence and maximal greatness of REL-01-a (Definitions~2.2 to~2.4). Write $M(x,w)$ for: $x$ exists and is maximally excellent at $w$. Then $x$ is maximally great at $w$, $\MG(x,w)$, if $M(x,v)$ for every $v \in R(w)$.

\begin{definition}\label{d:mgstar}
A being $x$ is \emph{ancestrally maximally great} at $w$, $\MGs(x,w)$, if $M(x,v)$ for every $v \in \Rp(w)$: $x$ exists and is maximally excellent at every world reachable from $w$ by a finite chain of accessibility steps.
\end{definition}

Since $R(w) \subseteq \Rp(w)$, ancestral maximal greatness entails maximal greatness. On a transitive frame $\Rp = R$ and the two coincide. Ancestral maximal greatness is the natural reading of the informal idea behind Definition~2.4 of REL-01-a when accessibility may fail to be transitive: a being that is maximally great in the intended sense exists and is maximally excellent not only at the worlds possible relative to here, but at the worlds possible relative to those, and so on.

\begin{principle}[Strong possibility premise]\label{pr:strong}
At the actual world, it is possible that some being is ancestrally maximally great: $@ \in \pos\,\exists x\,\MGs(x,\cdot)$.
\end{principle}

Principle~\ref{pr:strong} is the stronger possibility premise that REL-01-a announces for this manuscript. It entails Assumption~3.1 of REL-01-a on every frame, and is entailed by it on every transitive frame.

\subsection{The argument schemas}

The steps and conclusions of the argument are captured by eight schemas. In each, $p$ stands for the proposition that a given being exists and is maximally excellent; $\nec p$ then says that the being is maximally great, and $\nplus p$ that it is ancestrally maximally great.

\begin{center}\footnotesize
\begin{tabular}{@{}llll@{}}
\toprule
& Schema & Role in the argument & Condition \\
\midrule
S1 & $\pos\nec p \to p$ & Assumption~3.1 yields Corollary~3.5 & B \\
S2 & $\pos\nec p \to \nec p$ & Lemma~3.3 of REL-01-a & 5 \\
S3 & $\pos\nec p \to \nec\nec p$ & Lemma~3.2; Assumption~3.1 yields Theorem~3.4 & U \\
S4 & $\nec p \to \nec\nec p$ & Anselm's principle (REL-01-g) for $\MG$ & 4 \\
S5 & $\pos\nplus p \to p$ & Principle~\ref{pr:strong} yields Corollary~3.5 & C \\
S6 & $\pos\nplus p \to \nec\nplus p$ & Principle~\ref{pr:strong} yields $\nec\exists x\,\MGs$ & U$^+$ \\
S7 & $\pos(\nec p \wedge \nec\nec p) \to p \wedge \nec\nec p$ & two-step version (Section~\ref{s:finite}) & C$_2$ \\
S8 & $\pos(\nec p \wedge \nec\nec p \wedge \nec\nec\nec p) \to p \wedge \nec\nec p$ & three-step version (Section~\ref{s:finite}) & C$_3$ \\
\bottomrule
\end{tabular}
\end{center}

The last column anticipates the results of Section~\ref{s:corr}. Schema S4 is the form that Anselm's principle of REL-01-g, $\nec(G \to \nec G)$, takes when $G$ is analysed by Definition~2.4; it is also the principle behind the converse argument of REL-01-a (Proposition~4.1), since $\pos\neg\nec p \to \neg\nec p$ is equivalent to S4.

The quantified forms of the schemas replace $\pos\nec p$ by $\pos\exists x\,\MG(x,\cdot)$, $\nec\nec p$ by $\nec\exists x\,\MG(x,\cdot)$, $p$ in S1 and S5 by $\exists x\, M(x,\cdot)$, and similarly for the ancestral notions. The next lemma shows that nothing is lost by working with the propositional forms.

\begin{lemma}[Reduction]\label{l:reduction}
On every frame, and for every domain size, the quantified form of each of S1--S8 is valid if and only if its propositional form is valid.
\end{lemma}

\begin{proof}
Suppose the propositional form fails at a world $w$ under a valuation $V$. Give every being of the domain the same profile, $M(x,v)$ if and only if $v \in V$; each quantified clause then reduces to the corresponding propositional clause, and the quantified form fails at $w$. Conversely, suppose the quantified form fails at $w$. Its antecedent supplies a being $b$ that has the relevant property (maximal greatness, ancestral or $k$-step maximal greatness, or in S4 maximal greatness at $w$ itself) at the relevant world, and its consequent fails for every being, in particular for $b$. Let $V = \{v : M(b,v)\}$. Under $V$ the propositional antecedent holds at $w$, because $b$ witnesses it, and the propositional consequent fails at $w$, because it fails for $b$.
\end{proof}

The proof uses the constant domain of REL-01-a (Section~2.2) only in that the quantifiers range over the same beings at every world. For variable domains the reduction needs the premise that REL-01-a, Chapter~6, adds; we do not pursue that case.

\section{Exact frame conditions}\label{s:corr}

Each schema corresponds to a first-order condition on frames. The results for S1, S2 and S4 are standard \citep{HughesCresswell1996, Blackburn2001}; we include the short proofs because the other three follow the same pattern, and because the pattern is what makes the counterexamples in the listing intelligible.

\begin{proposition}\label{p:s1}
S1 is valid on a frame if and only if the frame is symmetric.
\end{proposition}

\begin{proof}
If the frame is symmetric and $\nec p$ holds at some $v \in R(w)$, then $w \in R(v)$, so $p$ holds at $w$. If $wRv$ but not $vRw$, let $p$ be $R(v)$: $\nec p$ holds at $v$, so $\pos\nec p$ holds at $w$, but $p$ fails at $w$.
\end{proof}

\begin{proposition}\label{p:s2}
S2 is valid on a frame if and only if the frame is Euclidean.
\end{proposition}

\begin{proof}
If the frame is Euclidean, $\nec p$ holds at $v \in R(w)$ and $u \in R(w)$, then $vRu$, so $p$ holds at $u$. If $wRv$, $wRu$ and not $vRu$, let $p$ be $R(v)$: $\pos\nec p$ holds at $w$ and $p$ fails at $u \in R(w)$.
\end{proof}

\begin{proposition}\label{p:s3}
S3 is valid on a frame if and only if the frame is uniform.
\end{proposition}

\begin{proof}
If the frame is uniform, $\nec p$ holds at $v \in R(w)$, $u \in R(w)$ and $z \in R(u)$, then $z \in R(u) = R(v)$, so $p$ holds at $z$; hence $\nec\nec p$ holds at $w$. If $u, v \in R(w)$ and some $z \in R(u)$ is not in $R(v)$, let $p$ be $R(v)$: $\pos\nec p$ holds at $w$, and $\nec p$ fails at $u$ because $p$ fails at $z$.
\end{proof}

\begin{proposition}\label{p:s4}
S4 is valid on a frame if and only if the frame is transitive.
\end{proposition}

\begin{proof}
Standard: if $wRv$, $vRu$ and not $wRu$, let $p$ be $R(w)$; then $\nec p$ holds at $w$ and $\nec p$ fails at $v$. The converse is immediate.
\end{proof}

\begin{proposition}\label{p:s5}
S5 is valid on a frame if and only if the frame is retraceable.
\end{proposition}

\begin{proof}
If the frame is retraceable and $\nplus p$ holds at $v \in R(w)$, then $w \in \Rp(v)$, so $p$ holds at $w$. If $wRv$ but $w \notin \Rp(v)$, let $p$ be $\Rp(v)$: $\nplus p$ holds at $v$, so $\pos\nplus p$ holds at $w$, and $p$ fails at $w$.
\end{proof}

\begin{proposition}\label{p:s6}
S6 is valid on a frame if and only if the frame is ancestrally uniform.
\end{proposition}

\begin{proof}
As for Proposition~\ref{p:s3}, with $\Rp(u)$ and $\Rp(v)$ in place of $R(u)$ and $R(v)$, using $\Rp(u) \subseteq \Rp(w')$ whenever $w' R u$.
\end{proof}

\begin{lemma}\label{l:inclusions}
Every symmetric frame is retraceable; every retraceable frame is ancestrally uniform; every Euclidean frame is uniform. None of the converses holds.
\end{lemma}

\begin{proof}
The first claim is immediate. For the second, let $u, v \in R(w)$. Then $\Rp(u) \subseteq \Rp(w)$, and since $w \in \Rp(u)$ by retraceability, $\Rp(w) \subseteq \Rp(u)$; so $\Rp(u) = \Rp(w) = \Rp(v)$. For the third, let $u, v \in R(w)$ and $z \in R(u)$. By the Euclidean condition applied at $w$, $uRv$; applied at $u$ to $v$ and $z$, $vRz$. So $R(u) \subseteq R(v)$, and by symmetry of the roles, equality. Counterexamples to the converses are the frames F3.27, F2.5 and F2.6 of the listing, discussed in Section~\ref{s:computation}: a three-cycle, retraceable but not symmetric; a frame in which two worlds both see one of them, ancestrally uniform but not retraceable; and two worlds that see each other and not themselves, uniform but not Euclidean.
\end{proof}

\begin{corollary}\label{c:orig}
The argument of REL-01-a, from Assumption~3.1 to both Theorem~3.4 and Corollary~3.5, is valid on a frame if and only if the frame is symmetric and uniform. Every equivalence relation is symmetric and uniform, but not conversely: the symmetric uniform frames include frames that are neither reflexive nor transitive.
\end{corollary}

\begin{proof}
By Lemma~\ref{l:reduction} and Propositions~\ref{p:s1} and~\ref{p:s3}. An equivalence relation is Euclidean and hence uniform by Lemma~\ref{l:inclusions}. The frame F2.6, two worlds each seeing only the other, is symmetric and uniform, and is neither reflexive nor transitive.
\end{proof}

Corollary~\ref{c:orig} sharpens the remark of REL-01-a that its proof does not survive the move to a weaker logic. The proof does not, because Lemma~3.3 is the S5 principle S2 and fails on every non-Euclidean frame. The conclusions survive on a class of frames strictly larger than the equivalence relations, but only on frames that satisfy uniformity, a condition that is implied neither by the conditions of S4 nor by those of the Brouwerian system KTB. In particular they fail on some reflexive transitive frames, so the argument from Assumption~3.1 is not valid in S4, which is the logic \citet{Salmon1989} defends against; and they fail on some reflexive symmetric frames, so it is not valid in the Brouwerian system KTB either.

\section{The argument without the S5 axiom}\label{s:main}

\begin{theorem}\label{t:main}
On every retraceable frame, Principle~\ref{pr:strong} entails:
\begin{enumerate}[label=(\roman*)]
\item necessarily, some being is ancestrally maximally great: $@ \in \nec\,\exists x\, \MGs(x,\cdot)$;
\item necessarily, some being is maximally great: $@ \in \nec\,\exists x\, \MG(x,\cdot)$, which is the conclusion of Theorem~3.4 of REL-01-a;
\item some being actually exists and is maximally excellent, which is the conclusion of Corollary~3.5 of REL-01-a;
\item some being is ancestrally maximally great at the actual world.
\end{enumerate}
\end{theorem}

\begin{proof}
By Principle~\ref{pr:strong} there are a world $v \in R(@)$ and a being $b$ such that $M(b,z)$ for every $z \in \Rp(v)$. Since $@ R v$, retraceability gives $v \Rp @$. Hence $@ \in \Rp(v)$, and every world reachable from $@$ is reachable from $v$: $\Rp(@) \subseteq \Rp(v)$. So $M(b,@)$, which is (iii), and $M(b,z)$ for every $z \in \Rp(@)$, which is (iv). Let $u \in R(@)$. Then $\Rp(u) \subseteq \Rp(@)$, so $b$ is ancestrally maximally great at $u$, which gives (i); and since $R(u) \subseteq \Rp(u)$, $b$ is maximally great at $u$, which gives (ii).
\end{proof}

The proof uses neither reflexivity nor transitivity, and it uses symmetry only in the weakened form of retraceability. It does not use Lemma~3.3 of REL-01-a or any other form of the S5 principle.

\begin{corollary}\label{c:systems}
The conclusions (i)--(iv) of Theorem~\ref{t:main} follow from Principle~\ref{pr:strong} in the Brouwerian systems KB, KDB and KTB, whose frames are the symmetric, the symmetric serial and the symmetric reflexive frames.
\end{corollary}

\begin{proof}
Every symmetric frame is retraceable (Lemma~\ref{l:inclusions}).
\end{proof}

\begin{proposition}\label{p:exact}
Retraceability is necessary as well as sufficient: if a frame is not retraceable, there is a model on it in which Principle~\ref{pr:strong} holds at some world and (iii) fails there. Ancestral uniformity is necessary and sufficient for (i) alone.
\end{proposition}

\begin{proof}
By Lemma~\ref{l:reduction}, Proposition~\ref{p:s5} for (iii) and Proposition~\ref{p:s6} for (i).
\end{proof}

\begin{proposition}\label{p:coincide}
On every transitive frame, and in particular under Assumption~2.1 of REL-01-a, Principle~\ref{pr:strong} is equivalent to Assumption~3.1 of REL-01-a.
\end{proposition}

\begin{proof}
On a transitive frame $\Rp = R$, so $\MGs = \MG$.
\end{proof}

\begin{proposition}\label{p:anselm}
Anselm's principle for ancestral maximal greatness,
\[\nec(\exists x\,\MGs(x,\cdot) \to \nec\,\exists x\,\MGs(x,\cdot)),\]
is valid on every frame. So is the corresponding converse argument: if it is possible that no being is ancestrally maximally great, then no being is ancestrally maximally great.
\end{proposition}

\begin{proof}
If $b$ is ancestrally maximally great at $w$ and $u \in R(w)$, then $\Rp(u) \subseteq \Rp(w)$, so $b$ is ancestrally maximally great at $u$. The converse argument is the contrapositive of the same fact.
\end{proof}

Proposition~\ref{p:anselm} explains where the frame conditions of the main argument go in the present version. In REL-01-a the work done by symmetry and transitivity in Lemma~3.2, and in REL-01-g the work done by Anselm's principle (Assumption~2.2 of REL-01-g), is to make maximal greatness persist along accessibility. Ancestral maximal greatness persists along accessibility by its definition, so the substitute principle needs no frame condition at all. What remains is to return from the world at which the possibility premise is witnessed to the actual world, and that is exactly retraceability.

\subsection{Finite strengthenings}\label{s:finite}

The reply to the third objection below notes that ancestral maximal greatness is not expressible by a single formula of the basic modal language. On symmetric frames that expressive cost can be avoided, because a bounded strengthening of the premise suffices.

\begin{definition}\label{d:ck}
For $k \geq 1$, let $R^{\leq k}(w)$ be the set of worlds reachable from $w$ by a chain of between $1$ and $k$ accessibility steps. A being is \emph{$k$-step maximally great} at $w$ if it exists and is maximally excellent at every world in $R^{\leq k}(w)$. A frame satisfies C$_k$ if (a) $wRv$ implies $w \in R^{\leq k}(v)$, and (b) $u, v \in R(w)$ implies $R(u) \subseteq R^{\leq k}(v)$.
\end{definition}

One-step maximal greatness is maximal greatness, and $k$-step maximal greatness is expressed by $\nec p \wedge \nec\nec p \wedge \cdots \wedge \nec^k p$, which is a formula of the basic modal language for each finite $k$. Clause (a) of C$_k$ says that every step can be retraced in at most $k$ steps; clause (b) says that whatever is seen from one successor of a world is reachable in at most $k$ steps from every other successor.

\begin{proposition}\label{p:finite}
The possibility that some being is $k$-step maximally great entails, at every world of a frame, both that some being exists and is maximally excellent there and that necessarily some being is maximally great there, if and only if the frame satisfies C$_k$. Every symmetric frame satisfies C$_3$. Some reflexive symmetric frames do not satisfy C$_2$.
\end{proposition}

\begin{proof}
Sufficiency: let $b$ be $k$-step maximally great at $v \in R(w)$. By (a), $w \in R^{\leq k}(v)$, so $b$ exists and is maximally excellent at $w$; for $u \in R(w)$, $R(u) \subseteq R^{\leq k}(v)$ by (b), so $b$ is maximally great at $u$. Necessity: if (a) fails for $wRv$, or (b) fails for $u, v \in R(w)$, let $p$ be $R^{\leq k}(v)$; in either case the antecedent of the schema (S7 for $k=2$, S8 for $k=3$, and its analogue for other $k$) holds at $w$ and the consequent fails, and Lemma~\ref{l:reduction} lifts the countermodel to the quantified form. On a symmetric frame, $wRv$ gives $vRw$, a one-step return, and for $u, v \in R(w)$ and $z \in R(u)$ the chain $v R w R u R z$ has three steps. The frame F4.2683 of the listing is reflexive and symmetric and fails C$_2$, as explained in Section~\ref{s:computation}.
\end{proof}

On symmetric frames, then, the strengthened premise can be stated in the object language as $\pos\exists x(\nec Mx \wedge \nec\nec Mx \wedge \nec\nec\nec Mx)$, and three iterations are the least number that works on all of them. This gives a second version of the argument without the S5 axiom, valid in KB, KDB and KTB, whose premise is weaker than Principle~\ref{pr:strong} but whose frame condition is stronger than retraceability: of the \stNumC{} retraceable frames in the enumeration, \stNumCnotSEight{} do not satisfy C$_3$.

\section{The computation}\label{s:computation}

The general results are checked by an exhaustive computation, whose method is described in Appendix~\ref{app:method}. Every frame with at most four worlds is generated up to isomorphism; there are \stNOne, \stNTwo, \stNThree{} and \stNFour{} such frames with one to four worlds, \stTotal{} in all. For each frame, the program computes the eight conditions T, B, 4, 5, D, U, U$^+$ and C, and decides the validity of S1--S8 by evaluating each schema at every world under every valuation. It also decides the quantified forms over a domain of two beings under every assignment of the profile $M$. It verifies that every verdict agrees with the condition that Propositions~\ref{p:s1}--\ref{p:s6} and~\ref{p:finite} assign to the schema, and that the quantified and propositional verdicts agree, as Lemma~\ref{l:reduction} requires. There is no disagreement among the \stChecks{} verdicts.

Table~\ref{tab:systems} gives the result by system: for each class of frames defined by a combination of standard conditions, the number of frames of the class (up to four worlds) and the number on which each schema is valid; \textbf{all} marks a schema valid on every frame of the class.

\begin{table}[ht]
\caption{Validity of the argument schemas by system, frames with at most four worlds.}\label{tab:systems}
% generated by gen/generate.py -- do not edit
\begin{center}\small\begin{tabular}{@{}lrrrrrrrrr@{}}\toprule
System & frames & S1 & S2 & S3 & S4 & S5 & S6 & S7 & S8\\\midrule
K & 3160 & 118 & 50 & 127 & 291 & 1296 & 1597 & 609 & 1264\\
KD & 2417 & 89 & 30 & 71 & 115 & 1247 & 1506 & 564 & 1215\\
KT & 238 & 18 & 11 & 11 & 46 & 101 & 101 & 61 & 101\\
KB & 118 & \textbf{all} & 25 & 38 & 25 & \textbf{all} & \textbf{all} & 108 & \textbf{all}\\
K4 & 291 & 25 & 42 & 68 & \textbf{all} & 25 & 68 & 25 & 25\\
K5 & 50 & 25 & \textbf{all} & \textbf{all} & 42 & 25 & \textbf{all} & 25 & 25\\
KDB & 89 & \textbf{all} & 11 & 20 & 11 & \textbf{all} & \textbf{all} & 79 & \textbf{all}\\
KTB & 18 & \textbf{all} & 11 & 11 & 11 & \textbf{all} & \textbf{all} & 17 & \textbf{all}\\
KD4 & 115 & 11 & 23 & 29 & \textbf{all} & 11 & 29 & 11 & 11\\
S4 & 46 & 11 & 11 & 11 & \textbf{all} & 11 & 11 & 11 & 11\\
K45 & 42 & 25 & \textbf{all} & \textbf{all} & \textbf{all} & 25 & \textbf{all} & 25 & 25\\
KD5 & 30 & 11 & \textbf{all} & \textbf{all} & 23 & 11 & \textbf{all} & 11 & 11\\
KD45 & 23 & 11 & \textbf{all} & \textbf{all} & \textbf{all} & 11 & \textbf{all} & 11 & 11\\
KB4 & 25 & \textbf{all} & \textbf{all} & \textbf{all} & \textbf{all} & \textbf{all} & \textbf{all} & \textbf{all} & \textbf{all}\\
S5 & 11 & \textbf{all} & \textbf{all} & \textbf{all} & \textbf{all} & \textbf{all} & \textbf{all} & \textbf{all} & \textbf{all}\\
\bottomrule\end{tabular}\end{center}

\end{table}

The table displays the results of Sections~\ref{s:corr} and~\ref{s:main} in miniature. S5 and S6, the two schemas of the strengthened argument, are valid on every frame of KB, KDB and KTB. S1 is valid on exactly the frames of the B systems. S3, the step to Theorem~3.4 from the original premise, is valid on all frames of the systems containing 5 and on only \stNumArgOrig{} of the \stNumB{} symmetric frames, the symmetric uniform ones of Corollary~\ref{c:orig}. In the S4 class, neither the original nor the strengthened argument is valid on the frames that are not equivalence relations. Of the \stTotal{} frames, \stNumC{} are retraceable and only \stNumB{} are symmetric: retraceability is a much weaker requirement than symmetry, and \stNumCnotB{} frames validate the strengthened argument without being symmetric.

Seven frames illustrate the results. The listing writes a frame as the successor sets of worlds $0, 1, \ldots$ separated by dots, so that \texttt{01.0} is the frame in which world $0$ sees itself and world $1$, and world $1$ sees world $0$.

\emph{F2.6} (\texttt{1.0}). Two worlds, each seeing only the other. The frame is symmetric and uniform but neither reflexive, transitive nor Euclidean, and all six schemas except S2 and S4 are valid on it. The original argument of REL-01-a is therefore valid on F2.6 although Lemma~3.3 fails there: the conclusions are reached without the S5 principle, through Lemma~3.2, which holds on F2.6 because the frame is uniform. This frame is the witness for the second part of Corollary~\ref{c:orig}.

\emph{F2.8} (\texttt{01.0}). World 0 sees itself and world 1; world 1 sees world 0. The frame is symmetric but not uniform. S1 is valid: from the original premise, a maximally excellent being actually exists. S3 fails, under the valuation that makes $p$ true at world 0 alone, evaluated at world 0: a being that exists and is maximally excellent at world 0 only is maximally great at world 1, so the original premise holds at world 0, but it is not maximally great at world 0, because world 0 sees world 1. The strengthened premise is false at world 0 under this valuation, since world 1 reaches itself through world 0, and S5 and S6 are valid on the frame. F2.8 is the smallest frame on which the strengthened premise does what the original one cannot.

\emph{F2.5} (\texttt{0.0}). Both worlds see world 0 only. The frame is serial, transitive and Euclidean (KD45), hence uniform, and S2, S3, S4 and S6 are valid. But S1 and S5 fail at world 1 under the valuation that makes $p$ true at world 0 only, and the listing records this witness as \texttt{1:10}. At world 1 it is necessary that some being is maximally great, and nothing exists and is maximally excellent at world 1. This is the doxastic pattern familiar from KD45: necessity without truth. It shows that Theorem~3.4 does not by itself yield Corollary~3.5, which REL-01-a derives with the help of reflexivity.

\emph{F3.27} (\texttt{1.2.0}). A three-cycle. It is retraceable and uniform but not symmetric, and S1 fails: from the original premise, nothing need actually exist. S3, S5 and S6 are valid. The strengthened argument succeeds here without symmetry, which shows that retraceability, not symmetry, is the exact condition of Theorem~\ref{t:main}.

\emph{F3.93} (\texttt{01.12.02}). A reflexive three-cycle: each world sees itself and its successor. It is reflexive and retraceable, but neither symmetric, transitive nor uniform. Both steps of the original argument fail (S1 and S3), and both steps of the strengthened argument hold (S5 and S6). Of the \stNThree{} three-world frames, this is the simplest reflexive frame on which the two versions come apart completely.

\emph{F2.9} (\texttt{01.1}). World 0 sees both worlds, world 1 sees only itself: the smallest reflexive transitive frame that is not an equivalence relation. Only S4 is valid. Both versions of the argument fail, and the witness \texttt{0:01} makes $p$ true at world 1 alone: a being that exists and is maximally excellent at world 1 only is ancestrally maximally great at world 1, so the strengthened premise holds at world 0, but the being does not exist at world 0. This is the shape of the cases of \citet{Chandler1976} and \citet{Salmon1989}: a world possible relative to the actual world from which the actual world is not possible. Neither version of the argument survives the move to S4, and no version can, since the frame is a countermodel to S5 and S6 alike.

\emph{F4.559 and F4.2683} (\texttt{13.02.1.0} and \texttt{013.012.12.03}). The first is the path $3$--$0$--$1$--$2$ with accessibility in both directions; the second adds reflexivity. Both are symmetric, so S5, S6 and S8 are valid on them, but S7 fails at world~1 when $p$ is false at world~3 alone. A being that exists and is maximally excellent everywhere except at world~3 is two-step maximally great at world~2, which is accessible from world~1; but world~3 is accessible from world~0, which is accessible from world~1, so the being is not maximally great at world~0, and Theorem~3.4 fails at world~1. The two frames are among the smallest witnesses for the last clause of Proposition~\ref{p:finite}: among symmetric frames with at most four worlds, \stNumBnotSSeven{} fail C$_2$ and none fails C$_3$.

\section{Objections and replies}\label{s:objections}

\begin{objection}
The strengthened premise is stronger, and a stronger premise is less credible. The version without the S5 axiom buys validity with plausibility.
\end{objection}

\begin{reply}
That is the cost REL-01-a announces, and we do not deny it. Two points limit it. First, on transitive frames the premises coincide (Proposition~\ref{p:coincide}); the strengthening is felt only by someone who already denies transitivity, and for such a person the original premise does not express the idea of maximal greatness well, since a being that is maximally great in the sense of Definition~2.4 may fail to exist at a world that is possible relative to a possible world. Second, the support that Volume~III gives the possibility premise transfers. REL-01-f shows that what Volume~III needs from its Lemma~2417 is that the great-making properties be jointly instantiable, and that the proof uses neither the possibility premise nor S5. The same proof gives Principle~\ref{pr:strong} if the great-making properties include the ancestral forms of necessary existence and necessary maximal excellence, since by the argument of REL-01-f their conjunction is ancestral maximal greatness.
\end{reply}

\begin{objection}
The converse argument gains more than the original. Possibly no being is ancestrally maximally great is a weaker premise than possibly no being is maximally great, and by Proposition~\ref{p:anselm} the converse argument from it needs no frame condition at all. The dialectical balance of REL-01-a (Proposition~4.1) is upset in favour of the atheist.
\end{objection}

\begin{reply}
The converse argument does become frame-free, but the balance is not upset on the frames on which the strengthened argument is valid. On a retraceable frame the worlds reachable from the actual world form a single cluster under $\Rp$, and on that cluster the strengthened premise and its converse are contradictories at every world, exactly as Assumption~3.1 and the hypothesis of Proposition~4.1 are in S5. The dispute again reduces to a choice between two possibility premises. On frames that are not retraceable the converse argument succeeds and the original one fails, but on those frames, as F2.9 shows, no version of the argument succeeds, so the asymmetry there is not a new cost of the present version.
\end{reply}

\begin{objection}
Ancestral maximal greatness is not expressible in the language of propositional modal logic: $\nplus$ is an infinitary operator. A premise stated with it is not a modal premise at all.
\end{objection}

\begin{reply}
The premise is stated in the same metalanguage of worlds and accessibility in which REL-01-a states Definition~2.4, and nothing in the argument requires it to be the translation of a modal formula. If one wants an object-language version, $\nplus p$ is the infinite conjunction of $\nec p, \nec\nec p, \nec\nec\nec p, \ldots$, and on a frame with $n$ worlds the first $n$ conjuncts suffice. The enumeration in Appendix~\ref{app:listing} uses this fact: the program computes $\Rp$ by closure, which on frames with at most four worlds coincides with the fourfold iteration.
\end{reply}

\begin{objection}
The S5 objection was never about frames. Those who reject S5 for metaphysical modality, such as \citet{Salmon1989}, reject transitivity and keep reflexivity; the relevant system is a reflexive non-transitive one, and the question is whether the argument survives there.
\end{objection}

\begin{reply}
The results answer that question. None of the six conditions of Section~\ref{s:corr} is implied by reflexivity, and reflexivity plays no role in any of the proofs. Among reflexive non-transitive frames, the strengthened argument is valid on exactly the retraceable ones (F3.93 is an example), and the original argument on none that are not symmetric and uniform. Whether metaphysical accessibility is retraceable when it is not transitive is a question about tolerance in origins that our results do not settle; what they settle is that this, and nothing else, is what the argument needs.
\end{reply}

\begin{objection}
The reduction to a single being (Lemma~\ref{l:reduction}) depends on the constant domain of REL-01-a, which REL-01-g criticizes. On variable domains the frame conditions might differ.
\end{objection}

\begin{reply}
They might, and we have not computed them. Two remarks. The sentential version of REL-01-g avoids the domain question by treating \enquote{a maximally great being exists} as an unanalysed sentence $G$; for that version the frame conditions are those of the propositional schemas with Anselm's principle as an explicit premise, and our Propositions~\ref{p:s1}--\ref{p:s6} apply directly. And on variable domains the existence predicate already appears inside $M$, so that the reduction fails only if quantifiers range over different beings at different worlds; the additional premise that REL-01-a, Chapter~6, introduces for variable domains is designed to close exactly that gap.
\end{reply}

\section{What the results do not establish}\label{s:not}

The results concern validity, not soundness. They do not show that Principle~\ref{pr:strong}, or Assumption~3.1 of REL-01-a, is true; that remains the task of Volume~III. They do not show that metaphysical accessibility is symmetric, retraceable or uniform. They establish frame correspondence for six schemas, not completeness: we have not axiomatized the logic of uniform frames or of retraceable frames, and the latter is not finitely axiomatizable in the basic modal language for the reason given in the reply to the third objection. The enumeration is bounded at four worlds, but every claim the manuscript makes is proved for all frames; the enumeration serves as a check on the proofs and as a catalogue of counterexamples, not as their ground.

\section{Conclusion}

The conclusions of the main argument of REL-01-a do not need the S5 axiom. With the original possibility premise they follow on exactly the symmetric uniform frames, a class strictly larger than the equivalence relations and containing frames that are neither reflexive nor transitive. With the stronger premise that possibly some being is ancestrally maximally great, they follow on exactly the retraceable frames, and so in KB, KDB and KTB. The substitute principle that does the work of transitivity, Anselm's principle for ancestral maximal greatness, holds on every frame. What no version can do without is the possibility of returning from the world that witnesses the possibility premise to the actual world.

\begin{thebibliography}{99}\small
\bibitem[Benzmüller and Woltzenlogel~Paleo(2014)]{Benzmuller2014} Benzmüller, C. and Woltzenlogel Paleo, B. (2014). Automating Gödel's ontological proof of God's existence with higher-order automated theorem provers. In \emph{Proceedings of the 21st European Conference on Artificial Intelligence (ECAI 2014)} (pp.~93--98). IOS Press.
\bibitem[Blackburn et~al.(2001)]{Blackburn2001} Blackburn, P., de~Rijke, M. and Venema, Y. (2001). \emph{Modal Logic}. Cambridge University Press.
\bibitem[Chandler(1976)]{Chandler1976} Chandler, H.~S. (1976). Plantinga and the contingently possible. \emph{Analysis}, 36(2), 106--109.
\bibitem[Hartshorne(1962)]{Hartshorne1962} Hartshorne, C. (1962). \emph{The Logic of Perfection}. Open Court.
\bibitem[Hughes and Cresswell(1996)]{HughesCresswell1996} Hughes, G.~E. and Cresswell, M.~J. (1996). \emph{A New Introduction to Modal Logic}. Routledge.
\bibitem[Kripke(1963)]{Kripke1963} Kripke, S.~A. (1963). Semantical analysis of modal logic I: Normal modal propositional calculi. \emph{Zeitschrift für mathematische Logik und Grundlagen der Mathematik}, 9, 67--96.
\bibitem[Malcolm(1960)]{Malcolm1960} Malcolm, N. (1960). Anselm's ontological arguments. \emph{The Philosophical Review}, 69(1), 41--62.
\bibitem[Oppy(1995)]{Oppy1995} Oppy, G. (1995). \emph{Ontological Arguments and Belief in God}. Cambridge University Press.
\bibitem[Plantinga(1974)]{Plantinga1974} Plantinga, A. (1974). \emph{The Nature of Necessity}. Clarendon Press.
\bibitem[Salmon(1989)]{Salmon1989} Salmon, N. (1989). The logic of what might have been. \emph{The Philosophical Review}, 98(1), 3--34.
\bibitem[Williamson(2013)]{Williamson2013} Williamson, T. (2013). \emph{Modal Logic as Metaphysics}. Oxford University Press.
\end{thebibliography}

\clearpage
\appendix
\section{The enumeration of frames}\label{app:method}

\subsection{What is enumerated}

The program \texttt{gen/generate.py} (Python~3, standard library only) generates every binary relation on $\{0, \ldots, n-1\}$ for $n = 1, \ldots, 4$, encodes it as an $n^2$-bit mask, and keeps exactly one representative of each isomorphism class: the relation whose mask is least among the masks of its images under all $n!$ permutations of the worlds. Representatives are numbered by size and then by number of arrows; the identifier F$n$.$k$ is the $k$-th class with $n$ worlds. Since every relation is generated and every class is represented by its least member, the enumeration is exhaustive up to isomorphism; the counts \stNOne, \stNTwo, \stNThree, \stNFour{} agree with the known numbers of unlabelled directed graphs with loops allowed on one to four vertices.

\subsection{How each entry is checked}

For each frame the program computes $R$, its transitive closure $\Rp$ (by Warshall's algorithm), the sets of worlds reachable in at most two and three steps, and the ten conditions of Definitions~\ref{d:std}, \ref{d:conds} and~\ref{d:ck}. It then evaluates each schema S1--S8 at every world under every valuation of $p$, $2^n$ valuations per frame and \stValuations{} in all, and records the schema as valid if no world and valuation falsify it, together with the first falsifying world and valuation otherwise. The quantified forms are evaluated over a constant domain of two beings under all $2^{2n}$ assignments of the profile $M$, \stQValuations{} assignments in all. The program halts with an error if any verdict disagrees with the condition paired with the schema in Sections~\ref{s:corr} and~\ref{s:finite}, if a quantified verdict disagrees with the propositional one, or if one of the implications of Lemma~\ref{l:inclusions} fails. It also verifies that ancestral uniformity coincides with the condition that every successor of a successor of $w$ is reachable from every successor of $w$, an equivalent form used in checking Proposition~\ref{p:s6}. The run completes without error.

\subsection{Summary statistics}

Of the \stTotal{} frames, \stNumB{} are symmetric, \stNumC{} retraceable, \stNumU{} uniform, \stNumFive{} Euclidean, \stNumFour{} transitive and \stNumSFive{} are equivalence relations. There are \stNumCnotB{} retraceable frames that are not symmetric, \stNumUnotFive{} uniform frames that are not Euclidean, \stNumBnotU{} symmetric frames that are not uniform, and \stNumBnotFour{} symmetric frames that are not transitive. On \stNumSOneNotSThree{} frames the original premise yields actual existence but not necessary maximal greatness, and on \stNumSThreeNotSOne{} the reverse. The original argument (S1 and S3) is valid on \stNumArgOrig{} frames and the strengthened argument (S5 and S6) on \stNumArgStar, exactly the retraceable frames. Table~A.1 breaks the counts down by the combination of standard conditions a frame satisfies.

\begin{center}\small Table A.1: Frames by combination of standard conditions (a dash marks a failed condition), with the number satisfying U and C and the number validating each schema.\end{center}
% generated by gen/generate.py -- do not edit
{\footnotesize\setlength{\tabcolsep}{3.5pt}\begin{longtable}{@{}lrrrrrrrrrrrrrrr@{}}\toprule
TB45D & $n{=}1$ & $2$ & $3$ & $4$ & all & U & C & S1 & S2 & S3 & S4 & S5 & S6 & S7 & S8\\\midrule\endhead
\texttt{TB45D} & 1 & 2 & 3 & 5 & 11 & 11 & 11 & 11 & 11 & 11 & 11 & 11 & 11 & 11 & 11\\
\texttt{-{}-{}45D} & 0 & 1 & 3 & 8 & 12 & 12 & 0 & 0 & 12 & 12 & 12 & 0 & 12 & 0 & 0\\
\texttt{-{}B45-{}} & 1 & 2 & 4 & 7 & 14 & 14 & 14 & 14 & 14 & 14 & 14 & 14 & 14 & 14 & 14\\
\texttt{T-{}4-{}D} & 0 & 1 & 6 & 28 & 35 & 0 & 0 & 0 & 0 & 0 & 35 & 0 & 0 & 0 & 0\\
\texttt{TB-{}-{}D} & 0 & 0 & 1 & 6 & 7 & 0 & 7 & 7 & 0 & 0 & 0 & 7 & 7 & 6 & 7\\
\texttt{-{}-{}-{}5D} & 0 & 0 & 1 & 6 & 7 & 7 & 0 & 0 & 7 & 7 & 0 & 0 & 7 & 0 & 0\\
\texttt{-{}-{}4-{}D} & 0 & 0 & 5 & 52 & 57 & 6 & 0 & 0 & 0 & 6 & 57 & 0 & 6 & 0 & 0\\
\texttt{-{}-{}45-{}} & 0 & 0 & 1 & 4 & 5 & 5 & 0 & 0 & 5 & 5 & 5 & 0 & 5 & 0 & 0\\
\texttt{-{}B-{}-{}D} & 0 & 2 & 10 & 59 & 71 & 9 & 71 & 71 & 0 & 9 & 0 & 71 & 71 & 62 & 71\\
\texttt{T-{}-{}-{}D} & 0 & 0 & 6 & 179 & 185 & 0 & 83 & 0 & 0 & 0 & 0 & 83 & 83 & 44 & 83\\
\texttt{-{}-{}-{}-{}D} & 0 & 0 & 35 & 1997 & 2032 & 26 & 1075 & 0 & 0 & 26 & 0 & 1075 & 1309 & 441 & 1043\\
\texttt{-{}-{}-{}5-{}} & 0 & 0 & 0 & 1 & 1 & 1 & 0 & 0 & 1 & 1 & 0 & 0 & 1 & 0 & 0\\
\texttt{-{}-{}4-{}-{}} & 0 & 2 & 17 & 138 & 157 & 20 & 0 & 0 & 0 & 20 & 157 & 0 & 20 & 0 & 0\\
\texttt{-{}B-{}-{}-{}} & 0 & 0 & 2 & 13 & 15 & 4 & 15 & 15 & 0 & 4 & 0 & 15 & 15 & 15 & 15\\
\texttt{-{}-{}-{}-{}-{}} & 0 & 0 & 10 & 541 & 551 & 12 & 20 & 0 & 0 & 12 & 0 & 20 & 36 & 16 & 20\\
\bottomrule\end{longtable}}


\section{Listing of all frames with at most four worlds}\label{app:listing}

Each line gives: the identifier; the frame as successor sets of worlds $0, 1, \ldots$ separated by dots (\texttt{-} for the empty set); the standard conditions T, B, 4, 5, D, with a dash for each that fails; the conditions U, U$^+$ and C (header \texttt{UPC}, with U$^+$ written \texttt{P}), with a dash for each that fails; the verdicts on S1--S8 in order (\texttt{+} valid, \texttt{-} not valid); and, in the last column, \texttt{ok} if the strengthened argument is valid on the frame, or otherwise a falsifying witness $w$\texttt{:}$b_0b_1\ldots$, where $w$ is a world and $b_i$ is $1$ if $p$ is true at world $i$. An unprimed witness falsifies S5 and a primed one (where S5 holds) falsifies S6.

% generated by gen/generate.py -- do not edit
\begingroup\fontsize{6.6}{7.7}\selectfont\setlength{\columnsep}{12pt}\setlength{\parindent}{0pt}
\subsection*{Frames with 1 world (2 isomorphism classes)}
\begin{multicols}{2}\raggedright
\texttt{id~~~~~~successors~TB45D~UPC~S1..S8~~~w*}\par
\texttt{F1.1~~~~-{}~~~~~~~~~~-{}B45-{}~UPC~++++++++~ok}\par
\texttt{F1.2~~~~0~~~~~~~~~~TB45D~UPC~++++++++~ok}\par
\end{multicols}
\subsection*{Frames with 2 worlds (10 isomorphism classes)}
\begin{multicols}{2}\raggedright
\texttt{id~~~~~~successors~TB45D~UPC~S1..S8~~~w*}\par
\texttt{F2.1~~~~-{}.-{}~~~~~~~~-{}B45-{}~UPC~++++++++~ok}\par
\texttt{F2.2~~~~0.-{}~~~~~~~~-{}B45-{}~UPC~++++++++~ok}\par
\texttt{F2.3~~~~1.-{}~~~~~~~~-{}-{}4-{}-{}~UP-{}~-{}-{}++-{}+-{}-{}~0:00}\par
\texttt{F2.4~~~~01.-{}~~~~~~~-{}-{}4-{}-{}~-{}-{}-{}~-{}-{}-{}+-{}-{}-{}-{}~0:00}\par
\texttt{F2.5~~~~0.0~~~~~~~~-{}-{}45D~UP-{}~-{}+++-{}+-{}-{}~1:10}\par
\texttt{F2.6~~~~1.0~~~~~~~~-{}B-{}-{}D~UPC~+-{}+-{}++++~ok}\par
\texttt{F2.7~~~~0.1~~~~~~~~TB45D~UPC~++++++++~ok}\par
\texttt{F2.8~~~~01.0~~~~~~~-{}B-{}-{}D~-{}PC~+-{}-{}-{}++++~ok}\par
\texttt{F2.9~~~~01.1~~~~~~~T-{}4-{}D~-{}-{}-{}~-{}-{}-{}+-{}-{}-{}-{}~0:01}\par
\texttt{F2.10~~~01.01~~~~~~TB45D~UPC~++++++++~ok}\par
\end{multicols}
\subsection*{Frames with 3 worlds (104 isomorphism classes)}
\begin{multicols}{2}\raggedright
\texttt{id~~~~~~successors~~TB45D~UPC~S1..S8~~~w*}\par
\texttt{F3.1~~~~-{}.-{}.-{}~~~~~~~-{}B45-{}~UPC~++++++++~ok}\par
\texttt{F3.2~~~~0.-{}.-{}~~~~~~~-{}B45-{}~UPC~++++++++~ok}\par
\texttt{F3.3~~~~1.-{}.-{}~~~~~~~-{}-{}4-{}-{}~UP-{}~-{}-{}++-{}+-{}-{}~0:000}\par
\texttt{F3.4~~~~01.-{}.-{}~~~~~~-{}-{}4-{}-{}~-{}-{}-{}~-{}-{}-{}+-{}-{}-{}-{}~0:000}\par
\texttt{F3.5~~~~12.-{}.-{}~~~~~~-{}-{}4-{}-{}~UP-{}~-{}-{}++-{}+-{}-{}~0:000}\par
\texttt{F3.6~~~~0.0.-{}~~~~~~~-{}-{}45-{}~UP-{}~-{}+++-{}+-{}-{}~1:100}\par
\texttt{F3.7~~~~1.0.-{}~~~~~~~-{}B-{}-{}-{}~UPC~+-{}+-{}++++~ok}\par
\texttt{F3.8~~~~2.0.-{}~~~~~~~-{}-{}-{}-{}-{}~UP-{}~-{}-{}+-{}-{}+-{}-{}~0:000}\par
\texttt{F3.9~~~~0.1.-{}~~~~~~~-{}B45-{}~UPC~++++++++~ok}\par
\texttt{F3.10~~~2.1.-{}~~~~~~~-{}-{}4-{}-{}~UP-{}~-{}-{}++-{}+-{}-{}~0:000}\par
\texttt{F3.11~~~2.2.-{}~~~~~~~-{}-{}4-{}-{}~UP-{}~-{}-{}++-{}+-{}-{}~0:000}\par
\texttt{F3.12~~~012.-{}.-{}~~~~~-{}-{}4-{}-{}~-{}-{}-{}~-{}-{}-{}+-{}-{}-{}-{}~0:000}\par
\texttt{F3.13~~~01.0.-{}~~~~~~-{}B-{}-{}-{}~-{}PC~+-{}-{}-{}++++~ok}\par
\texttt{F3.14~~~02.0.-{}~~~~~~-{}-{}-{}-{}-{}~-{}-{}-{}~-{}-{}-{}-{}-{}-{}-{}-{}~0:000}\par
\texttt{F3.15~~~12.0.-{}~~~~~~-{}-{}-{}-{}-{}~-{}-{}-{}~-{}-{}-{}-{}-{}-{}-{}-{}~0:000}\par
\texttt{F3.16~~~01.1.-{}~~~~~~-{}-{}4-{}-{}~-{}-{}-{}~-{}-{}-{}+-{}-{}-{}-{}~0:010}\par
\texttt{F3.17~~~02.1.-{}~~~~~~-{}-{}4-{}-{}~-{}-{}-{}~-{}-{}-{}+-{}-{}-{}-{}~0:000}\par
\texttt{F3.18~~~12.1.-{}~~~~~~-{}-{}4-{}-{}~-{}-{}-{}~-{}-{}-{}+-{}-{}-{}-{}~0:000}\par
\texttt{F3.19~~~2.01.-{}~~~~~~-{}-{}-{}-{}-{}~-{}-{}-{}~-{}-{}-{}-{}-{}-{}-{}-{}~0:000}\par
\texttt{F3.20~~~02.2.-{}~~~~~~-{}-{}4-{}-{}~-{}-{}-{}~-{}-{}-{}+-{}-{}-{}-{}~0:000}\par
\texttt{F3.21~~~12.2.-{}~~~~~~-{}-{}4-{}-{}~-{}-{}-{}~-{}-{}-{}+-{}-{}-{}-{}~0:000}\par
\texttt{F3.22~~~0.0.0~~~~~~~-{}-{}45D~UP-{}~-{}+++-{}+-{}-{}~1:100}\par
\texttt{F3.23~~~1.0.0~~~~~~~-{}-{}-{}-{}D~UP-{}~-{}-{}+-{}-{}+-{}-{}~2:110}\par
\texttt{F3.24~~~0.1.0~~~~~~~-{}-{}45D~UP-{}~-{}+++-{}+-{}-{}~2:100}\par
\texttt{F3.25~~~1.1.0~~~~~~~-{}-{}-{}-{}D~UP-{}~-{}-{}+-{}-{}+-{}-{}~0:010}\par
\texttt{F3.26~~~2.1.0~~~~~~~-{}B-{}-{}D~UPC~+-{}+-{}++++~ok}\par
\texttt{F3.27~~~1.2.0~~~~~~~-{}-{}-{}-{}D~UPC~-{}-{}+-{}++++~ok}\par
\texttt{F3.28~~~0.1.2~~~~~~~TB45D~UPC~++++++++~ok}\par
\texttt{F3.29~~~012.0.-{}~~~~~-{}-{}-{}-{}-{}~-{}-{}-{}~-{}-{}-{}-{}-{}-{}-{}-{}~0:000}\par
\texttt{F3.30~~~012.1.-{}~~~~~-{}-{}4-{}-{}~-{}-{}-{}~-{}-{}-{}+-{}-{}-{}-{}~0:000}\par
\texttt{F3.31~~~01.01.-{}~~~~~-{}B45-{}~UPC~++++++++~ok}\par
\texttt{F3.32~~~02.01.-{}~~~~~-{}-{}-{}-{}-{}~-{}-{}-{}~-{}-{}-{}-{}-{}-{}-{}-{}~0:000}\par
\texttt{F3.33~~~12.01.-{}~~~~~-{}-{}-{}-{}-{}~-{}-{}-{}~-{}-{}-{}-{}-{}-{}-{}-{}~0:000}\par
\texttt{F3.34~~~012.2.-{}~~~~~-{}-{}4-{}-{}~-{}-{}-{}~-{}-{}-{}+-{}-{}-{}-{}~0:000}\par
\texttt{F3.35~~~02.02.-{}~~~~~-{}-{}4-{}-{}~-{}-{}-{}~-{}-{}-{}+-{}-{}-{}-{}~0:000}\par
\texttt{F3.36~~~12.02.-{}~~~~~-{}-{}-{}-{}-{}~-{}-{}-{}~-{}-{}-{}-{}-{}-{}-{}-{}~0:000}\par
\texttt{F3.37~~~02.12.-{}~~~~~-{}-{}4-{}-{}~-{}-{}-{}~-{}-{}-{}+-{}-{}-{}-{}~0:000}\par
\texttt{F3.38~~~01.0.0~~~~~~-{}-{}-{}-{}D~-{}P-{}~-{}-{}-{}-{}-{}+-{}-{}~2:110}\par
\texttt{F3.39~~~12.0.0~~~~~~-{}B-{}-{}D~UPC~+-{}+-{}++++~ok}\par
\texttt{F3.40~~~01.1.0~~~~~~-{}-{}-{}-{}D~-{}-{}-{}~-{}-{}-{}-{}-{}-{}-{}-{}~0:010}\par
\texttt{F3.41~~~02.1.0~~~~~~-{}B-{}-{}D~-{}PC~+-{}-{}-{}++++~ok}\par
\texttt{F3.42~~~12.1.0~~~~~~-{}-{}-{}-{}D~-{}-{}-{}~-{}-{}-{}-{}-{}-{}-{}-{}~0:010}\par
\texttt{F3.43~~~0.01.0~~~~~~-{}-{}4-{}D~-{}-{}-{}~-{}-{}-{}+-{}-{}-{}-{}~1:100}\par
\texttt{F3.44~~~1.01.0~~~~~~-{}-{}-{}-{}D~-{}P-{}~-{}-{}-{}-{}-{}+-{}-{}~2:110}\par
\texttt{F3.45~~~2.01.0~~~~~~-{}-{}-{}-{}D~-{}-{}-{}~-{}-{}-{}-{}-{}-{}-{}-{}~1:101}\par
\texttt{F3.46~~~01.2.0~~~~~~-{}-{}-{}-{}D~-{}PC~-{}-{}-{}-{}++-{}+~ok}\par
\texttt{F3.47~~~12.2.0~~~~~~-{}-{}-{}-{}D~-{}PC~-{}-{}-{}-{}++++~ok}\par
\texttt{F3.48~~~0.02.0~~~~~~-{}-{}4-{}D~UP-{}~-{}-{}++-{}+-{}-{}~1:100}\par
\texttt{F3.49~~~2.02.0~~~~~~-{}-{}-{}-{}D~-{}P-{}~-{}-{}-{}-{}-{}+-{}-{}~1:101}\par
\texttt{F3.50~~~0.12.0~~~~~~-{}-{}-{}-{}D~-{}-{}-{}~-{}-{}-{}-{}-{}-{}-{}-{}~1:100}\par
\texttt{F3.51~~~0.1.01~~~~~~-{}-{}4-{}D~-{}-{}-{}~-{}-{}-{}+-{}-{}-{}-{}~2:100}\par
\texttt{F3.52~~~01.1.2~~~~~~T-{}4-{}D~-{}-{}-{}~-{}-{}-{}+-{}-{}-{}-{}~0:010}\par
\texttt{F3.53~~~012.01.-{}~~~~-{}-{}-{}-{}-{}~-{}-{}-{}~-{}-{}-{}-{}-{}-{}-{}-{}~0:000}\par
\texttt{F3.54~~~012.02.-{}~~~~-{}-{}-{}-{}-{}~-{}-{}-{}~-{}-{}-{}-{}-{}-{}-{}-{}~0:000}\par
\texttt{F3.55~~~012.12.-{}~~~~-{}-{}4-{}-{}~-{}-{}-{}~-{}-{}-{}+-{}-{}-{}-{}~0:000}\par
\texttt{F3.56~~~012.0.0~~~~~-{}B-{}-{}D~-{}PC~+-{}-{}-{}++++~ok}\par
\texttt{F3.57~~~012.1.0~~~~~-{}-{}-{}-{}D~-{}-{}-{}~-{}-{}-{}-{}-{}-{}-{}-{}~0:010}\par
\texttt{F3.58~~~01.01.0~~~~~-{}-{}-{}5D~UP-{}~-{}++-{}-{}+-{}-{}~2:110}\par
\texttt{F3.59~~~02.01.0~~~~~-{}-{}-{}-{}D~-{}-{}-{}~-{}-{}-{}-{}-{}-{}-{}-{}~1:101}\par
\texttt{F3.60~~~12.01.0~~~~~-{}B-{}-{}D~-{}PC~+-{}-{}-{}++++~ok}\par
\texttt{F3.61~~~012.2.0~~~~~-{}-{}-{}-{}D~-{}PC~-{}-{}-{}-{}++-{}+~ok}\par
\texttt{F3.62~~~01.02.0~~~~~-{}-{}-{}-{}D~-{}PC~-{}-{}-{}-{}++++~ok}\par
\texttt{F3.63~~~02.02.0~~~~~-{}-{}-{}-{}D~-{}P-{}~-{}-{}-{}-{}-{}+-{}-{}~1:101}\par
\texttt{F3.64~~~12.02.0~~~~~-{}-{}-{}-{}D~-{}PC~-{}-{}-{}-{}++++~ok}\par
\texttt{F3.65~~~01.12.0~~~~~-{}-{}-{}-{}D~-{}PC~-{}-{}-{}-{}++-{}+~ok}\par
\texttt{F3.66~~~02.12.0~~~~~-{}-{}-{}-{}D~-{}-{}-{}~-{}-{}-{}-{}-{}-{}-{}-{}~1:101}\par
\texttt{F3.67~~~12.12.0~~~~~-{}-{}-{}-{}D~-{}PC~-{}-{}-{}-{}++++~ok}\par
\texttt{F3.68~~~0.012.0~~~~~-{}-{}4-{}D~-{}-{}-{}~-{}-{}-{}+-{}-{}-{}-{}~1:100}\par
\texttt{F3.69~~~2.012.0~~~~~-{}-{}-{}-{}D~-{}-{}-{}~-{}-{}-{}-{}-{}-{}-{}-{}~1:101}\par
\texttt{F3.70~~~01.1.01~~~~~-{}-{}4-{}D~-{}-{}-{}~-{}-{}-{}+-{}-{}-{}-{}~0:010}\par
\texttt{F3.71~~~02.1.01~~~~~-{}-{}-{}-{}D~-{}-{}-{}~-{}-{}-{}-{}-{}-{}-{}-{}~2:010}\par
\texttt{F3.72~~~12.1.01~~~~~-{}-{}-{}-{}D~-{}-{}-{}~-{}-{}-{}-{}-{}-{}-{}-{}~0:010}\par
\texttt{F3.73~~~012.1.2~~~~~T-{}4-{}D~-{}-{}-{}~-{}-{}-{}+-{}-{}-{}-{}~0:010}\par
\texttt{F3.74~~~01.01.2~~~~~TB45D~UPC~++++++++~ok}\par
\texttt{F3.75~~~02.01.2~~~~~T-{}-{}-{}D~-{}-{}-{}~-{}-{}-{}-{}-{}-{}-{}-{}~0:001}\par
\texttt{F3.76~~~02.12.2~~~~~T-{}4-{}D~-{}-{}-{}~-{}-{}-{}+-{}-{}-{}-{}~0:001}\par
\texttt{F3.77~~~012.012.-{}~~~-{}-{}4-{}-{}~-{}-{}-{}~-{}-{}-{}+-{}-{}-{}-{}~0:000}\par
\texttt{F3.78~~~012.01.0~~~~-{}B-{}-{}D~-{}PC~+-{}-{}-{}++++~ok}\par
\texttt{F3.79~~~012.02.0~~~~-{}-{}-{}-{}D~-{}PC~-{}-{}-{}-{}++++~ok}\par
\texttt{F3.80~~~012.12.0~~~~-{}-{}-{}-{}D~-{}PC~-{}-{}-{}-{}++++~ok}\par
\texttt{F3.81~~~01.012.0~~~~-{}-{}-{}-{}D~-{}PC~-{}-{}-{}-{}++-{}+~ok}\par
\texttt{F3.82~~~02.012.0~~~~-{}-{}-{}-{}D~-{}-{}-{}~-{}-{}-{}-{}-{}-{}-{}-{}~1:101}\par
\texttt{F3.83~~~12.012.0~~~~-{}-{}-{}-{}D~-{}PC~-{}-{}-{}-{}++++~ok}\par
\texttt{F3.84~~~012.1.01~~~~-{}-{}-{}-{}D~-{}-{}-{}~-{}-{}-{}-{}-{}-{}-{}-{}~0:010}\par
\texttt{F3.85~~~01.01.01~~~~-{}-{}45D~UP-{}~-{}+++-{}+-{}-{}~2:110}\par
\texttt{F3.86~~~02.01.01~~~~-{}-{}-{}-{}D~-{}PC~-{}-{}-{}-{}++++~ok}\par
\texttt{F3.87~~~12.01.01~~~~-{}-{}-{}-{}D~-{}PC~-{}-{}-{}-{}++++~ok}\par
\texttt{F3.88~~~12.02.01~~~~-{}B-{}-{}D~-{}PC~+-{}-{}-{}++++~ok}\par
\texttt{F3.89~~~02.12.01~~~~-{}B-{}-{}D~-{}PC~+-{}-{}-{}++++~ok}\par
\texttt{F3.90~~~012.01.2~~~~T-{}-{}-{}D~-{}-{}-{}~-{}-{}-{}-{}-{}-{}-{}-{}~0:001}\par
\texttt{F3.91~~~012.12.2~~~~T-{}4-{}D~-{}-{}-{}~-{}-{}-{}+-{}-{}-{}-{}~0:001}\par
\texttt{F3.92~~~01.01.02~~~~T-{}-{}-{}D~-{}-{}-{}~-{}-{}-{}-{}-{}-{}-{}-{}~2:110}\par
\texttt{F3.93~~~01.12.02~~~~T-{}-{}-{}D~-{}PC~-{}-{}-{}-{}++++~ok}\par
\texttt{F3.94~~~012.012.0~~~-{}-{}-{}-{}D~-{}PC~-{}-{}-{}-{}++++~ok}\par
\texttt{F3.95~~~012.01.01~~~-{}-{}-{}-{}D~-{}PC~-{}-{}-{}-{}++++~ok}\par
\texttt{F3.96~~~012.02.01~~~-{}B-{}-{}D~-{}PC~+-{}-{}-{}++++~ok}\par
\texttt{F3.97~~~012.12.01~~~-{}-{}-{}-{}D~-{}PC~-{}-{}-{}-{}++++~ok}\par
\texttt{F3.98~~~012.012.2~~~T-{}4-{}D~-{}-{}-{}~-{}-{}-{}+-{}-{}-{}-{}~0:001}\par
\texttt{F3.99~~~012.01.02~~~TB-{}-{}D~-{}PC~+-{}-{}-{}++++~ok}\par
\texttt{F3.100~~012.12.02~~~T-{}-{}-{}D~-{}PC~-{}-{}-{}-{}++++~ok}\par
\texttt{F3.101~~02.012.02~~~T-{}4-{}D~-{}-{}-{}~-{}-{}-{}+-{}-{}-{}-{}~1:101}\par
\texttt{F3.102~~012.012.01~~-{}B-{}-{}D~-{}PC~+-{}-{}-{}++++~ok}\par
\texttt{F3.103~~012.012.02~~T-{}-{}-{}D~-{}PC~-{}-{}-{}-{}++++~ok}\par
\texttt{F3.104~~012.012.012~TB45D~UPC~++++++++~ok}\par
\end{multicols}
\subsection*{Frames with 4 worlds (3044 isomorphism classes)}
\begin{multicols}{2}\raggedright
\texttt{id~~~~~~successors~~~~~~~~~~TB45D~UPC~S1..S8~~~w*}\par
\texttt{F4.1~~~~-{}.-{}.-{}.-{}~~~~~~~~~~~~~-{}B45-{}~UPC~++++++++~ok}\par
\texttt{F4.2~~~~0.-{}.-{}.-{}~~~~~~~~~~~~~-{}B45-{}~UPC~++++++++~ok}\par
\texttt{F4.3~~~~1.-{}.-{}.-{}~~~~~~~~~~~~~-{}-{}4-{}-{}~UP-{}~-{}-{}++-{}+-{}-{}~0:0000}\par
\texttt{F4.4~~~~01.-{}.-{}.-{}~~~~~~~~~~~~-{}-{}4-{}-{}~-{}-{}-{}~-{}-{}-{}+-{}-{}-{}-{}~0:0000}\par
\texttt{F4.5~~~~12.-{}.-{}.-{}~~~~~~~~~~~~-{}-{}4-{}-{}~UP-{}~-{}-{}++-{}+-{}-{}~0:0000}\par
\texttt{F4.6~~~~0.0.-{}.-{}~~~~~~~~~~~~~-{}-{}45-{}~UP-{}~-{}+++-{}+-{}-{}~1:1000}\par
\texttt{F4.7~~~~1.0.-{}.-{}~~~~~~~~~~~~~-{}B-{}-{}-{}~UPC~+-{}+-{}++++~ok}\par
\texttt{F4.8~~~~2.0.-{}.-{}~~~~~~~~~~~~~-{}-{}-{}-{}-{}~UP-{}~-{}-{}+-{}-{}+-{}-{}~0:0000}\par
\texttt{F4.9~~~~0.1.-{}.-{}~~~~~~~~~~~~~-{}B45-{}~UPC~++++++++~ok}\par
\texttt{F4.10~~~2.1.-{}.-{}~~~~~~~~~~~~~-{}-{}4-{}-{}~UP-{}~-{}-{}++-{}+-{}-{}~0:0000}\par
\texttt{F4.11~~~2.2.-{}.-{}~~~~~~~~~~~~~-{}-{}4-{}-{}~UP-{}~-{}-{}++-{}+-{}-{}~0:0000}\par
\texttt{F4.12~~~3.2.-{}.-{}~~~~~~~~~~~~~-{}-{}4-{}-{}~UP-{}~-{}-{}++-{}+-{}-{}~0:0000}\par
\texttt{F4.13~~~012.-{}.-{}.-{}~~~~~~~~~~~-{}-{}4-{}-{}~-{}-{}-{}~-{}-{}-{}+-{}-{}-{}-{}~0:0000}\par
\texttt{F4.14~~~123.-{}.-{}.-{}~~~~~~~~~~~-{}-{}4-{}-{}~UP-{}~-{}-{}++-{}+-{}-{}~0:0000}\par
\texttt{F4.15~~~01.0.-{}.-{}~~~~~~~~~~~~-{}B-{}-{}-{}~-{}PC~+-{}-{}-{}++++~ok}\par
\texttt{F4.16~~~02.0.-{}.-{}~~~~~~~~~~~~-{}-{}-{}-{}-{}~-{}-{}-{}~-{}-{}-{}-{}-{}-{}-{}-{}~0:0000}\par
\texttt{F4.17~~~12.0.-{}.-{}~~~~~~~~~~~~-{}-{}-{}-{}-{}~-{}-{}-{}~-{}-{}-{}-{}-{}-{}-{}-{}~0:0000}\par
\texttt{F4.18~~~23.0.-{}.-{}~~~~~~~~~~~~-{}-{}-{}-{}-{}~UP-{}~-{}-{}+-{}-{}+-{}-{}~0:0000}\par
\texttt{F4.19~~~01.1.-{}.-{}~~~~~~~~~~~~-{}-{}4-{}-{}~-{}-{}-{}~-{}-{}-{}+-{}-{}-{}-{}~0:0100}\par
\texttt{F4.20~~~02.1.-{}.-{}~~~~~~~~~~~~-{}-{}4-{}-{}~-{}-{}-{}~-{}-{}-{}+-{}-{}-{}-{}~0:0000}\par
\texttt{F4.21~~~12.1.-{}.-{}~~~~~~~~~~~~-{}-{}4-{}-{}~-{}-{}-{}~-{}-{}-{}+-{}-{}-{}-{}~0:0000}\par
\texttt{F4.22~~~23.1.-{}.-{}~~~~~~~~~~~~-{}-{}4-{}-{}~UP-{}~-{}-{}++-{}+-{}-{}~0:0000}\par
\texttt{F4.23~~~2.01.-{}.-{}~~~~~~~~~~~~-{}-{}-{}-{}-{}~-{}-{}-{}~-{}-{}-{}-{}-{}-{}-{}-{}~0:0000}\par
\texttt{F4.24~~~02.2.-{}.-{}~~~~~~~~~~~~-{}-{}4-{}-{}~-{}-{}-{}~-{}-{}-{}+-{}-{}-{}-{}~0:0000}\par
\texttt{F4.25~~~12.2.-{}.-{}~~~~~~~~~~~~-{}-{}4-{}-{}~-{}-{}-{}~-{}-{}-{}+-{}-{}-{}-{}~0:0000}\par
\texttt{F4.26~~~03.2.-{}.-{}~~~~~~~~~~~~-{}-{}4-{}-{}~-{}-{}-{}~-{}-{}-{}+-{}-{}-{}-{}~0:0000}\par
\texttt{F4.27~~~13.2.-{}.-{}~~~~~~~~~~~~-{}-{}-{}-{}-{}~-{}-{}-{}~-{}-{}-{}-{}-{}-{}-{}-{}~0:0000}\par
\texttt{F4.28~~~23.2.-{}.-{}~~~~~~~~~~~~-{}-{}4-{}-{}~UP-{}~-{}-{}++-{}+-{}-{}~0:0000}\par
\texttt{F4.29~~~0.0.0.-{}~~~~~~~~~~~~~-{}-{}45-{}~UP-{}~-{}+++-{}+-{}-{}~1:1000}\par
\texttt{F4.30~~~1.0.0.-{}~~~~~~~~~~~~~-{}-{}-{}-{}-{}~UP-{}~-{}-{}+-{}-{}+-{}-{}~2:1100}\par
\texttt{F4.31~~~3.0.0.-{}~~~~~~~~~~~~~-{}-{}-{}-{}-{}~UP-{}~-{}-{}+-{}-{}+-{}-{}~0:0000}\par
\texttt{F4.32~~~0.1.0.-{}~~~~~~~~~~~~~-{}-{}45-{}~UP-{}~-{}+++-{}+-{}-{}~2:1000}\par
\texttt{F4.33~~~1.1.0.-{}~~~~~~~~~~~~~-{}-{}-{}-{}-{}~UP-{}~-{}-{}+-{}-{}+-{}-{}~0:0100}\par
\texttt{F4.34~~~2.1.0.-{}~~~~~~~~~~~~~-{}B-{}-{}-{}~UPC~+-{}+-{}++++~ok}\par
\texttt{F4.35~~~3.1.0.-{}~~~~~~~~~~~~~-{}-{}-{}-{}-{}~UP-{}~-{}-{}+-{}-{}+-{}-{}~0:0000}\par
\texttt{F4.36~~~1.2.0.-{}~~~~~~~~~~~~~-{}-{}-{}-{}-{}~UPC~-{}-{}+-{}++++~ok}\par
\texttt{F4.37~~~3.2.0.-{}~~~~~~~~~~~~~-{}-{}-{}-{}-{}~UP-{}~-{}-{}+-{}-{}+-{}-{}~0:0000}\par
\texttt{F4.38~~~0.3.0.-{}~~~~~~~~~~~~~-{}-{}4-{}-{}~UP-{}~-{}-{}++-{}+-{}-{}~1:0000}\par
\texttt{F4.39~~~2.3.0.-{}~~~~~~~~~~~~~-{}-{}-{}-{}-{}~UP-{}~-{}-{}+-{}-{}+-{}-{}~1:0000}\par
\texttt{F4.40~~~3.3.0.-{}~~~~~~~~~~~~~-{}-{}-{}-{}-{}~UP-{}~-{}-{}+-{}-{}+-{}-{}~0:0000}\par
\texttt{F4.41~~~0.1.2.-{}~~~~~~~~~~~~~-{}B45-{}~UPC~++++++++~ok}\par
\texttt{F4.42~~~3.1.2.-{}~~~~~~~~~~~~~-{}-{}4-{}-{}~UP-{}~-{}-{}++-{}+-{}-{}~0:0000}\par
\texttt{F4.43~~~3.3.2.-{}~~~~~~~~~~~~~-{}-{}4-{}-{}~UP-{}~-{}-{}++-{}+-{}-{}~0:0000}\par
\texttt{F4.44~~~3.3.3.-{}~~~~~~~~~~~~~-{}-{}4-{}-{}~UP-{}~-{}-{}++-{}+-{}-{}~0:0000}\par
\texttt{F4.45~~~0123.-{}.-{}.-{}~~~~~~~~~~-{}-{}4-{}-{}~-{}-{}-{}~-{}-{}-{}+-{}-{}-{}-{}~0:0000}\par
\texttt{F4.46~~~012.0.-{}.-{}~~~~~~~~~~~-{}-{}-{}-{}-{}~-{}-{}-{}~-{}-{}-{}-{}-{}-{}-{}-{}~0:0000}\par
\texttt{F4.47~~~023.0.-{}.-{}~~~~~~~~~~~-{}-{}-{}-{}-{}~-{}-{}-{}~-{}-{}-{}-{}-{}-{}-{}-{}~0:0000}\par
\texttt{F4.48~~~123.0.-{}.-{}~~~~~~~~~~~-{}-{}-{}-{}-{}~-{}-{}-{}~-{}-{}-{}-{}-{}-{}-{}-{}~0:0000}\par
\texttt{F4.49~~~012.1.-{}.-{}~~~~~~~~~~~-{}-{}4-{}-{}~-{}-{}-{}~-{}-{}-{}+-{}-{}-{}-{}~0:0000}\par
\texttt{F4.50~~~023.1.-{}.-{}~~~~~~~~~~~-{}-{}4-{}-{}~-{}-{}-{}~-{}-{}-{}+-{}-{}-{}-{}~0:0000}\par
\texttt{F4.51~~~123.1.-{}.-{}~~~~~~~~~~~-{}-{}4-{}-{}~-{}-{}-{}~-{}-{}-{}+-{}-{}-{}-{}~0:0000}\par
\texttt{F4.52~~~01.01.-{}.-{}~~~~~~~~~~~-{}B45-{}~UPC~++++++++~ok}\par
\texttt{F4.53~~~02.01.-{}.-{}~~~~~~~~~~~-{}-{}-{}-{}-{}~-{}-{}-{}~-{}-{}-{}-{}-{}-{}-{}-{}~0:0000}\par
\texttt{F4.54~~~12.01.-{}.-{}~~~~~~~~~~~-{}-{}-{}-{}-{}~-{}-{}-{}~-{}-{}-{}-{}-{}-{}-{}-{}~0:0000}\par
\texttt{F4.55~~~23.01.-{}.-{}~~~~~~~~~~~-{}-{}-{}-{}-{}~-{}-{}-{}~-{}-{}-{}-{}-{}-{}-{}-{}~0:0000}\par
\texttt{F4.56~~~012.2.-{}.-{}~~~~~~~~~~~-{}-{}4-{}-{}~-{}-{}-{}~-{}-{}-{}+-{}-{}-{}-{}~0:0000}\par
\texttt{F4.57~~~013.2.-{}.-{}~~~~~~~~~~~-{}-{}-{}-{}-{}~-{}-{}-{}~-{}-{}-{}-{}-{}-{}-{}-{}~0:0000}\par
\texttt{F4.58~~~023.2.-{}.-{}~~~~~~~~~~~-{}-{}4-{}-{}~-{}-{}-{}~-{}-{}-{}+-{}-{}-{}-{}~0:0000}\par
\texttt{F4.59~~~123.2.-{}.-{}~~~~~~~~~~~-{}-{}4-{}-{}~-{}-{}-{}~-{}-{}-{}+-{}-{}-{}-{}~0:0000}\par
\texttt{F4.60~~~02.02.-{}.-{}~~~~~~~~~~~-{}-{}4-{}-{}~-{}-{}-{}~-{}-{}-{}+-{}-{}-{}-{}~0:0000}\par
\texttt{F4.61~~~12.02.-{}.-{}~~~~~~~~~~~-{}-{}-{}-{}-{}~-{}-{}-{}~-{}-{}-{}-{}-{}-{}-{}-{}~0:0000}\par
\texttt{F4.62~~~03.02.-{}.-{}~~~~~~~~~~~-{}-{}-{}-{}-{}~-{}-{}-{}~-{}-{}-{}-{}-{}-{}-{}-{}~0:0000}\par
\texttt{F4.63~~~13.02.-{}.-{}~~~~~~~~~~~-{}-{}-{}-{}-{}~-{}-{}-{}~-{}-{}-{}-{}-{}-{}-{}-{}~0:0000}\par
\texttt{F4.64~~~23.02.-{}.-{}~~~~~~~~~~~-{}-{}-{}-{}-{}~-{}-{}-{}~-{}-{}-{}-{}-{}-{}-{}-{}~0:0000}\par
\texttt{F4.65~~~02.12.-{}.-{}~~~~~~~~~~~-{}-{}4-{}-{}~-{}-{}-{}~-{}-{}-{}+-{}-{}-{}-{}~0:0000}\par
\texttt{F4.66~~~03.12.-{}.-{}~~~~~~~~~~~-{}-{}4-{}-{}~-{}-{}-{}~-{}-{}-{}+-{}-{}-{}-{}~0:0000}\par
\texttt{F4.67~~~23.12.-{}.-{}~~~~~~~~~~~-{}-{}4-{}-{}~-{}-{}-{}~-{}-{}-{}+-{}-{}-{}-{}~0:0000}\par
\texttt{F4.68~~~23.23.-{}.-{}~~~~~~~~~~~-{}-{}4-{}-{}~UP-{}~-{}-{}++-{}+-{}-{}~0:0000}\par
\texttt{F4.69~~~01.0.0.-{}~~~~~~~~~~~~-{}-{}-{}-{}-{}~-{}P-{}~-{}-{}-{}-{}-{}+-{}-{}~2:1100}\par
\texttt{F4.70~~~12.0.0.-{}~~~~~~~~~~~~-{}B-{}-{}-{}~UPC~+-{}+-{}++++~ok}\par
\texttt{F4.71~~~03.0.0.-{}~~~~~~~~~~~~-{}-{}-{}-{}-{}~-{}-{}-{}~-{}-{}-{}-{}-{}-{}-{}-{}~0:0000}\par
\texttt{F4.72~~~13.0.0.-{}~~~~~~~~~~~~-{}-{}-{}-{}-{}~-{}-{}-{}~-{}-{}-{}-{}-{}-{}-{}-{}~0:0000}\par
\texttt{F4.73~~~01.1.0.-{}~~~~~~~~~~~~-{}-{}-{}-{}-{}~-{}-{}-{}~-{}-{}-{}-{}-{}-{}-{}-{}~0:0100}\par
\texttt{F4.74~~~02.1.0.-{}~~~~~~~~~~~~-{}B-{}-{}-{}~-{}PC~+-{}-{}-{}++++~ok}\par
\texttt{F4.75~~~12.1.0.-{}~~~~~~~~~~~~-{}-{}-{}-{}-{}~-{}-{}-{}~-{}-{}-{}-{}-{}-{}-{}-{}~0:0100}\par
\texttt{F4.76~~~03.1.0.-{}~~~~~~~~~~~~-{}-{}-{}-{}-{}~-{}-{}-{}~-{}-{}-{}-{}-{}-{}-{}-{}~0:0000}\par
\texttt{F4.77~~~13.1.0.-{}~~~~~~~~~~~~-{}-{}-{}-{}-{}~-{}-{}-{}~-{}-{}-{}-{}-{}-{}-{}-{}~0:0000}\par
\texttt{F4.78~~~23.1.0.-{}~~~~~~~~~~~~-{}-{}-{}-{}-{}~-{}-{}-{}~-{}-{}-{}-{}-{}-{}-{}-{}~0:0000}\par
\texttt{F4.79~~~0.01.0.-{}~~~~~~~~~~~~-{}-{}4-{}-{}~-{}-{}-{}~-{}-{}-{}+-{}-{}-{}-{}~1:1000}\par
\texttt{F4.80~~~1.01.0.-{}~~~~~~~~~~~~-{}-{}-{}-{}-{}~-{}P-{}~-{}-{}-{}-{}-{}+-{}-{}~2:1100}\par
\texttt{F4.81~~~2.01.0.-{}~~~~~~~~~~~~-{}-{}-{}-{}-{}~-{}-{}-{}~-{}-{}-{}-{}-{}-{}-{}-{}~1:1010}\par
\texttt{F4.82~~~3.01.0.-{}~~~~~~~~~~~~-{}-{}-{}-{}-{}~-{}-{}-{}~-{}-{}-{}-{}-{}-{}-{}-{}~0:0000}\par
\texttt{F4.83~~~01.2.0.-{}~~~~~~~~~~~~-{}-{}-{}-{}-{}~-{}PC~-{}-{}-{}-{}++-{}+~ok}\par
\texttt{F4.84~~~12.2.0.-{}~~~~~~~~~~~~-{}-{}-{}-{}-{}~-{}PC~-{}-{}-{}-{}++++~ok}\par
\texttt{F4.85~~~03.2.0.-{}~~~~~~~~~~~~-{}-{}-{}-{}-{}~-{}-{}-{}~-{}-{}-{}-{}-{}-{}-{}-{}~0:0000}\par
\texttt{F4.86~~~13.2.0.-{}~~~~~~~~~~~~-{}-{}-{}-{}-{}~-{}-{}-{}~-{}-{}-{}-{}-{}-{}-{}-{}~0:0000}\par
\texttt{F4.87~~~23.2.0.-{}~~~~~~~~~~~~-{}-{}-{}-{}-{}~-{}-{}-{}~-{}-{}-{}-{}-{}-{}-{}-{}~0:0000}\par
\texttt{F4.88~~~0.02.0.-{}~~~~~~~~~~~~-{}-{}4-{}-{}~UP-{}~-{}-{}++-{}+-{}-{}~1:1000}\par
\texttt{F4.89~~~2.02.0.-{}~~~~~~~~~~~~-{}-{}-{}-{}-{}~-{}P-{}~-{}-{}-{}-{}-{}+-{}-{}~1:1010}\par
\texttt{F4.90~~~3.02.0.-{}~~~~~~~~~~~~-{}-{}-{}-{}-{}~-{}-{}-{}~-{}-{}-{}-{}-{}-{}-{}-{}~0:0000}\par
\texttt{F4.91~~~0.12.0.-{}~~~~~~~~~~~~-{}-{}-{}-{}-{}~-{}-{}-{}~-{}-{}-{}-{}-{}-{}-{}-{}~1:1000}\par
\texttt{F4.92~~~3.12.0.-{}~~~~~~~~~~~~-{}-{}-{}-{}-{}~-{}-{}-{}~-{}-{}-{}-{}-{}-{}-{}-{}~0:0000}\par
\texttt{F4.93~~~01.3.0.-{}~~~~~~~~~~~~-{}-{}-{}-{}-{}~-{}-{}-{}~-{}-{}-{}-{}-{}-{}-{}-{}~1:0000}\par
\texttt{F4.94~~~02.3.0.-{}~~~~~~~~~~~~-{}-{}-{}-{}-{}~-{}P-{}~-{}-{}-{}-{}-{}+-{}-{}~1:0000}\par
\texttt{F4.95~~~12.3.0.-{}~~~~~~~~~~~~-{}-{}-{}-{}-{}~-{}-{}-{}~-{}-{}-{}-{}-{}-{}-{}-{}~1:0000}\par
\texttt{F4.96~~~03.3.0.-{}~~~~~~~~~~~~-{}-{}-{}-{}-{}~-{}-{}-{}~-{}-{}-{}-{}-{}-{}-{}-{}~0:0000}\par
\texttt{F4.97~~~13.3.0.-{}~~~~~~~~~~~~-{}-{}-{}-{}-{}~-{}-{}-{}~-{}-{}-{}-{}-{}-{}-{}-{}~0:0000}\par
\texttt{F4.98~~~23.3.0.-{}~~~~~~~~~~~~-{}-{}-{}-{}-{}~-{}-{}-{}~-{}-{}-{}-{}-{}-{}-{}-{}~0:0000}\par
\texttt{F4.99~~~0.03.0.-{}~~~~~~~~~~~~-{}-{}4-{}-{}~-{}-{}-{}~-{}-{}-{}+-{}-{}-{}-{}~1:0000}\par
\texttt{F4.100~~2.03.0.-{}~~~~~~~~~~~~-{}-{}-{}-{}-{}~-{}-{}-{}~-{}-{}-{}-{}-{}-{}-{}-{}~1:0000}\par
\texttt{F4.101~~3.03.0.-{}~~~~~~~~~~~~-{}-{}-{}-{}-{}~-{}-{}-{}~-{}-{}-{}-{}-{}-{}-{}-{}~0:0000}\par
\texttt{F4.102~~0.13.0.-{}~~~~~~~~~~~~-{}-{}4-{}-{}~-{}-{}-{}~-{}-{}-{}+-{}-{}-{}-{}~1:0000}\par
\texttt{F4.103~~2.13.0.-{}~~~~~~~~~~~~-{}-{}-{}-{}-{}~-{}-{}-{}~-{}-{}-{}-{}-{}-{}-{}-{}~1:0000}\par
\texttt{F4.104~~3.13.0.-{}~~~~~~~~~~~~-{}-{}-{}-{}-{}~-{}-{}-{}~-{}-{}-{}-{}-{}-{}-{}-{}~0:0000}\par
\texttt{F4.105~~0.23.0.-{}~~~~~~~~~~~~-{}-{}-{}-{}-{}~-{}-{}-{}~-{}-{}-{}-{}-{}-{}-{}-{}~1:0000}\par
\texttt{F4.106~~3.23.0.-{}~~~~~~~~~~~~-{}-{}-{}-{}-{}~-{}-{}-{}~-{}-{}-{}-{}-{}-{}-{}-{}~0:0000}\par
\texttt{F4.107~~0.1.01.-{}~~~~~~~~~~~~-{}-{}4-{}-{}~-{}-{}-{}~-{}-{}-{}+-{}-{}-{}-{}~2:1000}\par
\texttt{F4.108~~3.1.01.-{}~~~~~~~~~~~~-{}-{}-{}-{}-{}~-{}-{}-{}~-{}-{}-{}-{}-{}-{}-{}-{}~0:0000}\par
\texttt{F4.109~~3.3.01.-{}~~~~~~~~~~~~-{}-{}-{}-{}-{}~UP-{}~-{}-{}+-{}-{}+-{}-{}~0:0000}\par
\texttt{F4.110~~01.1.2.-{}~~~~~~~~~~~~-{}-{}4-{}-{}~-{}-{}-{}~-{}-{}-{}+-{}-{}-{}-{}~0:0100}\par
\texttt{F4.111~~03.1.2.-{}~~~~~~~~~~~~-{}-{}4-{}-{}~-{}-{}-{}~-{}-{}-{}+-{}-{}-{}-{}~0:0000}\par
\texttt{F4.112~~13.1.2.-{}~~~~~~~~~~~~-{}-{}4-{}-{}~-{}-{}-{}~-{}-{}-{}+-{}-{}-{}-{}~0:0000}\par
\texttt{F4.113~~3.01.2.-{}~~~~~~~~~~~~-{}-{}-{}-{}-{}~-{}-{}-{}~-{}-{}-{}-{}-{}-{}-{}-{}~0:0000}\par
\texttt{F4.114~~3.12.2.-{}~~~~~~~~~~~~-{}-{}4-{}-{}~-{}-{}-{}~-{}-{}-{}+-{}-{}-{}-{}~0:0000}\par
\texttt{F4.115~~03.3.2.-{}~~~~~~~~~~~~-{}-{}4-{}-{}~-{}-{}-{}~-{}-{}-{}+-{}-{}-{}-{}~0:0000}\par
\texttt{F4.116~~13.3.2.-{}~~~~~~~~~~~~-{}-{}4-{}-{}~-{}-{}-{}~-{}-{}-{}+-{}-{}-{}-{}~0:0000}\par
\texttt{F4.117~~23.3.2.-{}~~~~~~~~~~~~-{}-{}4-{}-{}~-{}-{}-{}~-{}-{}-{}+-{}-{}-{}-{}~0:0000}\par
\texttt{F4.118~~3.3.02.-{}~~~~~~~~~~~~-{}-{}-{}-{}-{}~-{}-{}-{}~-{}-{}-{}-{}-{}-{}-{}-{}~0:0000}\par
\texttt{F4.119~~03.3.3.-{}~~~~~~~~~~~~-{}-{}4-{}-{}~-{}-{}-{}~-{}-{}-{}+-{}-{}-{}-{}~0:0000}\par
\texttt{F4.120~~13.3.3.-{}~~~~~~~~~~~~-{}-{}4-{}-{}~-{}-{}-{}~-{}-{}-{}+-{}-{}-{}-{}~0:0000}\par
\texttt{F4.121~~0.0.0.0~~~~~~~~~~~~~-{}-{}45D~UP-{}~-{}+++-{}+-{}-{}~1:1000}\par
\texttt{F4.122~~1.0.0.0~~~~~~~~~~~~~-{}-{}-{}-{}D~UP-{}~-{}-{}+-{}-{}+-{}-{}~2:1100}\par
\texttt{F4.123~~0.1.0.0~~~~~~~~~~~~~-{}-{}45D~UP-{}~-{}+++-{}+-{}-{}~2:1000}\par
\texttt{F4.124~~1.1.0.0~~~~~~~~~~~~~-{}-{}-{}-{}D~UP-{}~-{}-{}+-{}-{}+-{}-{}~0:0100}\par
\texttt{F4.125~~2.1.0.0~~~~~~~~~~~~~-{}-{}-{}-{}D~UP-{}~-{}-{}+-{}-{}+-{}-{}~3:1010}\par
\texttt{F4.126~~0.2.0.0~~~~~~~~~~~~~-{}-{}-{}-{}D~UP-{}~-{}-{}+-{}-{}+-{}-{}~1:1000}\par
\texttt{F4.127~~1.2.0.0~~~~~~~~~~~~~-{}-{}-{}-{}D~UP-{}~-{}-{}+-{}-{}+-{}-{}~3:1110}\par
\texttt{F4.128~~2.2.0.0~~~~~~~~~~~~~-{}-{}-{}-{}D~UP-{}~-{}-{}+-{}-{}+-{}-{}~1:1010}\par
\texttt{F4.129~~3.2.0.0~~~~~~~~~~~~~-{}-{}-{}-{}D~UP-{}~-{}-{}+-{}-{}+-{}-{}~1:1001}\par
\texttt{F4.130~~0.1.1.0~~~~~~~~~~~~~-{}-{}45D~UP-{}~-{}+++-{}+-{}-{}~3:1000}\par
\texttt{F4.131~~2.1.1.0~~~~~~~~~~~~~-{}-{}-{}-{}D~UP-{}~-{}-{}+-{}-{}+-{}-{}~0:0100}\par
\texttt{F4.132~~3.1.1.0~~~~~~~~~~~~~-{}-{}-{}-{}D~UP-{}~-{}-{}+-{}-{}+-{}-{}~2:0100}\par
\texttt{F4.133~~3.2.1.0~~~~~~~~~~~~~-{}B-{}-{}D~UPC~+-{}+-{}++++~ok}\par
\texttt{F4.134~~2.3.1.0~~~~~~~~~~~~~-{}-{}-{}-{}D~UPC~-{}-{}+-{}++-{}+~ok}\par
\texttt{F4.135~~0.1.2.0~~~~~~~~~~~~~-{}-{}45D~UP-{}~-{}+++-{}+-{}-{}~3:1000}\par
\texttt{F4.136~~1.1.2.0~~~~~~~~~~~~~-{}-{}-{}-{}D~UP-{}~-{}-{}+-{}-{}+-{}-{}~0:0100}\par
\texttt{F4.137~~3.1.2.0~~~~~~~~~~~~~-{}B-{}-{}D~UPC~+-{}+-{}++++~ok}\par
\texttt{F4.138~~1.3.2.0~~~~~~~~~~~~~-{}-{}-{}-{}D~UPC~-{}-{}+-{}++++~ok}\par
\texttt{F4.139~~0.1.2.3~~~~~~~~~~~~~TB45D~UPC~++++++++~ok}\par
\texttt{F4.140~~0123.0.-{}.-{}~~~~~~~~~~-{}-{}-{}-{}-{}~-{}-{}-{}~-{}-{}-{}-{}-{}-{}-{}-{}~0:0000}\par
\texttt{F4.141~~0123.1.-{}.-{}~~~~~~~~~~-{}-{}4-{}-{}~-{}-{}-{}~-{}-{}-{}+-{}-{}-{}-{}~0:0000}\par
\texttt{F4.142~~012.01.-{}.-{}~~~~~~~~~~-{}-{}-{}-{}-{}~-{}-{}-{}~-{}-{}-{}-{}-{}-{}-{}-{}~0:0000}\par
\texttt{F4.143~~023.01.-{}.-{}~~~~~~~~~~-{}-{}-{}-{}-{}~-{}-{}-{}~-{}-{}-{}-{}-{}-{}-{}-{}~0:0000}\par
\texttt{F4.144~~123.01.-{}.-{}~~~~~~~~~~-{}-{}-{}-{}-{}~-{}-{}-{}~-{}-{}-{}-{}-{}-{}-{}-{}~0:0000}\par
\texttt{F4.145~~0123.2.-{}.-{}~~~~~~~~~~-{}-{}4-{}-{}~-{}-{}-{}~-{}-{}-{}+-{}-{}-{}-{}~0:0000}\par
\texttt{F4.146~~012.02.-{}.-{}~~~~~~~~~~-{}-{}-{}-{}-{}~-{}-{}-{}~-{}-{}-{}-{}-{}-{}-{}-{}~0:0000}\par
\texttt{F4.147~~013.02.-{}.-{}~~~~~~~~~~-{}-{}-{}-{}-{}~-{}-{}-{}~-{}-{}-{}-{}-{}-{}-{}-{}~0:0000}\par
\texttt{F4.148~~023.02.-{}.-{}~~~~~~~~~~-{}-{}-{}-{}-{}~-{}-{}-{}~-{}-{}-{}-{}-{}-{}-{}-{}~0:0000}\par
\texttt{F4.149~~123.02.-{}.-{}~~~~~~~~~~-{}-{}-{}-{}-{}~-{}-{}-{}~-{}-{}-{}-{}-{}-{}-{}-{}~0:0000}\par
\texttt{F4.150~~012.12.-{}.-{}~~~~~~~~~~-{}-{}4-{}-{}~-{}-{}-{}~-{}-{}-{}+-{}-{}-{}-{}~0:0000}\par
\texttt{F4.151~~013.12.-{}.-{}~~~~~~~~~~-{}-{}-{}-{}-{}~-{}-{}-{}~-{}-{}-{}-{}-{}-{}-{}-{}~0:0000}\par
\texttt{F4.152~~023.12.-{}.-{}~~~~~~~~~~-{}-{}4-{}-{}~-{}-{}-{}~-{}-{}-{}+-{}-{}-{}-{}~0:0000}\par
\texttt{F4.153~~123.12.-{}.-{}~~~~~~~~~~-{}-{}4-{}-{}~-{}-{}-{}~-{}-{}-{}+-{}-{}-{}-{}~0:0000}\par
\texttt{F4.154~~23.012.-{}.-{}~~~~~~~~~~-{}-{}-{}-{}-{}~-{}-{}-{}~-{}-{}-{}-{}-{}-{}-{}-{}~0:0000}\par
\texttt{F4.155~~023.23.-{}.-{}~~~~~~~~~~-{}-{}4-{}-{}~-{}-{}-{}~-{}-{}-{}+-{}-{}-{}-{}~0:0000}\par
\texttt{F4.156~~123.23.-{}.-{}~~~~~~~~~~-{}-{}4-{}-{}~-{}-{}-{}~-{}-{}-{}+-{}-{}-{}-{}~0:0000}\par
\texttt{F4.157~~012.0.0.-{}~~~~~~~~~~~-{}B-{}-{}-{}~-{}PC~+-{}-{}-{}++++~ok}\par
\texttt{F4.158~~013.0.0.-{}~~~~~~~~~~~-{}-{}-{}-{}-{}~-{}-{}-{}~-{}-{}-{}-{}-{}-{}-{}-{}~0:0000}\par
\texttt{F4.159~~123.0.0.-{}~~~~~~~~~~~-{}-{}-{}-{}-{}~-{}-{}-{}~-{}-{}-{}-{}-{}-{}-{}-{}~0:0000}\par
\texttt{F4.160~~012.1.0.-{}~~~~~~~~~~~-{}-{}-{}-{}-{}~-{}-{}-{}~-{}-{}-{}-{}-{}-{}-{}-{}~0:0100}\par
\texttt{F4.161~~013.1.0.-{}~~~~~~~~~~~-{}-{}-{}-{}-{}~-{}-{}-{}~-{}-{}-{}-{}-{}-{}-{}-{}~0:0000}\par
\texttt{F4.162~~023.1.0.-{}~~~~~~~~~~~-{}-{}-{}-{}-{}~-{}-{}-{}~-{}-{}-{}-{}-{}-{}-{}-{}~0:0000}\par
\texttt{F4.163~~123.1.0.-{}~~~~~~~~~~~-{}-{}-{}-{}-{}~-{}-{}-{}~-{}-{}-{}-{}-{}-{}-{}-{}~0:0000}\par
\texttt{F4.164~~01.01.0.-{}~~~~~~~~~~~-{}-{}-{}5-{}~UP-{}~-{}++-{}-{}+-{}-{}~2:1100}\par
\texttt{F4.165~~02.01.0.-{}~~~~~~~~~~~-{}-{}-{}-{}-{}~-{}-{}-{}~-{}-{}-{}-{}-{}-{}-{}-{}~1:1010}\par
\texttt{F4.166~~12.01.0.-{}~~~~~~~~~~~-{}B-{}-{}-{}~-{}PC~+-{}-{}-{}++++~ok}\par
\texttt{F4.167~~03.01.0.-{}~~~~~~~~~~~-{}-{}-{}-{}-{}~-{}-{}-{}~-{}-{}-{}-{}-{}-{}-{}-{}~0:0000}\par
\texttt{F4.168~~13.01.0.-{}~~~~~~~~~~~-{}-{}-{}-{}-{}~-{}-{}-{}~-{}-{}-{}-{}-{}-{}-{}-{}~0:0000}\par
\texttt{F4.169~~23.01.0.-{}~~~~~~~~~~~-{}-{}-{}-{}-{}~-{}-{}-{}~-{}-{}-{}-{}-{}-{}-{}-{}~0:0000}\par
\texttt{F4.170~~012.2.0.-{}~~~~~~~~~~~-{}-{}-{}-{}-{}~-{}PC~-{}-{}-{}-{}++-{}+~ok}\par
\texttt{F4.171~~013.2.0.-{}~~~~~~~~~~~-{}-{}-{}-{}-{}~-{}-{}-{}~-{}-{}-{}-{}-{}-{}-{}-{}~0:0000}\par
\texttt{F4.172~~023.2.0.-{}~~~~~~~~~~~-{}-{}-{}-{}-{}~-{}-{}-{}~-{}-{}-{}-{}-{}-{}-{}-{}~0:0000}\par
\texttt{F4.173~~123.2.0.-{}~~~~~~~~~~~-{}-{}-{}-{}-{}~-{}-{}-{}~-{}-{}-{}-{}-{}-{}-{}-{}~0:0000}\par
\texttt{F4.174~~01.02.0.-{}~~~~~~~~~~~-{}-{}-{}-{}-{}~-{}PC~-{}-{}-{}-{}++++~ok}\par
\texttt{F4.175~~02.02.0.-{}~~~~~~~~~~~-{}-{}-{}-{}-{}~-{}P-{}~-{}-{}-{}-{}-{}+-{}-{}~1:1010}\par
\texttt{F4.176~~12.02.0.-{}~~~~~~~~~~~-{}-{}-{}-{}-{}~-{}PC~-{}-{}-{}-{}++++~ok}\par
\texttt{F4.177~~03.02.0.-{}~~~~~~~~~~~-{}-{}-{}-{}-{}~-{}-{}-{}~-{}-{}-{}-{}-{}-{}-{}-{}~0:0000}\par
\texttt{F4.178~~13.02.0.-{}~~~~~~~~~~~-{}-{}-{}-{}-{}~-{}-{}-{}~-{}-{}-{}-{}-{}-{}-{}-{}~0:0000}\par
\texttt{F4.179~~23.02.0.-{}~~~~~~~~~~~-{}-{}-{}-{}-{}~-{}-{}-{}~-{}-{}-{}-{}-{}-{}-{}-{}~0:0000}\par
\texttt{F4.180~~01.12.0.-{}~~~~~~~~~~~-{}-{}-{}-{}-{}~-{}PC~-{}-{}-{}-{}++-{}+~ok}\par
\texttt{F4.181~~02.12.0.-{}~~~~~~~~~~~-{}-{}-{}-{}-{}~-{}-{}-{}~-{}-{}-{}-{}-{}-{}-{}-{}~1:1010}\par
\texttt{F4.182~~12.12.0.-{}~~~~~~~~~~~-{}-{}-{}-{}-{}~-{}PC~-{}-{}-{}-{}++++~ok}\par
\texttt{F4.183~~03.12.0.-{}~~~~~~~~~~~-{}-{}-{}-{}-{}~-{}-{}-{}~-{}-{}-{}-{}-{}-{}-{}-{}~0:0000}\par
\texttt{F4.184~~13.12.0.-{}~~~~~~~~~~~-{}-{}-{}-{}-{}~-{}-{}-{}~-{}-{}-{}-{}-{}-{}-{}-{}~0:0000}\par
\texttt{F4.185~~23.12.0.-{}~~~~~~~~~~~-{}-{}-{}-{}-{}~-{}-{}-{}~-{}-{}-{}-{}-{}-{}-{}-{}~0:0000}\par
\texttt{F4.186~~0.012.0.-{}~~~~~~~~~~~-{}-{}4-{}-{}~-{}-{}-{}~-{}-{}-{}+-{}-{}-{}-{}~1:1000}\par
\texttt{F4.187~~2.012.0.-{}~~~~~~~~~~~-{}-{}-{}-{}-{}~-{}-{}-{}~-{}-{}-{}-{}-{}-{}-{}-{}~1:1010}\par
\texttt{F4.188~~3.012.0.-{}~~~~~~~~~~~-{}-{}-{}-{}-{}~-{}-{}-{}~-{}-{}-{}-{}-{}-{}-{}-{}~0:0000}\par
\texttt{F4.189~~012.3.0.-{}~~~~~~~~~~~-{}-{}-{}-{}-{}~-{}-{}-{}~-{}-{}-{}-{}-{}-{}-{}-{}~1:0000}\par
\texttt{F4.190~~013.3.0.-{}~~~~~~~~~~~-{}-{}-{}-{}-{}~-{}-{}-{}~-{}-{}-{}-{}-{}-{}-{}-{}~0:0000}\par
\texttt{F4.191~~023.3.0.-{}~~~~~~~~~~~-{}-{}-{}-{}-{}~-{}-{}-{}~-{}-{}-{}-{}-{}-{}-{}-{}~0:0000}\par
\texttt{F4.192~~123.3.0.-{}~~~~~~~~~~~-{}-{}-{}-{}-{}~-{}-{}-{}~-{}-{}-{}-{}-{}-{}-{}-{}~0:0000}\par
\texttt{F4.193~~01.03.0.-{}~~~~~~~~~~~-{}-{}-{}-{}-{}~-{}-{}-{}~-{}-{}-{}-{}-{}-{}-{}-{}~1:0000}\par
\texttt{F4.194~~02.03.0.-{}~~~~~~~~~~~-{}-{}-{}-{}-{}~-{}-{}-{}~-{}-{}-{}-{}-{}-{}-{}-{}~1:0000}\par
\texttt{F4.195~~12.03.0.-{}~~~~~~~~~~~-{}-{}-{}-{}-{}~-{}-{}-{}~-{}-{}-{}-{}-{}-{}-{}-{}~1:0000}\par
\texttt{F4.196~~03.03.0.-{}~~~~~~~~~~~-{}-{}-{}-{}-{}~-{}-{}-{}~-{}-{}-{}-{}-{}-{}-{}-{}~0:0000}\par
\texttt{F4.197~~13.03.0.-{}~~~~~~~~~~~-{}-{}-{}-{}-{}~-{}-{}-{}~-{}-{}-{}-{}-{}-{}-{}-{}~0:0000}\par
\texttt{F4.198~~23.03.0.-{}~~~~~~~~~~~-{}-{}-{}-{}-{}~-{}-{}-{}~-{}-{}-{}-{}-{}-{}-{}-{}~0:0000}\par
\texttt{F4.199~~01.13.0.-{}~~~~~~~~~~~-{}-{}-{}-{}-{}~-{}-{}-{}~-{}-{}-{}-{}-{}-{}-{}-{}~1:0000}\par
\texttt{F4.200~~02.13.0.-{}~~~~~~~~~~~-{}-{}-{}-{}-{}~-{}-{}-{}~-{}-{}-{}-{}-{}-{}-{}-{}~1:0000}\par
\texttt{F4.201~~12.13.0.-{}~~~~~~~~~~~-{}-{}-{}-{}-{}~-{}-{}-{}~-{}-{}-{}-{}-{}-{}-{}-{}~1:0000}\par
\texttt{F4.202~~03.13.0.-{}~~~~~~~~~~~-{}-{}-{}-{}-{}~-{}-{}-{}~-{}-{}-{}-{}-{}-{}-{}-{}~0:0000}\par
\texttt{F4.203~~13.13.0.-{}~~~~~~~~~~~-{}-{}-{}-{}-{}~-{}-{}-{}~-{}-{}-{}-{}-{}-{}-{}-{}~0:0000}\par
\texttt{F4.204~~23.13.0.-{}~~~~~~~~~~~-{}-{}-{}-{}-{}~-{}-{}-{}~-{}-{}-{}-{}-{}-{}-{}-{}~0:0000}\par
\texttt{F4.205~~0.013.0.-{}~~~~~~~~~~~-{}-{}4-{}-{}~-{}-{}-{}~-{}-{}-{}+-{}-{}-{}-{}~1:0000}\par
\texttt{F4.206~~2.013.0.-{}~~~~~~~~~~~-{}-{}-{}-{}-{}~-{}-{}-{}~-{}-{}-{}-{}-{}-{}-{}-{}~1:0000}\par
\texttt{F4.207~~3.013.0.-{}~~~~~~~~~~~-{}-{}-{}-{}-{}~-{}-{}-{}~-{}-{}-{}-{}-{}-{}-{}-{}~0:0000}\par
\texttt{F4.208~~01.23.0.-{}~~~~~~~~~~~-{}-{}-{}-{}-{}~-{}-{}-{}~-{}-{}-{}-{}-{}-{}-{}-{}~1:0000}\par
\texttt{F4.209~~02.23.0.-{}~~~~~~~~~~~-{}-{}-{}-{}-{}~-{}-{}-{}~-{}-{}-{}-{}-{}-{}-{}-{}~1:0000}\par
\texttt{F4.210~~12.23.0.-{}~~~~~~~~~~~-{}-{}-{}-{}-{}~-{}-{}-{}~-{}-{}-{}-{}-{}-{}-{}-{}~1:0000}\par
\texttt{F4.211~~03.23.0.-{}~~~~~~~~~~~-{}-{}-{}-{}-{}~-{}-{}-{}~-{}-{}-{}-{}-{}-{}-{}-{}~0:0000}\par
\texttt{F4.212~~13.23.0.-{}~~~~~~~~~~~-{}-{}-{}-{}-{}~-{}-{}-{}~-{}-{}-{}-{}-{}-{}-{}-{}~0:0000}\par
\texttt{F4.213~~23.23.0.-{}~~~~~~~~~~~-{}-{}-{}-{}-{}~-{}-{}-{}~-{}-{}-{}-{}-{}-{}-{}-{}~0:0000}\par
\texttt{F4.214~~0.023.0.-{}~~~~~~~~~~~-{}-{}4-{}-{}~-{}-{}-{}~-{}-{}-{}+-{}-{}-{}-{}~1:0000}\par
\texttt{F4.215~~2.023.0.-{}~~~~~~~~~~~-{}-{}-{}-{}-{}~-{}-{}-{}~-{}-{}-{}-{}-{}-{}-{}-{}~1:0000}\par
\texttt{F4.216~~3.023.0.-{}~~~~~~~~~~~-{}-{}-{}-{}-{}~-{}-{}-{}~-{}-{}-{}-{}-{}-{}-{}-{}~0:0000}\par
\texttt{F4.217~~0.123.0.-{}~~~~~~~~~~~-{}-{}-{}-{}-{}~-{}-{}-{}~-{}-{}-{}-{}-{}-{}-{}-{}~1:0000}\par
\texttt{F4.218~~3.123.0.-{}~~~~~~~~~~~-{}-{}-{}-{}-{}~-{}-{}-{}~-{}-{}-{}-{}-{}-{}-{}-{}~0:0000}\par
\texttt{F4.219~~01.1.01.-{}~~~~~~~~~~~-{}-{}4-{}-{}~-{}-{}-{}~-{}-{}-{}+-{}-{}-{}-{}~0:0100}\par
\texttt{F4.220~~02.1.01.-{}~~~~~~~~~~~-{}-{}-{}-{}-{}~-{}-{}-{}~-{}-{}-{}-{}-{}-{}-{}-{}~2:0100}\par
\texttt{F4.221~~12.1.01.-{}~~~~~~~~~~~-{}-{}-{}-{}-{}~-{}-{}-{}~-{}-{}-{}-{}-{}-{}-{}-{}~0:0100}\par
\texttt{F4.222~~03.1.01.-{}~~~~~~~~~~~-{}-{}-{}-{}-{}~-{}-{}-{}~-{}-{}-{}-{}-{}-{}-{}-{}~0:0000}\par
\texttt{F4.223~~13.1.01.-{}~~~~~~~~~~~-{}-{}-{}-{}-{}~-{}-{}-{}~-{}-{}-{}-{}-{}-{}-{}-{}~0:0000}\par
\texttt{F4.224~~23.1.01.-{}~~~~~~~~~~~-{}-{}-{}-{}-{}~-{}-{}-{}~-{}-{}-{}-{}-{}-{}-{}-{}~0:0000}\par
\texttt{F4.225~~3.01.01.-{}~~~~~~~~~~~-{}-{}-{}-{}-{}~-{}-{}-{}~-{}-{}-{}-{}-{}-{}-{}-{}~0:0000}\par
\texttt{F4.226~~3.02.01.-{}~~~~~~~~~~~-{}-{}-{}-{}-{}~-{}-{}-{}~-{}-{}-{}-{}-{}-{}-{}-{}~0:0000}\par
\texttt{F4.227~~3.12.01.-{}~~~~~~~~~~~-{}-{}-{}-{}-{}~-{}-{}-{}~-{}-{}-{}-{}-{}-{}-{}-{}~0:0000}\par
\texttt{F4.228~~03.3.01.-{}~~~~~~~~~~~-{}-{}-{}-{}-{}~-{}-{}-{}~-{}-{}-{}-{}-{}-{}-{}-{}~0:0000}\par
\texttt{F4.229~~13.3.01.-{}~~~~~~~~~~~-{}-{}-{}-{}-{}~-{}-{}-{}~-{}-{}-{}-{}-{}-{}-{}-{}~0:0000}\par
\texttt{F4.230~~23.3.01.-{}~~~~~~~~~~~-{}-{}-{}-{}-{}~-{}-{}-{}~-{}-{}-{}-{}-{}-{}-{}-{}~0:0000}\par
\texttt{F4.231~~012.1.2.-{}~~~~~~~~~~~-{}-{}4-{}-{}~-{}-{}-{}~-{}-{}-{}+-{}-{}-{}-{}~0:0100}\par
\texttt{F4.232~~013.1.2.-{}~~~~~~~~~~~-{}-{}4-{}-{}~-{}-{}-{}~-{}-{}-{}+-{}-{}-{}-{}~0:0000}\par
\texttt{F4.233~~123.1.2.-{}~~~~~~~~~~~-{}-{}4-{}-{}~-{}-{}-{}~-{}-{}-{}+-{}-{}-{}-{}~0:0000}\par
\texttt{F4.234~~01.01.2.-{}~~~~~~~~~~~-{}B45-{}~UPC~++++++++~ok}\par
\texttt{F4.235~~02.01.2.-{}~~~~~~~~~~~-{}-{}-{}-{}-{}~-{}-{}-{}~-{}-{}-{}-{}-{}-{}-{}-{}~0:0010}\par
\texttt{F4.236~~03.01.2.-{}~~~~~~~~~~~-{}-{}-{}-{}-{}~-{}-{}-{}~-{}-{}-{}-{}-{}-{}-{}-{}~0:0000}\par
\texttt{F4.237~~13.01.2.-{}~~~~~~~~~~~-{}-{}-{}-{}-{}~-{}-{}-{}~-{}-{}-{}-{}-{}-{}-{}-{}~0:0000}\par
\texttt{F4.238~~23.01.2.-{}~~~~~~~~~~~-{}-{}-{}-{}-{}~-{}-{}-{}~-{}-{}-{}-{}-{}-{}-{}-{}~0:0000}\par
\texttt{F4.239~~02.12.2.-{}~~~~~~~~~~~-{}-{}4-{}-{}~-{}-{}-{}~-{}-{}-{}+-{}-{}-{}-{}~0:0010}\par
\texttt{F4.240~~03.12.2.-{}~~~~~~~~~~~-{}-{}4-{}-{}~-{}-{}-{}~-{}-{}-{}+-{}-{}-{}-{}~0:0000}\par
\texttt{F4.241~~13.12.2.-{}~~~~~~~~~~~-{}-{}-{}-{}-{}~-{}-{}-{}~-{}-{}-{}-{}-{}-{}-{}-{}~0:0000}\par
\texttt{F4.242~~23.12.2.-{}~~~~~~~~~~~-{}-{}4-{}-{}~-{}-{}-{}~-{}-{}-{}+-{}-{}-{}-{}~0:0000}\par
\texttt{F4.243~~3.012.2.-{}~~~~~~~~~~~-{}-{}-{}-{}-{}~-{}-{}-{}~-{}-{}-{}-{}-{}-{}-{}-{}~0:0000}\par
\texttt{F4.244~~013.3.2.-{}~~~~~~~~~~~-{}-{}4-{}-{}~-{}-{}-{}~-{}-{}-{}+-{}-{}-{}-{}~0:0000}\par
\texttt{F4.245~~023.3.2.-{}~~~~~~~~~~~-{}-{}4-{}-{}~-{}-{}-{}~-{}-{}-{}+-{}-{}-{}-{}~0:0000}\par
\texttt{F4.246~~123.3.2.-{}~~~~~~~~~~~-{}-{}4-{}-{}~-{}-{}-{}~-{}-{}-{}+-{}-{}-{}-{}~0:0000}\par
\texttt{F4.247~~03.03.2.-{}~~~~~~~~~~~-{}-{}4-{}-{}~-{}-{}-{}~-{}-{}-{}+-{}-{}-{}-{}~0:0000}\par
\texttt{F4.248~~13.03.2.-{}~~~~~~~~~~~-{}-{}-{}-{}-{}~-{}-{}-{}~-{}-{}-{}-{}-{}-{}-{}-{}~0:0000}\par
\texttt{F4.249~~23.03.2.-{}~~~~~~~~~~~-{}-{}-{}-{}-{}~-{}-{}-{}~-{}-{}-{}-{}-{}-{}-{}-{}~0:0000}\par
\texttt{F4.250~~03.13.2.-{}~~~~~~~~~~~-{}-{}4-{}-{}~-{}-{}-{}~-{}-{}-{}+-{}-{}-{}-{}~0:0000}\par
\texttt{F4.251~~23.13.2.-{}~~~~~~~~~~~-{}-{}4-{}-{}~-{}-{}-{}~-{}-{}-{}+-{}-{}-{}-{}~0:0000}\par
\texttt{F4.252~~23.23.2.-{}~~~~~~~~~~~-{}-{}4-{}-{}~-{}-{}-{}~-{}-{}-{}+-{}-{}-{}-{}~0:0000}\par
\texttt{F4.253~~3.01.02.-{}~~~~~~~~~~~-{}-{}-{}-{}-{}~-{}-{}-{}~-{}-{}-{}-{}-{}-{}-{}-{}~0:0000}\par
\texttt{F4.254~~3.12.02.-{}~~~~~~~~~~~-{}-{}-{}-{}-{}~-{}-{}-{}~-{}-{}-{}-{}-{}-{}-{}-{}~0:0000}\par
\texttt{F4.255~~02.3.02.-{}~~~~~~~~~~~-{}-{}4-{}-{}~UP-{}~-{}-{}++-{}+-{}-{}~1:0000}\par
\texttt{F4.256~~03.3.02.-{}~~~~~~~~~~~-{}-{}-{}-{}-{}~-{}-{}-{}~-{}-{}-{}-{}-{}-{}-{}-{}~0:0000}\par
\texttt{F4.257~~13.3.02.-{}~~~~~~~~~~~-{}-{}-{}-{}-{}~-{}-{}-{}~-{}-{}-{}-{}-{}-{}-{}-{}~0:0000}\par
\texttt{F4.258~~23.3.02.-{}~~~~~~~~~~~-{}-{}-{}-{}-{}~-{}-{}-{}~-{}-{}-{}-{}-{}-{}-{}-{}~0:0000}\par
\texttt{F4.259~~3.03.02.-{}~~~~~~~~~~~-{}-{}-{}-{}-{}~-{}-{}-{}~-{}-{}-{}-{}-{}-{}-{}-{}~0:0000}\par
\texttt{F4.260~~3.13.02.-{}~~~~~~~~~~~-{}-{}-{}-{}-{}~-{}-{}-{}~-{}-{}-{}-{}-{}-{}-{}-{}~0:0000}\par
\texttt{F4.261~~3.23.02.-{}~~~~~~~~~~~-{}-{}-{}-{}-{}~-{}-{}-{}~-{}-{}-{}-{}-{}-{}-{}-{}~0:0000}\par
\texttt{F4.262~~3.3.012.-{}~~~~~~~~~~~-{}-{}-{}-{}-{}~-{}-{}-{}~-{}-{}-{}-{}-{}-{}-{}-{}~0:0000}\par
\texttt{F4.263~~013.3.3.-{}~~~~~~~~~~~-{}-{}4-{}-{}~-{}-{}-{}~-{}-{}-{}+-{}-{}-{}-{}~0:0000}\par
\texttt{F4.264~~123.3.3.-{}~~~~~~~~~~~-{}-{}4-{}-{}~-{}-{}-{}~-{}-{}-{}+-{}-{}-{}-{}~0:0000}\par
\texttt{F4.265~~03.03.3.-{}~~~~~~~~~~~-{}-{}4-{}-{}~-{}-{}-{}~-{}-{}-{}+-{}-{}-{}-{}~0:0000}\par
\texttt{F4.266~~13.03.3.-{}~~~~~~~~~~~-{}-{}-{}-{}-{}~-{}-{}-{}~-{}-{}-{}-{}-{}-{}-{}-{}~0:0000}\par
\texttt{F4.267~~23.03.3.-{}~~~~~~~~~~~-{}-{}-{}-{}-{}~-{}-{}-{}~-{}-{}-{}-{}-{}-{}-{}-{}~0:0000}\par
\texttt{F4.268~~03.13.3.-{}~~~~~~~~~~~-{}-{}4-{}-{}~-{}-{}-{}~-{}-{}-{}+-{}-{}-{}-{}~0:0000}\par
\texttt{F4.269~~23.13.3.-{}~~~~~~~~~~~-{}-{}4-{}-{}~-{}-{}-{}~-{}-{}-{}+-{}-{}-{}-{}~0:0000}\par
\texttt{F4.270~~23.23.3.-{}~~~~~~~~~~~-{}-{}4-{}-{}~-{}-{}-{}~-{}-{}-{}+-{}-{}-{}-{}~0:0000}\par
\texttt{F4.271~~01.0.0.0~~~~~~~~~~~~-{}-{}-{}-{}D~-{}P-{}~-{}-{}-{}-{}-{}+-{}-{}~2:1100}\par
\texttt{F4.272~~12.0.0.0~~~~~~~~~~~~-{}-{}-{}-{}D~UP-{}~-{}-{}+-{}-{}+-{}-{}~3:1110}\par
\texttt{F4.273~~01.1.0.0~~~~~~~~~~~~-{}-{}-{}-{}D~-{}-{}-{}~-{}-{}-{}-{}-{}-{}-{}-{}~0:0100}\par
\texttt{F4.274~~02.1.0.0~~~~~~~~~~~~-{}-{}-{}-{}D~-{}P-{}~-{}-{}-{}-{}-{}+-{}-{}~3:1010}\par
\texttt{F4.275~~12.1.0.0~~~~~~~~~~~~-{}-{}-{}-{}D~-{}-{}-{}~-{}-{}-{}-{}-{}-{}-{}-{}~0:0100}\par
\texttt{F4.276~~23.1.0.0~~~~~~~~~~~~-{}B-{}-{}D~UPC~+-{}+-{}++++~ok}\par
\texttt{F4.277~~0.01.0.0~~~~~~~~~~~~-{}-{}4-{}D~-{}-{}-{}~-{}-{}-{}+-{}-{}-{}-{}~1:1000}\par
\texttt{F4.278~~1.01.0.0~~~~~~~~~~~~-{}-{}-{}-{}D~-{}P-{}~-{}-{}-{}-{}-{}+-{}-{}~2:1100}\par
\texttt{F4.279~~2.01.0.0~~~~~~~~~~~~-{}-{}-{}-{}D~-{}-{}-{}~-{}-{}-{}-{}-{}-{}-{}-{}~1:1010}\par
\texttt{F4.280~~01.2.0.0~~~~~~~~~~~~-{}-{}-{}-{}D~-{}P-{}~-{}-{}-{}-{}-{}+-{}-{}~3:1110}\par
\texttt{F4.281~~02.2.0.0~~~~~~~~~~~~-{}-{}-{}-{}D~-{}P-{}~-{}-{}-{}-{}-{}+-{}-{}~1:1010}\par
\texttt{F4.282~~12.2.0.0~~~~~~~~~~~~-{}-{}-{}-{}D~-{}P-{}~-{}-{}-{}-{}-{}+-{}-{}~3:1110}\par
\texttt{F4.283~~03.2.0.0~~~~~~~~~~~~-{}-{}-{}-{}D~-{}P-{}~-{}-{}-{}-{}-{}+-{}-{}~1:1001}\par
\texttt{F4.284~~13.2.0.0~~~~~~~~~~~~-{}-{}-{}-{}D~-{}PC~-{}-{}-{}-{}++-{}+~ok}\par
\texttt{F4.285~~23.2.0.0~~~~~~~~~~~~-{}-{}-{}-{}D~UP-{}~-{}-{}+-{}-{}+-{}-{}~1:1011}\par
\texttt{F4.286~~0.02.0.0~~~~~~~~~~~~-{}-{}4-{}D~UP-{}~-{}-{}++-{}+-{}-{}~1:1000}\par
\texttt{F4.287~~1.02.0.0~~~~~~~~~~~~-{}-{}-{}-{}D~-{}P-{}~-{}-{}-{}-{}-{}+-{}-{}~3:1110}\par
\texttt{F4.288~~2.02.0.0~~~~~~~~~~~~-{}-{}-{}-{}D~-{}P-{}~-{}-{}-{}-{}-{}+-{}-{}~1:1010}\par
\texttt{F4.289~~3.02.0.0~~~~~~~~~~~~-{}-{}-{}-{}D~-{}P-{}~-{}-{}-{}-{}-{}+-{}-{}~1:1001}\par
\texttt{F4.290~~0.12.0.0~~~~~~~~~~~~-{}-{}-{}-{}D~-{}-{}-{}~-{}-{}-{}-{}-{}-{}-{}-{}~1:1000}\par
\texttt{F4.291~~1.12.0.0~~~~~~~~~~~~-{}-{}-{}-{}D~-{}P-{}~-{}-{}-{}-{}-{}+-{}-{}~3:1110}\par
\texttt{F4.292~~2.12.0.0~~~~~~~~~~~~-{}-{}-{}-{}D~-{}-{}-{}~-{}-{}-{}-{}-{}-{}-{}-{}~1:1010}\par
\texttt{F4.293~~3.12.0.0~~~~~~~~~~~~-{}-{}-{}-{}D~-{}-{}-{}~-{}-{}-{}-{}-{}-{}-{}-{}~1:1001}\par
\texttt{F4.294~~0.23.0.0~~~~~~~~~~~~-{}-{}-{}-{}D~UP-{}~-{}-{}+-{}-{}+-{}-{}~1:1000}\par
\texttt{F4.295~~1.23.0.0~~~~~~~~~~~~-{}-{}-{}-{}D~UPC~-{}-{}+-{}++++~ok}\par
\texttt{F4.296~~2.23.0.0~~~~~~~~~~~~-{}-{}-{}-{}D~UP-{}~-{}-{}+-{}-{}+-{}-{}~1:1010}\par
\texttt{F4.297~~01.1.1.0~~~~~~~~~~~~-{}-{}-{}-{}D~-{}-{}-{}~-{}-{}-{}-{}-{}-{}-{}-{}~0:0100}\par
\texttt{F4.298~~02.1.1.0~~~~~~~~~~~~-{}-{}-{}-{}D~-{}-{}-{}~-{}-{}-{}-{}-{}-{}-{}-{}~0:0100}\par
\texttt{F4.299~~12.1.1.0~~~~~~~~~~~~-{}-{}-{}-{}D~UP-{}~-{}-{}+-{}-{}+-{}-{}~0:0100}\par
\texttt{F4.300~~03.1.1.0~~~~~~~~~~~~-{}-{}-{}-{}D~-{}P-{}~-{}-{}-{}-{}-{}+-{}-{}~2:0100}\par
\texttt{F4.301~~13.1.1.0~~~~~~~~~~~~-{}-{}-{}-{}D~-{}-{}-{}~-{}-{}-{}-{}-{}-{}-{}-{}~0:0100}\par
\texttt{F4.302~~23.1.1.0~~~~~~~~~~~~-{}-{}-{}-{}D~-{}-{}-{}~-{}-{}-{}-{}-{}-{}-{}-{}~0:0100}\par
\texttt{F4.303~~2.01.1.0~~~~~~~~~~~~-{}-{}-{}-{}D~-{}P-{}~-{}-{}-{}-{}-{}+-{}-{}~3:1110}\par
\texttt{F4.304~~3.01.1.0~~~~~~~~~~~~-{}-{}-{}-{}D~-{}-{}-{}~-{}-{}-{}-{}-{}-{}-{}-{}~1:1001}\par
\texttt{F4.305~~12.2.1.0~~~~~~~~~~~~-{}-{}-{}-{}D~-{}P-{}~-{}-{}-{}-{}-{}+-{}-{}~0:0110}\par
\texttt{F4.306~~03.2.1.0~~~~~~~~~~~~-{}B-{}-{}D~-{}PC~+-{}-{}-{}++++~ok}\par
\texttt{F4.307~~13.2.1.0~~~~~~~~~~~~-{}-{}-{}-{}D~-{}-{}-{}~-{}-{}-{}-{}-{}-{}-{}-{}~0:0110}\par
\texttt{F4.308~~2.02.1.0~~~~~~~~~~~~-{}-{}-{}-{}D~-{}P-{}~-{}-{}-{}-{}-{}+-{}-{}~3:1110}\par
\texttt{F4.309~~2.12.1.0~~~~~~~~~~~~-{}-{}-{}-{}D~-{}P-{}~-{}-{}-{}-{}-{}+-{}-{}~0:0110}\par
\texttt{F4.310~~02.3.1.0~~~~~~~~~~~~-{}-{}-{}-{}D~-{}PC~-{}-{}-{}-{}++-{}-{}~ok}\par
\texttt{F4.311~~12.3.1.0~~~~~~~~~~~~-{}-{}-{}-{}D~-{}PC~-{}-{}-{}-{}++-{}+~ok}\par
\texttt{F4.312~~23.3.1.0~~~~~~~~~~~~-{}-{}-{}-{}D~-{}PC~-{}-{}-{}-{}++-{}+~ok}\par
\texttt{F4.313~~0.1.01.0~~~~~~~~~~~~-{}-{}4-{}D~-{}-{}-{}~-{}-{}-{}+-{}-{}-{}-{}~2:1000}\par
\texttt{F4.314~~1.1.01.0~~~~~~~~~~~~-{}-{}-{}-{}D~UP-{}~-{}-{}+-{}-{}+-{}-{}~0:0100}\par
\texttt{F4.315~~2.1.01.0~~~~~~~~~~~~-{}-{}-{}-{}D~-{}-{}-{}~-{}-{}-{}-{}-{}-{}-{}-{}~2:0100}\par
\texttt{F4.316~~3.1.01.0~~~~~~~~~~~~-{}-{}-{}-{}D~-{}-{}-{}~-{}-{}-{}-{}-{}-{}-{}-{}~2:0100}\par
\texttt{F4.317~~0.3.01.0~~~~~~~~~~~~-{}-{}-{}-{}D~-{}-{}-{}~-{}-{}-{}-{}-{}-{}-{}-{}~1:1000}\par
\texttt{F4.318~~1.3.01.0~~~~~~~~~~~~-{}-{}-{}-{}D~-{}P-{}~-{}-{}-{}-{}-{}+-{}-{}~2:1101}\par
\texttt{F4.319~~01.1.2.0~~~~~~~~~~~~-{}-{}-{}-{}D~-{}-{}-{}~-{}-{}-{}-{}-{}-{}-{}-{}~0:0100}\par
\texttt{F4.320~~12.1.2.0~~~~~~~~~~~~-{}-{}-{}-{}D~-{}-{}-{}~-{}-{}-{}-{}-{}-{}-{}-{}~0:0100}\par
\texttt{F4.321~~03.1.2.0~~~~~~~~~~~~-{}B-{}-{}D~-{}PC~+-{}-{}-{}++++~ok}\par
\texttt{F4.322~~13.1.2.0~~~~~~~~~~~~-{}-{}-{}-{}D~-{}-{}-{}~-{}-{}-{}-{}-{}-{}-{}-{}~0:0100}\par
\texttt{F4.323~~0.01.2.0~~~~~~~~~~~~-{}-{}4-{}D~-{}-{}-{}~-{}-{}-{}+-{}-{}-{}-{}~1:1000}\par
\texttt{F4.324~~1.01.2.0~~~~~~~~~~~~-{}-{}-{}-{}D~-{}P-{}~-{}-{}-{}-{}-{}+-{}-{}~3:1100}\par
\texttt{F4.325~~2.01.2.0~~~~~~~~~~~~-{}-{}-{}-{}D~-{}-{}-{}~-{}-{}-{}-{}-{}-{}-{}-{}~0:0010}\par
\texttt{F4.326~~3.01.2.0~~~~~~~~~~~~-{}-{}-{}-{}D~-{}-{}-{}~-{}-{}-{}-{}-{}-{}-{}-{}~1:1001}\par
\texttt{F4.327~~0.12.2.0~~~~~~~~~~~~-{}-{}4-{}D~-{}-{}-{}~-{}-{}-{}+-{}-{}-{}-{}~3:1000}\par
\texttt{F4.328~~1.12.2.0~~~~~~~~~~~~-{}-{}-{}-{}D~-{}-{}-{}~-{}-{}-{}-{}-{}-{}-{}-{}~1:0010}\par
\texttt{F4.329~~2.12.2.0~~~~~~~~~~~~-{}-{}-{}-{}D~-{}-{}-{}~-{}-{}-{}-{}-{}-{}-{}-{}~0:0010}\par
\texttt{F4.330~~3.12.2.0~~~~~~~~~~~~-{}-{}-{}-{}D~-{}-{}-{}~-{}-{}-{}-{}-{}-{}-{}-{}~1:0010}\par
\texttt{F4.331~~01.3.2.0~~~~~~~~~~~~-{}-{}-{}-{}D~-{}PC~-{}-{}-{}-{}++-{}+~ok}\par
\texttt{F4.332~~12.3.2.0~~~~~~~~~~~~-{}-{}-{}-{}D~-{}-{}-{}~-{}-{}-{}-{}-{}-{}-{}-{}~0:0010}\par
\texttt{F4.333~~13.3.2.0~~~~~~~~~~~~-{}-{}-{}-{}D~-{}PC~-{}-{}-{}-{}++++~ok}\par
\texttt{F4.334~~0.03.2.0~~~~~~~~~~~~-{}-{}4-{}D~UP-{}~-{}-{}++-{}+-{}-{}~1:1000}\par
\texttt{F4.335~~2.03.2.0~~~~~~~~~~~~-{}-{}-{}-{}D~-{}-{}-{}~-{}-{}-{}-{}-{}-{}-{}-{}~0:0010}\par
\texttt{F4.336~~3.03.2.0~~~~~~~~~~~~-{}-{}-{}-{}D~-{}P-{}~-{}-{}-{}-{}-{}+-{}-{}~1:1001}\par
\texttt{F4.337~~0.13.2.0~~~~~~~~~~~~-{}-{}-{}-{}D~-{}-{}-{}~-{}-{}-{}-{}-{}-{}-{}-{}~1:1000}\par
\texttt{F4.338~~2.13.2.0~~~~~~~~~~~~-{}-{}-{}-{}D~-{}-{}-{}~-{}-{}-{}-{}-{}-{}-{}-{}~0:0010}\par
\texttt{F4.339~~0.23.2.0~~~~~~~~~~~~-{}-{}-{}-{}D~-{}-{}-{}~-{}-{}-{}-{}-{}-{}-{}-{}~1:1000}\par
\texttt{F4.340~~1.3.02.0~~~~~~~~~~~~-{}-{}-{}-{}D~-{}-{}-{}~-{}-{}-{}-{}-{}-{}-{}-{}~2:1101}\par
\texttt{F4.341~~0.1.2.01~~~~~~~~~~~~-{}-{}4-{}D~-{}-{}-{}~-{}-{}-{}+-{}-{}-{}-{}~3:1000}\par
\texttt{F4.342~~01.1.2.3~~~~~~~~~~~~T-{}4-{}D~-{}-{}-{}~-{}-{}-{}+-{}-{}-{}-{}~0:0100}\par
\texttt{F4.343~~0123.01.-{}.-{}~~~~~~~~~-{}-{}-{}-{}-{}~-{}-{}-{}~-{}-{}-{}-{}-{}-{}-{}-{}~0:0000}\par
\texttt{F4.344~~0123.02.-{}.-{}~~~~~~~~~-{}-{}-{}-{}-{}~-{}-{}-{}~-{}-{}-{}-{}-{}-{}-{}-{}~0:0000}\par
\texttt{F4.345~~0123.12.-{}.-{}~~~~~~~~~-{}-{}4-{}-{}~-{}-{}-{}~-{}-{}-{}+-{}-{}-{}-{}~0:0000}\par
\texttt{F4.346~~012.012.-{}.-{}~~~~~~~~~-{}-{}4-{}-{}~-{}-{}-{}~-{}-{}-{}+-{}-{}-{}-{}~0:0000}\par
\texttt{F4.347~~013.012.-{}.-{}~~~~~~~~~-{}-{}-{}-{}-{}~-{}-{}-{}~-{}-{}-{}-{}-{}-{}-{}-{}~0:0000}\par
\texttt{F4.348~~023.012.-{}.-{}~~~~~~~~~-{}-{}-{}-{}-{}~-{}-{}-{}~-{}-{}-{}-{}-{}-{}-{}-{}~0:0000}\par
\texttt{F4.349~~123.012.-{}.-{}~~~~~~~~~-{}-{}-{}-{}-{}~-{}-{}-{}~-{}-{}-{}-{}-{}-{}-{}-{}~0:0000}\par
\texttt{F4.350~~0123.23.-{}.-{}~~~~~~~~~-{}-{}4-{}-{}~-{}-{}-{}~-{}-{}-{}+-{}-{}-{}-{}~0:0000}\par
\texttt{F4.351~~023.023.-{}.-{}~~~~~~~~~-{}-{}4-{}-{}~-{}-{}-{}~-{}-{}-{}+-{}-{}-{}-{}~0:0000}\par
\texttt{F4.352~~123.023.-{}.-{}~~~~~~~~~-{}-{}-{}-{}-{}~-{}-{}-{}~-{}-{}-{}-{}-{}-{}-{}-{}~0:0000}\par
\texttt{F4.353~~023.123.-{}.-{}~~~~~~~~~-{}-{}4-{}-{}~-{}-{}-{}~-{}-{}-{}+-{}-{}-{}-{}~0:0000}\par
\texttt{F4.354~~0123.0.0.-{}~~~~~~~~~~-{}-{}-{}-{}-{}~-{}-{}-{}~-{}-{}-{}-{}-{}-{}-{}-{}~0:0000}\par
\texttt{F4.355~~0123.1.0.-{}~~~~~~~~~~-{}-{}-{}-{}-{}~-{}-{}-{}~-{}-{}-{}-{}-{}-{}-{}-{}~0:0000}\par
\texttt{F4.356~~012.01.0.-{}~~~~~~~~~~-{}B-{}-{}-{}~-{}PC~+-{}-{}-{}++++~ok}\par
\texttt{F4.357~~013.01.0.-{}~~~~~~~~~~-{}-{}-{}-{}-{}~-{}-{}-{}~-{}-{}-{}-{}-{}-{}-{}-{}~0:0000}\par
\texttt{F4.358~~023.01.0.-{}~~~~~~~~~~-{}-{}-{}-{}-{}~-{}-{}-{}~-{}-{}-{}-{}-{}-{}-{}-{}~0:0000}\par
\texttt{F4.359~~123.01.0.-{}~~~~~~~~~~-{}-{}-{}-{}-{}~-{}-{}-{}~-{}-{}-{}-{}-{}-{}-{}-{}~0:0000}\par
\texttt{F4.360~~0123.2.0.-{}~~~~~~~~~~-{}-{}-{}-{}-{}~-{}-{}-{}~-{}-{}-{}-{}-{}-{}-{}-{}~0:0000}\par
\texttt{F4.361~~012.02.0.-{}~~~~~~~~~~-{}-{}-{}-{}-{}~-{}PC~-{}-{}-{}-{}++++~ok}\par
\texttt{F4.362~~013.02.0.-{}~~~~~~~~~~-{}-{}-{}-{}-{}~-{}-{}-{}~-{}-{}-{}-{}-{}-{}-{}-{}~0:0000}\par
\texttt{F4.363~~023.02.0.-{}~~~~~~~~~~-{}-{}-{}-{}-{}~-{}-{}-{}~-{}-{}-{}-{}-{}-{}-{}-{}~0:0000}\par
\texttt{F4.364~~123.02.0.-{}~~~~~~~~~~-{}-{}-{}-{}-{}~-{}-{}-{}~-{}-{}-{}-{}-{}-{}-{}-{}~0:0000}\par
\texttt{F4.365~~012.12.0.-{}~~~~~~~~~~-{}-{}-{}-{}-{}~-{}PC~-{}-{}-{}-{}++++~ok}\par
\texttt{F4.366~~013.12.0.-{}~~~~~~~~~~-{}-{}-{}-{}-{}~-{}-{}-{}~-{}-{}-{}-{}-{}-{}-{}-{}~0:0000}\par
\texttt{F4.367~~023.12.0.-{}~~~~~~~~~~-{}-{}-{}-{}-{}~-{}-{}-{}~-{}-{}-{}-{}-{}-{}-{}-{}~0:0000}\par
\texttt{F4.368~~123.12.0.-{}~~~~~~~~~~-{}-{}-{}-{}-{}~-{}-{}-{}~-{}-{}-{}-{}-{}-{}-{}-{}~0:0000}\par
\texttt{F4.369~~01.012.0.-{}~~~~~~~~~~-{}-{}-{}-{}-{}~-{}PC~-{}-{}-{}-{}++-{}+~ok}\par
\texttt{F4.370~~02.012.0.-{}~~~~~~~~~~-{}-{}-{}-{}-{}~-{}-{}-{}~-{}-{}-{}-{}-{}-{}-{}-{}~1:1010}\par
\texttt{F4.371~~12.012.0.-{}~~~~~~~~~~-{}-{}-{}-{}-{}~-{}PC~-{}-{}-{}-{}++++~ok}\par
\texttt{F4.372~~03.012.0.-{}~~~~~~~~~~-{}-{}-{}-{}-{}~-{}-{}-{}~-{}-{}-{}-{}-{}-{}-{}-{}~0:0000}\par
\texttt{F4.373~~13.012.0.-{}~~~~~~~~~~-{}-{}-{}-{}-{}~-{}-{}-{}~-{}-{}-{}-{}-{}-{}-{}-{}~0:0000}\par
\texttt{F4.374~~23.012.0.-{}~~~~~~~~~~-{}-{}-{}-{}-{}~-{}-{}-{}~-{}-{}-{}-{}-{}-{}-{}-{}~0:0000}\par
\texttt{F4.375~~0123.3.0.-{}~~~~~~~~~~-{}-{}-{}-{}-{}~-{}-{}-{}~-{}-{}-{}-{}-{}-{}-{}-{}~0:0000}\par
\texttt{F4.376~~012.03.0.-{}~~~~~~~~~~-{}-{}-{}-{}-{}~-{}-{}-{}~-{}-{}-{}-{}-{}-{}-{}-{}~1:0000}\par
\texttt{F4.377~~013.03.0.-{}~~~~~~~~~~-{}-{}-{}-{}-{}~-{}-{}-{}~-{}-{}-{}-{}-{}-{}-{}-{}~0:0000}\par
\texttt{F4.378~~023.03.0.-{}~~~~~~~~~~-{}-{}-{}-{}-{}~-{}-{}-{}~-{}-{}-{}-{}-{}-{}-{}-{}~0:0000}\par
\texttt{F4.379~~123.03.0.-{}~~~~~~~~~~-{}-{}-{}-{}-{}~-{}-{}-{}~-{}-{}-{}-{}-{}-{}-{}-{}~0:0000}\par
\texttt{F4.380~~012.13.0.-{}~~~~~~~~~~-{}-{}-{}-{}-{}~-{}-{}-{}~-{}-{}-{}-{}-{}-{}-{}-{}~1:0000}\par
\texttt{F4.381~~013.13.0.-{}~~~~~~~~~~-{}-{}-{}-{}-{}~-{}-{}-{}~-{}-{}-{}-{}-{}-{}-{}-{}~0:0000}\par
\texttt{F4.382~~023.13.0.-{}~~~~~~~~~~-{}-{}-{}-{}-{}~-{}-{}-{}~-{}-{}-{}-{}-{}-{}-{}-{}~0:0000}\par
\texttt{F4.383~~123.13.0.-{}~~~~~~~~~~-{}-{}-{}-{}-{}~-{}-{}-{}~-{}-{}-{}-{}-{}-{}-{}-{}~0:0000}\par
\texttt{F4.384~~01.013.0.-{}~~~~~~~~~~-{}-{}-{}-{}-{}~-{}-{}-{}~-{}-{}-{}-{}-{}-{}-{}-{}~1:0000}\par
\texttt{F4.385~~02.013.0.-{}~~~~~~~~~~-{}-{}-{}-{}-{}~-{}-{}-{}~-{}-{}-{}-{}-{}-{}-{}-{}~1:0000}\par
\texttt{F4.386~~12.013.0.-{}~~~~~~~~~~-{}-{}-{}-{}-{}~-{}-{}-{}~-{}-{}-{}-{}-{}-{}-{}-{}~1:0000}\par
\texttt{F4.387~~03.013.0.-{}~~~~~~~~~~-{}-{}-{}-{}-{}~-{}-{}-{}~-{}-{}-{}-{}-{}-{}-{}-{}~0:0000}\par
\texttt{F4.388~~13.013.0.-{}~~~~~~~~~~-{}-{}-{}-{}-{}~-{}-{}-{}~-{}-{}-{}-{}-{}-{}-{}-{}~0:0000}\par
\texttt{F4.389~~23.013.0.-{}~~~~~~~~~~-{}-{}-{}-{}-{}~-{}-{}-{}~-{}-{}-{}-{}-{}-{}-{}-{}~0:0000}\par
\texttt{F4.390~~012.23.0.-{}~~~~~~~~~~-{}-{}-{}-{}-{}~-{}-{}-{}~-{}-{}-{}-{}-{}-{}-{}-{}~1:0000}\par
\texttt{F4.391~~013.23.0.-{}~~~~~~~~~~-{}-{}-{}-{}-{}~-{}-{}-{}~-{}-{}-{}-{}-{}-{}-{}-{}~0:0000}\par
\texttt{F4.392~~023.23.0.-{}~~~~~~~~~~-{}-{}-{}-{}-{}~-{}-{}-{}~-{}-{}-{}-{}-{}-{}-{}-{}~0:0000}\par
\texttt{F4.393~~123.23.0.-{}~~~~~~~~~~-{}-{}-{}-{}-{}~-{}-{}-{}~-{}-{}-{}-{}-{}-{}-{}-{}~0:0000}\par
\texttt{F4.394~~01.023.0.-{}~~~~~~~~~~-{}-{}-{}-{}-{}~-{}-{}-{}~-{}-{}-{}-{}-{}-{}-{}-{}~1:0000}\par
\texttt{F4.395~~02.023.0.-{}~~~~~~~~~~-{}-{}-{}-{}-{}~-{}-{}-{}~-{}-{}-{}-{}-{}-{}-{}-{}~1:0000}\par
\texttt{F4.396~~12.023.0.-{}~~~~~~~~~~-{}-{}-{}-{}-{}~-{}-{}-{}~-{}-{}-{}-{}-{}-{}-{}-{}~1:0000}\par
\texttt{F4.397~~03.023.0.-{}~~~~~~~~~~-{}-{}-{}-{}-{}~-{}-{}-{}~-{}-{}-{}-{}-{}-{}-{}-{}~0:0000}\par
\texttt{F4.398~~13.023.0.-{}~~~~~~~~~~-{}-{}-{}-{}-{}~-{}-{}-{}~-{}-{}-{}-{}-{}-{}-{}-{}~0:0000}\par
\texttt{F4.399~~23.023.0.-{}~~~~~~~~~~-{}-{}-{}-{}-{}~-{}-{}-{}~-{}-{}-{}-{}-{}-{}-{}-{}~0:0000}\par
\texttt{F4.400~~01.123.0.-{}~~~~~~~~~~-{}-{}-{}-{}-{}~-{}-{}-{}~-{}-{}-{}-{}-{}-{}-{}-{}~1:0000}\par
\texttt{F4.401~~02.123.0.-{}~~~~~~~~~~-{}-{}-{}-{}-{}~-{}-{}-{}~-{}-{}-{}-{}-{}-{}-{}-{}~1:0000}\par
\texttt{F4.402~~12.123.0.-{}~~~~~~~~~~-{}-{}-{}-{}-{}~-{}-{}-{}~-{}-{}-{}-{}-{}-{}-{}-{}~1:0000}\par
\texttt{F4.403~~03.123.0.-{}~~~~~~~~~~-{}-{}-{}-{}-{}~-{}-{}-{}~-{}-{}-{}-{}-{}-{}-{}-{}~0:0000}\par
\texttt{F4.404~~13.123.0.-{}~~~~~~~~~~-{}-{}-{}-{}-{}~-{}-{}-{}~-{}-{}-{}-{}-{}-{}-{}-{}~0:0000}\par
\texttt{F4.405~~23.123.0.-{}~~~~~~~~~~-{}-{}-{}-{}-{}~-{}-{}-{}~-{}-{}-{}-{}-{}-{}-{}-{}~0:0000}\par
\texttt{F4.406~~0.0123.0.-{}~~~~~~~~~~-{}-{}4-{}-{}~-{}-{}-{}~-{}-{}-{}+-{}-{}-{}-{}~1:0000}\par
\texttt{F4.407~~2.0123.0.-{}~~~~~~~~~~-{}-{}-{}-{}-{}~-{}-{}-{}~-{}-{}-{}-{}-{}-{}-{}-{}~1:0000}\par
\texttt{F4.408~~3.0123.0.-{}~~~~~~~~~~-{}-{}-{}-{}-{}~-{}-{}-{}~-{}-{}-{}-{}-{}-{}-{}-{}~0:0000}\par
\texttt{F4.409~~012.1.01.-{}~~~~~~~~~~-{}-{}-{}-{}-{}~-{}-{}-{}~-{}-{}-{}-{}-{}-{}-{}-{}~0:0100}\par
\texttt{F4.410~~013.1.01.-{}~~~~~~~~~~-{}-{}-{}-{}-{}~-{}-{}-{}~-{}-{}-{}-{}-{}-{}-{}-{}~0:0000}\par
\texttt{F4.411~~023.1.01.-{}~~~~~~~~~~-{}-{}-{}-{}-{}~-{}-{}-{}~-{}-{}-{}-{}-{}-{}-{}-{}~0:0000}\par
\texttt{F4.412~~123.1.01.-{}~~~~~~~~~~-{}-{}-{}-{}-{}~-{}-{}-{}~-{}-{}-{}-{}-{}-{}-{}-{}~0:0000}\par
\texttt{F4.413~~01.01.01.-{}~~~~~~~~~~-{}-{}45-{}~UP-{}~-{}+++-{}+-{}-{}~2:1100}\par
\texttt{F4.414~~02.01.01.-{}~~~~~~~~~~-{}-{}-{}-{}-{}~-{}PC~-{}-{}-{}-{}++++~ok}\par
\texttt{F4.415~~12.01.01.-{}~~~~~~~~~~-{}-{}-{}-{}-{}~-{}PC~-{}-{}-{}-{}++++~ok}\par
\texttt{F4.416~~03.01.01.-{}~~~~~~~~~~-{}-{}-{}-{}-{}~-{}-{}-{}~-{}-{}-{}-{}-{}-{}-{}-{}~0:0000}\par
\texttt{F4.417~~13.01.01.-{}~~~~~~~~~~-{}-{}-{}-{}-{}~-{}-{}-{}~-{}-{}-{}-{}-{}-{}-{}-{}~0:0000}\par
\texttt{F4.418~~23.01.01.-{}~~~~~~~~~~-{}-{}-{}-{}-{}~-{}-{}-{}~-{}-{}-{}-{}-{}-{}-{}-{}~0:0000}\par
\texttt{F4.419~~12.02.01.-{}~~~~~~~~~~-{}B-{}-{}-{}~-{}PC~+-{}-{}-{}++++~ok}\par
\texttt{F4.420~~03.02.01.-{}~~~~~~~~~~-{}-{}-{}-{}-{}~-{}-{}-{}~-{}-{}-{}-{}-{}-{}-{}-{}~0:0000}\par
\texttt{F4.421~~13.02.01.-{}~~~~~~~~~~-{}-{}-{}-{}-{}~-{}-{}-{}~-{}-{}-{}-{}-{}-{}-{}-{}~0:0000}\par
\texttt{F4.422~~02.12.01.-{}~~~~~~~~~~-{}B-{}-{}-{}~-{}PC~+-{}-{}-{}++++~ok}\par
\texttt{F4.423~~03.12.01.-{}~~~~~~~~~~-{}-{}-{}-{}-{}~-{}-{}-{}~-{}-{}-{}-{}-{}-{}-{}-{}~0:0000}\par
\texttt{F4.424~~13.12.01.-{}~~~~~~~~~~-{}-{}-{}-{}-{}~-{}-{}-{}~-{}-{}-{}-{}-{}-{}-{}-{}~0:0000}\par
\texttt{F4.425~~23.12.01.-{}~~~~~~~~~~-{}-{}-{}-{}-{}~-{}-{}-{}~-{}-{}-{}-{}-{}-{}-{}-{}~0:0000}\par
\texttt{F4.426~~3.012.01.-{}~~~~~~~~~~-{}-{}-{}-{}-{}~-{}-{}-{}~-{}-{}-{}-{}-{}-{}-{}-{}~0:0000}\par
\texttt{F4.427~~013.3.01.-{}~~~~~~~~~~-{}-{}-{}-{}-{}~-{}-{}-{}~-{}-{}-{}-{}-{}-{}-{}-{}~0:0000}\par
\texttt{F4.428~~023.3.01.-{}~~~~~~~~~~-{}-{}-{}-{}-{}~-{}-{}-{}~-{}-{}-{}-{}-{}-{}-{}-{}~0:0000}\par
\texttt{F4.429~~123.3.01.-{}~~~~~~~~~~-{}-{}-{}-{}-{}~-{}-{}-{}~-{}-{}-{}-{}-{}-{}-{}-{}~0:0000}\par
\texttt{F4.430~~03.03.01.-{}~~~~~~~~~~-{}-{}-{}-{}-{}~-{}-{}-{}~-{}-{}-{}-{}-{}-{}-{}-{}~0:0000}\par
\texttt{F4.431~~13.03.01.-{}~~~~~~~~~~-{}-{}-{}-{}-{}~-{}-{}-{}~-{}-{}-{}-{}-{}-{}-{}-{}~0:0000}\par
\texttt{F4.432~~23.03.01.-{}~~~~~~~~~~-{}-{}-{}-{}-{}~-{}-{}-{}~-{}-{}-{}-{}-{}-{}-{}-{}~0:0000}\par
\texttt{F4.433~~03.13.01.-{}~~~~~~~~~~-{}-{}-{}-{}-{}~-{}-{}-{}~-{}-{}-{}-{}-{}-{}-{}-{}~0:0000}\par
\texttt{F4.434~~23.13.01.-{}~~~~~~~~~~-{}-{}-{}-{}-{}~-{}-{}-{}~-{}-{}-{}-{}-{}-{}-{}-{}~0:0000}\par
\texttt{F4.435~~23.23.01.-{}~~~~~~~~~~-{}-{}-{}-{}-{}~-{}-{}-{}~-{}-{}-{}-{}-{}-{}-{}-{}~0:0000}\par
\texttt{F4.436~~0123.1.2.-{}~~~~~~~~~~-{}-{}4-{}-{}~-{}-{}-{}~-{}-{}-{}+-{}-{}-{}-{}~0:0000}\par
\texttt{F4.437~~012.01.2.-{}~~~~~~~~~~-{}-{}-{}-{}-{}~-{}-{}-{}~-{}-{}-{}-{}-{}-{}-{}-{}~0:0010}\par
\texttt{F4.438~~013.01.2.-{}~~~~~~~~~~-{}-{}-{}-{}-{}~-{}-{}-{}~-{}-{}-{}-{}-{}-{}-{}-{}~0:0000}\par
\texttt{F4.439~~023.01.2.-{}~~~~~~~~~~-{}-{}-{}-{}-{}~-{}-{}-{}~-{}-{}-{}-{}-{}-{}-{}-{}~0:0000}\par
\texttt{F4.440~~123.01.2.-{}~~~~~~~~~~-{}-{}-{}-{}-{}~-{}-{}-{}~-{}-{}-{}-{}-{}-{}-{}-{}~0:0000}\par
\texttt{F4.441~~012.12.2.-{}~~~~~~~~~~-{}-{}4-{}-{}~-{}-{}-{}~-{}-{}-{}+-{}-{}-{}-{}~0:0010}\par
\texttt{F4.442~~013.12.2.-{}~~~~~~~~~~-{}-{}-{}-{}-{}~-{}-{}-{}~-{}-{}-{}-{}-{}-{}-{}-{}~0:0000}\par
\texttt{F4.443~~023.12.2.-{}~~~~~~~~~~-{}-{}4-{}-{}~-{}-{}-{}~-{}-{}-{}+-{}-{}-{}-{}~0:0000}\par
\texttt{F4.444~~123.12.2.-{}~~~~~~~~~~-{}-{}4-{}-{}~-{}-{}-{}~-{}-{}-{}+-{}-{}-{}-{}~0:0000}\par
\texttt{F4.445~~03.012.2.-{}~~~~~~~~~~-{}-{}-{}-{}-{}~-{}-{}-{}~-{}-{}-{}-{}-{}-{}-{}-{}~0:0000}\par
\texttt{F4.446~~13.012.2.-{}~~~~~~~~~~-{}-{}-{}-{}-{}~-{}-{}-{}~-{}-{}-{}-{}-{}-{}-{}-{}~0:0000}\par
\texttt{F4.447~~23.012.2.-{}~~~~~~~~~~-{}-{}-{}-{}-{}~-{}-{}-{}~-{}-{}-{}-{}-{}-{}-{}-{}~0:0000}\par
\texttt{F4.448~~0123.3.2.-{}~~~~~~~~~~-{}-{}4-{}-{}~-{}-{}-{}~-{}-{}-{}+-{}-{}-{}-{}~0:0000}\par
\texttt{F4.449~~013.03.2.-{}~~~~~~~~~~-{}-{}-{}-{}-{}~-{}-{}-{}~-{}-{}-{}-{}-{}-{}-{}-{}~0:0000}\par
\texttt{F4.450~~023.03.2.-{}~~~~~~~~~~-{}-{}-{}-{}-{}~-{}-{}-{}~-{}-{}-{}-{}-{}-{}-{}-{}~0:0000}\par
\texttt{F4.451~~123.03.2.-{}~~~~~~~~~~-{}-{}-{}-{}-{}~-{}-{}-{}~-{}-{}-{}-{}-{}-{}-{}-{}~0:0000}\par
\texttt{F4.452~~013.13.2.-{}~~~~~~~~~~-{}-{}4-{}-{}~-{}-{}-{}~-{}-{}-{}+-{}-{}-{}-{}~0:0000}\par
\texttt{F4.453~~023.13.2.-{}~~~~~~~~~~-{}-{}4-{}-{}~-{}-{}-{}~-{}-{}-{}+-{}-{}-{}-{}~0:0000}\par
\texttt{F4.454~~123.13.2.-{}~~~~~~~~~~-{}-{}4-{}-{}~-{}-{}-{}~-{}-{}-{}+-{}-{}-{}-{}~0:0000}\par
\texttt{F4.455~~23.013.2.-{}~~~~~~~~~~-{}-{}-{}-{}-{}~-{}-{}-{}~-{}-{}-{}-{}-{}-{}-{}-{}~0:0000}\par
\texttt{F4.456~~023.23.2.-{}~~~~~~~~~~-{}-{}4-{}-{}~-{}-{}-{}~-{}-{}-{}+-{}-{}-{}-{}~0:0000}\par
\texttt{F4.457~~123.23.2.-{}~~~~~~~~~~-{}-{}4-{}-{}~-{}-{}-{}~-{}-{}-{}+-{}-{}-{}-{}~0:0000}\par
\texttt{F4.458~~01.01.02.-{}~~~~~~~~~~-{}-{}-{}-{}-{}~-{}-{}-{}~-{}-{}-{}-{}-{}-{}-{}-{}~2:1100}\par
\texttt{F4.459~~03.01.02.-{}~~~~~~~~~~-{}-{}-{}-{}-{}~-{}-{}-{}~-{}-{}-{}-{}-{}-{}-{}-{}~0:0000}\par
\texttt{F4.460~~13.01.02.-{}~~~~~~~~~~-{}-{}-{}-{}-{}~-{}-{}-{}~-{}-{}-{}-{}-{}-{}-{}-{}~0:0000}\par
\texttt{F4.461~~01.12.02.-{}~~~~~~~~~~-{}-{}-{}-{}-{}~-{}PC~-{}-{}-{}-{}++++~ok}\par
\texttt{F4.462~~03.12.02.-{}~~~~~~~~~~-{}-{}-{}-{}-{}~-{}-{}-{}~-{}-{}-{}-{}-{}-{}-{}-{}~0:0000}\par
\texttt{F4.463~~13.12.02.-{}~~~~~~~~~~-{}-{}-{}-{}-{}~-{}-{}-{}~-{}-{}-{}-{}-{}-{}-{}-{}~0:0000}\par
\texttt{F4.464~~23.12.02.-{}~~~~~~~~~~-{}-{}-{}-{}-{}~-{}-{}-{}~-{}-{}-{}-{}-{}-{}-{}-{}~0:0000}\par
\texttt{F4.465~~3.012.02.-{}~~~~~~~~~~-{}-{}-{}-{}-{}~-{}-{}-{}~-{}-{}-{}-{}-{}-{}-{}-{}~0:0000}\par
\texttt{F4.466~~012.3.02.-{}~~~~~~~~~~-{}-{}-{}-{}-{}~-{}-{}-{}~-{}-{}-{}-{}-{}-{}-{}-{}~1:0000}\par
\texttt{F4.467~~013.3.02.-{}~~~~~~~~~~-{}-{}-{}-{}-{}~-{}-{}-{}~-{}-{}-{}-{}-{}-{}-{}-{}~0:0000}\par
\texttt{F4.468~~023.3.02.-{}~~~~~~~~~~-{}-{}-{}-{}-{}~-{}-{}-{}~-{}-{}-{}-{}-{}-{}-{}-{}~0:0000}\par
\texttt{F4.469~~123.3.02.-{}~~~~~~~~~~-{}-{}-{}-{}-{}~-{}-{}-{}~-{}-{}-{}-{}-{}-{}-{}-{}~0:0000}\par
\texttt{F4.470~~02.03.02.-{}~~~~~~~~~~-{}-{}-{}-{}-{}~-{}-{}-{}~-{}-{}-{}-{}-{}-{}-{}-{}~1:0000}\par
\texttt{F4.471~~03.03.02.-{}~~~~~~~~~~-{}-{}-{}-{}-{}~-{}-{}-{}~-{}-{}-{}-{}-{}-{}-{}-{}~0:0000}\par
\texttt{F4.472~~13.03.02.-{}~~~~~~~~~~-{}-{}-{}-{}-{}~-{}-{}-{}~-{}-{}-{}-{}-{}-{}-{}-{}~0:0000}\par
\texttt{F4.473~~23.03.02.-{}~~~~~~~~~~-{}-{}-{}-{}-{}~-{}-{}-{}~-{}-{}-{}-{}-{}-{}-{}-{}~0:0000}\par
\texttt{F4.474~~02.13.02.-{}~~~~~~~~~~-{}-{}4-{}-{}~-{}-{}-{}~-{}-{}-{}+-{}-{}-{}-{}~1:0000}\par
\texttt{F4.475~~03.13.02.-{}~~~~~~~~~~-{}-{}-{}-{}-{}~-{}-{}-{}~-{}-{}-{}-{}-{}-{}-{}-{}~0:0000}\par
\texttt{F4.476~~13.13.02.-{}~~~~~~~~~~-{}-{}-{}-{}-{}~-{}-{}-{}~-{}-{}-{}-{}-{}-{}-{}-{}~0:0000}\par
\texttt{F4.477~~23.13.02.-{}~~~~~~~~~~-{}-{}-{}-{}-{}~-{}-{}-{}~-{}-{}-{}-{}-{}-{}-{}-{}~0:0000}\par
\texttt{F4.478~~3.013.02.-{}~~~~~~~~~~-{}-{}-{}-{}-{}~-{}-{}-{}~-{}-{}-{}-{}-{}-{}-{}-{}~0:0000}\par
\texttt{F4.479~~03.23.02.-{}~~~~~~~~~~-{}-{}-{}-{}-{}~-{}-{}-{}~-{}-{}-{}-{}-{}-{}-{}-{}~0:0000}\par
\texttt{F4.480~~13.23.02.-{}~~~~~~~~~~-{}-{}-{}-{}-{}~-{}-{}-{}~-{}-{}-{}-{}-{}-{}-{}-{}~0:0000}\par
\texttt{F4.481~~23.23.02.-{}~~~~~~~~~~-{}-{}-{}-{}-{}~-{}-{}-{}~-{}-{}-{}-{}-{}-{}-{}-{}~0:0000}\par
\texttt{F4.482~~3.023.02.-{}~~~~~~~~~~-{}-{}-{}-{}-{}~-{}-{}-{}~-{}-{}-{}-{}-{}-{}-{}-{}~0:0000}\par
\texttt{F4.483~~3.123.02.-{}~~~~~~~~~~-{}-{}-{}-{}-{}~-{}-{}-{}~-{}-{}-{}-{}-{}-{}-{}-{}~0:0000}\par
\texttt{F4.484~~03.3.012.-{}~~~~~~~~~~-{}-{}-{}-{}-{}~-{}-{}-{}~-{}-{}-{}-{}-{}-{}-{}-{}~0:0000}\par
\texttt{F4.485~~13.3.012.-{}~~~~~~~~~~-{}-{}-{}-{}-{}~-{}-{}-{}~-{}-{}-{}-{}-{}-{}-{}-{}~0:0000}\par
\texttt{F4.486~~23.3.012.-{}~~~~~~~~~~-{}-{}-{}-{}-{}~-{}-{}-{}~-{}-{}-{}-{}-{}-{}-{}-{}~0:0000}\par
\texttt{F4.487~~0123.3.3.-{}~~~~~~~~~~-{}-{}4-{}-{}~-{}-{}-{}~-{}-{}-{}+-{}-{}-{}-{}~0:0000}\par
\texttt{F4.488~~013.03.3.-{}~~~~~~~~~~-{}-{}-{}-{}-{}~-{}-{}-{}~-{}-{}-{}-{}-{}-{}-{}-{}~0:0000}\par
\texttt{F4.489~~023.03.3.-{}~~~~~~~~~~-{}-{}-{}-{}-{}~-{}-{}-{}~-{}-{}-{}-{}-{}-{}-{}-{}~0:0000}\par
\texttt{F4.490~~123.03.3.-{}~~~~~~~~~~-{}-{}-{}-{}-{}~-{}-{}-{}~-{}-{}-{}-{}-{}-{}-{}-{}~0:0000}\par
\texttt{F4.491~~013.13.3.-{}~~~~~~~~~~-{}-{}4-{}-{}~-{}-{}-{}~-{}-{}-{}+-{}-{}-{}-{}~0:0000}\par
\texttt{F4.492~~023.13.3.-{}~~~~~~~~~~-{}-{}4-{}-{}~-{}-{}-{}~-{}-{}-{}+-{}-{}-{}-{}~0:0000}\par
\texttt{F4.493~~123.13.3.-{}~~~~~~~~~~-{}-{}4-{}-{}~-{}-{}-{}~-{}-{}-{}+-{}-{}-{}-{}~0:0000}\par
\texttt{F4.494~~23.013.3.-{}~~~~~~~~~~-{}-{}-{}-{}-{}~-{}-{}-{}~-{}-{}-{}-{}-{}-{}-{}-{}~0:0000}\par
\texttt{F4.495~~023.23.3.-{}~~~~~~~~~~-{}-{}4-{}-{}~-{}-{}-{}~-{}-{}-{}+-{}-{}-{}-{}~0:0000}\par
\texttt{F4.496~~123.23.3.-{}~~~~~~~~~~-{}-{}4-{}-{}~-{}-{}-{}~-{}-{}-{}+-{}-{}-{}-{}~0:0000}\par
\texttt{F4.497~~03.03.03.-{}~~~~~~~~~~-{}-{}4-{}-{}~-{}-{}-{}~-{}-{}-{}+-{}-{}-{}-{}~0:0000}\par
\texttt{F4.498~~13.03.03.-{}~~~~~~~~~~-{}-{}-{}-{}-{}~-{}-{}-{}~-{}-{}-{}-{}-{}-{}-{}-{}~0:0000}\par
\texttt{F4.499~~03.13.03.-{}~~~~~~~~~~-{}-{}4-{}-{}~-{}-{}-{}~-{}-{}-{}+-{}-{}-{}-{}~0:0000}\par
\texttt{F4.500~~13.13.03.-{}~~~~~~~~~~-{}-{}-{}-{}-{}~-{}-{}-{}~-{}-{}-{}-{}-{}-{}-{}-{}~0:0000}\par
\texttt{F4.501~~23.13.03.-{}~~~~~~~~~~-{}-{}-{}-{}-{}~-{}-{}-{}~-{}-{}-{}-{}-{}-{}-{}-{}~0:0000}\par
\texttt{F4.502~~13.23.03.-{}~~~~~~~~~~-{}-{}-{}-{}-{}~-{}-{}-{}~-{}-{}-{}-{}-{}-{}-{}-{}~0:0000}\par
\texttt{F4.503~~03.13.23.-{}~~~~~~~~~~-{}-{}4-{}-{}~-{}-{}-{}~-{}-{}-{}+-{}-{}-{}-{}~0:0000}\par
\texttt{F4.504~~012.0.0.0~~~~~~~~~~~-{}-{}-{}-{}D~-{}P-{}~-{}-{}-{}-{}-{}+-{}-{}~3:1110}\par
\texttt{F4.505~~123.0.0.0~~~~~~~~~~~-{}B-{}-{}D~UPC~+-{}+-{}++++~ok}\par
\texttt{F4.506~~012.1.0.0~~~~~~~~~~~-{}-{}-{}-{}D~-{}-{}-{}~-{}-{}-{}-{}-{}-{}-{}-{}~0:0100}\par
\texttt{F4.507~~023.1.0.0~~~~~~~~~~~-{}B-{}-{}D~-{}PC~+-{}-{}-{}++++~ok}\par
\texttt{F4.508~~123.1.0.0~~~~~~~~~~~-{}-{}-{}-{}D~-{}-{}-{}~-{}-{}-{}-{}-{}-{}-{}-{}~0:0100}\par
\texttt{F4.509~~01.01.0.0~~~~~~~~~~~-{}-{}-{}5D~UP-{}~-{}++-{}-{}+-{}-{}~2:1100}\par
\texttt{F4.510~~02.01.0.0~~~~~~~~~~~-{}-{}-{}-{}D~-{}-{}-{}~-{}-{}-{}-{}-{}-{}-{}-{}~1:1010}\par
\texttt{F4.511~~12.01.0.0~~~~~~~~~~~-{}-{}-{}-{}D~-{}P-{}~-{}-{}-{}-{}-{}+-{}-{}~3:1110}\par
\texttt{F4.512~~23.01.0.0~~~~~~~~~~~-{}-{}-{}-{}D~-{}-{}-{}~-{}-{}-{}-{}-{}-{}-{}-{}~1:1011}\par
\texttt{F4.513~~012.2.0.0~~~~~~~~~~~-{}-{}-{}-{}D~-{}P-{}~-{}-{}-{}-{}-{}+-{}-{}~3:1110}\par
\texttt{F4.514~~013.2.0.0~~~~~~~~~~~-{}-{}-{}-{}D~-{}PC~-{}-{}-{}-{}++-{}+~ok}\par
\texttt{F4.515~~023.2.0.0~~~~~~~~~~~-{}-{}-{}-{}D~-{}P-{}~-{}-{}-{}-{}-{}+-{}-{}~1:1011}\par
\texttt{F4.516~~123.2.0.0~~~~~~~~~~~-{}-{}-{}-{}D~-{}PC~-{}-{}-{}-{}++++~ok}\par
\texttt{F4.517~~01.02.0.0~~~~~~~~~~~-{}-{}-{}-{}D~-{}P-{}~-{}-{}-{}-{}-{}+-{}-{}~3:1110}\par
\texttt{F4.518~~02.02.0.0~~~~~~~~~~~-{}-{}-{}-{}D~-{}P-{}~-{}-{}-{}-{}-{}+-{}-{}~1:1010}\par
\texttt{F4.519~~12.02.0.0~~~~~~~~~~~-{}-{}-{}-{}D~-{}P-{}~-{}-{}-{}-{}-{}+-{}-{}~3:1110}\par
\texttt{F4.520~~03.02.0.0~~~~~~~~~~~-{}-{}-{}-{}D~-{}P-{}~-{}-{}-{}-{}-{}+-{}-{}~1:1001}\par
\texttt{F4.521~~13.02.0.0~~~~~~~~~~~-{}-{}-{}-{}D~-{}PC~-{}-{}-{}-{}++-{}+~ok}\par
\texttt{F4.522~~23.02.0.0~~~~~~~~~~~-{}-{}-{}-{}D~-{}P-{}~-{}-{}-{}-{}-{}+-{}-{}~1:1011}\par
\texttt{F4.523~~01.12.0.0~~~~~~~~~~~-{}-{}-{}-{}D~-{}P-{}~-{}-{}-{}-{}-{}+-{}-{}~3:1110}\par
\texttt{F4.524~~02.12.0.0~~~~~~~~~~~-{}-{}-{}-{}D~-{}-{}-{}~-{}-{}-{}-{}-{}-{}-{}-{}~1:1010}\par
\texttt{F4.525~~12.12.0.0~~~~~~~~~~~-{}-{}-{}-{}D~-{}P-{}~-{}-{}-{}-{}-{}+-{}-{}~3:1110}\par
\texttt{F4.526~~03.12.0.0~~~~~~~~~~~-{}-{}-{}-{}D~-{}-{}-{}~-{}-{}-{}-{}-{}-{}-{}-{}~1:1001}\par
\texttt{F4.527~~13.12.0.0~~~~~~~~~~~-{}-{}-{}-{}D~-{}PC~-{}-{}-{}-{}++-{}+~ok}\par
\texttt{F4.528~~23.12.0.0~~~~~~~~~~~-{}-{}-{}-{}D~-{}-{}-{}~-{}-{}-{}-{}-{}-{}-{}-{}~1:1011}\par
\texttt{F4.529~~0.012.0.0~~~~~~~~~~~-{}-{}4-{}D~-{}-{}-{}~-{}-{}-{}+-{}-{}-{}-{}~1:1000}\par
\texttt{F4.530~~1.012.0.0~~~~~~~~~~~-{}-{}-{}-{}D~-{}P-{}~-{}-{}-{}-{}-{}+-{}-{}~3:1110}\par
\texttt{F4.531~~2.012.0.0~~~~~~~~~~~-{}-{}-{}-{}D~-{}-{}-{}~-{}-{}-{}-{}-{}-{}-{}-{}~1:1010}\par
\texttt{F4.532~~3.012.0.0~~~~~~~~~~~-{}-{}-{}-{}D~-{}-{}-{}~-{}-{}-{}-{}-{}-{}-{}-{}~1:1001}\par
\texttt{F4.533~~01.23.0.0~~~~~~~~~~~-{}-{}-{}-{}D~-{}PC~-{}-{}-{}-{}++-{}+~ok}\par
\texttt{F4.534~~02.23.0.0~~~~~~~~~~~-{}-{}-{}-{}D~-{}P-{}~-{}-{}-{}-{}-{}+-{}-{}~1:1010}\par
\texttt{F4.535~~12.23.0.0~~~~~~~~~~~-{}-{}-{}-{}D~-{}PC~-{}-{}-{}-{}++-{}+~ok}\par
\texttt{F4.536~~23.23.0.0~~~~~~~~~~~-{}-{}-{}-{}D~UP-{}~-{}-{}+-{}-{}+-{}-{}~1:1011}\par
\texttt{F4.537~~0.023.0.0~~~~~~~~~~~-{}-{}4-{}D~UP-{}~-{}-{}++-{}+-{}-{}~1:1000}\par
\texttt{F4.538~~1.023.0.0~~~~~~~~~~~-{}-{}-{}-{}D~-{}PC~-{}-{}-{}-{}++++~ok}\par
\texttt{F4.539~~2.023.0.0~~~~~~~~~~~-{}-{}-{}-{}D~-{}P-{}~-{}-{}-{}-{}-{}+-{}-{}~1:1010}\par
\texttt{F4.540~~0.123.0.0~~~~~~~~~~~-{}-{}-{}-{}D~-{}-{}-{}~-{}-{}-{}-{}-{}-{}-{}-{}~1:1000}\par
\texttt{F4.541~~1.123.0.0~~~~~~~~~~~-{}-{}-{}-{}D~-{}PC~-{}-{}-{}-{}++-{}+~ok}\par
\texttt{F4.542~~2.123.0.0~~~~~~~~~~~-{}-{}-{}-{}D~-{}-{}-{}~-{}-{}-{}-{}-{}-{}-{}-{}~1:1010}\par
\texttt{F4.543~~012.1.1.0~~~~~~~~~~~-{}-{}-{}-{}D~-{}-{}-{}~-{}-{}-{}-{}-{}-{}-{}-{}~0:0100}\par
\texttt{F4.544~~013.1.1.0~~~~~~~~~~~-{}-{}-{}-{}D~-{}-{}-{}~-{}-{}-{}-{}-{}-{}-{}-{}~0:0100}\par
\texttt{F4.545~~023.1.1.0~~~~~~~~~~~-{}-{}-{}-{}D~-{}-{}-{}~-{}-{}-{}-{}-{}-{}-{}-{}~0:0100}\par
\texttt{F4.546~~123.1.1.0~~~~~~~~~~~-{}-{}-{}-{}D~-{}-{}-{}~-{}-{}-{}-{}-{}-{}-{}-{}~0:0100}\par
\texttt{F4.547~~01.01.1.0~~~~~~~~~~~-{}-{}-{}5D~UP-{}~-{}++-{}-{}+-{}-{}~2:1100}\par
\texttt{F4.548~~02.01.1.0~~~~~~~~~~~-{}-{}-{}-{}D~-{}P-{}~-{}-{}-{}-{}-{}+-{}-{}~3:1110}\par
\texttt{F4.549~~12.01.1.0~~~~~~~~~~~-{}-{}-{}-{}D~-{}P-{}~-{}-{}-{}-{}-{}+-{}-{}~3:1110}\par
\texttt{F4.550~~03.01.1.0~~~~~~~~~~~-{}-{}-{}-{}D~-{}-{}-{}~-{}-{}-{}-{}-{}-{}-{}-{}~1:1001}\par
\texttt{F4.551~~13.01.1.0~~~~~~~~~~~-{}-{}-{}-{}D~-{}P-{}~-{}-{}-{}-{}-{}+-{}-{}~2:1101}\par
\texttt{F4.552~~23.01.1.0~~~~~~~~~~~-{}-{}-{}-{}D~-{}PC~-{}-{}-{}-{}++-{}+~ok}\par
\texttt{F4.553~~012.2.1.0~~~~~~~~~~~-{}-{}-{}-{}D~-{}-{}-{}~-{}-{}-{}-{}-{}-{}-{}-{}~0:0110}\par
\texttt{F4.554~~013.2.1.0~~~~~~~~~~~-{}-{}-{}-{}D~-{}-{}-{}~-{}-{}-{}-{}-{}-{}-{}-{}~0:0110}\par
\texttt{F4.555~~123.2.1.0~~~~~~~~~~~-{}-{}-{}-{}D~-{}-{}-{}~-{}-{}-{}-{}-{}-{}-{}-{}~0:0110}\par
\texttt{F4.556~~02.02.1.0~~~~~~~~~~~-{}-{}-{}-{}D~-{}P-{}~-{}-{}-{}-{}-{}+-{}-{}~3:1110}\par
\texttt{F4.557~~12.02.1.0~~~~~~~~~~~-{}-{}-{}-{}D~-{}P-{}~-{}-{}-{}-{}-{}+-{}-{}~3:1110}\par
\texttt{F4.558~~03.02.1.0~~~~~~~~~~~-{}-{}-{}-{}D~-{}-{}-{}~-{}-{}-{}-{}-{}-{}-{}-{}~1:1001}\par
\texttt{F4.559~~13.02.1.0~~~~~~~~~~~-{}B-{}-{}D~-{}PC~+-{}-{}-{}++-{}+~ok}\par
\texttt{F4.560~~23.02.1.0~~~~~~~~~~~-{}-{}-{}-{}D~-{}PC~-{}-{}-{}-{}++-{}+~ok}\par
\texttt{F4.561~~02.12.1.0~~~~~~~~~~~-{}-{}-{}-{}D~-{}-{}-{}~-{}-{}-{}-{}-{}-{}-{}-{}~0:0110}\par
\texttt{F4.562~~12.12.1.0~~~~~~~~~~~-{}-{}-{}-{}D~-{}P-{}~-{}-{}-{}-{}-{}+-{}-{}~0:0110}\par
\texttt{F4.563~~03.12.1.0~~~~~~~~~~~-{}B-{}-{}D~-{}PC~+-{}-{}-{}++++~ok}\par
\texttt{F4.564~~23.12.1.0~~~~~~~~~~~-{}-{}-{}-{}D~-{}-{}-{}~-{}-{}-{}-{}-{}-{}-{}-{}~0:0110}\par
\texttt{F4.565~~2.012.1.0~~~~~~~~~~~-{}-{}-{}-{}D~-{}P-{}~-{}-{}-{}-{}-{}+-{}-{}~3:1110}\par
\texttt{F4.566~~012.3.1.0~~~~~~~~~~~-{}-{}-{}-{}D~-{}PC~-{}-{}-{}-{}++-{}-{}~ok}\par
\texttt{F4.567~~023.3.1.0~~~~~~~~~~~-{}-{}-{}-{}D~-{}PC~-{}-{}-{}-{}++-{}-{}~ok}\par
\texttt{F4.568~~123.3.1.0~~~~~~~~~~~-{}-{}-{}-{}D~-{}PC~-{}-{}-{}-{}++-{}+~ok}\par
\texttt{F4.569~~02.03.1.0~~~~~~~~~~~-{}-{}-{}-{}D~-{}PC~-{}-{}-{}-{}++-{}+~ok}\par
\texttt{F4.570~~12.03.1.0~~~~~~~~~~~-{}-{}-{}-{}D~-{}PC~-{}-{}-{}-{}++++~ok}\par
\texttt{F4.571~~23.03.1.0~~~~~~~~~~~-{}-{}-{}-{}D~-{}PC~-{}-{}-{}-{}++-{}+~ok}\par
\texttt{F4.572~~02.13.1.0~~~~~~~~~~~-{}-{}-{}-{}D~-{}PC~-{}-{}-{}-{}++-{}-{}~ok}\par
\texttt{F4.573~~23.13.1.0~~~~~~~~~~~-{}-{}-{}-{}D~-{}PC~-{}-{}-{}-{}++-{}+~ok}\par
\texttt{F4.574~~23.23.1.0~~~~~~~~~~~-{}-{}-{}-{}D~-{}PC~-{}-{}-{}-{}++-{}+~ok}\par
\texttt{F4.575~~01.1.01.0~~~~~~~~~~~-{}-{}-{}-{}D~-{}-{}-{}~-{}-{}-{}-{}-{}-{}-{}-{}~0:0100}\par
\texttt{F4.576~~02.1.01.0~~~~~~~~~~~-{}-{}-{}-{}D~-{}-{}-{}~-{}-{}-{}-{}-{}-{}-{}-{}~2:0100}\par
\texttt{F4.577~~12.1.01.0~~~~~~~~~~~-{}-{}-{}-{}D~-{}-{}-{}~-{}-{}-{}-{}-{}-{}-{}-{}~0:0100}\par
\texttt{F4.578~~03.1.01.0~~~~~~~~~~~-{}-{}-{}-{}D~-{}-{}-{}~-{}-{}-{}-{}-{}-{}-{}-{}~2:0100}\par
\texttt{F4.579~~13.1.01.0~~~~~~~~~~~-{}-{}-{}-{}D~-{}-{}-{}~-{}-{}-{}-{}-{}-{}-{}-{}~0:0100}\par
\texttt{F4.580~~23.1.01.0~~~~~~~~~~~-{}-{}-{}-{}D~-{}-{}-{}~-{}-{}-{}-{}-{}-{}-{}-{}~2:0100}\par
\texttt{F4.581~~0.01.01.0~~~~~~~~~~~-{}-{}4-{}D~-{}-{}-{}~-{}-{}-{}+-{}-{}-{}-{}~1:1000}\par
\texttt{F4.582~~1.01.01.0~~~~~~~~~~~-{}-{}-{}-{}D~-{}P-{}~-{}-{}-{}-{}-{}+-{}-{}~2:1100}\par
\texttt{F4.583~~2.01.01.0~~~~~~~~~~~-{}-{}-{}-{}D~-{}P-{}~-{}-{}-{}-{}-{}+-{}-{}~3:1110}\par
\texttt{F4.584~~3.01.01.0~~~~~~~~~~~-{}-{}-{}-{}D~-{}-{}-{}~-{}-{}-{}-{}-{}-{}-{}-{}~1:1001}\par
\texttt{F4.585~~0.02.01.0~~~~~~~~~~~-{}-{}-{}-{}D~-{}-{}-{}~-{}-{}-{}-{}-{}-{}-{}-{}~1:1000}\par
\texttt{F4.586~~1.02.01.0~~~~~~~~~~~-{}-{}-{}-{}D~-{}P-{}~-{}-{}-{}-{}-{}+-{}-{}~3:1110}\par
\texttt{F4.587~~3.02.01.0~~~~~~~~~~~-{}-{}-{}-{}D~-{}-{}-{}~-{}-{}-{}-{}-{}-{}-{}-{}~1:1001}\par
\texttt{F4.588~~0.12.01.0~~~~~~~~~~~-{}-{}-{}-{}D~-{}-{}-{}~-{}-{}-{}-{}-{}-{}-{}-{}~2:1000}\par
\texttt{F4.589~~1.12.01.0~~~~~~~~~~~-{}-{}-{}-{}D~-{}P-{}~-{}-{}-{}-{}-{}+-{}-{}~3:1110}\par
\texttt{F4.590~~2.12.01.0~~~~~~~~~~~-{}-{}-{}-{}D~-{}P-{}~-{}-{}-{}-{}-{}+-{}-{}~3:1110}\par
\texttt{F4.591~~3.12.01.0~~~~~~~~~~~-{}-{}-{}-{}D~-{}-{}-{}~-{}-{}-{}-{}-{}-{}-{}-{}~2:1001}\par
\texttt{F4.592~~01.3.01.0~~~~~~~~~~~-{}-{}-{}-{}D~-{}P-{}~-{}-{}-{}-{}-{}+-{}-{}~2:1101}\par
\texttt{F4.593~~02.3.01.0~~~~~~~~~~~-{}-{}-{}-{}D~-{}PC~-{}-{}-{}-{}++-{}+~ok}\par
\texttt{F4.594~~12.3.01.0~~~~~~~~~~~-{}-{}-{}-{}D~-{}PC~-{}-{}-{}-{}++-{}+~ok}\par
\texttt{F4.595~~03.3.01.0~~~~~~~~~~~-{}-{}-{}-{}D~-{}P-{}~-{}-{}-{}-{}-{}+-{}-{}~1:1001}\par
\texttt{F4.596~~13.3.01.0~~~~~~~~~~~-{}-{}-{}-{}D~-{}P-{}~-{}-{}-{}-{}-{}+-{}-{}~2:1101}\par
\texttt{F4.597~~23.3.01.0~~~~~~~~~~~-{}-{}-{}-{}D~-{}PC~-{}-{}-{}-{}++-{}+~ok}\par
\texttt{F4.598~~0.03.01.0~~~~~~~~~~~-{}-{}-{}-{}D~-{}-{}-{}~-{}-{}-{}-{}-{}-{}-{}-{}~1:1000}\par
\texttt{F4.599~~1.03.01.0~~~~~~~~~~~-{}-{}-{}-{}D~-{}P-{}~-{}-{}-{}-{}-{}+-{}-{}~2:1101}\par
\texttt{F4.600~~2.03.01.0~~~~~~~~~~~-{}-{}-{}-{}D~-{}PC~-{}-{}-{}-{}++-{}+~ok}\par
\texttt{F4.601~~3.03.01.0~~~~~~~~~~~-{}-{}-{}-{}D~-{}P-{}~-{}-{}-{}-{}-{}+-{}-{}~1:1001}\par
\texttt{F4.602~~0.13.01.0~~~~~~~~~~~-{}-{}-{}-{}D~-{}-{}-{}~-{}-{}-{}-{}-{}-{}-{}-{}~1:1000}\par
\texttt{F4.603~~1.13.01.0~~~~~~~~~~~-{}-{}-{}-{}D~-{}P-{}~-{}-{}-{}-{}-{}+-{}-{}~2:1101}\par
\texttt{F4.604~~2.13.01.0~~~~~~~~~~~-{}-{}-{}-{}D~-{}PC~-{}-{}-{}-{}++-{}-{}~ok}\par
\texttt{F4.605~~3.13.01.0~~~~~~~~~~~-{}-{}-{}-{}D~-{}-{}-{}~-{}-{}-{}-{}-{}-{}-{}-{}~1:1001}\par
\texttt{F4.606~~0.23.01.0~~~~~~~~~~~-{}-{}-{}-{}D~-{}-{}-{}~-{}-{}-{}-{}-{}-{}-{}-{}~1:1000}\par
\texttt{F4.607~~1.23.01.0~~~~~~~~~~~-{}-{}-{}-{}D~-{}PC~-{}-{}-{}-{}++++~ok}\par
\texttt{F4.608~~3.23.01.0~~~~~~~~~~~-{}-{}-{}-{}D~-{}-{}-{}~-{}-{}-{}-{}-{}-{}-{}-{}~1:1001}\par
\texttt{F4.609~~012.1.2.0~~~~~~~~~~~-{}-{}-{}-{}D~-{}-{}-{}~-{}-{}-{}-{}-{}-{}-{}-{}~0:0100}\par
\texttt{F4.610~~013.1.2.0~~~~~~~~~~~-{}-{}-{}-{}D~-{}-{}-{}~-{}-{}-{}-{}-{}-{}-{}-{}~0:0100}\par
\texttt{F4.611~~123.1.2.0~~~~~~~~~~~-{}-{}-{}-{}D~-{}-{}-{}~-{}-{}-{}-{}-{}-{}-{}-{}~0:0100}\par
\texttt{F4.612~~01.01.2.0~~~~~~~~~~~-{}-{}-{}5D~UP-{}~-{}++-{}-{}+-{}-{}~3:1100}\par
\texttt{F4.613~~02.01.2.0~~~~~~~~~~~-{}-{}-{}-{}D~-{}-{}-{}~-{}-{}-{}-{}-{}-{}-{}-{}~0:0010}\par
\texttt{F4.614~~12.01.2.0~~~~~~~~~~~-{}-{}-{}-{}D~-{}-{}-{}~-{}-{}-{}-{}-{}-{}-{}-{}~0:0010}\par
\texttt{F4.615~~03.01.2.0~~~~~~~~~~~-{}-{}-{}-{}D~-{}-{}-{}~-{}-{}-{}-{}-{}-{}-{}-{}~1:1001}\par
\texttt{F4.616~~13.01.2.0~~~~~~~~~~~-{}B-{}-{}D~-{}PC~+-{}-{}-{}++++~ok}\par
\texttt{F4.617~~23.01.2.0~~~~~~~~~~~-{}-{}-{}-{}D~-{}-{}-{}~-{}-{}-{}-{}-{}-{}-{}-{}~0:0010}\par
\texttt{F4.618~~01.12.2.0~~~~~~~~~~~-{}-{}-{}-{}D~-{}-{}-{}~-{}-{}-{}-{}-{}-{}-{}-{}~1:0010}\par
\texttt{F4.619~~02.12.2.0~~~~~~~~~~~-{}-{}-{}-{}D~-{}-{}-{}~-{}-{}-{}-{}-{}-{}-{}-{}~0:0010}\par
\texttt{F4.620~~12.12.2.0~~~~~~~~~~~-{}-{}-{}-{}D~-{}-{}-{}~-{}-{}-{}-{}-{}-{}-{}-{}~0:0010}\par
\texttt{F4.621~~03.12.2.0~~~~~~~~~~~-{}-{}-{}-{}D~-{}-{}-{}~-{}-{}-{}-{}-{}-{}-{}-{}~1:0010}\par
\texttt{F4.622~~13.12.2.0~~~~~~~~~~~-{}-{}-{}-{}D~-{}-{}-{}~-{}-{}-{}-{}-{}-{}-{}-{}~1:0010}\par
\texttt{F4.623~~23.12.2.0~~~~~~~~~~~-{}-{}-{}-{}D~-{}-{}-{}~-{}-{}-{}-{}-{}-{}-{}-{}~0:0010}\par
\texttt{F4.624~~0.012.2.0~~~~~~~~~~~-{}-{}4-{}D~-{}-{}-{}~-{}-{}-{}+-{}-{}-{}-{}~1:1000}\par
\texttt{F4.625~~1.012.2.0~~~~~~~~~~~-{}-{}-{}-{}D~-{}-{}-{}~-{}-{}-{}-{}-{}-{}-{}-{}~1:0010}\par
\texttt{F4.626~~2.012.2.0~~~~~~~~~~~-{}-{}-{}-{}D~-{}-{}-{}~-{}-{}-{}-{}-{}-{}-{}-{}~0:0010}\par
\texttt{F4.627~~3.012.2.0~~~~~~~~~~~-{}-{}-{}-{}D~-{}-{}-{}~-{}-{}-{}-{}-{}-{}-{}-{}~1:0010}\par
\texttt{F4.628~~012.3.2.0~~~~~~~~~~~-{}-{}-{}-{}D~-{}-{}-{}~-{}-{}-{}-{}-{}-{}-{}-{}~0:0010}\par
\texttt{F4.629~~013.3.2.0~~~~~~~~~~~-{}-{}-{}-{}D~-{}PC~-{}-{}-{}-{}++-{}+~ok}\par
\texttt{F4.630~~123.3.2.0~~~~~~~~~~~-{}-{}-{}-{}D~-{}-{}-{}~-{}-{}-{}-{}-{}-{}-{}-{}~0:0010}\par
\texttt{F4.631~~01.03.2.0~~~~~~~~~~~-{}-{}-{}-{}D~-{}PC~-{}-{}-{}-{}++++~ok}\par
\texttt{F4.632~~02.03.2.0~~~~~~~~~~~-{}-{}-{}-{}D~-{}-{}-{}~-{}-{}-{}-{}-{}-{}-{}-{}~0:0010}\par
\texttt{F4.633~~12.03.2.0~~~~~~~~~~~-{}-{}-{}-{}D~-{}-{}-{}~-{}-{}-{}-{}-{}-{}-{}-{}~0:0010}\par
\texttt{F4.634~~03.03.2.0~~~~~~~~~~~-{}-{}-{}-{}D~-{}P-{}~-{}-{}-{}-{}-{}+-{}-{}~1:1001}\par
\texttt{F4.635~~13.03.2.0~~~~~~~~~~~-{}-{}-{}-{}D~-{}PC~-{}-{}-{}-{}++++~ok}\par
\texttt{F4.636~~23.03.2.0~~~~~~~~~~~-{}-{}-{}-{}D~-{}-{}-{}~-{}-{}-{}-{}-{}-{}-{}-{}~0:0010}\par
\texttt{F4.637~~01.13.2.0~~~~~~~~~~~-{}-{}-{}-{}D~-{}PC~-{}-{}-{}-{}++-{}+~ok}\par
\texttt{F4.638~~02.13.2.0~~~~~~~~~~~-{}-{}-{}-{}D~-{}-{}-{}~-{}-{}-{}-{}-{}-{}-{}-{}~0:0010}\par
\texttt{F4.639~~12.13.2.0~~~~~~~~~~~-{}-{}-{}-{}D~-{}-{}-{}~-{}-{}-{}-{}-{}-{}-{}-{}~0:0010}\par
\texttt{F4.640~~03.13.2.0~~~~~~~~~~~-{}-{}-{}-{}D~-{}-{}-{}~-{}-{}-{}-{}-{}-{}-{}-{}~1:1001}\par
\texttt{F4.641~~13.13.2.0~~~~~~~~~~~-{}-{}-{}-{}D~-{}PC~-{}-{}-{}-{}++++~ok}\par
\texttt{F4.642~~23.13.2.0~~~~~~~~~~~-{}-{}-{}-{}D~-{}-{}-{}~-{}-{}-{}-{}-{}-{}-{}-{}~0:0010}\par
\texttt{F4.643~~0.013.2.0~~~~~~~~~~~-{}-{}4-{}D~-{}-{}-{}~-{}-{}-{}+-{}-{}-{}-{}~1:1000}\par
\texttt{F4.644~~2.013.2.0~~~~~~~~~~~-{}-{}-{}-{}D~-{}-{}-{}~-{}-{}-{}-{}-{}-{}-{}-{}~0:0010}\par
\texttt{F4.645~~3.013.2.0~~~~~~~~~~~-{}-{}-{}-{}D~-{}-{}-{}~-{}-{}-{}-{}-{}-{}-{}-{}~1:1001}\par
\texttt{F4.646~~01.23.2.0~~~~~~~~~~~-{}-{}-{}-{}D~-{}-{}-{}~-{}-{}-{}-{}-{}-{}-{}-{}~1:0010}\par
\texttt{F4.647~~02.23.2.0~~~~~~~~~~~-{}-{}-{}-{}D~-{}-{}-{}~-{}-{}-{}-{}-{}-{}-{}-{}~0:0010}\par
\texttt{F4.648~~12.23.2.0~~~~~~~~~~~-{}-{}-{}-{}D~-{}-{}-{}~-{}-{}-{}-{}-{}-{}-{}-{}~0:0010}\par
\texttt{F4.649~~03.23.2.0~~~~~~~~~~~-{}-{}-{}-{}D~-{}-{}-{}~-{}-{}-{}-{}-{}-{}-{}-{}~1:0010}\par
\texttt{F4.650~~13.23.2.0~~~~~~~~~~~-{}-{}-{}-{}D~-{}-{}-{}~-{}-{}-{}-{}-{}-{}-{}-{}~1:0010}\par
\texttt{F4.651~~23.23.2.0~~~~~~~~~~~-{}-{}-{}-{}D~-{}-{}-{}~-{}-{}-{}-{}-{}-{}-{}-{}~0:0010}\par
\texttt{F4.652~~0.023.2.0~~~~~~~~~~~-{}-{}4-{}D~-{}-{}-{}~-{}-{}-{}+-{}-{}-{}-{}~1:1000}\par
\texttt{F4.653~~2.023.2.0~~~~~~~~~~~-{}-{}-{}-{}D~-{}-{}-{}~-{}-{}-{}-{}-{}-{}-{}-{}~0:0010}\par
\texttt{F4.654~~3.023.2.0~~~~~~~~~~~-{}-{}-{}-{}D~-{}-{}-{}~-{}-{}-{}-{}-{}-{}-{}-{}~1:0010}\par
\texttt{F4.655~~0.123.2.0~~~~~~~~~~~-{}-{}-{}-{}D~-{}-{}-{}~-{}-{}-{}-{}-{}-{}-{}-{}~1:1000}\par
\texttt{F4.656~~2.123.2.0~~~~~~~~~~~-{}-{}-{}-{}D~-{}-{}-{}~-{}-{}-{}-{}-{}-{}-{}-{}~0:0010}\par
\texttt{F4.657~~0.01.02.0~~~~~~~~~~~-{}-{}4-{}D~-{}-{}-{}~-{}-{}-{}+-{}-{}-{}-{}~1:1000}\par
\texttt{F4.658~~1.01.02.0~~~~~~~~~~~-{}-{}-{}-{}D~-{}-{}-{}~-{}-{}-{}-{}-{}-{}-{}-{}~2:1100}\par
\texttt{F4.659~~3.01.02.0~~~~~~~~~~~-{}-{}-{}-{}D~-{}-{}-{}~-{}-{}-{}-{}-{}-{}-{}-{}~1:1001}\par
\texttt{F4.660~~0.12.02.0~~~~~~~~~~~-{}-{}-{}-{}D~-{}-{}-{}~-{}-{}-{}-{}-{}-{}-{}-{}~2:1000}\par
\texttt{F4.661~~1.12.02.0~~~~~~~~~~~-{}-{}-{}-{}D~-{}P-{}~-{}-{}-{}-{}-{}+-{}-{}~3:1110}\par
\texttt{F4.662~~2.12.02.0~~~~~~~~~~~-{}-{}-{}-{}D~-{}-{}-{}~-{}-{}-{}-{}-{}-{}-{}-{}~1:1010}\par
\texttt{F4.663~~3.12.02.0~~~~~~~~~~~-{}-{}-{}-{}D~-{}-{}-{}~-{}-{}-{}-{}-{}-{}-{}-{}~2:1001}\par
\texttt{F4.664~~01.3.02.0~~~~~~~~~~~-{}-{}-{}-{}D~-{}-{}-{}~-{}-{}-{}-{}-{}-{}-{}-{}~2:1101}\par
\texttt{F4.665~~02.3.02.0~~~~~~~~~~~-{}-{}-{}-{}D~UP-{}~-{}-{}+-{}-{}+-{}-{}~1:1010}\par
\texttt{F4.666~~12.3.02.0~~~~~~~~~~~-{}-{}-{}-{}D~-{}PC~-{}-{}-{}-{}++-{}+~ok}\par
\texttt{F4.667~~13.3.02.0~~~~~~~~~~~-{}-{}-{}-{}D~-{}-{}-{}~-{}-{}-{}-{}-{}-{}-{}-{}~2:1101}\par
\texttt{F4.668~~0.03.02.0~~~~~~~~~~~-{}-{}4-{}D~-{}-{}-{}~-{}-{}-{}+-{}-{}-{}-{}~1:1000}\par
\texttt{F4.669~~1.03.02.0~~~~~~~~~~~-{}-{}-{}-{}D~-{}-{}-{}~-{}-{}-{}-{}-{}-{}-{}-{}~2:1101}\par
\texttt{F4.670~~2.03.02.0~~~~~~~~~~~-{}-{}-{}-{}D~-{}P-{}~-{}-{}-{}-{}-{}+-{}-{}~1:1010}\par
\texttt{F4.671~~3.03.02.0~~~~~~~~~~~-{}-{}-{}-{}D~-{}-{}-{}~-{}-{}-{}-{}-{}-{}-{}-{}~1:1001}\par
\texttt{F4.672~~0.13.02.0~~~~~~~~~~~-{}-{}-{}-{}D~-{}-{}-{}~-{}-{}-{}-{}-{}-{}-{}-{}~1:1000}\par
\texttt{F4.673~~1.13.02.0~~~~~~~~~~~-{}-{}-{}-{}D~-{}-{}-{}~-{}-{}-{}-{}-{}-{}-{}-{}~2:1101}\par
\texttt{F4.674~~2.13.02.0~~~~~~~~~~~-{}-{}-{}-{}D~-{}-{}-{}~-{}-{}-{}-{}-{}-{}-{}-{}~1:1010}\par
\texttt{F4.675~~3.13.02.0~~~~~~~~~~~-{}-{}-{}-{}D~-{}-{}-{}~-{}-{}-{}-{}-{}-{}-{}-{}~1:1001}\par
\texttt{F4.676~~0.23.02.0~~~~~~~~~~~-{}-{}-{}-{}D~-{}-{}-{}~-{}-{}-{}-{}-{}-{}-{}-{}~1:1000}\par
\texttt{F4.677~~1.23.02.0~~~~~~~~~~~-{}-{}-{}-{}D~-{}PC~-{}-{}-{}-{}++-{}+~ok}\par
\texttt{F4.678~~0.12.12.0~~~~~~~~~~~-{}-{}45D~UP-{}~-{}+++-{}+-{}-{}~3:1000}\par
\texttt{F4.679~~3.12.12.0~~~~~~~~~~~-{}B-{}-{}D~UPC~+-{}+-{}++++~ok}\par
\texttt{F4.680~~01.3.12.0~~~~~~~~~~~-{}-{}-{}-{}D~-{}-{}-{}~-{}-{}-{}-{}-{}-{}-{}-{}~2:1101}\par
\texttt{F4.681~~02.3.12.0~~~~~~~~~~~-{}-{}-{}-{}D~-{}PC~-{}-{}-{}-{}++-{}-{}~ok}\par
\texttt{F4.682~~12.3.12.0~~~~~~~~~~~-{}-{}-{}-{}D~-{}PC~-{}-{}-{}-{}++-{}+~ok}\par
\texttt{F4.683~~13.3.12.0~~~~~~~~~~~-{}-{}-{}-{}D~-{}-{}-{}~-{}-{}-{}-{}-{}-{}-{}-{}~2:1101}\par
\texttt{F4.684~~0.03.12.0~~~~~~~~~~~-{}-{}-{}-{}D~-{}-{}-{}~-{}-{}-{}-{}-{}-{}-{}-{}~1:1000}\par
\texttt{F4.685~~2.03.12.0~~~~~~~~~~~-{}-{}-{}-{}D~-{}PC~-{}-{}-{}-{}++-{}+~ok}\par
\texttt{F4.686~~3.03.12.0~~~~~~~~~~~-{}-{}-{}-{}D~-{}-{}-{}~-{}-{}-{}-{}-{}-{}-{}-{}~1:1001}\par
\texttt{F4.687~~0.13.12.0~~~~~~~~~~~-{}-{}-{}-{}D~-{}-{}-{}~-{}-{}-{}-{}-{}-{}-{}-{}~1:1000}\par
\texttt{F4.688~~0.23.12.0~~~~~~~~~~~-{}-{}-{}-{}D~-{}-{}-{}~-{}-{}-{}-{}-{}-{}-{}-{}~1:1000}\par
\texttt{F4.689~~1.3.012.0~~~~~~~~~~~-{}-{}-{}-{}D~-{}-{}-{}~-{}-{}-{}-{}-{}-{}-{}-{}~2:1101}\par
\texttt{F4.690~~02.23.3.0~~~~~~~~~~~-{}-{}-{}-{}D~-{}P-{}~-{}-{}-{}-{}-{}+-{}-{}~1:1011}\par
\texttt{F4.691~~12.23.3.0~~~~~~~~~~~-{}-{}-{}-{}D~-{}PC~-{}-{}-{}-{}++-{}+~ok}\par
\texttt{F4.692~~23.23.3.0~~~~~~~~~~~-{}-{}-{}-{}D~-{}P-{}~-{}-{}-{}-{}-{}+-{}-{}~1:1011}\par
\texttt{F4.693~~2.023.3.0~~~~~~~~~~~-{}-{}-{}-{}D~-{}P-{}~-{}-{}-{}-{}-{}+-{}-{}~1:1011}\par
\texttt{F4.694~~0.03.03.0~~~~~~~~~~~-{}-{}4-{}D~UP-{}~-{}-{}++-{}+-{}-{}~1:1000}\par
\texttt{F4.695~~3.03.03.0~~~~~~~~~~~-{}-{}-{}-{}D~-{}P-{}~-{}-{}-{}-{}-{}+-{}-{}~1:1001}\par
\texttt{F4.696~~0.13.03.0~~~~~~~~~~~-{}-{}-{}-{}D~-{}-{}-{}~-{}-{}-{}-{}-{}-{}-{}-{}~1:1000}\par
\texttt{F4.697~~0.23.03.0~~~~~~~~~~~-{}-{}-{}-{}D~-{}-{}-{}~-{}-{}-{}-{}-{}-{}-{}-{}~1:1000}\par
\texttt{F4.698~~0.13.13.0~~~~~~~~~~~-{}-{}-{}-{}D~-{}-{}-{}~-{}-{}-{}-{}-{}-{}-{}-{}~1:1000}\par
\texttt{F4.699~~0.23.13.0~~~~~~~~~~~-{}-{}-{}-{}D~-{}-{}-{}~-{}-{}-{}-{}-{}-{}-{}-{}~1:1000}\par
\texttt{F4.700~~0.13.23.0~~~~~~~~~~~-{}-{}-{}-{}D~-{}-{}-{}~-{}-{}-{}-{}-{}-{}-{}-{}~1:1000}\par
\texttt{F4.701~~0.1.01.01~~~~~~~~~~~-{}-{}4-{}D~-{}-{}-{}~-{}-{}-{}+-{}-{}-{}-{}~2:1000}\par
\texttt{F4.702~~01.1.2.01~~~~~~~~~~~-{}-{}4-{}D~-{}-{}-{}~-{}-{}-{}+-{}-{}-{}-{}~0:0100}\par
\texttt{F4.703~~02.1.2.01~~~~~~~~~~~-{}-{}-{}-{}D~-{}-{}-{}~-{}-{}-{}-{}-{}-{}-{}-{}~3:0100}\par
\texttt{F4.704~~12.1.2.01~~~~~~~~~~~-{}-{}-{}-{}D~-{}-{}-{}~-{}-{}-{}-{}-{}-{}-{}-{}~0:0100}\par
\texttt{F4.705~~03.1.2.01~~~~~~~~~~~-{}-{}-{}-{}D~-{}-{}-{}~-{}-{}-{}-{}-{}-{}-{}-{}~3:0100}\par
\texttt{F4.706~~13.1.2.01~~~~~~~~~~~-{}-{}-{}-{}D~-{}-{}-{}~-{}-{}-{}-{}-{}-{}-{}-{}~0:0100}\par
\texttt{F4.707~~23.1.2.01~~~~~~~~~~~-{}-{}-{}-{}D~-{}-{}-{}~-{}-{}-{}-{}-{}-{}-{}-{}~3:0100}\par
\texttt{F4.708~~0.1.02.01~~~~~~~~~~~-{}-{}4-{}D~-{}-{}-{}~-{}-{}-{}+-{}-{}-{}-{}~2:1000}\par
\texttt{F4.709~~0.1.23.01~~~~~~~~~~~-{}-{}-{}-{}D~-{}-{}-{}~-{}-{}-{}-{}-{}-{}-{}-{}~3:1000}\par
\texttt{F4.710~~0.1.2.012~~~~~~~~~~~-{}-{}4-{}D~-{}-{}-{}~-{}-{}-{}+-{}-{}-{}-{}~3:1000}\par
\texttt{F4.711~~012.1.2.3~~~~~~~~~~~T-{}4-{}D~-{}-{}-{}~-{}-{}-{}+-{}-{}-{}-{}~0:0100}\par
\texttt{F4.712~~01.01.2.3~~~~~~~~~~~TB45D~UPC~++++++++~ok}\par
\texttt{F4.713~~02.01.2.3~~~~~~~~~~~T-{}-{}-{}D~-{}-{}-{}~-{}-{}-{}-{}-{}-{}-{}-{}~0:0010}\par
\texttt{F4.714~~02.12.2.3~~~~~~~~~~~T-{}4-{}D~-{}-{}-{}~-{}-{}-{}+-{}-{}-{}-{}~0:0010}\par
\texttt{F4.715~~03.12.2.3~~~~~~~~~~~T-{}4-{}D~-{}-{}-{}~-{}-{}-{}+-{}-{}-{}-{}~1:0010}\par
\texttt{F4.716~~0123.012.-{}.-{}~~~~~~~~-{}-{}-{}-{}-{}~-{}-{}-{}~-{}-{}-{}-{}-{}-{}-{}-{}~0:0000}\par
\texttt{F4.717~~0123.023.-{}.-{}~~~~~~~~-{}-{}-{}-{}-{}~-{}-{}-{}~-{}-{}-{}-{}-{}-{}-{}-{}~0:0000}\par
\texttt{F4.718~~0123.123.-{}.-{}~~~~~~~~-{}-{}4-{}-{}~-{}-{}-{}~-{}-{}-{}+-{}-{}-{}-{}~0:0000}\par
\texttt{F4.719~~0123.01.0.-{}~~~~~~~~~-{}-{}-{}-{}-{}~-{}-{}-{}~-{}-{}-{}-{}-{}-{}-{}-{}~0:0000}\par
\texttt{F4.720~~0123.02.0.-{}~~~~~~~~~-{}-{}-{}-{}-{}~-{}-{}-{}~-{}-{}-{}-{}-{}-{}-{}-{}~0:0000}\par
\texttt{F4.721~~0123.12.0.-{}~~~~~~~~~-{}-{}-{}-{}-{}~-{}-{}-{}~-{}-{}-{}-{}-{}-{}-{}-{}~0:0000}\par
\texttt{F4.722~~012.012.0.-{}~~~~~~~~~-{}-{}-{}-{}-{}~-{}PC~-{}-{}-{}-{}++++~ok}\par
\texttt{F4.723~~013.012.0.-{}~~~~~~~~~-{}-{}-{}-{}-{}~-{}-{}-{}~-{}-{}-{}-{}-{}-{}-{}-{}~0:0000}\par
\texttt{F4.724~~023.012.0.-{}~~~~~~~~~-{}-{}-{}-{}-{}~-{}-{}-{}~-{}-{}-{}-{}-{}-{}-{}-{}~0:0000}\par
\texttt{F4.725~~123.012.0.-{}~~~~~~~~~-{}-{}-{}-{}-{}~-{}-{}-{}~-{}-{}-{}-{}-{}-{}-{}-{}~0:0000}\par
\texttt{F4.726~~0123.03.0.-{}~~~~~~~~~-{}-{}-{}-{}-{}~-{}-{}-{}~-{}-{}-{}-{}-{}-{}-{}-{}~0:0000}\par
\texttt{F4.727~~0123.13.0.-{}~~~~~~~~~-{}-{}-{}-{}-{}~-{}-{}-{}~-{}-{}-{}-{}-{}-{}-{}-{}~0:0000}\par
\texttt{F4.728~~012.013.0.-{}~~~~~~~~~-{}-{}-{}-{}-{}~-{}-{}-{}~-{}-{}-{}-{}-{}-{}-{}-{}~1:0000}\par
\texttt{F4.729~~013.013.0.-{}~~~~~~~~~-{}-{}-{}-{}-{}~-{}-{}-{}~-{}-{}-{}-{}-{}-{}-{}-{}~0:0000}\par
\texttt{F4.730~~023.013.0.-{}~~~~~~~~~-{}-{}-{}-{}-{}~-{}-{}-{}~-{}-{}-{}-{}-{}-{}-{}-{}~0:0000}\par
\texttt{F4.731~~123.013.0.-{}~~~~~~~~~-{}-{}-{}-{}-{}~-{}-{}-{}~-{}-{}-{}-{}-{}-{}-{}-{}~0:0000}\par
\texttt{F4.732~~0123.23.0.-{}~~~~~~~~~-{}-{}-{}-{}-{}~-{}-{}-{}~-{}-{}-{}-{}-{}-{}-{}-{}~0:0000}\par
\texttt{F4.733~~012.023.0.-{}~~~~~~~~~-{}-{}-{}-{}-{}~-{}-{}-{}~-{}-{}-{}-{}-{}-{}-{}-{}~1:0000}\par
\texttt{F4.734~~013.023.0.-{}~~~~~~~~~-{}-{}-{}-{}-{}~-{}-{}-{}~-{}-{}-{}-{}-{}-{}-{}-{}~0:0000}\par
\texttt{F4.735~~023.023.0.-{}~~~~~~~~~-{}-{}-{}-{}-{}~-{}-{}-{}~-{}-{}-{}-{}-{}-{}-{}-{}~0:0000}\par
\texttt{F4.736~~123.023.0.-{}~~~~~~~~~-{}-{}-{}-{}-{}~-{}-{}-{}~-{}-{}-{}-{}-{}-{}-{}-{}~0:0000}\par
\texttt{F4.737~~012.123.0.-{}~~~~~~~~~-{}-{}-{}-{}-{}~-{}-{}-{}~-{}-{}-{}-{}-{}-{}-{}-{}~1:0000}\par
\texttt{F4.738~~013.123.0.-{}~~~~~~~~~-{}-{}-{}-{}-{}~-{}-{}-{}~-{}-{}-{}-{}-{}-{}-{}-{}~0:0000}\par
\texttt{F4.739~~023.123.0.-{}~~~~~~~~~-{}-{}-{}-{}-{}~-{}-{}-{}~-{}-{}-{}-{}-{}-{}-{}-{}~0:0000}\par
\texttt{F4.740~~123.123.0.-{}~~~~~~~~~-{}-{}-{}-{}-{}~-{}-{}-{}~-{}-{}-{}-{}-{}-{}-{}-{}~0:0000}\par
\texttt{F4.741~~01.0123.0.-{}~~~~~~~~~-{}-{}-{}-{}-{}~-{}-{}-{}~-{}-{}-{}-{}-{}-{}-{}-{}~1:0000}\par
\texttt{F4.742~~02.0123.0.-{}~~~~~~~~~-{}-{}-{}-{}-{}~-{}-{}-{}~-{}-{}-{}-{}-{}-{}-{}-{}~1:0000}\par
\texttt{F4.743~~12.0123.0.-{}~~~~~~~~~-{}-{}-{}-{}-{}~-{}-{}-{}~-{}-{}-{}-{}-{}-{}-{}-{}~1:0000}\par
\texttt{F4.744~~03.0123.0.-{}~~~~~~~~~-{}-{}-{}-{}-{}~-{}-{}-{}~-{}-{}-{}-{}-{}-{}-{}-{}~0:0000}\par
\texttt{F4.745~~13.0123.0.-{}~~~~~~~~~-{}-{}-{}-{}-{}~-{}-{}-{}~-{}-{}-{}-{}-{}-{}-{}-{}~0:0000}\par
\texttt{F4.746~~23.0123.0.-{}~~~~~~~~~-{}-{}-{}-{}-{}~-{}-{}-{}~-{}-{}-{}-{}-{}-{}-{}-{}~0:0000}\par
\texttt{F4.747~~0123.1.01.-{}~~~~~~~~~-{}-{}-{}-{}-{}~-{}-{}-{}~-{}-{}-{}-{}-{}-{}-{}-{}~0:0000}\par
\texttt{F4.748~~012.01.01.-{}~~~~~~~~~-{}-{}-{}-{}-{}~-{}PC~-{}-{}-{}-{}++++~ok}\par
\texttt{F4.749~~013.01.01.-{}~~~~~~~~~-{}-{}-{}-{}-{}~-{}-{}-{}~-{}-{}-{}-{}-{}-{}-{}-{}~0:0000}\par
\texttt{F4.750~~023.01.01.-{}~~~~~~~~~-{}-{}-{}-{}-{}~-{}-{}-{}~-{}-{}-{}-{}-{}-{}-{}-{}~0:0000}\par
\texttt{F4.751~~123.01.01.-{}~~~~~~~~~-{}-{}-{}-{}-{}~-{}-{}-{}~-{}-{}-{}-{}-{}-{}-{}-{}~0:0000}\par
\texttt{F4.752~~012.02.01.-{}~~~~~~~~~-{}B-{}-{}-{}~-{}PC~+-{}-{}-{}++++~ok}\par
\texttt{F4.753~~013.02.01.-{}~~~~~~~~~-{}-{}-{}-{}-{}~-{}-{}-{}~-{}-{}-{}-{}-{}-{}-{}-{}~0:0000}\par
\texttt{F4.754~~123.02.01.-{}~~~~~~~~~-{}-{}-{}-{}-{}~-{}-{}-{}~-{}-{}-{}-{}-{}-{}-{}-{}~0:0000}\par
\texttt{F4.755~~012.12.01.-{}~~~~~~~~~-{}-{}-{}-{}-{}~-{}PC~-{}-{}-{}-{}++++~ok}\par
\texttt{F4.756~~013.12.01.-{}~~~~~~~~~-{}-{}-{}-{}-{}~-{}-{}-{}~-{}-{}-{}-{}-{}-{}-{}-{}~0:0000}\par
\texttt{F4.757~~023.12.01.-{}~~~~~~~~~-{}-{}-{}-{}-{}~-{}-{}-{}~-{}-{}-{}-{}-{}-{}-{}-{}~0:0000}\par
\texttt{F4.758~~123.12.01.-{}~~~~~~~~~-{}-{}-{}-{}-{}~-{}-{}-{}~-{}-{}-{}-{}-{}-{}-{}-{}~0:0000}\par
\texttt{F4.759~~03.012.01.-{}~~~~~~~~~-{}-{}-{}-{}-{}~-{}-{}-{}~-{}-{}-{}-{}-{}-{}-{}-{}~0:0000}\par
\texttt{F4.760~~13.012.01.-{}~~~~~~~~~-{}-{}-{}-{}-{}~-{}-{}-{}~-{}-{}-{}-{}-{}-{}-{}-{}~0:0000}\par
\texttt{F4.761~~23.012.01.-{}~~~~~~~~~-{}-{}-{}-{}-{}~-{}-{}-{}~-{}-{}-{}-{}-{}-{}-{}-{}~0:0000}\par
\texttt{F4.762~~0123.3.01.-{}~~~~~~~~~-{}-{}-{}-{}-{}~-{}-{}-{}~-{}-{}-{}-{}-{}-{}-{}-{}~0:0000}\par
\texttt{F4.763~~013.03.01.-{}~~~~~~~~~-{}-{}-{}-{}-{}~-{}-{}-{}~-{}-{}-{}-{}-{}-{}-{}-{}~0:0000}\par
\texttt{F4.764~~023.03.01.-{}~~~~~~~~~-{}-{}-{}-{}-{}~-{}-{}-{}~-{}-{}-{}-{}-{}-{}-{}-{}~0:0000}\par
\texttt{F4.765~~123.03.01.-{}~~~~~~~~~-{}-{}-{}-{}-{}~-{}-{}-{}~-{}-{}-{}-{}-{}-{}-{}-{}~0:0000}\par
\texttt{F4.766~~013.13.01.-{}~~~~~~~~~-{}-{}-{}-{}-{}~-{}-{}-{}~-{}-{}-{}-{}-{}-{}-{}-{}~0:0000}\par
\texttt{F4.767~~023.13.01.-{}~~~~~~~~~-{}-{}-{}-{}-{}~-{}-{}-{}~-{}-{}-{}-{}-{}-{}-{}-{}~0:0000}\par
\texttt{F4.768~~123.13.01.-{}~~~~~~~~~-{}-{}-{}-{}-{}~-{}-{}-{}~-{}-{}-{}-{}-{}-{}-{}-{}~0:0000}\par
\texttt{F4.769~~23.013.01.-{}~~~~~~~~~-{}-{}-{}-{}-{}~-{}-{}-{}~-{}-{}-{}-{}-{}-{}-{}-{}~0:0000}\par
\texttt{F4.770~~023.23.01.-{}~~~~~~~~~-{}-{}-{}-{}-{}~-{}-{}-{}~-{}-{}-{}-{}-{}-{}-{}-{}~0:0000}\par
\texttt{F4.771~~123.23.01.-{}~~~~~~~~~-{}-{}-{}-{}-{}~-{}-{}-{}~-{}-{}-{}-{}-{}-{}-{}-{}~0:0000}\par
\texttt{F4.772~~0123.01.2.-{}~~~~~~~~~-{}-{}-{}-{}-{}~-{}-{}-{}~-{}-{}-{}-{}-{}-{}-{}-{}~0:0000}\par
\texttt{F4.773~~0123.12.2.-{}~~~~~~~~~-{}-{}4-{}-{}~-{}-{}-{}~-{}-{}-{}+-{}-{}-{}-{}~0:0000}\par
\texttt{F4.774~~012.012.2.-{}~~~~~~~~~-{}-{}4-{}-{}~-{}-{}-{}~-{}-{}-{}+-{}-{}-{}-{}~0:0010}\par
\texttt{F4.775~~013.012.2.-{}~~~~~~~~~-{}-{}-{}-{}-{}~-{}-{}-{}~-{}-{}-{}-{}-{}-{}-{}-{}~0:0000}\par
\texttt{F4.776~~023.012.2.-{}~~~~~~~~~-{}-{}-{}-{}-{}~-{}-{}-{}~-{}-{}-{}-{}-{}-{}-{}-{}~0:0000}\par
\texttt{F4.777~~123.012.2.-{}~~~~~~~~~-{}-{}-{}-{}-{}~-{}-{}-{}~-{}-{}-{}-{}-{}-{}-{}-{}~0:0000}\par
\texttt{F4.778~~0123.03.2.-{}~~~~~~~~~-{}-{}-{}-{}-{}~-{}-{}-{}~-{}-{}-{}-{}-{}-{}-{}-{}~0:0000}\par
\texttt{F4.779~~0123.13.2.-{}~~~~~~~~~-{}-{}4-{}-{}~-{}-{}-{}~-{}-{}-{}+-{}-{}-{}-{}~0:0000}\par
\texttt{F4.780~~013.013.2.-{}~~~~~~~~~-{}-{}4-{}-{}~-{}-{}-{}~-{}-{}-{}+-{}-{}-{}-{}~0:0000}\par
\texttt{F4.781~~023.013.2.-{}~~~~~~~~~-{}-{}-{}-{}-{}~-{}-{}-{}~-{}-{}-{}-{}-{}-{}-{}-{}~0:0000}\par
\texttt{F4.782~~123.013.2.-{}~~~~~~~~~-{}-{}-{}-{}-{}~-{}-{}-{}~-{}-{}-{}-{}-{}-{}-{}-{}~0:0000}\par
\texttt{F4.783~~0123.23.2.-{}~~~~~~~~~-{}-{}4-{}-{}~-{}-{}-{}~-{}-{}-{}+-{}-{}-{}-{}~0:0000}\par
\texttt{F4.784~~023.023.2.-{}~~~~~~~~~-{}-{}4-{}-{}~-{}-{}-{}~-{}-{}-{}+-{}-{}-{}-{}~0:0000}\par
\texttt{F4.785~~123.023.2.-{}~~~~~~~~~-{}-{}-{}-{}-{}~-{}-{}-{}~-{}-{}-{}-{}-{}-{}-{}-{}~0:0000}\par
\texttt{F4.786~~023.123.2.-{}~~~~~~~~~-{}-{}4-{}-{}~-{}-{}-{}~-{}-{}-{}+-{}-{}-{}-{}~0:0000}\par
\texttt{F4.787~~012.01.02.-{}~~~~~~~~~-{}B-{}-{}-{}~-{}PC~+-{}-{}-{}++++~ok}\par
\texttt{F4.788~~013.01.02.-{}~~~~~~~~~-{}-{}-{}-{}-{}~-{}-{}-{}~-{}-{}-{}-{}-{}-{}-{}-{}~0:0000}\par
\texttt{F4.789~~123.01.02.-{}~~~~~~~~~-{}-{}-{}-{}-{}~-{}-{}-{}~-{}-{}-{}-{}-{}-{}-{}-{}~0:0000}\par
\texttt{F4.790~~012.12.02.-{}~~~~~~~~~-{}-{}-{}-{}-{}~-{}PC~-{}-{}-{}-{}++++~ok}\par
\texttt{F4.791~~013.12.02.-{}~~~~~~~~~-{}-{}-{}-{}-{}~-{}-{}-{}~-{}-{}-{}-{}-{}-{}-{}-{}~0:0000}\par
\texttt{F4.792~~023.12.02.-{}~~~~~~~~~-{}-{}-{}-{}-{}~-{}-{}-{}~-{}-{}-{}-{}-{}-{}-{}-{}~0:0000}\par
\texttt{F4.793~~123.12.02.-{}~~~~~~~~~-{}-{}-{}-{}-{}~-{}-{}-{}~-{}-{}-{}-{}-{}-{}-{}-{}~0:0000}\par
\texttt{F4.794~~02.012.02.-{}~~~~~~~~~-{}-{}4-{}-{}~-{}-{}-{}~-{}-{}-{}+-{}-{}-{}-{}~1:1010}\par
\texttt{F4.795~~03.012.02.-{}~~~~~~~~~-{}-{}-{}-{}-{}~-{}-{}-{}~-{}-{}-{}-{}-{}-{}-{}-{}~0:0000}\par
\texttt{F4.796~~13.012.02.-{}~~~~~~~~~-{}-{}-{}-{}-{}~-{}-{}-{}~-{}-{}-{}-{}-{}-{}-{}-{}~0:0000}\par
\texttt{F4.797~~23.012.02.-{}~~~~~~~~~-{}-{}-{}-{}-{}~-{}-{}-{}~-{}-{}-{}-{}-{}-{}-{}-{}~0:0000}\par
\texttt{F4.798~~0123.3.02.-{}~~~~~~~~~-{}-{}-{}-{}-{}~-{}-{}-{}~-{}-{}-{}-{}-{}-{}-{}-{}~0:0000}\par
\texttt{F4.799~~012.03.02.-{}~~~~~~~~~-{}-{}-{}-{}-{}~-{}-{}-{}~-{}-{}-{}-{}-{}-{}-{}-{}~1:0000}\par
\texttt{F4.800~~013.03.02.-{}~~~~~~~~~-{}-{}-{}-{}-{}~-{}-{}-{}~-{}-{}-{}-{}-{}-{}-{}-{}~0:0000}\par
\texttt{F4.801~~023.03.02.-{}~~~~~~~~~-{}-{}-{}-{}-{}~-{}-{}-{}~-{}-{}-{}-{}-{}-{}-{}-{}~0:0000}\par
\texttt{F4.802~~123.03.02.-{}~~~~~~~~~-{}-{}-{}-{}-{}~-{}-{}-{}~-{}-{}-{}-{}-{}-{}-{}-{}~0:0000}\par
\texttt{F4.803~~012.13.02.-{}~~~~~~~~~-{}-{}-{}-{}-{}~-{}-{}-{}~-{}-{}-{}-{}-{}-{}-{}-{}~1:0000}\par
\texttt{F4.804~~013.13.02.-{}~~~~~~~~~-{}-{}-{}-{}-{}~-{}-{}-{}~-{}-{}-{}-{}-{}-{}-{}-{}~0:0000}\par
\texttt{F4.805~~023.13.02.-{}~~~~~~~~~-{}-{}-{}-{}-{}~-{}-{}-{}~-{}-{}-{}-{}-{}-{}-{}-{}~0:0000}\par
\texttt{F4.806~~123.13.02.-{}~~~~~~~~~-{}-{}-{}-{}-{}~-{}-{}-{}~-{}-{}-{}-{}-{}-{}-{}-{}~0:0000}\par
\texttt{F4.807~~02.013.02.-{}~~~~~~~~~-{}-{}-{}-{}-{}~-{}-{}-{}~-{}-{}-{}-{}-{}-{}-{}-{}~1:0000}\par
\texttt{F4.808~~03.013.02.-{}~~~~~~~~~-{}-{}-{}-{}-{}~-{}-{}-{}~-{}-{}-{}-{}-{}-{}-{}-{}~0:0000}\par
\texttt{F4.809~~13.013.02.-{}~~~~~~~~~-{}-{}-{}-{}-{}~-{}-{}-{}~-{}-{}-{}-{}-{}-{}-{}-{}~0:0000}\par
\texttt{F4.810~~23.013.02.-{}~~~~~~~~~-{}-{}-{}-{}-{}~-{}-{}-{}~-{}-{}-{}-{}-{}-{}-{}-{}~0:0000}\par
\texttt{F4.811~~012.23.02.-{}~~~~~~~~~-{}-{}-{}-{}-{}~-{}-{}-{}~-{}-{}-{}-{}-{}-{}-{}-{}~1:0000}\par
\texttt{F4.812~~013.23.02.-{}~~~~~~~~~-{}-{}-{}-{}-{}~-{}-{}-{}~-{}-{}-{}-{}-{}-{}-{}-{}~0:0000}\par
\texttt{F4.813~~023.23.02.-{}~~~~~~~~~-{}-{}-{}-{}-{}~-{}-{}-{}~-{}-{}-{}-{}-{}-{}-{}-{}~0:0000}\par
\texttt{F4.814~~123.23.02.-{}~~~~~~~~~-{}-{}-{}-{}-{}~-{}-{}-{}~-{}-{}-{}-{}-{}-{}-{}-{}~0:0000}\par
\texttt{F4.815~~02.023.02.-{}~~~~~~~~~-{}-{}4-{}-{}~-{}-{}-{}~-{}-{}-{}+-{}-{}-{}-{}~1:0000}\par
\texttt{F4.816~~03.023.02.-{}~~~~~~~~~-{}-{}-{}-{}-{}~-{}-{}-{}~-{}-{}-{}-{}-{}-{}-{}-{}~0:0000}\par
\texttt{F4.817~~13.023.02.-{}~~~~~~~~~-{}-{}-{}-{}-{}~-{}-{}-{}~-{}-{}-{}-{}-{}-{}-{}-{}~0:0000}\par
\texttt{F4.818~~23.023.02.-{}~~~~~~~~~-{}-{}-{}-{}-{}~-{}-{}-{}~-{}-{}-{}-{}-{}-{}-{}-{}~0:0000}\par
\texttt{F4.819~~03.123.02.-{}~~~~~~~~~-{}-{}-{}-{}-{}~-{}-{}-{}~-{}-{}-{}-{}-{}-{}-{}-{}~0:0000}\par
\texttt{F4.820~~13.123.02.-{}~~~~~~~~~-{}-{}-{}-{}-{}~-{}-{}-{}~-{}-{}-{}-{}-{}-{}-{}-{}~0:0000}\par
\texttt{F4.821~~23.123.02.-{}~~~~~~~~~-{}-{}-{}-{}-{}~-{}-{}-{}~-{}-{}-{}-{}-{}-{}-{}-{}~0:0000}\par
\texttt{F4.822~~3.0123.02.-{}~~~~~~~~~-{}-{}-{}-{}-{}~-{}-{}-{}~-{}-{}-{}-{}-{}-{}-{}-{}~0:0000}\par
\texttt{F4.823~~3.012.012.-{}~~~~~~~~~-{}-{}-{}-{}-{}~-{}-{}-{}~-{}-{}-{}-{}-{}-{}-{}-{}~0:0000}\par
\texttt{F4.824~~013.3.012.-{}~~~~~~~~~-{}-{}-{}-{}-{}~-{}-{}-{}~-{}-{}-{}-{}-{}-{}-{}-{}~0:0000}\par
\texttt{F4.825~~023.3.012.-{}~~~~~~~~~-{}-{}-{}-{}-{}~-{}-{}-{}~-{}-{}-{}-{}-{}-{}-{}-{}~0:0000}\par
\texttt{F4.826~~123.3.012.-{}~~~~~~~~~-{}-{}-{}-{}-{}~-{}-{}-{}~-{}-{}-{}-{}-{}-{}-{}-{}~0:0000}\par
\texttt{F4.827~~03.03.012.-{}~~~~~~~~~-{}-{}-{}-{}-{}~-{}-{}-{}~-{}-{}-{}-{}-{}-{}-{}-{}~0:0000}\par
\texttt{F4.828~~13.03.012.-{}~~~~~~~~~-{}-{}-{}-{}-{}~-{}-{}-{}~-{}-{}-{}-{}-{}-{}-{}-{}~0:0000}\par
\texttt{F4.829~~23.03.012.-{}~~~~~~~~~-{}-{}-{}-{}-{}~-{}-{}-{}~-{}-{}-{}-{}-{}-{}-{}-{}~0:0000}\par
\texttt{F4.830~~03.13.012.-{}~~~~~~~~~-{}-{}-{}-{}-{}~-{}-{}-{}~-{}-{}-{}-{}-{}-{}-{}-{}~0:0000}\par
\texttt{F4.831~~23.13.012.-{}~~~~~~~~~-{}-{}-{}-{}-{}~-{}-{}-{}~-{}-{}-{}-{}-{}-{}-{}-{}~0:0000}\par
\texttt{F4.832~~23.23.012.-{}~~~~~~~~~-{}-{}-{}-{}-{}~-{}-{}-{}~-{}-{}-{}-{}-{}-{}-{}-{}~0:0000}\par
\texttt{F4.833~~0123.03.3.-{}~~~~~~~~~-{}-{}-{}-{}-{}~-{}-{}-{}~-{}-{}-{}-{}-{}-{}-{}-{}~0:0000}\par
\texttt{F4.834~~0123.13.3.-{}~~~~~~~~~-{}-{}4-{}-{}~-{}-{}-{}~-{}-{}-{}+-{}-{}-{}-{}~0:0000}\par
\texttt{F4.835~~013.013.3.-{}~~~~~~~~~-{}-{}4-{}-{}~-{}-{}-{}~-{}-{}-{}+-{}-{}-{}-{}~0:0000}\par
\texttt{F4.836~~023.013.3.-{}~~~~~~~~~-{}-{}-{}-{}-{}~-{}-{}-{}~-{}-{}-{}-{}-{}-{}-{}-{}~0:0000}\par
\texttt{F4.837~~123.013.3.-{}~~~~~~~~~-{}-{}-{}-{}-{}~-{}-{}-{}~-{}-{}-{}-{}-{}-{}-{}-{}~0:0000}\par
\texttt{F4.838~~0123.23.3.-{}~~~~~~~~~-{}-{}4-{}-{}~-{}-{}-{}~-{}-{}-{}+-{}-{}-{}-{}~0:0000}\par
\texttt{F4.839~~023.023.3.-{}~~~~~~~~~-{}-{}4-{}-{}~-{}-{}-{}~-{}-{}-{}+-{}-{}-{}-{}~0:0000}\par
\texttt{F4.840~~123.023.3.-{}~~~~~~~~~-{}-{}-{}-{}-{}~-{}-{}-{}~-{}-{}-{}-{}-{}-{}-{}-{}~0:0000}\par
\texttt{F4.841~~023.123.3.-{}~~~~~~~~~-{}-{}4-{}-{}~-{}-{}-{}~-{}-{}-{}+-{}-{}-{}-{}~0:0000}\par
\texttt{F4.842~~013.03.03.-{}~~~~~~~~~-{}-{}-{}-{}-{}~-{}-{}-{}~-{}-{}-{}-{}-{}-{}-{}-{}~0:0000}\par
\texttt{F4.843~~123.03.03.-{}~~~~~~~~~-{}-{}-{}-{}-{}~-{}-{}-{}~-{}-{}-{}-{}-{}-{}-{}-{}~0:0000}\par
\texttt{F4.844~~013.13.03.-{}~~~~~~~~~-{}-{}-{}-{}-{}~-{}-{}-{}~-{}-{}-{}-{}-{}-{}-{}-{}~0:0000}\par
\texttt{F4.845~~023.13.03.-{}~~~~~~~~~-{}-{}-{}-{}-{}~-{}-{}-{}~-{}-{}-{}-{}-{}-{}-{}-{}~0:0000}\par
\texttt{F4.846~~123.13.03.-{}~~~~~~~~~-{}-{}-{}-{}-{}~-{}-{}-{}~-{}-{}-{}-{}-{}-{}-{}-{}~0:0000}\par
\texttt{F4.847~~03.013.03.-{}~~~~~~~~~-{}-{}4-{}-{}~-{}-{}-{}~-{}-{}-{}+-{}-{}-{}-{}~0:0000}\par
\texttt{F4.848~~13.013.03.-{}~~~~~~~~~-{}-{}-{}-{}-{}~-{}-{}-{}~-{}-{}-{}-{}-{}-{}-{}-{}~0:0000}\par
\texttt{F4.849~~23.013.03.-{}~~~~~~~~~-{}-{}-{}-{}-{}~-{}-{}-{}~-{}-{}-{}-{}-{}-{}-{}-{}~0:0000}\par
\texttt{F4.850~~013.23.03.-{}~~~~~~~~~-{}-{}-{}-{}-{}~-{}-{}-{}~-{}-{}-{}-{}-{}-{}-{}-{}~0:0000}\par
\texttt{F4.851~~123.23.03.-{}~~~~~~~~~-{}-{}-{}-{}-{}~-{}-{}-{}~-{}-{}-{}-{}-{}-{}-{}-{}~0:0000}\par
\texttt{F4.852~~03.023.03.-{}~~~~~~~~~-{}-{}4-{}-{}~-{}-{}-{}~-{}-{}-{}+-{}-{}-{}-{}~0:0000}\par
\texttt{F4.853~~23.023.03.-{}~~~~~~~~~-{}-{}-{}-{}-{}~-{}-{}-{}~-{}-{}-{}-{}-{}-{}-{}-{}~0:0000}\par
\texttt{F4.854~~03.123.03.-{}~~~~~~~~~-{}-{}-{}-{}-{}~-{}-{}-{}~-{}-{}-{}-{}-{}-{}-{}-{}~0:0000}\par
\texttt{F4.855~~03.13.013.-{}~~~~~~~~~-{}-{}4-{}-{}~-{}-{}-{}~-{}-{}-{}+-{}-{}-{}-{}~0:0000}\par
\texttt{F4.856~~013.13.23.-{}~~~~~~~~~-{}-{}4-{}-{}~-{}-{}-{}~-{}-{}-{}+-{}-{}-{}-{}~0:0000}\par
\texttt{F4.857~~0123.0.0.0~~~~~~~~~~-{}B-{}-{}D~-{}PC~+-{}-{}-{}++++~ok}\par
\texttt{F4.858~~0123.1.0.0~~~~~~~~~~-{}-{}-{}-{}D~-{}-{}-{}~-{}-{}-{}-{}-{}-{}-{}-{}~0:0100}\par
\texttt{F4.859~~012.01.0.0~~~~~~~~~~-{}-{}-{}-{}D~-{}P-{}~-{}-{}-{}-{}-{}+-{}-{}~3:1110}\par
\texttt{F4.860~~023.01.0.0~~~~~~~~~~-{}-{}-{}-{}D~-{}-{}-{}~-{}-{}-{}-{}-{}-{}-{}-{}~1:1011}\par
\texttt{F4.861~~123.01.0.0~~~~~~~~~~-{}B-{}-{}D~-{}PC~+-{}-{}-{}++++~ok}\par
\texttt{F4.862~~0123.2.0.0~~~~~~~~~~-{}-{}-{}-{}D~-{}PC~-{}-{}-{}-{}++-{}+~ok}\par
\texttt{F4.863~~012.02.0.0~~~~~~~~~~-{}-{}-{}-{}D~-{}P-{}~-{}-{}-{}-{}-{}+-{}-{}~3:1110}\par
\texttt{F4.864~~013.02.0.0~~~~~~~~~~-{}-{}-{}-{}D~-{}PC~-{}-{}-{}-{}++-{}+~ok}\par
\texttt{F4.865~~023.02.0.0~~~~~~~~~~-{}-{}-{}-{}D~-{}P-{}~-{}-{}-{}-{}-{}+-{}-{}~1:1011}\par
\texttt{F4.866~~123.02.0.0~~~~~~~~~~-{}-{}-{}-{}D~-{}PC~-{}-{}-{}-{}++++~ok}\par
\texttt{F4.867~~012.12.0.0~~~~~~~~~~-{}-{}-{}-{}D~-{}P-{}~-{}-{}-{}-{}-{}+-{}-{}~3:1110}\par
\texttt{F4.868~~013.12.0.0~~~~~~~~~~-{}-{}-{}-{}D~-{}PC~-{}-{}-{}-{}++-{}+~ok}\par
\texttt{F4.869~~023.12.0.0~~~~~~~~~~-{}-{}-{}-{}D~-{}-{}-{}~-{}-{}-{}-{}-{}-{}-{}-{}~1:1011}\par
\texttt{F4.870~~123.12.0.0~~~~~~~~~~-{}-{}-{}-{}D~-{}PC~-{}-{}-{}-{}++++~ok}\par
\texttt{F4.871~~01.012.0.0~~~~~~~~~~-{}-{}-{}-{}D~-{}P-{}~-{}-{}-{}-{}-{}+-{}-{}~3:1110}\par
\texttt{F4.872~~02.012.0.0~~~~~~~~~~-{}-{}-{}-{}D~-{}-{}-{}~-{}-{}-{}-{}-{}-{}-{}-{}~1:1010}\par
\texttt{F4.873~~12.012.0.0~~~~~~~~~~-{}-{}-{}-{}D~-{}P-{}~-{}-{}-{}-{}-{}+-{}-{}~3:1110}\par
\texttt{F4.874~~03.012.0.0~~~~~~~~~~-{}-{}-{}-{}D~-{}-{}-{}~-{}-{}-{}-{}-{}-{}-{}-{}~1:1001}\par
\texttt{F4.875~~13.012.0.0~~~~~~~~~~-{}-{}-{}-{}D~-{}PC~-{}-{}-{}-{}++-{}+~ok}\par
\texttt{F4.876~~23.012.0.0~~~~~~~~~~-{}-{}-{}-{}D~-{}-{}-{}~-{}-{}-{}-{}-{}-{}-{}-{}~1:1011}\par
\texttt{F4.877~~012.23.0.0~~~~~~~~~~-{}-{}-{}-{}D~-{}PC~-{}-{}-{}-{}++-{}+~ok}\par
\texttt{F4.878~~023.23.0.0~~~~~~~~~~-{}-{}-{}-{}D~-{}P-{}~-{}-{}-{}-{}-{}+-{}-{}~1:1011}\par
\texttt{F4.879~~123.23.0.0~~~~~~~~~~-{}-{}-{}-{}D~-{}PC~-{}-{}-{}-{}++++~ok}\par
\texttt{F4.880~~01.023.0.0~~~~~~~~~~-{}-{}-{}-{}D~-{}PC~-{}-{}-{}-{}++++~ok}\par
\texttt{F4.881~~02.023.0.0~~~~~~~~~~-{}-{}-{}-{}D~-{}P-{}~-{}-{}-{}-{}-{}+-{}-{}~1:1010}\par
\texttt{F4.882~~12.023.0.0~~~~~~~~~~-{}-{}-{}-{}D~-{}PC~-{}-{}-{}-{}++-{}+~ok}\par
\texttt{F4.883~~23.023.0.0~~~~~~~~~~-{}-{}-{}-{}D~-{}P-{}~-{}-{}-{}-{}-{}+-{}-{}~1:1011}\par
\texttt{F4.884~~01.123.0.0~~~~~~~~~~-{}-{}-{}-{}D~-{}PC~-{}-{}-{}-{}++-{}+~ok}\par
\texttt{F4.885~~02.123.0.0~~~~~~~~~~-{}-{}-{}-{}D~-{}-{}-{}~-{}-{}-{}-{}-{}-{}-{}-{}~1:1010}\par
\texttt{F4.886~~12.123.0.0~~~~~~~~~~-{}-{}-{}-{}D~-{}PC~-{}-{}-{}-{}++-{}+~ok}\par
\texttt{F4.887~~23.123.0.0~~~~~~~~~~-{}-{}-{}-{}D~-{}-{}-{}~-{}-{}-{}-{}-{}-{}-{}-{}~1:1011}\par
\texttt{F4.888~~0.0123.0.0~~~~~~~~~~-{}-{}4-{}D~-{}-{}-{}~-{}-{}-{}+-{}-{}-{}-{}~1:1000}\par
\texttt{F4.889~~1.0123.0.0~~~~~~~~~~-{}-{}-{}-{}D~-{}PC~-{}-{}-{}-{}++-{}+~ok}\par
\texttt{F4.890~~2.0123.0.0~~~~~~~~~~-{}-{}-{}-{}D~-{}-{}-{}~-{}-{}-{}-{}-{}-{}-{}-{}~1:1010}\par
\texttt{F4.891~~0123.1.1.0~~~~~~~~~~-{}-{}-{}-{}D~-{}-{}-{}~-{}-{}-{}-{}-{}-{}-{}-{}~0:0100}\par
\texttt{F4.892~~012.01.1.0~~~~~~~~~~-{}-{}-{}-{}D~-{}P-{}~-{}-{}-{}-{}-{}+-{}-{}~3:1110}\par
\texttt{F4.893~~013.01.1.0~~~~~~~~~~-{}-{}-{}-{}D~-{}P-{}~-{}-{}-{}-{}-{}+-{}-{}~2:1101}\par
\texttt{F4.894~~023.01.1.0~~~~~~~~~~-{}-{}-{}-{}D~-{}PC~-{}-{}-{}-{}++-{}+~ok}\par
\texttt{F4.895~~123.01.1.0~~~~~~~~~~-{}-{}-{}-{}D~-{}PC~-{}-{}-{}-{}++++~ok}\par
\texttt{F4.896~~0123.2.1.0~~~~~~~~~~-{}-{}-{}-{}D~-{}-{}-{}~-{}-{}-{}-{}-{}-{}-{}-{}~0:0110}\par
\texttt{F4.897~~012.02.1.0~~~~~~~~~~-{}-{}-{}-{}D~-{}P-{}~-{}-{}-{}-{}-{}+-{}-{}~3:1110}\par
\texttt{F4.898~~013.02.1.0~~~~~~~~~~-{}B-{}-{}D~-{}PC~+-{}-{}-{}++-{}+~ok}\par
\texttt{F4.899~~023.02.1.0~~~~~~~~~~-{}-{}-{}-{}D~-{}PC~-{}-{}-{}-{}++-{}+~ok}\par
\texttt{F4.900~~123.02.1.0~~~~~~~~~~-{}-{}-{}-{}D~-{}PC~-{}-{}-{}-{}++-{}+~ok}\par
\texttt{F4.901~~012.12.1.0~~~~~~~~~~-{}-{}-{}-{}D~-{}-{}-{}~-{}-{}-{}-{}-{}-{}-{}-{}~0:0110}\par
\texttt{F4.902~~013.12.1.0~~~~~~~~~~-{}-{}-{}-{}D~-{}-{}-{}~-{}-{}-{}-{}-{}-{}-{}-{}~0:0110}\par
\texttt{F4.903~~023.12.1.0~~~~~~~~~~-{}-{}-{}-{}D~-{}-{}-{}~-{}-{}-{}-{}-{}-{}-{}-{}~0:0110}\par
\texttt{F4.904~~123.12.1.0~~~~~~~~~~-{}-{}-{}-{}D~-{}-{}-{}~-{}-{}-{}-{}-{}-{}-{}-{}~0:0110}\par
\texttt{F4.905~~02.012.1.0~~~~~~~~~~-{}-{}-{}-{}D~-{}P-{}~-{}-{}-{}-{}-{}+-{}-{}~3:1110}\par
\texttt{F4.906~~12.012.1.0~~~~~~~~~~-{}-{}-{}-{}D~-{}P-{}~-{}-{}-{}-{}-{}+-{}-{}~3:1110}\par
\texttt{F4.907~~23.012.1.0~~~~~~~~~~-{}-{}-{}-{}D~-{}PC~-{}-{}-{}-{}++-{}+~ok}\par
\texttt{F4.908~~0123.3.1.0~~~~~~~~~~-{}-{}-{}-{}D~-{}PC~-{}-{}-{}-{}++-{}-{}~ok}\par
\texttt{F4.909~~012.03.1.0~~~~~~~~~~-{}-{}-{}-{}D~-{}PC~-{}-{}-{}-{}++-{}+~ok}\par
\texttt{F4.910~~023.03.1.0~~~~~~~~~~-{}-{}-{}-{}D~-{}PC~-{}-{}-{}-{}++-{}+~ok}\par
\texttt{F4.911~~123.03.1.0~~~~~~~~~~-{}-{}-{}-{}D~-{}PC~-{}-{}-{}-{}++++~ok}\par
\texttt{F4.912~~012.13.1.0~~~~~~~~~~-{}-{}-{}-{}D~-{}PC~-{}-{}-{}-{}++-{}-{}~ok}\par
\texttt{F4.913~~023.13.1.0~~~~~~~~~~-{}-{}-{}-{}D~-{}PC~-{}-{}-{}-{}++-{}-{}~ok}\par
\texttt{F4.914~~123.13.1.0~~~~~~~~~~-{}-{}-{}-{}D~-{}PC~-{}-{}-{}-{}++-{}+~ok}\par
\texttt{F4.915~~23.013.1.0~~~~~~~~~~-{}-{}-{}-{}D~-{}PC~-{}-{}-{}-{}++-{}+~ok}\par
\texttt{F4.916~~023.23.1.0~~~~~~~~~~-{}-{}-{}-{}D~-{}PC~-{}-{}-{}-{}++-{}+~ok}\par
\texttt{F4.917~~123.23.1.0~~~~~~~~~~-{}-{}-{}-{}D~-{}PC~-{}-{}-{}-{}++-{}+~ok}\par
\texttt{F4.918~~012.1.01.0~~~~~~~~~~-{}-{}-{}-{}D~-{}-{}-{}~-{}-{}-{}-{}-{}-{}-{}-{}~0:0100}\par
\texttt{F4.919~~013.1.01.0~~~~~~~~~~-{}-{}-{}-{}D~-{}-{}-{}~-{}-{}-{}-{}-{}-{}-{}-{}~0:0100}\par
\texttt{F4.920~~023.1.01.0~~~~~~~~~~-{}-{}-{}-{}D~-{}-{}-{}~-{}-{}-{}-{}-{}-{}-{}-{}~2:0100}\par
\texttt{F4.921~~123.1.01.0~~~~~~~~~~-{}-{}-{}-{}D~-{}-{}-{}~-{}-{}-{}-{}-{}-{}-{}-{}~0:0100}\par
\texttt{F4.922~~01.01.01.0~~~~~~~~~~-{}-{}-{}5D~UP-{}~-{}++-{}-{}+-{}-{}~2:1100}\par
\texttt{F4.923~~02.01.01.0~~~~~~~~~~-{}-{}-{}-{}D~-{}P-{}~-{}-{}-{}-{}-{}+-{}-{}~3:1110}\par
\texttt{F4.924~~12.01.01.0~~~~~~~~~~-{}-{}-{}-{}D~-{}P-{}~-{}-{}-{}-{}-{}+-{}-{}~3:1110}\par
\texttt{F4.925~~03.01.01.0~~~~~~~~~~-{}-{}-{}-{}D~-{}-{}-{}~-{}-{}-{}-{}-{}-{}-{}-{}~1:1001}\par
\texttt{F4.926~~13.01.01.0~~~~~~~~~~-{}-{}-{}-{}D~-{}P-{}~-{}-{}-{}-{}-{}+-{}-{}~2:1101}\par
\texttt{F4.927~~23.01.01.0~~~~~~~~~~-{}-{}-{}-{}D~-{}PC~-{}-{}-{}-{}++-{}+~ok}\par
\texttt{F4.928~~01.02.01.0~~~~~~~~~~-{}-{}-{}-{}D~-{}P-{}~-{}-{}-{}-{}-{}+-{}-{}~3:1110}\par
\texttt{F4.929~~12.02.01.0~~~~~~~~~~-{}-{}-{}-{}D~-{}P-{}~-{}-{}-{}-{}-{}+-{}-{}~3:1110}\par
\texttt{F4.930~~03.02.01.0~~~~~~~~~~-{}-{}-{}-{}D~-{}-{}-{}~-{}-{}-{}-{}-{}-{}-{}-{}~1:1001}\par
\texttt{F4.931~~13.02.01.0~~~~~~~~~~-{}-{}-{}-{}D~-{}PC~-{}-{}-{}-{}++-{}+~ok}\par
\texttt{F4.932~~01.12.01.0~~~~~~~~~~-{}-{}-{}-{}D~-{}P-{}~-{}-{}-{}-{}-{}+-{}-{}~3:1110}\par
\texttt{F4.933~~02.12.01.0~~~~~~~~~~-{}-{}-{}-{}D~-{}P-{}~-{}-{}-{}-{}-{}+-{}-{}~3:1110}\par
\texttt{F4.934~~12.12.01.0~~~~~~~~~~-{}-{}-{}-{}D~-{}P-{}~-{}-{}-{}-{}-{}+-{}-{}~3:1110}\par
\texttt{F4.935~~03.12.01.0~~~~~~~~~~-{}-{}-{}-{}D~-{}-{}-{}~-{}-{}-{}-{}-{}-{}-{}-{}~2:1001}\par
\texttt{F4.936~~13.12.01.0~~~~~~~~~~-{}-{}-{}-{}D~-{}PC~-{}-{}-{}-{}++-{}+~ok}\par
\texttt{F4.937~~23.12.01.0~~~~~~~~~~-{}B-{}-{}D~-{}PC~+-{}-{}-{}++-{}+~ok}\par
\texttt{F4.938~~0.012.01.0~~~~~~~~~~-{}-{}-{}-{}D~-{}-{}-{}~-{}-{}-{}-{}-{}-{}-{}-{}~1:1000}\par
\texttt{F4.939~~1.012.01.0~~~~~~~~~~-{}-{}-{}-{}D~-{}P-{}~-{}-{}-{}-{}-{}+-{}-{}~3:1110}\par
\texttt{F4.940~~2.012.01.0~~~~~~~~~~-{}-{}-{}-{}D~-{}P-{}~-{}-{}-{}-{}-{}+-{}-{}~3:1110}\par
\texttt{F4.941~~3.012.01.0~~~~~~~~~~-{}-{}-{}-{}D~-{}-{}-{}~-{}-{}-{}-{}-{}-{}-{}-{}~1:1001}\par
\texttt{F4.942~~012.3.01.0~~~~~~~~~~-{}-{}-{}-{}D~-{}PC~-{}-{}-{}-{}++-{}+~ok}\par
\texttt{F4.943~~013.3.01.0~~~~~~~~~~-{}-{}-{}-{}D~-{}P-{}~-{}-{}-{}-{}-{}+-{}-{}~2:1101}\par
\texttt{F4.944~~023.3.01.0~~~~~~~~~~-{}-{}-{}-{}D~-{}PC~-{}-{}-{}-{}++-{}+~ok}\par
\texttt{F4.945~~123.3.01.0~~~~~~~~~~-{}-{}-{}-{}D~-{}PC~-{}-{}-{}-{}++-{}+~ok}\par
\texttt{F4.946~~01.03.01.0~~~~~~~~~~-{}-{}-{}-{}D~-{}P-{}~-{}-{}-{}-{}-{}+-{}-{}~2:1101}\par
\texttt{F4.947~~02.03.01.0~~~~~~~~~~-{}-{}-{}-{}D~-{}PC~-{}-{}-{}-{}++-{}+~ok}\par
\texttt{F4.948~~12.03.01.0~~~~~~~~~~-{}-{}-{}-{}D~-{}PC~-{}-{}-{}-{}++++~ok}\par
\texttt{F4.949~~03.03.01.0~~~~~~~~~~-{}-{}-{}-{}D~-{}P-{}~-{}-{}-{}-{}-{}+-{}-{}~1:1001}\par
\texttt{F4.950~~13.03.01.0~~~~~~~~~~-{}-{}-{}-{}D~-{}P-{}~-{}-{}-{}-{}-{}+-{}-{}~2:1101}\par
\texttt{F4.951~~23.03.01.0~~~~~~~~~~-{}-{}-{}-{}D~-{}PC~-{}-{}-{}-{}++-{}+~ok}\par
\texttt{F4.952~~01.13.01.0~~~~~~~~~~-{}-{}-{}-{}D~-{}P-{}~-{}-{}-{}-{}-{}+-{}-{}~2:1101}\par
\texttt{F4.953~~02.13.01.0~~~~~~~~~~-{}-{}-{}-{}D~-{}PC~-{}-{}-{}-{}++-{}-{}~ok}\par
\texttt{F4.954~~12.13.01.0~~~~~~~~~~-{}-{}-{}-{}D~-{}PC~-{}-{}-{}-{}++-{}+~ok}\par
\texttt{F4.955~~03.13.01.0~~~~~~~~~~-{}-{}-{}-{}D~-{}-{}-{}~-{}-{}-{}-{}-{}-{}-{}-{}~1:1001}\par
\texttt{F4.956~~13.13.01.0~~~~~~~~~~-{}-{}-{}-{}D~-{}P-{}~-{}-{}-{}-{}-{}+-{}-{}~2:1101}\par
\texttt{F4.957~~23.13.01.0~~~~~~~~~~-{}-{}-{}-{}D~-{}PC~-{}-{}-{}-{}++-{}+~ok}\par
\texttt{F4.958~~0.013.01.0~~~~~~~~~~-{}-{}-{}-{}D~-{}-{}-{}~-{}-{}-{}-{}-{}-{}-{}-{}~1:1000}\par
\texttt{F4.959~~1.013.01.0~~~~~~~~~~-{}-{}-{}-{}D~-{}P-{}~-{}-{}-{}-{}-{}+-{}-{}~2:1101}\par
\texttt{F4.960~~2.013.01.0~~~~~~~~~~-{}-{}-{}-{}D~-{}PC~-{}-{}-{}-{}++-{}-{}~ok}\par
\texttt{F4.961~~3.013.01.0~~~~~~~~~~-{}-{}-{}-{}D~-{}-{}-{}~-{}-{}-{}-{}-{}-{}-{}-{}~1:1001}\par
\texttt{F4.962~~01.23.01.0~~~~~~~~~~-{}-{}-{}-{}D~-{}PC~-{}-{}-{}-{}++++~ok}\par
\texttt{F4.963~~02.23.01.0~~~~~~~~~~-{}-{}-{}-{}D~-{}PC~-{}-{}-{}-{}++-{}+~ok}\par
\texttt{F4.964~~12.23.01.0~~~~~~~~~~-{}-{}-{}-{}D~-{}PC~-{}-{}-{}-{}++++~ok}\par
\texttt{F4.965~~03.23.01.0~~~~~~~~~~-{}-{}-{}-{}D~-{}-{}-{}~-{}-{}-{}-{}-{}-{}-{}-{}~1:1001}\par
\texttt{F4.966~~13.23.01.0~~~~~~~~~~-{}-{}-{}-{}D~-{}PC~-{}-{}-{}-{}++-{}+~ok}\par
\texttt{F4.967~~23.23.01.0~~~~~~~~~~-{}-{}-{}-{}D~-{}PC~-{}-{}-{}-{}++-{}+~ok}\par
\texttt{F4.968~~0.023.01.0~~~~~~~~~~-{}-{}-{}-{}D~-{}-{}-{}~-{}-{}-{}-{}-{}-{}-{}-{}~1:1000}\par
\texttt{F4.969~~1.023.01.0~~~~~~~~~~-{}-{}-{}-{}D~-{}PC~-{}-{}-{}-{}++++~ok}\par
\texttt{F4.970~~2.023.01.0~~~~~~~~~~-{}-{}-{}-{}D~-{}PC~-{}-{}-{}-{}++-{}+~ok}\par
\texttt{F4.971~~3.023.01.0~~~~~~~~~~-{}-{}-{}-{}D~-{}-{}-{}~-{}-{}-{}-{}-{}-{}-{}-{}~1:1001}\par
\texttt{F4.972~~0.123.01.0~~~~~~~~~~-{}-{}-{}-{}D~-{}-{}-{}~-{}-{}-{}-{}-{}-{}-{}-{}~1:1000}\par
\texttt{F4.973~~1.123.01.0~~~~~~~~~~-{}-{}-{}-{}D~-{}PC~-{}-{}-{}-{}++-{}+~ok}\par
\texttt{F4.974~~2.123.01.0~~~~~~~~~~-{}-{}-{}-{}D~-{}PC~-{}-{}-{}-{}++-{}-{}~ok}\par
\texttt{F4.975~~3.123.01.0~~~~~~~~~~-{}-{}-{}-{}D~-{}-{}-{}~-{}-{}-{}-{}-{}-{}-{}-{}~1:1001}\par
\texttt{F4.976~~0123.1.2.0~~~~~~~~~~-{}-{}-{}-{}D~-{}-{}-{}~-{}-{}-{}-{}-{}-{}-{}-{}~0:0100}\par
\texttt{F4.977~~012.01.2.0~~~~~~~~~~-{}-{}-{}-{}D~-{}-{}-{}~-{}-{}-{}-{}-{}-{}-{}-{}~0:0010}\par
\texttt{F4.978~~013.01.2.0~~~~~~~~~~-{}B-{}-{}D~-{}PC~+-{}-{}-{}++++~ok}\par
\texttt{F4.979~~023.01.2.0~~~~~~~~~~-{}-{}-{}-{}D~-{}-{}-{}~-{}-{}-{}-{}-{}-{}-{}-{}~0:0010}\par
\texttt{F4.980~~123.01.2.0~~~~~~~~~~-{}-{}-{}-{}D~-{}-{}-{}~-{}-{}-{}-{}-{}-{}-{}-{}~0:0010}\par
\texttt{F4.981~~012.12.2.0~~~~~~~~~~-{}-{}-{}-{}D~-{}-{}-{}~-{}-{}-{}-{}-{}-{}-{}-{}~0:0010}\par
\texttt{F4.982~~013.12.2.0~~~~~~~~~~-{}-{}-{}-{}D~-{}-{}-{}~-{}-{}-{}-{}-{}-{}-{}-{}~1:0010}\par
\texttt{F4.983~~023.12.2.0~~~~~~~~~~-{}-{}-{}-{}D~-{}-{}-{}~-{}-{}-{}-{}-{}-{}-{}-{}~0:0010}\par
\texttt{F4.984~~123.12.2.0~~~~~~~~~~-{}-{}-{}-{}D~-{}-{}-{}~-{}-{}-{}-{}-{}-{}-{}-{}~0:0010}\par
\texttt{F4.985~~01.012.2.0~~~~~~~~~~-{}-{}-{}-{}D~-{}-{}-{}~-{}-{}-{}-{}-{}-{}-{}-{}~1:0010}\par
\texttt{F4.986~~02.012.2.0~~~~~~~~~~-{}-{}-{}-{}D~-{}-{}-{}~-{}-{}-{}-{}-{}-{}-{}-{}~0:0010}\par
\texttt{F4.987~~12.012.2.0~~~~~~~~~~-{}-{}-{}-{}D~-{}-{}-{}~-{}-{}-{}-{}-{}-{}-{}-{}~0:0010}\par
\texttt{F4.988~~03.012.2.0~~~~~~~~~~-{}-{}-{}-{}D~-{}-{}-{}~-{}-{}-{}-{}-{}-{}-{}-{}~1:0010}\par
\texttt{F4.989~~13.012.2.0~~~~~~~~~~-{}-{}-{}-{}D~-{}-{}-{}~-{}-{}-{}-{}-{}-{}-{}-{}~1:0010}\par
\texttt{F4.990~~23.012.2.0~~~~~~~~~~-{}-{}-{}-{}D~-{}-{}-{}~-{}-{}-{}-{}-{}-{}-{}-{}~0:0010}\par
\texttt{F4.991~~0123.3.2.0~~~~~~~~~~-{}-{}-{}-{}D~-{}-{}-{}~-{}-{}-{}-{}-{}-{}-{}-{}~0:0010}\par
\texttt{F4.992~~012.03.2.0~~~~~~~~~~-{}-{}-{}-{}D~-{}-{}-{}~-{}-{}-{}-{}-{}-{}-{}-{}~0:0010}\par
\texttt{F4.993~~013.03.2.0~~~~~~~~~~-{}-{}-{}-{}D~-{}PC~-{}-{}-{}-{}++++~ok}\par
\texttt{F4.994~~023.03.2.0~~~~~~~~~~-{}-{}-{}-{}D~-{}-{}-{}~-{}-{}-{}-{}-{}-{}-{}-{}~0:0010}\par
\texttt{F4.995~~123.03.2.0~~~~~~~~~~-{}-{}-{}-{}D~-{}-{}-{}~-{}-{}-{}-{}-{}-{}-{}-{}~0:0010}\par
\texttt{F4.996~~012.13.2.0~~~~~~~~~~-{}-{}-{}-{}D~-{}-{}-{}~-{}-{}-{}-{}-{}-{}-{}-{}~0:0010}\par
\texttt{F4.997~~013.13.2.0~~~~~~~~~~-{}-{}-{}-{}D~-{}PC~-{}-{}-{}-{}++++~ok}\par
\texttt{F4.998~~023.13.2.0~~~~~~~~~~-{}-{}-{}-{}D~-{}-{}-{}~-{}-{}-{}-{}-{}-{}-{}-{}~0:0010}\par
\texttt{F4.999~~123.13.2.0~~~~~~~~~~-{}-{}-{}-{}D~-{}-{}-{}~-{}-{}-{}-{}-{}-{}-{}-{}~0:0010}\par
\texttt{F4.1000~01.013.2.0~~~~~~~~~~-{}-{}-{}-{}D~-{}PC~-{}-{}-{}-{}++-{}+~ok}\par
\texttt{F4.1001~02.013.2.0~~~~~~~~~~-{}-{}-{}-{}D~-{}-{}-{}~-{}-{}-{}-{}-{}-{}-{}-{}~0:0010}\par
\texttt{F4.1002~12.013.2.0~~~~~~~~~~-{}-{}-{}-{}D~-{}-{}-{}~-{}-{}-{}-{}-{}-{}-{}-{}~0:0010}\par
\texttt{F4.1003~03.013.2.0~~~~~~~~~~-{}-{}-{}-{}D~-{}-{}-{}~-{}-{}-{}-{}-{}-{}-{}-{}~1:1001}\par
\texttt{F4.1004~13.013.2.0~~~~~~~~~~-{}-{}-{}-{}D~-{}PC~-{}-{}-{}-{}++++~ok}\par
\texttt{F4.1005~23.013.2.0~~~~~~~~~~-{}-{}-{}-{}D~-{}-{}-{}~-{}-{}-{}-{}-{}-{}-{}-{}~0:0010}\par
\texttt{F4.1006~012.23.2.0~~~~~~~~~~-{}-{}-{}-{}D~-{}-{}-{}~-{}-{}-{}-{}-{}-{}-{}-{}~0:0010}\par
\texttt{F4.1007~013.23.2.0~~~~~~~~~~-{}-{}-{}-{}D~-{}-{}-{}~-{}-{}-{}-{}-{}-{}-{}-{}~1:0010}\par
\texttt{F4.1008~023.23.2.0~~~~~~~~~~-{}-{}-{}-{}D~-{}-{}-{}~-{}-{}-{}-{}-{}-{}-{}-{}~0:0010}\par
\texttt{F4.1009~123.23.2.0~~~~~~~~~~-{}-{}-{}-{}D~-{}-{}-{}~-{}-{}-{}-{}-{}-{}-{}-{}~0:0010}\par
\texttt{F4.1010~01.023.2.0~~~~~~~~~~-{}-{}-{}-{}D~-{}-{}-{}~-{}-{}-{}-{}-{}-{}-{}-{}~1:0010}\par
\texttt{F4.1011~02.023.2.0~~~~~~~~~~-{}-{}-{}-{}D~-{}-{}-{}~-{}-{}-{}-{}-{}-{}-{}-{}~0:0010}\par
\texttt{F4.1012~12.023.2.0~~~~~~~~~~-{}-{}-{}-{}D~-{}-{}-{}~-{}-{}-{}-{}-{}-{}-{}-{}~0:0010}\par
\texttt{F4.1013~03.023.2.0~~~~~~~~~~-{}-{}-{}-{}D~-{}-{}-{}~-{}-{}-{}-{}-{}-{}-{}-{}~1:0010}\par
\texttt{F4.1014~13.023.2.0~~~~~~~~~~-{}-{}-{}-{}D~-{}-{}-{}~-{}-{}-{}-{}-{}-{}-{}-{}~1:0010}\par
\texttt{F4.1015~23.023.2.0~~~~~~~~~~-{}-{}-{}-{}D~-{}-{}-{}~-{}-{}-{}-{}-{}-{}-{}-{}~0:0010}\par
\texttt{F4.1016~01.123.2.0~~~~~~~~~~-{}-{}-{}-{}D~-{}-{}-{}~-{}-{}-{}-{}-{}-{}-{}-{}~1:0010}\par
\texttt{F4.1017~02.123.2.0~~~~~~~~~~-{}-{}-{}-{}D~-{}-{}-{}~-{}-{}-{}-{}-{}-{}-{}-{}~0:0010}\par
\texttt{F4.1018~12.123.2.0~~~~~~~~~~-{}-{}-{}-{}D~-{}-{}-{}~-{}-{}-{}-{}-{}-{}-{}-{}~0:0010}\par
\texttt{F4.1019~03.123.2.0~~~~~~~~~~-{}-{}-{}-{}D~-{}-{}-{}~-{}-{}-{}-{}-{}-{}-{}-{}~1:0010}\par
\texttt{F4.1020~13.123.2.0~~~~~~~~~~-{}-{}-{}-{}D~-{}-{}-{}~-{}-{}-{}-{}-{}-{}-{}-{}~1:0010}\par
\texttt{F4.1021~23.123.2.0~~~~~~~~~~-{}-{}-{}-{}D~-{}-{}-{}~-{}-{}-{}-{}-{}-{}-{}-{}~0:0010}\par
\texttt{F4.1022~0.0123.2.0~~~~~~~~~~-{}-{}4-{}D~-{}-{}-{}~-{}-{}-{}+-{}-{}-{}-{}~1:1000}\par
\texttt{F4.1023~2.0123.2.0~~~~~~~~~~-{}-{}-{}-{}D~-{}-{}-{}~-{}-{}-{}-{}-{}-{}-{}-{}~0:0010}\par
\texttt{F4.1024~3.0123.2.0~~~~~~~~~~-{}-{}-{}-{}D~-{}-{}-{}~-{}-{}-{}-{}-{}-{}-{}-{}~1:0010}\par
\texttt{F4.1025~01.01.02.0~~~~~~~~~~-{}-{}-{}-{}D~-{}-{}-{}~-{}-{}-{}-{}-{}-{}-{}-{}~2:1100}\par
\texttt{F4.1026~12.01.02.0~~~~~~~~~~-{}-{}-{}-{}D~-{}P-{}~-{}-{}-{}-{}-{}+-{}-{}~3:1110}\par
\texttt{F4.1027~03.01.02.0~~~~~~~~~~-{}-{}-{}-{}D~-{}-{}-{}~-{}-{}-{}-{}-{}-{}-{}-{}~1:1001}\par
\texttt{F4.1028~13.01.02.0~~~~~~~~~~-{}-{}-{}-{}D~-{}-{}-{}~-{}-{}-{}-{}-{}-{}-{}-{}~2:1101}\par
\texttt{F4.1029~01.12.02.0~~~~~~~~~~-{}-{}-{}-{}D~-{}P-{}~-{}-{}-{}-{}-{}+-{}-{}~3:1110}\par
\texttt{F4.1030~02.12.02.0~~~~~~~~~~-{}-{}-{}-{}D~-{}-{}-{}~-{}-{}-{}-{}-{}-{}-{}-{}~1:1010}\par
\texttt{F4.1031~12.12.02.0~~~~~~~~~~-{}-{}-{}-{}D~-{}P-{}~-{}-{}-{}-{}-{}+-{}-{}~3:1110}\par
\texttt{F4.1032~03.12.02.0~~~~~~~~~~-{}-{}-{}-{}D~-{}-{}-{}~-{}-{}-{}-{}-{}-{}-{}-{}~2:1001}\par
\texttt{F4.1033~13.12.02.0~~~~~~~~~~-{}-{}-{}-{}D~-{}PC~-{}-{}-{}-{}++-{}+~ok}\par
\texttt{F4.1034~23.12.02.0~~~~~~~~~~-{}-{}-{}-{}D~-{}-{}-{}~-{}-{}-{}-{}-{}-{}-{}-{}~1:1011}\par
\texttt{F4.1035~0.012.02.0~~~~~~~~~~-{}-{}4-{}D~-{}-{}-{}~-{}-{}-{}+-{}-{}-{}-{}~1:1000}\par
\texttt{F4.1036~1.012.02.0~~~~~~~~~~-{}-{}-{}-{}D~-{}P-{}~-{}-{}-{}-{}-{}+-{}-{}~3:1110}\par
\texttt{F4.1037~2.012.02.0~~~~~~~~~~-{}-{}-{}-{}D~-{}-{}-{}~-{}-{}-{}-{}-{}-{}-{}-{}~1:1010}\par
\texttt{F4.1038~3.012.02.0~~~~~~~~~~-{}-{}-{}-{}D~-{}-{}-{}~-{}-{}-{}-{}-{}-{}-{}-{}~1:1001}\par
\texttt{F4.1039~012.3.02.0~~~~~~~~~~-{}-{}-{}-{}D~-{}PC~-{}-{}-{}-{}++-{}+~ok}\par
\texttt{F4.1040~013.3.02.0~~~~~~~~~~-{}-{}-{}-{}D~-{}-{}-{}~-{}-{}-{}-{}-{}-{}-{}-{}~2:1101}\par
\texttt{F4.1041~023.3.02.0~~~~~~~~~~-{}-{}-{}-{}D~-{}P-{}~-{}-{}-{}-{}-{}+-{}-{}~1:1011}\par
\texttt{F4.1042~123.3.02.0~~~~~~~~~~-{}-{}-{}-{}D~-{}PC~-{}-{}-{}-{}++-{}+~ok}\par
\texttt{F4.1043~01.03.02.0~~~~~~~~~~-{}-{}-{}-{}D~-{}-{}-{}~-{}-{}-{}-{}-{}-{}-{}-{}~2:1101}\par
\texttt{F4.1044~02.03.02.0~~~~~~~~~~-{}-{}-{}-{}D~-{}P-{}~-{}-{}-{}-{}-{}+-{}-{}~1:1010}\par
\texttt{F4.1045~12.03.02.0~~~~~~~~~~-{}-{}-{}-{}D~-{}PC~-{}-{}-{}-{}++-{}+~ok}\par
\texttt{F4.1046~03.03.02.0~~~~~~~~~~-{}-{}-{}-{}D~-{}-{}-{}~-{}-{}-{}-{}-{}-{}-{}-{}~1:1001}\par
\texttt{F4.1047~13.03.02.0~~~~~~~~~~-{}-{}-{}-{}D~-{}-{}-{}~-{}-{}-{}-{}-{}-{}-{}-{}~2:1101}\par
\texttt{F4.1048~23.03.02.0~~~~~~~~~~-{}-{}-{}-{}D~-{}P-{}~-{}-{}-{}-{}-{}+-{}-{}~1:1011}\par
\texttt{F4.1049~01.13.02.0~~~~~~~~~~-{}-{}-{}-{}D~-{}-{}-{}~-{}-{}-{}-{}-{}-{}-{}-{}~2:1101}\par
\texttt{F4.1050~02.13.02.0~~~~~~~~~~-{}-{}-{}-{}D~-{}-{}-{}~-{}-{}-{}-{}-{}-{}-{}-{}~1:1010}\par
\texttt{F4.1051~12.13.02.0~~~~~~~~~~-{}-{}-{}-{}D~-{}PC~-{}-{}-{}-{}++-{}+~ok}\par
\texttt{F4.1052~03.13.02.0~~~~~~~~~~-{}-{}-{}-{}D~-{}-{}-{}~-{}-{}-{}-{}-{}-{}-{}-{}~1:1001}\par
\texttt{F4.1053~13.13.02.0~~~~~~~~~~-{}-{}-{}-{}D~-{}-{}-{}~-{}-{}-{}-{}-{}-{}-{}-{}~2:1101}\par
\texttt{F4.1054~23.13.02.0~~~~~~~~~~-{}-{}-{}-{}D~-{}-{}-{}~-{}-{}-{}-{}-{}-{}-{}-{}~1:1011}\par
\texttt{F4.1055~0.013.02.0~~~~~~~~~~-{}-{}4-{}D~-{}-{}-{}~-{}-{}-{}+-{}-{}-{}-{}~1:1000}\par
\texttt{F4.1056~1.013.02.0~~~~~~~~~~-{}-{}-{}-{}D~-{}-{}-{}~-{}-{}-{}-{}-{}-{}-{}-{}~2:1101}\par
\texttt{F4.1057~2.013.02.0~~~~~~~~~~-{}-{}-{}-{}D~-{}-{}-{}~-{}-{}-{}-{}-{}-{}-{}-{}~1:1010}\par
\texttt{F4.1058~3.013.02.0~~~~~~~~~~-{}-{}-{}-{}D~-{}-{}-{}~-{}-{}-{}-{}-{}-{}-{}-{}~1:1001}\par
\texttt{F4.1059~01.23.02.0~~~~~~~~~~-{}-{}-{}-{}D~-{}PC~-{}-{}-{}-{}++-{}+~ok}\par
\texttt{F4.1060~02.23.02.0~~~~~~~~~~-{}-{}-{}-{}D~-{}P-{}~-{}-{}-{}-{}-{}+-{}-{}~1:1010}\par
\texttt{F4.1061~12.23.02.0~~~~~~~~~~-{}-{}-{}-{}D~-{}PC~-{}-{}-{}-{}++-{}+~ok}\par
\texttt{F4.1062~03.23.02.0~~~~~~~~~~-{}-{}-{}-{}D~-{}-{}-{}~-{}-{}-{}-{}-{}-{}-{}-{}~1:1001}\par
\texttt{F4.1063~13.23.02.0~~~~~~~~~~-{}-{}-{}-{}D~-{}PC~-{}-{}-{}-{}++-{}+~ok}\par
\texttt{F4.1064~23.23.02.0~~~~~~~~~~-{}-{}-{}-{}D~-{}P-{}~-{}-{}-{}-{}-{}+-{}-{}~1:1011}\par
\texttt{F4.1065~0.023.02.0~~~~~~~~~~-{}-{}4-{}D~-{}-{}-{}~-{}-{}-{}+-{}-{}-{}-{}~1:1000}\par
\texttt{F4.1066~1.023.02.0~~~~~~~~~~-{}-{}-{}-{}D~-{}PC~-{}-{}-{}-{}++-{}+~ok}\par
\texttt{F4.1067~2.023.02.0~~~~~~~~~~-{}-{}-{}-{}D~-{}P-{}~-{}-{}-{}-{}-{}+-{}-{}~1:1010}\par
\texttt{F4.1068~3.023.02.0~~~~~~~~~~-{}-{}-{}-{}D~-{}-{}-{}~-{}-{}-{}-{}-{}-{}-{}-{}~1:1001}\par
\texttt{F4.1069~0.123.02.0~~~~~~~~~~-{}-{}-{}-{}D~-{}-{}-{}~-{}-{}-{}-{}-{}-{}-{}-{}~1:1000}\par
\texttt{F4.1070~1.123.02.0~~~~~~~~~~-{}-{}-{}-{}D~-{}PC~-{}-{}-{}-{}++-{}+~ok}\par
\texttt{F4.1071~2.123.02.0~~~~~~~~~~-{}-{}-{}-{}D~-{}-{}-{}~-{}-{}-{}-{}-{}-{}-{}-{}~1:1010}\par
\texttt{F4.1072~3.123.02.0~~~~~~~~~~-{}-{}-{}-{}D~-{}-{}-{}~-{}-{}-{}-{}-{}-{}-{}-{}~1:1001}\par
\texttt{F4.1073~01.12.12.0~~~~~~~~~~-{}-{}-{}-{}D~-{}-{}-{}~-{}-{}-{}-{}-{}-{}-{}-{}~0:0110}\par
\texttt{F4.1074~12.12.12.0~~~~~~~~~~-{}-{}-{}-{}D~UP-{}~-{}-{}+-{}-{}+-{}-{}~0:0110}\par
\texttt{F4.1075~03.12.12.0~~~~~~~~~~-{}B-{}-{}D~-{}PC~+-{}-{}-{}++++~ok}\par
\texttt{F4.1076~13.12.12.0~~~~~~~~~~-{}-{}-{}-{}D~-{}-{}-{}~-{}-{}-{}-{}-{}-{}-{}-{}~0:0110}\par
\texttt{F4.1077~0.012.12.0~~~~~~~~~~-{}-{}-{}-{}D~-{}-{}-{}~-{}-{}-{}-{}-{}-{}-{}-{}~1:1000}\par
\texttt{F4.1078~2.012.12.0~~~~~~~~~~-{}-{}-{}-{}D~-{}P-{}~-{}-{}-{}-{}-{}+-{}-{}~3:1110}\par
\texttt{F4.1079~3.012.12.0~~~~~~~~~~-{}-{}-{}-{}D~-{}-{}-{}~-{}-{}-{}-{}-{}-{}-{}-{}~1:1001}\par
\texttt{F4.1080~012.3.12.0~~~~~~~~~~-{}-{}-{}-{}D~-{}PC~-{}-{}-{}-{}++-{}+~ok}\par
\texttt{F4.1081~013.3.12.0~~~~~~~~~~-{}-{}-{}-{}D~-{}-{}-{}~-{}-{}-{}-{}-{}-{}-{}-{}~2:1101}\par
\texttt{F4.1082~023.3.12.0~~~~~~~~~~-{}-{}-{}-{}D~-{}PC~-{}-{}-{}-{}++-{}-{}~ok}\par
\texttt{F4.1083~123.3.12.0~~~~~~~~~~-{}-{}-{}-{}D~-{}PC~-{}-{}-{}-{}++-{}+~ok}\par
\texttt{F4.1084~01.03.12.0~~~~~~~~~~-{}-{}-{}-{}D~-{}-{}-{}~-{}-{}-{}-{}-{}-{}-{}-{}~2:1101}\par
\texttt{F4.1085~02.03.12.0~~~~~~~~~~-{}-{}-{}-{}D~-{}PC~-{}-{}-{}-{}++-{}+~ok}\par
\texttt{F4.1086~12.03.12.0~~~~~~~~~~-{}-{}-{}-{}D~-{}PC~-{}-{}-{}-{}++++~ok}\par
\texttt{F4.1087~03.03.12.0~~~~~~~~~~-{}-{}-{}-{}D~-{}-{}-{}~-{}-{}-{}-{}-{}-{}-{}-{}~1:1001}\par
\texttt{F4.1088~13.03.12.0~~~~~~~~~~-{}-{}-{}-{}D~-{}-{}-{}~-{}-{}-{}-{}-{}-{}-{}-{}~2:1101}\par
\texttt{F4.1089~23.03.12.0~~~~~~~~~~-{}-{}-{}-{}D~-{}PC~-{}-{}-{}-{}++-{}+~ok}\par
\texttt{F4.1090~01.13.12.0~~~~~~~~~~-{}-{}-{}-{}D~-{}-{}-{}~-{}-{}-{}-{}-{}-{}-{}-{}~2:1101}\par
\texttt{F4.1091~02.13.12.0~~~~~~~~~~-{}-{}-{}-{}D~-{}PC~-{}-{}-{}-{}++-{}-{}~ok}\par
\texttt{F4.1092~12.13.12.0~~~~~~~~~~-{}-{}-{}-{}D~-{}PC~-{}-{}-{}-{}++-{}+~ok}\par
\texttt{F4.1093~03.13.12.0~~~~~~~~~~-{}-{}-{}-{}D~-{}-{}-{}~-{}-{}-{}-{}-{}-{}-{}-{}~1:1001}\par
\texttt{F4.1094~13.13.12.0~~~~~~~~~~-{}-{}-{}-{}D~-{}-{}-{}~-{}-{}-{}-{}-{}-{}-{}-{}~2:1101}\par
\texttt{F4.1095~23.13.12.0~~~~~~~~~~-{}-{}-{}-{}D~-{}PC~-{}-{}-{}-{}++-{}+~ok}\par
\texttt{F4.1096~0.013.12.0~~~~~~~~~~-{}-{}-{}-{}D~-{}-{}-{}~-{}-{}-{}-{}-{}-{}-{}-{}~1:1000}\par
\texttt{F4.1097~2.013.12.0~~~~~~~~~~-{}-{}-{}-{}D~-{}PC~-{}-{}-{}-{}++-{}-{}~ok}\par
\texttt{F4.1098~3.013.12.0~~~~~~~~~~-{}-{}-{}-{}D~-{}-{}-{}~-{}-{}-{}-{}-{}-{}-{}-{}~1:1001}\par
\texttt{F4.1099~01.23.12.0~~~~~~~~~~-{}-{}-{}-{}D~-{}PC~-{}-{}-{}-{}++-{}+~ok}\par
\texttt{F4.1100~02.23.12.0~~~~~~~~~~-{}-{}-{}-{}D~-{}PC~-{}-{}-{}-{}++-{}+~ok}\par
\texttt{F4.1101~12.23.12.0~~~~~~~~~~-{}-{}-{}-{}D~-{}PC~-{}-{}-{}-{}++-{}+~ok}\par
\texttt{F4.1102~03.23.12.0~~~~~~~~~~-{}-{}-{}-{}D~-{}-{}-{}~-{}-{}-{}-{}-{}-{}-{}-{}~1:1001}\par
\texttt{F4.1103~13.23.12.0~~~~~~~~~~-{}-{}-{}-{}D~-{}PC~-{}-{}-{}-{}++-{}+~ok}\par
\texttt{F4.1104~23.23.12.0~~~~~~~~~~-{}-{}-{}-{}D~-{}PC~-{}-{}-{}-{}++-{}+~ok}\par
\texttt{F4.1105~0.023.12.0~~~~~~~~~~-{}-{}-{}-{}D~-{}-{}-{}~-{}-{}-{}-{}-{}-{}-{}-{}~1:1000}\par
\texttt{F4.1106~2.023.12.0~~~~~~~~~~-{}-{}-{}-{}D~-{}PC~-{}-{}-{}-{}++-{}+~ok}\par
\texttt{F4.1107~3.023.12.0~~~~~~~~~~-{}-{}-{}-{}D~-{}-{}-{}~-{}-{}-{}-{}-{}-{}-{}-{}~1:1001}\par
\texttt{F4.1108~0.123.12.0~~~~~~~~~~-{}-{}-{}-{}D~-{}-{}-{}~-{}-{}-{}-{}-{}-{}-{}-{}~1:1000}\par
\texttt{F4.1109~2.123.12.0~~~~~~~~~~-{}-{}-{}-{}D~-{}PC~-{}-{}-{}-{}++-{}-{}~ok}\par
\texttt{F4.1110~01.3.012.0~~~~~~~~~~-{}-{}-{}-{}D~-{}-{}-{}~-{}-{}-{}-{}-{}-{}-{}-{}~2:1101}\par
\texttt{F4.1111~12.3.012.0~~~~~~~~~~-{}-{}-{}-{}D~-{}PC~-{}-{}-{}-{}++-{}+~ok}\par
\texttt{F4.1112~13.3.012.0~~~~~~~~~~-{}-{}-{}-{}D~-{}-{}-{}~-{}-{}-{}-{}-{}-{}-{}-{}~2:1101}\par
\texttt{F4.1113~0.03.012.0~~~~~~~~~~-{}-{}-{}-{}D~-{}-{}-{}~-{}-{}-{}-{}-{}-{}-{}-{}~1:1000}\par
\texttt{F4.1114~1.03.012.0~~~~~~~~~~-{}-{}-{}-{}D~-{}-{}-{}~-{}-{}-{}-{}-{}-{}-{}-{}~2:1101}\par
\texttt{F4.1115~2.03.012.0~~~~~~~~~~-{}-{}-{}-{}D~-{}PC~-{}-{}-{}-{}++-{}+~ok}\par
\texttt{F4.1116~3.03.012.0~~~~~~~~~~-{}-{}-{}-{}D~-{}-{}-{}~-{}-{}-{}-{}-{}-{}-{}-{}~1:1001}\par
\texttt{F4.1117~0.13.012.0~~~~~~~~~~-{}-{}-{}-{}D~-{}-{}-{}~-{}-{}-{}-{}-{}-{}-{}-{}~1:1000}\par
\texttt{F4.1118~1.13.012.0~~~~~~~~~~-{}-{}-{}-{}D~-{}-{}-{}~-{}-{}-{}-{}-{}-{}-{}-{}~2:1101}\par
\texttt{F4.1119~0.23.012.0~~~~~~~~~~-{}-{}-{}-{}D~-{}-{}-{}~-{}-{}-{}-{}-{}-{}-{}-{}~1:1000}\par
\texttt{F4.1120~1.23.012.0~~~~~~~~~~-{}-{}-{}-{}D~-{}PC~-{}-{}-{}-{}++-{}+~ok}\par
\texttt{F4.1121~012.23.3.0~~~~~~~~~~-{}-{}-{}-{}D~-{}PC~-{}-{}-{}-{}++-{}+~ok}\par
\texttt{F4.1122~023.23.3.0~~~~~~~~~~-{}-{}-{}-{}D~-{}P-{}~-{}-{}-{}-{}-{}+-{}-{}~1:1011}\par
\texttt{F4.1123~123.23.3.0~~~~~~~~~~-{}-{}-{}-{}D~-{}PC~-{}-{}-{}-{}++-{}+~ok}\par
\texttt{F4.1124~02.023.3.0~~~~~~~~~~-{}-{}-{}-{}D~-{}P-{}~-{}-{}-{}-{}-{}+-{}-{}~1:1011}\par
\texttt{F4.1125~12.023.3.0~~~~~~~~~~-{}-{}-{}-{}D~-{}PC~-{}-{}-{}-{}++-{}+~ok}\par
\texttt{F4.1126~23.023.3.0~~~~~~~~~~-{}-{}-{}-{}D~-{}P-{}~-{}-{}-{}-{}-{}+-{}-{}~1:1011}\par
\texttt{F4.1127~02.123.3.0~~~~~~~~~~-{}-{}-{}-{}D~-{}-{}-{}~-{}-{}-{}-{}-{}-{}-{}-{}~1:1011}\par
\texttt{F4.1128~12.123.3.0~~~~~~~~~~-{}-{}-{}-{}D~-{}PC~-{}-{}-{}-{}++-{}+~ok}\par
\texttt{F4.1129~23.123.3.0~~~~~~~~~~-{}-{}-{}-{}D~-{}-{}-{}~-{}-{}-{}-{}-{}-{}-{}-{}~1:1011}\par
\texttt{F4.1130~2.0123.3.0~~~~~~~~~~-{}-{}-{}-{}D~-{}-{}-{}~-{}-{}-{}-{}-{}-{}-{}-{}~1:1011}\par
\texttt{F4.1131~01.03.03.0~~~~~~~~~~-{}-{}-{}-{}D~-{}P-{}~-{}-{}-{}-{}-{}+-{}-{}~2:1101}\par
\texttt{F4.1132~12.03.03.0~~~~~~~~~~-{}-{}-{}-{}D~-{}PC~-{}-{}-{}-{}++++~ok}\par
\texttt{F4.1133~03.03.03.0~~~~~~~~~~-{}-{}-{}-{}D~-{}P-{}~-{}-{}-{}-{}-{}+-{}-{}~1:1001}\par
\texttt{F4.1134~13.03.03.0~~~~~~~~~~-{}-{}-{}-{}D~-{}P-{}~-{}-{}-{}-{}-{}+-{}-{}~2:1101}\par
\texttt{F4.1135~01.13.03.0~~~~~~~~~~-{}-{}-{}-{}D~-{}P-{}~-{}-{}-{}-{}-{}+-{}-{}~2:1101}\par
\texttt{F4.1136~02.13.03.0~~~~~~~~~~-{}-{}-{}-{}D~-{}-{}-{}~-{}-{}-{}-{}-{}-{}-{}-{}~1:1011}\par
\texttt{F4.1137~12.13.03.0~~~~~~~~~~-{}-{}-{}-{}D~-{}PC~-{}-{}-{}-{}++-{}+~ok}\par
\texttt{F4.1138~03.13.03.0~~~~~~~~~~-{}-{}-{}-{}D~-{}-{}-{}~-{}-{}-{}-{}-{}-{}-{}-{}~1:1001}\par
\texttt{F4.1139~13.13.03.0~~~~~~~~~~-{}-{}-{}-{}D~-{}P-{}~-{}-{}-{}-{}-{}+-{}-{}~2:1101}\par
\texttt{F4.1140~23.13.03.0~~~~~~~~~~-{}-{}-{}-{}D~-{}-{}-{}~-{}-{}-{}-{}-{}-{}-{}-{}~1:1011}\par
\texttt{F4.1141~0.013.03.0~~~~~~~~~~-{}-{}4-{}D~-{}-{}-{}~-{}-{}-{}+-{}-{}-{}-{}~1:1000}\par
\texttt{F4.1142~3.013.03.0~~~~~~~~~~-{}-{}-{}-{}D~-{}-{}-{}~-{}-{}-{}-{}-{}-{}-{}-{}~1:1001}\par
\texttt{F4.1143~01.23.03.0~~~~~~~~~~-{}-{}-{}-{}D~-{}PC~-{}-{}-{}-{}++-{}+~ok}\par
\texttt{F4.1144~02.23.03.0~~~~~~~~~~-{}-{}-{}-{}D~-{}P-{}~-{}-{}-{}-{}-{}+-{}-{}~1:1011}\par
\texttt{F4.1145~12.23.03.0~~~~~~~~~~-{}-{}-{}-{}D~-{}PC~-{}-{}-{}-{}++-{}+~ok}\par
\texttt{F4.1146~03.23.03.0~~~~~~~~~~-{}-{}-{}-{}D~-{}P-{}~-{}-{}-{}-{}-{}+-{}-{}~1:1001}\par
\texttt{F4.1147~13.23.03.0~~~~~~~~~~-{}-{}-{}-{}D~-{}PC~-{}-{}-{}-{}++-{}+~ok}\par
\texttt{F4.1148~23.23.03.0~~~~~~~~~~-{}-{}-{}-{}D~-{}P-{}~-{}-{}-{}-{}-{}+-{}-{}~1:1011}\par
\texttt{F4.1149~0.023.03.0~~~~~~~~~~-{}-{}4-{}D~-{}-{}-{}~-{}-{}-{}+-{}-{}-{}-{}~1:1000}\par
\texttt{F4.1150~3.023.03.0~~~~~~~~~~-{}-{}-{}-{}D~-{}P-{}~-{}-{}-{}-{}-{}+-{}-{}~1:1001}\par
\texttt{F4.1151~0.123.03.0~~~~~~~~~~-{}-{}-{}-{}D~-{}-{}-{}~-{}-{}-{}-{}-{}-{}-{}-{}~1:1000}\par
\texttt{F4.1152~01.13.13.0~~~~~~~~~~-{}-{}-{}-{}D~-{}P-{}~-{}-{}-{}-{}-{}+-{}-{}~2:1101}\par
\texttt{F4.1153~02.13.13.0~~~~~~~~~~-{}-{}-{}-{}D~-{}PC~-{}-{}-{}-{}++-{}+~ok}\par
\texttt{F4.1154~12.13.13.0~~~~~~~~~~-{}-{}-{}-{}D~-{}PC~-{}-{}-{}-{}++-{}+~ok}\par
\texttt{F4.1155~03.13.13.0~~~~~~~~~~-{}-{}-{}-{}D~-{}-{}-{}~-{}-{}-{}-{}-{}-{}-{}-{}~1:1001}\par
\texttt{F4.1156~13.13.13.0~~~~~~~~~~-{}-{}-{}-{}D~-{}P-{}~-{}-{}-{}-{}-{}+-{}-{}~2:1101}\par
\texttt{F4.1157~23.13.13.0~~~~~~~~~~-{}-{}-{}-{}D~-{}PC~-{}-{}-{}-{}++-{}+~ok}\par
\texttt{F4.1158~0.013.13.0~~~~~~~~~~-{}-{}-{}-{}D~-{}-{}-{}~-{}-{}-{}-{}-{}-{}-{}-{}~1:1000}\par
\texttt{F4.1159~01.23.13.0~~~~~~~~~~-{}-{}-{}-{}D~-{}PC~-{}-{}-{}-{}++-{}+~ok}\par
\texttt{F4.1160~12.23.13.0~~~~~~~~~~-{}-{}-{}-{}D~-{}PC~-{}-{}-{}-{}++-{}+~ok}\par
\texttt{F4.1161~03.23.13.0~~~~~~~~~~-{}-{}-{}-{}D~-{}-{}-{}~-{}-{}-{}-{}-{}-{}-{}-{}~1:1001}\par
\texttt{F4.1162~13.23.13.0~~~~~~~~~~-{}-{}-{}-{}D~-{}PC~-{}-{}-{}-{}++-{}+~ok}\par
\texttt{F4.1163~0.023.13.0~~~~~~~~~~-{}-{}-{}-{}D~-{}-{}-{}~-{}-{}-{}-{}-{}-{}-{}-{}~1:1000}\par
\texttt{F4.1164~0.123.13.0~~~~~~~~~~-{}-{}-{}-{}D~-{}-{}-{}~-{}-{}-{}-{}-{}-{}-{}-{}~1:1000}\par
\texttt{F4.1165~0.13.013.0~~~~~~~~~~-{}-{}-{}-{}D~-{}-{}-{}~-{}-{}-{}-{}-{}-{}-{}-{}~1:1000}\par
\texttt{F4.1166~01.13.23.0~~~~~~~~~~-{}-{}-{}-{}D~-{}-{}-{}~-{}-{}-{}-{}-{}-{}-{}-{}~2:1101}\par
\texttt{F4.1167~12.13.23.0~~~~~~~~~~-{}-{}-{}-{}D~-{}PC~-{}-{}-{}-{}++-{}+~ok}\par
\texttt{F4.1168~03.13.23.0~~~~~~~~~~-{}-{}-{}-{}D~-{}-{}-{}~-{}-{}-{}-{}-{}-{}-{}-{}~1:1001}\par
\texttt{F4.1169~13.13.23.0~~~~~~~~~~-{}-{}-{}-{}D~-{}-{}-{}~-{}-{}-{}-{}-{}-{}-{}-{}~2:1101}\par
\texttt{F4.1170~0.013.23.0~~~~~~~~~~-{}-{}-{}-{}D~-{}-{}-{}~-{}-{}-{}-{}-{}-{}-{}-{}~1:1000}\par
\texttt{F4.1171~0.123.23.0~~~~~~~~~~-{}-{}-{}-{}D~-{}-{}-{}~-{}-{}-{}-{}-{}-{}-{}-{}~1:1000}\par
\texttt{F4.1172~01.1.01.01~~~~~~~~~~-{}-{}4-{}D~-{}-{}-{}~-{}-{}-{}+-{}-{}-{}-{}~0:0100}\par
\texttt{F4.1173~02.1.01.01~~~~~~~~~~-{}-{}-{}-{}D~-{}-{}-{}~-{}-{}-{}-{}-{}-{}-{}-{}~2:0100}\par
\texttt{F4.1174~12.1.01.01~~~~~~~~~~-{}-{}-{}-{}D~-{}-{}-{}~-{}-{}-{}-{}-{}-{}-{}-{}~0:0100}\par
\texttt{F4.1175~23.1.01.01~~~~~~~~~~-{}-{}-{}-{}D~-{}-{}-{}~-{}-{}-{}-{}-{}-{}-{}-{}~2:0100}\par
\texttt{F4.1176~012.1.2.01~~~~~~~~~~-{}-{}-{}-{}D~-{}-{}-{}~-{}-{}-{}-{}-{}-{}-{}-{}~0:0100}\par
\texttt{F4.1177~013.1.2.01~~~~~~~~~~-{}-{}-{}-{}D~-{}-{}-{}~-{}-{}-{}-{}-{}-{}-{}-{}~0:0100}\par
\texttt{F4.1178~023.1.2.01~~~~~~~~~~-{}-{}-{}-{}D~-{}-{}-{}~-{}-{}-{}-{}-{}-{}-{}-{}~3:0100}\par
\texttt{F4.1179~123.1.2.01~~~~~~~~~~-{}-{}-{}-{}D~-{}-{}-{}~-{}-{}-{}-{}-{}-{}-{}-{}~0:0100}\par
\texttt{F4.1180~01.01.2.01~~~~~~~~~~-{}-{}45D~UP-{}~-{}+++-{}+-{}-{}~3:1100}\par
\texttt{F4.1181~02.01.2.01~~~~~~~~~~-{}-{}-{}-{}D~-{}-{}-{}~-{}-{}-{}-{}-{}-{}-{}-{}~0:0010}\par
\texttt{F4.1182~12.01.2.01~~~~~~~~~~-{}-{}-{}-{}D~-{}-{}-{}~-{}-{}-{}-{}-{}-{}-{}-{}~0:0010}\par
\texttt{F4.1183~03.01.2.01~~~~~~~~~~-{}-{}-{}-{}D~-{}PC~-{}-{}-{}-{}++++~ok}\par
\texttt{F4.1184~13.01.2.01~~~~~~~~~~-{}-{}-{}-{}D~-{}PC~-{}-{}-{}-{}++++~ok}\par
\texttt{F4.1185~23.01.2.01~~~~~~~~~~-{}-{}-{}-{}D~-{}-{}-{}~-{}-{}-{}-{}-{}-{}-{}-{}~0:0010}\par
\texttt{F4.1186~02.02.2.01~~~~~~~~~~-{}-{}-{}-{}D~-{}-{}-{}~-{}-{}-{}-{}-{}-{}-{}-{}~0:0010}\par
\texttt{F4.1187~12.02.2.01~~~~~~~~~~-{}-{}-{}-{}D~-{}-{}-{}~-{}-{}-{}-{}-{}-{}-{}-{}~0:0010}\par
\texttt{F4.1188~03.02.2.01~~~~~~~~~~-{}-{}-{}-{}D~-{}-{}-{}~-{}-{}-{}-{}-{}-{}-{}-{}~1:0010}\par
\texttt{F4.1189~13.02.2.01~~~~~~~~~~-{}-{}-{}-{}D~-{}-{}-{}~-{}-{}-{}-{}-{}-{}-{}-{}~1:0010}\par
\texttt{F4.1190~23.02.2.01~~~~~~~~~~-{}-{}-{}-{}D~-{}-{}-{}~-{}-{}-{}-{}-{}-{}-{}-{}~0:0010}\par
\texttt{F4.1191~02.12.2.01~~~~~~~~~~-{}-{}-{}-{}D~-{}-{}-{}~-{}-{}-{}-{}-{}-{}-{}-{}~0:0010}\par
\texttt{F4.1192~03.12.2.01~~~~~~~~~~-{}-{}-{}-{}D~-{}-{}-{}~-{}-{}-{}-{}-{}-{}-{}-{}~1:0010}\par
\texttt{F4.1193~13.12.2.01~~~~~~~~~~-{}-{}-{}-{}D~-{}-{}-{}~-{}-{}-{}-{}-{}-{}-{}-{}~1:0010}\par
\texttt{F4.1194~23.12.2.01~~~~~~~~~~-{}-{}-{}-{}D~-{}-{}-{}~-{}-{}-{}-{}-{}-{}-{}-{}~0:0010}\par
\texttt{F4.1195~13.03.2.01~~~~~~~~~~-{}B-{}-{}D~-{}PC~+-{}-{}-{}++++~ok}\par
\texttt{F4.1196~03.13.2.01~~~~~~~~~~-{}B-{}-{}D~-{}PC~+-{}-{}-{}++++~ok}\par
\texttt{F4.1197~23.13.2.01~~~~~~~~~~-{}-{}-{}-{}D~-{}-{}-{}~-{}-{}-{}-{}-{}-{}-{}-{}~0:0010}\par
\texttt{F4.1198~01.1.02.01~~~~~~~~~~-{}-{}-{}-{}D~-{}-{}-{}~-{}-{}-{}-{}-{}-{}-{}-{}~0:0100}\par
\texttt{F4.1199~02.1.02.01~~~~~~~~~~-{}-{}-{}-{}D~-{}-{}-{}~-{}-{}-{}-{}-{}-{}-{}-{}~3:0100}\par
\texttt{F4.1200~12.1.02.01~~~~~~~~~~-{}-{}-{}-{}D~-{}-{}-{}~-{}-{}-{}-{}-{}-{}-{}-{}~0:0100}\par
\texttt{F4.1201~03.1.02.01~~~~~~~~~~-{}-{}-{}-{}D~-{}-{}-{}~-{}-{}-{}-{}-{}-{}-{}-{}~3:0100}\par
\texttt{F4.1202~13.1.02.01~~~~~~~~~~-{}-{}-{}-{}D~-{}-{}-{}~-{}-{}-{}-{}-{}-{}-{}-{}~0:0100}\par
\texttt{F4.1203~0.01.02.01~~~~~~~~~~-{}-{}4-{}D~-{}-{}-{}~-{}-{}-{}+-{}-{}-{}-{}~1:1000}\par
\texttt{F4.1204~0.02.02.01~~~~~~~~~~-{}-{}-{}-{}D~-{}-{}-{}~-{}-{}-{}-{}-{}-{}-{}-{}~1:1000}\par
\texttt{F4.1205~0.12.02.01~~~~~~~~~~-{}-{}-{}-{}D~-{}-{}-{}~-{}-{}-{}-{}-{}-{}-{}-{}~2:1000}\par
\texttt{F4.1206~0.03.02.01~~~~~~~~~~-{}-{}-{}-{}D~-{}-{}-{}~-{}-{}-{}-{}-{}-{}-{}-{}~1:1000}\par
\texttt{F4.1207~0.13.02.01~~~~~~~~~~-{}-{}-{}-{}D~-{}-{}-{}~-{}-{}-{}-{}-{}-{}-{}-{}~2:1000}\par
\texttt{F4.1208~0.1.012.01~~~~~~~~~~-{}-{}4-{}D~-{}-{}-{}~-{}-{}-{}+-{}-{}-{}-{}~2:1000}\par
\texttt{F4.1209~0.02.03.01~~~~~~~~~~-{}-{}-{}-{}D~-{}-{}-{}~-{}-{}-{}-{}-{}-{}-{}-{}~1:1000}\par
\texttt{F4.1210~0.12.03.01~~~~~~~~~~-{}-{}-{}-{}D~-{}-{}-{}~-{}-{}-{}-{}-{}-{}-{}-{}~2:1000}\par
\texttt{F4.1211~0.1.013.01~~~~~~~~~~-{}-{}4-{}D~-{}-{}-{}~-{}-{}-{}+-{}-{}-{}-{}~2:1000}\par
\texttt{F4.1212~01.1.23.01~~~~~~~~~~-{}-{}-{}-{}D~-{}-{}-{}~-{}-{}-{}-{}-{}-{}-{}-{}~0:0100}\par
\texttt{F4.1213~02.1.23.01~~~~~~~~~~-{}-{}-{}-{}D~-{}-{}-{}~-{}-{}-{}-{}-{}-{}-{}-{}~3:0100}\par
\texttt{F4.1214~03.1.23.01~~~~~~~~~~-{}-{}-{}-{}D~-{}-{}-{}~-{}-{}-{}-{}-{}-{}-{}-{}~3:0100}\par
\texttt{F4.1215~0.1.023.01~~~~~~~~~~-{}-{}-{}-{}D~-{}-{}-{}~-{}-{}-{}-{}-{}-{}-{}-{}~2:1000}\par
\texttt{F4.1216~01.1.2.012~~~~~~~~~~-{}-{}4-{}D~-{}-{}-{}~-{}-{}-{}+-{}-{}-{}-{}~0:0100}\par
\texttt{F4.1217~03.1.2.012~~~~~~~~~~-{}-{}-{}-{}D~-{}-{}-{}~-{}-{}-{}-{}-{}-{}-{}-{}~3:0100}\par
\texttt{F4.1218~0123.1.2.3~~~~~~~~~~T-{}4-{}D~-{}-{}-{}~-{}-{}-{}+-{}-{}-{}-{}~0:0100}\par
\texttt{F4.1219~012.01.2.3~~~~~~~~~~T-{}-{}-{}D~-{}-{}-{}~-{}-{}-{}-{}-{}-{}-{}-{}~0:0010}\par
\texttt{F4.1220~023.01.2.3~~~~~~~~~~T-{}-{}-{}D~-{}-{}-{}~-{}-{}-{}-{}-{}-{}-{}-{}~0:0010}\par
\texttt{F4.1221~012.12.2.3~~~~~~~~~~T-{}4-{}D~-{}-{}-{}~-{}-{}-{}+-{}-{}-{}-{}~0:0010}\par
\texttt{F4.1222~013.12.2.3~~~~~~~~~~T-{}-{}-{}D~-{}-{}-{}~-{}-{}-{}-{}-{}-{}-{}-{}~1:0010}\par
\texttt{F4.1223~023.12.2.3~~~~~~~~~~T-{}4-{}D~-{}-{}-{}~-{}-{}-{}+-{}-{}-{}-{}~0:0010}\par
\texttt{F4.1224~01.01.02.3~~~~~~~~~~T-{}-{}-{}D~-{}-{}-{}~-{}-{}-{}-{}-{}-{}-{}-{}~2:1100}\par
\texttt{F4.1225~03.01.02.3~~~~~~~~~~T-{}-{}-{}D~-{}-{}-{}~-{}-{}-{}-{}-{}-{}-{}-{}~0:0001}\par
\texttt{F4.1226~01.12.02.3~~~~~~~~~~T-{}-{}-{}D~-{}PC~-{}-{}-{}-{}++++~ok}\par
\texttt{F4.1227~03.12.02.3~~~~~~~~~~T-{}-{}-{}D~-{}-{}-{}~-{}-{}-{}-{}-{}-{}-{}-{}~0:0001}\par
\texttt{F4.1228~02.13.02.3~~~~~~~~~~T-{}4-{}D~-{}-{}-{}~-{}-{}-{}+-{}-{}-{}-{}~1:0001}\par
\texttt{F4.1229~03.13.02.3~~~~~~~~~~T-{}-{}-{}D~-{}-{}-{}~-{}-{}-{}-{}-{}-{}-{}-{}~0:0001}\par
\texttt{F4.1230~03.13.23.3~~~~~~~~~~T-{}4-{}D~-{}-{}-{}~-{}-{}-{}+-{}-{}-{}-{}~0:0001}\par
\texttt{F4.1231~0123.0123.-{}.-{}~~~~~~~-{}-{}4-{}-{}~-{}-{}-{}~-{}-{}-{}+-{}-{}-{}-{}~0:0000}\par
\texttt{F4.1232~0123.012.0.-{}~~~~~~~~-{}-{}-{}-{}-{}~-{}-{}-{}~-{}-{}-{}-{}-{}-{}-{}-{}~0:0000}\par
\texttt{F4.1233~0123.013.0.-{}~~~~~~~~-{}-{}-{}-{}-{}~-{}-{}-{}~-{}-{}-{}-{}-{}-{}-{}-{}~0:0000}\par
\texttt{F4.1234~0123.023.0.-{}~~~~~~~~-{}-{}-{}-{}-{}~-{}-{}-{}~-{}-{}-{}-{}-{}-{}-{}-{}~0:0000}\par
\texttt{F4.1235~0123.123.0.-{}~~~~~~~~-{}-{}-{}-{}-{}~-{}-{}-{}~-{}-{}-{}-{}-{}-{}-{}-{}~0:0000}\par
\texttt{F4.1236~012.0123.0.-{}~~~~~~~~-{}-{}-{}-{}-{}~-{}-{}-{}~-{}-{}-{}-{}-{}-{}-{}-{}~1:0000}\par
\texttt{F4.1237~013.0123.0.-{}~~~~~~~~-{}-{}-{}-{}-{}~-{}-{}-{}~-{}-{}-{}-{}-{}-{}-{}-{}~0:0000}\par
\texttt{F4.1238~023.0123.0.-{}~~~~~~~~-{}-{}-{}-{}-{}~-{}-{}-{}~-{}-{}-{}-{}-{}-{}-{}-{}~0:0000}\par
\texttt{F4.1239~123.0123.0.-{}~~~~~~~~-{}-{}-{}-{}-{}~-{}-{}-{}~-{}-{}-{}-{}-{}-{}-{}-{}~0:0000}\par
\texttt{F4.1240~0123.01.01.-{}~~~~~~~~-{}-{}-{}-{}-{}~-{}-{}-{}~-{}-{}-{}-{}-{}-{}-{}-{}~0:0000}\par
\texttt{F4.1241~0123.02.01.-{}~~~~~~~~-{}-{}-{}-{}-{}~-{}-{}-{}~-{}-{}-{}-{}-{}-{}-{}-{}~0:0000}\par
\texttt{F4.1242~0123.12.01.-{}~~~~~~~~-{}-{}-{}-{}-{}~-{}-{}-{}~-{}-{}-{}-{}-{}-{}-{}-{}~0:0000}\par
\texttt{F4.1243~012.012.01.-{}~~~~~~~~-{}B-{}-{}-{}~-{}PC~+-{}-{}-{}++++~ok}\par
\texttt{F4.1244~013.012.01.-{}~~~~~~~~-{}-{}-{}-{}-{}~-{}-{}-{}~-{}-{}-{}-{}-{}-{}-{}-{}~0:0000}\par
\texttt{F4.1245~023.012.01.-{}~~~~~~~~-{}-{}-{}-{}-{}~-{}-{}-{}~-{}-{}-{}-{}-{}-{}-{}-{}~0:0000}\par
\texttt{F4.1246~123.012.01.-{}~~~~~~~~-{}-{}-{}-{}-{}~-{}-{}-{}~-{}-{}-{}-{}-{}-{}-{}-{}~0:0000}\par
\texttt{F4.1247~0123.03.01.-{}~~~~~~~~-{}-{}-{}-{}-{}~-{}-{}-{}~-{}-{}-{}-{}-{}-{}-{}-{}~0:0000}\par
\texttt{F4.1248~0123.13.01.-{}~~~~~~~~-{}-{}-{}-{}-{}~-{}-{}-{}~-{}-{}-{}-{}-{}-{}-{}-{}~0:0000}\par
\texttt{F4.1249~013.013.01.-{}~~~~~~~~-{}-{}-{}-{}-{}~-{}-{}-{}~-{}-{}-{}-{}-{}-{}-{}-{}~0:0000}\par
\texttt{F4.1250~023.013.01.-{}~~~~~~~~-{}-{}-{}-{}-{}~-{}-{}-{}~-{}-{}-{}-{}-{}-{}-{}-{}~0:0000}\par
\texttt{F4.1251~123.013.01.-{}~~~~~~~~-{}-{}-{}-{}-{}~-{}-{}-{}~-{}-{}-{}-{}-{}-{}-{}-{}~0:0000}\par
\texttt{F4.1252~0123.23.01.-{}~~~~~~~~-{}-{}-{}-{}-{}~-{}-{}-{}~-{}-{}-{}-{}-{}-{}-{}-{}~0:0000}\par
\texttt{F4.1253~023.023.01.-{}~~~~~~~~-{}-{}-{}-{}-{}~-{}-{}-{}~-{}-{}-{}-{}-{}-{}-{}-{}~0:0000}\par
\texttt{F4.1254~123.023.01.-{}~~~~~~~~-{}-{}-{}-{}-{}~-{}-{}-{}~-{}-{}-{}-{}-{}-{}-{}-{}~0:0000}\par
\texttt{F4.1255~023.123.01.-{}~~~~~~~~-{}-{}-{}-{}-{}~-{}-{}-{}~-{}-{}-{}-{}-{}-{}-{}-{}~0:0000}\par
\texttt{F4.1256~0123.012.2.-{}~~~~~~~~-{}-{}-{}-{}-{}~-{}-{}-{}~-{}-{}-{}-{}-{}-{}-{}-{}~0:0000}\par
\texttt{F4.1257~0123.013.2.-{}~~~~~~~~-{}-{}-{}-{}-{}~-{}-{}-{}~-{}-{}-{}-{}-{}-{}-{}-{}~0:0000}\par
\texttt{F4.1258~0123.023.2.-{}~~~~~~~~-{}-{}-{}-{}-{}~-{}-{}-{}~-{}-{}-{}-{}-{}-{}-{}-{}~0:0000}\par
\texttt{F4.1259~0123.123.2.-{}~~~~~~~~-{}-{}4-{}-{}~-{}-{}-{}~-{}-{}-{}+-{}-{}-{}-{}~0:0000}\par
\texttt{F4.1260~0123.01.02.-{}~~~~~~~~-{}-{}-{}-{}-{}~-{}-{}-{}~-{}-{}-{}-{}-{}-{}-{}-{}~0:0000}\par
\texttt{F4.1261~0123.12.02.-{}~~~~~~~~-{}-{}-{}-{}-{}~-{}-{}-{}~-{}-{}-{}-{}-{}-{}-{}-{}~0:0000}\par
\texttt{F4.1262~012.012.02.-{}~~~~~~~~-{}-{}-{}-{}-{}~-{}PC~-{}-{}-{}-{}++++~ok}\par
\texttt{F4.1263~013.012.02.-{}~~~~~~~~-{}-{}-{}-{}-{}~-{}-{}-{}~-{}-{}-{}-{}-{}-{}-{}-{}~0:0000}\par
\texttt{F4.1264~023.012.02.-{}~~~~~~~~-{}-{}-{}-{}-{}~-{}-{}-{}~-{}-{}-{}-{}-{}-{}-{}-{}~0:0000}\par
\texttt{F4.1265~123.012.02.-{}~~~~~~~~-{}-{}-{}-{}-{}~-{}-{}-{}~-{}-{}-{}-{}-{}-{}-{}-{}~0:0000}\par
\texttt{F4.1266~0123.03.02.-{}~~~~~~~~-{}-{}-{}-{}-{}~-{}-{}-{}~-{}-{}-{}-{}-{}-{}-{}-{}~0:0000}\par
\texttt{F4.1267~0123.13.02.-{}~~~~~~~~-{}-{}-{}-{}-{}~-{}-{}-{}~-{}-{}-{}-{}-{}-{}-{}-{}~0:0000}\par
\texttt{F4.1268~012.013.02.-{}~~~~~~~~-{}-{}-{}-{}-{}~-{}-{}-{}~-{}-{}-{}-{}-{}-{}-{}-{}~1:0000}\par
\texttt{F4.1269~013.013.02.-{}~~~~~~~~-{}-{}-{}-{}-{}~-{}-{}-{}~-{}-{}-{}-{}-{}-{}-{}-{}~0:0000}\par
\texttt{F4.1270~023.013.02.-{}~~~~~~~~-{}-{}-{}-{}-{}~-{}-{}-{}~-{}-{}-{}-{}-{}-{}-{}-{}~0:0000}\par
\texttt{F4.1271~123.013.02.-{}~~~~~~~~-{}-{}-{}-{}-{}~-{}-{}-{}~-{}-{}-{}-{}-{}-{}-{}-{}~0:0000}\par
\texttt{F4.1272~0123.23.02.-{}~~~~~~~~-{}-{}-{}-{}-{}~-{}-{}-{}~-{}-{}-{}-{}-{}-{}-{}-{}~0:0000}\par
\texttt{F4.1273~012.023.02.-{}~~~~~~~~-{}-{}-{}-{}-{}~-{}-{}-{}~-{}-{}-{}-{}-{}-{}-{}-{}~1:0000}\par
\texttt{F4.1274~013.023.02.-{}~~~~~~~~-{}-{}-{}-{}-{}~-{}-{}-{}~-{}-{}-{}-{}-{}-{}-{}-{}~0:0000}\par
\texttt{F4.1275~023.023.02.-{}~~~~~~~~-{}-{}-{}-{}-{}~-{}-{}-{}~-{}-{}-{}-{}-{}-{}-{}-{}~0:0000}\par
\texttt{F4.1276~123.023.02.-{}~~~~~~~~-{}-{}-{}-{}-{}~-{}-{}-{}~-{}-{}-{}-{}-{}-{}-{}-{}~0:0000}\par
\texttt{F4.1277~012.123.02.-{}~~~~~~~~-{}-{}-{}-{}-{}~-{}-{}-{}~-{}-{}-{}-{}-{}-{}-{}-{}~1:0000}\par
\texttt{F4.1278~013.123.02.-{}~~~~~~~~-{}-{}-{}-{}-{}~-{}-{}-{}~-{}-{}-{}-{}-{}-{}-{}-{}~0:0000}\par
\texttt{F4.1279~023.123.02.-{}~~~~~~~~-{}-{}-{}-{}-{}~-{}-{}-{}~-{}-{}-{}-{}-{}-{}-{}-{}~0:0000}\par
\texttt{F4.1280~123.123.02.-{}~~~~~~~~-{}-{}-{}-{}-{}~-{}-{}-{}~-{}-{}-{}-{}-{}-{}-{}-{}~0:0000}\par
\texttt{F4.1281~02.0123.02.-{}~~~~~~~~-{}-{}4-{}-{}~-{}-{}-{}~-{}-{}-{}+-{}-{}-{}-{}~1:0000}\par
\texttt{F4.1282~03.0123.02.-{}~~~~~~~~-{}-{}-{}-{}-{}~-{}-{}-{}~-{}-{}-{}-{}-{}-{}-{}-{}~0:0000}\par
\texttt{F4.1283~13.0123.02.-{}~~~~~~~~-{}-{}-{}-{}-{}~-{}-{}-{}~-{}-{}-{}-{}-{}-{}-{}-{}~0:0000}\par
\texttt{F4.1284~23.0123.02.-{}~~~~~~~~-{}-{}-{}-{}-{}~-{}-{}-{}~-{}-{}-{}-{}-{}-{}-{}-{}~0:0000}\par
\texttt{F4.1285~03.012.012.-{}~~~~~~~~-{}-{}-{}-{}-{}~-{}-{}-{}~-{}-{}-{}-{}-{}-{}-{}-{}~0:0000}\par
\texttt{F4.1286~13.012.012.-{}~~~~~~~~-{}-{}-{}-{}-{}~-{}-{}-{}~-{}-{}-{}-{}-{}-{}-{}-{}~0:0000}\par
\texttt{F4.1287~0123.3.012.-{}~~~~~~~~-{}-{}-{}-{}-{}~-{}-{}-{}~-{}-{}-{}-{}-{}-{}-{}-{}~0:0000}\par
\texttt{F4.1288~013.03.012.-{}~~~~~~~~-{}-{}-{}-{}-{}~-{}-{}-{}~-{}-{}-{}-{}-{}-{}-{}-{}~0:0000}\par
\texttt{F4.1289~023.03.012.-{}~~~~~~~~-{}-{}-{}-{}-{}~-{}-{}-{}~-{}-{}-{}-{}-{}-{}-{}-{}~0:0000}\par
\texttt{F4.1290~123.03.012.-{}~~~~~~~~-{}-{}-{}-{}-{}~-{}-{}-{}~-{}-{}-{}-{}-{}-{}-{}-{}~0:0000}\par
\texttt{F4.1291~013.13.012.-{}~~~~~~~~-{}-{}-{}-{}-{}~-{}-{}-{}~-{}-{}-{}-{}-{}-{}-{}-{}~0:0000}\par
\texttt{F4.1292~023.13.012.-{}~~~~~~~~-{}-{}-{}-{}-{}~-{}-{}-{}~-{}-{}-{}-{}-{}-{}-{}-{}~0:0000}\par
\texttt{F4.1293~123.13.012.-{}~~~~~~~~-{}-{}-{}-{}-{}~-{}-{}-{}~-{}-{}-{}-{}-{}-{}-{}-{}~0:0000}\par
\texttt{F4.1294~23.013.012.-{}~~~~~~~~-{}-{}-{}-{}-{}~-{}-{}-{}~-{}-{}-{}-{}-{}-{}-{}-{}~0:0000}\par
\texttt{F4.1295~023.23.012.-{}~~~~~~~~-{}-{}-{}-{}-{}~-{}-{}-{}~-{}-{}-{}-{}-{}-{}-{}-{}~0:0000}\par
\texttt{F4.1296~123.23.012.-{}~~~~~~~~-{}-{}-{}-{}-{}~-{}-{}-{}~-{}-{}-{}-{}-{}-{}-{}-{}~0:0000}\par
\texttt{F4.1297~0123.013.3.-{}~~~~~~~~-{}-{}-{}-{}-{}~-{}-{}-{}~-{}-{}-{}-{}-{}-{}-{}-{}~0:0000}\par
\texttt{F4.1298~0123.023.3.-{}~~~~~~~~-{}-{}-{}-{}-{}~-{}-{}-{}~-{}-{}-{}-{}-{}-{}-{}-{}~0:0000}\par
\texttt{F4.1299~0123.123.3.-{}~~~~~~~~-{}-{}4-{}-{}~-{}-{}-{}~-{}-{}-{}+-{}-{}-{}-{}~0:0000}\par
\texttt{F4.1300~0123.03.03.-{}~~~~~~~~-{}-{}-{}-{}-{}~-{}-{}-{}~-{}-{}-{}-{}-{}-{}-{}-{}~0:0000}\par
\texttt{F4.1301~0123.13.03.-{}~~~~~~~~-{}-{}-{}-{}-{}~-{}-{}-{}~-{}-{}-{}-{}-{}-{}-{}-{}~0:0000}\par
\texttt{F4.1302~013.013.03.-{}~~~~~~~~-{}-{}-{}-{}-{}~-{}-{}-{}~-{}-{}-{}-{}-{}-{}-{}-{}~0:0000}\par
\texttt{F4.1303~023.013.03.-{}~~~~~~~~-{}-{}-{}-{}-{}~-{}-{}-{}~-{}-{}-{}-{}-{}-{}-{}-{}~0:0000}\par
\texttt{F4.1304~123.013.03.-{}~~~~~~~~-{}-{}-{}-{}-{}~-{}-{}-{}~-{}-{}-{}-{}-{}-{}-{}-{}~0:0000}\par
\texttt{F4.1305~0123.23.03.-{}~~~~~~~~-{}-{}-{}-{}-{}~-{}-{}-{}~-{}-{}-{}-{}-{}-{}-{}-{}~0:0000}\par
\texttt{F4.1306~013.023.03.-{}~~~~~~~~-{}-{}-{}-{}-{}~-{}-{}-{}~-{}-{}-{}-{}-{}-{}-{}-{}~0:0000}\par
\texttt{F4.1307~023.023.03.-{}~~~~~~~~-{}-{}-{}-{}-{}~-{}-{}-{}~-{}-{}-{}-{}-{}-{}-{}-{}~0:0000}\par
\texttt{F4.1308~123.023.03.-{}~~~~~~~~-{}-{}-{}-{}-{}~-{}-{}-{}~-{}-{}-{}-{}-{}-{}-{}-{}~0:0000}\par
\texttt{F4.1309~013.123.03.-{}~~~~~~~~-{}-{}-{}-{}-{}~-{}-{}-{}~-{}-{}-{}-{}-{}-{}-{}-{}~0:0000}\par
\texttt{F4.1310~023.123.03.-{}~~~~~~~~-{}-{}-{}-{}-{}~-{}-{}-{}~-{}-{}-{}-{}-{}-{}-{}-{}~0:0000}\par
\texttt{F4.1311~123.123.03.-{}~~~~~~~~-{}-{}-{}-{}-{}~-{}-{}-{}~-{}-{}-{}-{}-{}-{}-{}-{}~0:0000}\par
\texttt{F4.1312~03.0123.03.-{}~~~~~~~~-{}-{}4-{}-{}~-{}-{}-{}~-{}-{}-{}+-{}-{}-{}-{}~0:0000}\par
\texttt{F4.1313~23.0123.03.-{}~~~~~~~~-{}-{}-{}-{}-{}~-{}-{}-{}~-{}-{}-{}-{}-{}-{}-{}-{}~0:0000}\par
\texttt{F4.1314~013.13.013.-{}~~~~~~~~-{}-{}4-{}-{}~-{}-{}-{}~-{}-{}-{}+-{}-{}-{}-{}~0:0000}\par
\texttt{F4.1315~023.13.013.-{}~~~~~~~~-{}-{}-{}-{}-{}~-{}-{}-{}~-{}-{}-{}-{}-{}-{}-{}-{}~0:0000}\par
\texttt{F4.1316~123.13.013.-{}~~~~~~~~-{}-{}-{}-{}-{}~-{}-{}-{}~-{}-{}-{}-{}-{}-{}-{}-{}~0:0000}\par
\texttt{F4.1317~0123.13.23.-{}~~~~~~~~-{}-{}4-{}-{}~-{}-{}-{}~-{}-{}-{}+-{}-{}-{}-{}~0:0000}\par
\texttt{F4.1318~013.013.23.-{}~~~~~~~~-{}-{}4-{}-{}~-{}-{}-{}~-{}-{}-{}+-{}-{}-{}-{}~0:0000}\par
\texttt{F4.1319~023.013.23.-{}~~~~~~~~-{}-{}-{}-{}-{}~-{}-{}-{}~-{}-{}-{}-{}-{}-{}-{}-{}~0:0000}\par
\texttt{F4.1320~023.123.23.-{}~~~~~~~~-{}-{}4-{}-{}~-{}-{}-{}~-{}-{}-{}+-{}-{}-{}-{}~0:0000}\par
\texttt{F4.1321~0123.01.0.0~~~~~~~~~-{}B-{}-{}D~-{}PC~+-{}-{}-{}++++~ok}\par
\texttt{F4.1322~0123.02.0.0~~~~~~~~~-{}-{}-{}-{}D~-{}PC~-{}-{}-{}-{}++++~ok}\par
\texttt{F4.1323~0123.12.0.0~~~~~~~~~-{}-{}-{}-{}D~-{}PC~-{}-{}-{}-{}++-{}+~ok}\par
\texttt{F4.1324~012.012.0.0~~~~~~~~~-{}-{}-{}-{}D~-{}P-{}~-{}-{}-{}-{}-{}+-{}-{}~3:1110}\par
\texttt{F4.1325~013.012.0.0~~~~~~~~~-{}-{}-{}-{}D~-{}PC~-{}-{}-{}-{}++-{}+~ok}\par
\texttt{F4.1326~023.012.0.0~~~~~~~~~-{}-{}-{}-{}D~-{}-{}-{}~-{}-{}-{}-{}-{}-{}-{}-{}~1:1011}\par
\texttt{F4.1327~123.012.0.0~~~~~~~~~-{}-{}-{}-{}D~-{}PC~-{}-{}-{}-{}++++~ok}\par
\texttt{F4.1328~0123.23.0.0~~~~~~~~~-{}-{}-{}-{}D~-{}PC~-{}-{}-{}-{}++-{}+~ok}\par
\texttt{F4.1329~012.023.0.0~~~~~~~~~-{}-{}-{}-{}D~-{}PC~-{}-{}-{}-{}++-{}+~ok}\par
\texttt{F4.1330~023.023.0.0~~~~~~~~~-{}-{}-{}-{}D~-{}P-{}~-{}-{}-{}-{}-{}+-{}-{}~1:1011}\par
\texttt{F4.1331~123.023.0.0~~~~~~~~~-{}-{}-{}-{}D~-{}PC~-{}-{}-{}-{}++++~ok}\par
\texttt{F4.1332~012.123.0.0~~~~~~~~~-{}-{}-{}-{}D~-{}PC~-{}-{}-{}-{}++-{}+~ok}\par
\texttt{F4.1333~023.123.0.0~~~~~~~~~-{}-{}-{}-{}D~-{}-{}-{}~-{}-{}-{}-{}-{}-{}-{}-{}~1:1011}\par
\texttt{F4.1334~123.123.0.0~~~~~~~~~-{}-{}-{}-{}D~-{}PC~-{}-{}-{}-{}++++~ok}\par
\texttt{F4.1335~01.0123.0.0~~~~~~~~~-{}-{}-{}-{}D~-{}PC~-{}-{}-{}-{}++-{}+~ok}\par
\texttt{F4.1336~02.0123.0.0~~~~~~~~~-{}-{}-{}-{}D~-{}-{}-{}~-{}-{}-{}-{}-{}-{}-{}-{}~1:1010}\par
\texttt{F4.1337~12.0123.0.0~~~~~~~~~-{}-{}-{}-{}D~-{}PC~-{}-{}-{}-{}++-{}+~ok}\par
\texttt{F4.1338~23.0123.0.0~~~~~~~~~-{}-{}-{}-{}D~-{}-{}-{}~-{}-{}-{}-{}-{}-{}-{}-{}~1:1011}\par
\texttt{F4.1339~0123.01.1.0~~~~~~~~~-{}-{}-{}-{}D~-{}PC~-{}-{}-{}-{}++-{}+~ok}\par
\texttt{F4.1340~0123.02.1.0~~~~~~~~~-{}-{}-{}-{}D~-{}PC~-{}-{}-{}-{}++-{}+~ok}\par
\texttt{F4.1341~0123.12.1.0~~~~~~~~~-{}-{}-{}-{}D~-{}-{}-{}~-{}-{}-{}-{}-{}-{}-{}-{}~0:0110}\par
\texttt{F4.1342~012.012.1.0~~~~~~~~~-{}-{}-{}-{}D~-{}P-{}~-{}-{}-{}-{}-{}+-{}-{}~3:1110}\par
\texttt{F4.1343~013.012.1.0~~~~~~~~~-{}B-{}-{}D~-{}PC~+-{}-{}-{}++-{}+~ok}\par
\texttt{F4.1344~023.012.1.0~~~~~~~~~-{}-{}-{}-{}D~-{}PC~-{}-{}-{}-{}++-{}+~ok}\par
\texttt{F4.1345~123.012.1.0~~~~~~~~~-{}-{}-{}-{}D~-{}PC~-{}-{}-{}-{}++-{}+~ok}\par
\texttt{F4.1346~0123.03.1.0~~~~~~~~~-{}-{}-{}-{}D~-{}PC~-{}-{}-{}-{}++-{}+~ok}\par
\texttt{F4.1347~0123.13.1.0~~~~~~~~~-{}-{}-{}-{}D~-{}PC~-{}-{}-{}-{}++-{}-{}~ok}\par
\texttt{F4.1348~012.013.1.0~~~~~~~~~-{}-{}-{}-{}D~-{}PC~-{}-{}-{}-{}++-{}+~ok}\par
\texttt{F4.1349~023.013.1.0~~~~~~~~~-{}-{}-{}-{}D~-{}PC~-{}-{}-{}-{}++-{}+~ok}\par
\texttt{F4.1350~123.013.1.0~~~~~~~~~-{}-{}-{}-{}D~-{}PC~-{}-{}-{}-{}++++~ok}\par
\texttt{F4.1351~0123.23.1.0~~~~~~~~~-{}-{}-{}-{}D~-{}PC~-{}-{}-{}-{}++-{}+~ok}\par
\texttt{F4.1352~023.023.1.0~~~~~~~~~-{}-{}-{}-{}D~-{}PC~-{}-{}-{}-{}++-{}+~ok}\par
\texttt{F4.1353~123.023.1.0~~~~~~~~~-{}-{}-{}-{}D~-{}PC~-{}-{}-{}-{}++++~ok}\par
\texttt{F4.1354~023.123.1.0~~~~~~~~~-{}-{}-{}-{}D~-{}PC~-{}-{}-{}-{}++-{}+~ok}\par
\texttt{F4.1355~0123.1.01.0~~~~~~~~~-{}-{}-{}-{}D~-{}-{}-{}~-{}-{}-{}-{}-{}-{}-{}-{}~0:0100}\par
\texttt{F4.1356~012.01.01.0~~~~~~~~~-{}-{}-{}-{}D~-{}P-{}~-{}-{}-{}-{}-{}+-{}-{}~3:1110}\par
\texttt{F4.1357~013.01.01.0~~~~~~~~~-{}-{}-{}-{}D~-{}P-{}~-{}-{}-{}-{}-{}+-{}-{}~2:1101}\par
\texttt{F4.1358~023.01.01.0~~~~~~~~~-{}-{}-{}-{}D~-{}PC~-{}-{}-{}-{}++-{}+~ok}\par
\texttt{F4.1359~123.01.01.0~~~~~~~~~-{}-{}-{}-{}D~-{}PC~-{}-{}-{}-{}++++~ok}\par
\texttt{F4.1360~012.02.01.0~~~~~~~~~-{}-{}-{}-{}D~-{}P-{}~-{}-{}-{}-{}-{}+-{}-{}~3:1110}\par
\texttt{F4.1361~013.02.01.0~~~~~~~~~-{}-{}-{}-{}D~-{}PC~-{}-{}-{}-{}++-{}+~ok}\par
\texttt{F4.1362~123.02.01.0~~~~~~~~~-{}B-{}-{}D~-{}PC~+-{}-{}-{}++++~ok}\par
\texttt{F4.1363~012.12.01.0~~~~~~~~~-{}-{}-{}-{}D~-{}P-{}~-{}-{}-{}-{}-{}+-{}-{}~3:1110}\par
\texttt{F4.1364~013.12.01.0~~~~~~~~~-{}-{}-{}-{}D~-{}PC~-{}-{}-{}-{}++-{}+~ok}\par
\texttt{F4.1365~023.12.01.0~~~~~~~~~-{}B-{}-{}D~-{}PC~+-{}-{}-{}++-{}+~ok}\par
\texttt{F4.1366~123.12.01.0~~~~~~~~~-{}-{}-{}-{}D~-{}PC~-{}-{}-{}-{}++-{}+~ok}\par
\texttt{F4.1367~01.012.01.0~~~~~~~~~-{}-{}-{}-{}D~-{}P-{}~-{}-{}-{}-{}-{}+-{}-{}~3:1110}\par
\texttt{F4.1368~02.012.01.0~~~~~~~~~-{}-{}-{}-{}D~-{}P-{}~-{}-{}-{}-{}-{}+-{}-{}~3:1110}\par
\texttt{F4.1369~12.012.01.0~~~~~~~~~-{}-{}-{}-{}D~-{}P-{}~-{}-{}-{}-{}-{}+-{}-{}~3:1110}\par
\texttt{F4.1370~03.012.01.0~~~~~~~~~-{}-{}-{}-{}D~-{}-{}-{}~-{}-{}-{}-{}-{}-{}-{}-{}~1:1001}\par
\texttt{F4.1371~13.012.01.0~~~~~~~~~-{}-{}-{}-{}D~-{}PC~-{}-{}-{}-{}++-{}+~ok}\par
\texttt{F4.1372~23.012.01.0~~~~~~~~~-{}-{}-{}-{}D~-{}PC~-{}-{}-{}-{}++-{}+~ok}\par
\texttt{F4.1373~0123.3.01.0~~~~~~~~~-{}-{}-{}-{}D~-{}PC~-{}-{}-{}-{}++-{}+~ok}\par
\texttt{F4.1374~012.03.01.0~~~~~~~~~-{}-{}-{}-{}D~-{}PC~-{}-{}-{}-{}++++~ok}\par
\texttt{F4.1375~013.03.01.0~~~~~~~~~-{}-{}-{}-{}D~-{}P-{}~-{}-{}-{}-{}-{}+-{}-{}~2:1101}\par
\texttt{F4.1376~023.03.01.0~~~~~~~~~-{}-{}-{}-{}D~-{}PC~-{}-{}-{}-{}++-{}+~ok}\par
\texttt{F4.1377~123.03.01.0~~~~~~~~~-{}-{}-{}-{}D~-{}PC~-{}-{}-{}-{}++++~ok}\par
\texttt{F4.1378~012.13.01.0~~~~~~~~~-{}-{}-{}-{}D~-{}PC~-{}-{}-{}-{}++-{}+~ok}\par
\texttt{F4.1379~013.13.01.0~~~~~~~~~-{}-{}-{}-{}D~-{}P-{}~-{}-{}-{}-{}-{}+-{}-{}~2:1101}\par
\texttt{F4.1380~023.13.01.0~~~~~~~~~-{}-{}-{}-{}D~-{}PC~-{}-{}-{}-{}++-{}+~ok}\par
\texttt{F4.1381~123.13.01.0~~~~~~~~~-{}-{}-{}-{}D~-{}PC~-{}-{}-{}-{}++-{}+~ok}\par
\texttt{F4.1382~01.013.01.0~~~~~~~~~-{}-{}-{}-{}D~-{}P-{}~-{}-{}-{}-{}-{}+-{}-{}~2:1101}\par
\texttt{F4.1383~02.013.01.0~~~~~~~~~-{}-{}-{}-{}D~-{}PC~-{}-{}-{}-{}++-{}-{}~ok}\par
\texttt{F4.1384~12.013.01.0~~~~~~~~~-{}-{}-{}-{}D~-{}PC~-{}-{}-{}-{}++-{}+~ok}\par
\texttt{F4.1385~03.013.01.0~~~~~~~~~-{}-{}-{}-{}D~-{}-{}-{}~-{}-{}-{}-{}-{}-{}-{}-{}~1:1001}\par
\texttt{F4.1386~13.013.01.0~~~~~~~~~-{}-{}-{}-{}D~-{}P-{}~-{}-{}-{}-{}-{}+-{}-{}~2:1101}\par
\texttt{F4.1387~23.013.01.0~~~~~~~~~-{}-{}-{}-{}D~-{}PC~-{}-{}-{}-{}++-{}+~ok}\par
\texttt{F4.1388~012.23.01.0~~~~~~~~~-{}-{}-{}-{}D~-{}PC~-{}-{}-{}-{}++++~ok}\par
\texttt{F4.1389~013.23.01.0~~~~~~~~~-{}-{}-{}-{}D~-{}PC~-{}-{}-{}-{}++-{}+~ok}\par
\texttt{F4.1390~023.23.01.0~~~~~~~~~-{}-{}-{}-{}D~-{}PC~-{}-{}-{}-{}++-{}+~ok}\par
\texttt{F4.1391~123.23.01.0~~~~~~~~~-{}-{}-{}-{}D~-{}PC~-{}-{}-{}-{}++++~ok}\par
\texttt{F4.1392~01.023.01.0~~~~~~~~~-{}-{}-{}-{}D~-{}PC~-{}-{}-{}-{}++++~ok}\par
\texttt{F4.1393~02.023.01.0~~~~~~~~~-{}-{}-{}-{}D~-{}PC~-{}-{}-{}-{}++-{}+~ok}\par
\texttt{F4.1394~12.023.01.0~~~~~~~~~-{}-{}-{}-{}D~-{}PC~-{}-{}-{}-{}++++~ok}\par
\texttt{F4.1395~03.023.01.0~~~~~~~~~-{}-{}-{}-{}D~-{}-{}-{}~-{}-{}-{}-{}-{}-{}-{}-{}~1:1001}\par
\texttt{F4.1396~13.023.01.0~~~~~~~~~-{}-{}-{}-{}D~-{}PC~-{}-{}-{}-{}++-{}+~ok}\par
\texttt{F4.1397~23.023.01.0~~~~~~~~~-{}-{}-{}-{}D~-{}PC~-{}-{}-{}-{}++-{}+~ok}\par
\texttt{F4.1398~01.123.01.0~~~~~~~~~-{}-{}-{}-{}D~-{}PC~-{}-{}-{}-{}++-{}+~ok}\par
\texttt{F4.1399~02.123.01.0~~~~~~~~~-{}-{}-{}-{}D~-{}PC~-{}-{}-{}-{}++-{}-{}~ok}\par
\texttt{F4.1400~12.123.01.0~~~~~~~~~-{}-{}-{}-{}D~-{}PC~-{}-{}-{}-{}++-{}+~ok}\par
\texttt{F4.1401~03.123.01.0~~~~~~~~~-{}-{}-{}-{}D~-{}-{}-{}~-{}-{}-{}-{}-{}-{}-{}-{}~1:1001}\par
\texttt{F4.1402~13.123.01.0~~~~~~~~~-{}-{}-{}-{}D~-{}PC~-{}-{}-{}-{}++-{}+~ok}\par
\texttt{F4.1403~23.123.01.0~~~~~~~~~-{}-{}-{}-{}D~-{}PC~-{}-{}-{}-{}++-{}+~ok}\par
\texttt{F4.1404~0.0123.01.0~~~~~~~~~-{}-{}-{}-{}D~-{}-{}-{}~-{}-{}-{}-{}-{}-{}-{}-{}~1:1000}\par
\texttt{F4.1405~1.0123.01.0~~~~~~~~~-{}-{}-{}-{}D~-{}PC~-{}-{}-{}-{}++-{}+~ok}\par
\texttt{F4.1406~2.0123.01.0~~~~~~~~~-{}-{}-{}-{}D~-{}PC~-{}-{}-{}-{}++-{}-{}~ok}\par
\texttt{F4.1407~3.0123.01.0~~~~~~~~~-{}-{}-{}-{}D~-{}-{}-{}~-{}-{}-{}-{}-{}-{}-{}-{}~1:1001}\par
\texttt{F4.1408~0123.01.2.0~~~~~~~~~-{}-{}-{}-{}D~-{}-{}-{}~-{}-{}-{}-{}-{}-{}-{}-{}~0:0010}\par
\texttt{F4.1409~0123.12.2.0~~~~~~~~~-{}-{}-{}-{}D~-{}-{}-{}~-{}-{}-{}-{}-{}-{}-{}-{}~0:0010}\par
\texttt{F4.1410~012.012.2.0~~~~~~~~~-{}-{}-{}-{}D~-{}-{}-{}~-{}-{}-{}-{}-{}-{}-{}-{}~0:0010}\par
\texttt{F4.1411~013.012.2.0~~~~~~~~~-{}-{}-{}-{}D~-{}-{}-{}~-{}-{}-{}-{}-{}-{}-{}-{}~1:0010}\par
\texttt{F4.1412~023.012.2.0~~~~~~~~~-{}-{}-{}-{}D~-{}-{}-{}~-{}-{}-{}-{}-{}-{}-{}-{}~0:0010}\par
\texttt{F4.1413~123.012.2.0~~~~~~~~~-{}-{}-{}-{}D~-{}-{}-{}~-{}-{}-{}-{}-{}-{}-{}-{}~0:0010}\par
\texttt{F4.1414~0123.03.2.0~~~~~~~~~-{}-{}-{}-{}D~-{}-{}-{}~-{}-{}-{}-{}-{}-{}-{}-{}~0:0010}\par
\texttt{F4.1415~0123.13.2.0~~~~~~~~~-{}-{}-{}-{}D~-{}-{}-{}~-{}-{}-{}-{}-{}-{}-{}-{}~0:0010}\par
\texttt{F4.1416~012.013.2.0~~~~~~~~~-{}-{}-{}-{}D~-{}-{}-{}~-{}-{}-{}-{}-{}-{}-{}-{}~0:0010}\par
\texttt{F4.1417~013.013.2.0~~~~~~~~~-{}-{}-{}-{}D~-{}PC~-{}-{}-{}-{}++++~ok}\par
\texttt{F4.1418~023.013.2.0~~~~~~~~~-{}-{}-{}-{}D~-{}-{}-{}~-{}-{}-{}-{}-{}-{}-{}-{}~0:0010}\par
\texttt{F4.1419~123.013.2.0~~~~~~~~~-{}-{}-{}-{}D~-{}-{}-{}~-{}-{}-{}-{}-{}-{}-{}-{}~0:0010}\par
\texttt{F4.1420~0123.23.2.0~~~~~~~~~-{}-{}-{}-{}D~-{}-{}-{}~-{}-{}-{}-{}-{}-{}-{}-{}~0:0010}\par
\texttt{F4.1421~012.023.2.0~~~~~~~~~-{}-{}-{}-{}D~-{}-{}-{}~-{}-{}-{}-{}-{}-{}-{}-{}~0:0010}\par
\texttt{F4.1422~013.023.2.0~~~~~~~~~-{}-{}-{}-{}D~-{}-{}-{}~-{}-{}-{}-{}-{}-{}-{}-{}~1:0010}\par
\texttt{F4.1423~023.023.2.0~~~~~~~~~-{}-{}-{}-{}D~-{}-{}-{}~-{}-{}-{}-{}-{}-{}-{}-{}~0:0010}\par
\texttt{F4.1424~123.023.2.0~~~~~~~~~-{}-{}-{}-{}D~-{}-{}-{}~-{}-{}-{}-{}-{}-{}-{}-{}~0:0010}\par
\texttt{F4.1425~012.123.2.0~~~~~~~~~-{}-{}-{}-{}D~-{}-{}-{}~-{}-{}-{}-{}-{}-{}-{}-{}~0:0010}\par
\texttt{F4.1426~013.123.2.0~~~~~~~~~-{}-{}-{}-{}D~-{}-{}-{}~-{}-{}-{}-{}-{}-{}-{}-{}~1:0010}\par
\texttt{F4.1427~023.123.2.0~~~~~~~~~-{}-{}-{}-{}D~-{}-{}-{}~-{}-{}-{}-{}-{}-{}-{}-{}~0:0010}\par
\texttt{F4.1428~123.123.2.0~~~~~~~~~-{}-{}-{}-{}D~-{}-{}-{}~-{}-{}-{}-{}-{}-{}-{}-{}~0:0010}\par
\texttt{F4.1429~01.0123.2.0~~~~~~~~~-{}-{}-{}-{}D~-{}-{}-{}~-{}-{}-{}-{}-{}-{}-{}-{}~1:0010}\par
\texttt{F4.1430~02.0123.2.0~~~~~~~~~-{}-{}-{}-{}D~-{}-{}-{}~-{}-{}-{}-{}-{}-{}-{}-{}~0:0010}\par
\texttt{F4.1431~12.0123.2.0~~~~~~~~~-{}-{}-{}-{}D~-{}-{}-{}~-{}-{}-{}-{}-{}-{}-{}-{}~0:0010}\par
\texttt{F4.1432~03.0123.2.0~~~~~~~~~-{}-{}-{}-{}D~-{}-{}-{}~-{}-{}-{}-{}-{}-{}-{}-{}~1:0010}\par
\texttt{F4.1433~13.0123.2.0~~~~~~~~~-{}-{}-{}-{}D~-{}-{}-{}~-{}-{}-{}-{}-{}-{}-{}-{}~1:0010}\par
\texttt{F4.1434~23.0123.2.0~~~~~~~~~-{}-{}-{}-{}D~-{}-{}-{}~-{}-{}-{}-{}-{}-{}-{}-{}~0:0010}\par
\texttt{F4.1435~012.01.02.0~~~~~~~~~-{}-{}-{}-{}D~-{}P-{}~-{}-{}-{}-{}-{}+-{}-{}~3:1110}\par
\texttt{F4.1436~013.01.02.0~~~~~~~~~-{}-{}-{}-{}D~-{}-{}-{}~-{}-{}-{}-{}-{}-{}-{}-{}~2:1101}\par
\texttt{F4.1437~123.01.02.0~~~~~~~~~-{}B-{}-{}D~-{}PC~+-{}-{}-{}++++~ok}\par
\texttt{F4.1438~012.12.02.0~~~~~~~~~-{}-{}-{}-{}D~-{}P-{}~-{}-{}-{}-{}-{}+-{}-{}~3:1110}\par
\texttt{F4.1439~013.12.02.0~~~~~~~~~-{}-{}-{}-{}D~-{}PC~-{}-{}-{}-{}++-{}+~ok}\par
\texttt{F4.1440~023.12.02.0~~~~~~~~~-{}-{}-{}-{}D~-{}-{}-{}~-{}-{}-{}-{}-{}-{}-{}-{}~1:1011}\par
\texttt{F4.1441~123.12.02.0~~~~~~~~~-{}-{}-{}-{}D~-{}PC~-{}-{}-{}-{}++++~ok}\par
\texttt{F4.1442~01.012.02.0~~~~~~~~~-{}-{}-{}-{}D~-{}P-{}~-{}-{}-{}-{}-{}+-{}-{}~3:1110}\par
\texttt{F4.1443~02.012.02.0~~~~~~~~~-{}-{}-{}-{}D~-{}-{}-{}~-{}-{}-{}-{}-{}-{}-{}-{}~1:1010}\par
\texttt{F4.1444~12.012.02.0~~~~~~~~~-{}-{}-{}-{}D~-{}P-{}~-{}-{}-{}-{}-{}+-{}-{}~3:1110}\par
\texttt{F4.1445~03.012.02.0~~~~~~~~~-{}-{}-{}-{}D~-{}-{}-{}~-{}-{}-{}-{}-{}-{}-{}-{}~1:1001}\par
\texttt{F4.1446~13.012.02.0~~~~~~~~~-{}-{}-{}-{}D~-{}PC~-{}-{}-{}-{}++-{}+~ok}\par
\texttt{F4.1447~23.012.02.0~~~~~~~~~-{}-{}-{}-{}D~-{}-{}-{}~-{}-{}-{}-{}-{}-{}-{}-{}~1:1011}\par
\texttt{F4.1448~0123.3.02.0~~~~~~~~~-{}-{}-{}-{}D~-{}PC~-{}-{}-{}-{}++-{}+~ok}\par
\texttt{F4.1449~012.03.02.0~~~~~~~~~-{}-{}-{}-{}D~-{}PC~-{}-{}-{}-{}++-{}+~ok}\par
\texttt{F4.1450~013.03.02.0~~~~~~~~~-{}-{}-{}-{}D~-{}-{}-{}~-{}-{}-{}-{}-{}-{}-{}-{}~2:1101}\par
\texttt{F4.1451~023.03.02.0~~~~~~~~~-{}-{}-{}-{}D~-{}P-{}~-{}-{}-{}-{}-{}+-{}-{}~1:1011}\par
\texttt{F4.1452~123.03.02.0~~~~~~~~~-{}-{}-{}-{}D~-{}PC~-{}-{}-{}-{}++++~ok}\par
\texttt{F4.1453~012.13.02.0~~~~~~~~~-{}-{}-{}-{}D~-{}PC~-{}-{}-{}-{}++-{}+~ok}\par
\texttt{F4.1454~013.13.02.0~~~~~~~~~-{}-{}-{}-{}D~-{}-{}-{}~-{}-{}-{}-{}-{}-{}-{}-{}~2:1101}\par
\texttt{F4.1455~023.13.02.0~~~~~~~~~-{}-{}-{}-{}D~-{}-{}-{}~-{}-{}-{}-{}-{}-{}-{}-{}~1:1011}\par
\texttt{F4.1456~123.13.02.0~~~~~~~~~-{}-{}-{}-{}D~-{}PC~-{}-{}-{}-{}++-{}+~ok}\par
\texttt{F4.1457~01.013.02.0~~~~~~~~~-{}-{}-{}-{}D~-{}-{}-{}~-{}-{}-{}-{}-{}-{}-{}-{}~2:1101}\par
\texttt{F4.1458~02.013.02.0~~~~~~~~~-{}-{}-{}-{}D~-{}-{}-{}~-{}-{}-{}-{}-{}-{}-{}-{}~1:1010}\par
\texttt{F4.1459~12.013.02.0~~~~~~~~~-{}-{}-{}-{}D~-{}PC~-{}-{}-{}-{}++-{}+~ok}\par
\texttt{F4.1460~03.013.02.0~~~~~~~~~-{}-{}-{}-{}D~-{}-{}-{}~-{}-{}-{}-{}-{}-{}-{}-{}~1:1001}\par
\texttt{F4.1461~13.013.02.0~~~~~~~~~-{}-{}-{}-{}D~-{}-{}-{}~-{}-{}-{}-{}-{}-{}-{}-{}~2:1101}\par
\texttt{F4.1462~23.013.02.0~~~~~~~~~-{}-{}-{}-{}D~-{}-{}-{}~-{}-{}-{}-{}-{}-{}-{}-{}~1:1011}\par
\texttt{F4.1463~012.23.02.0~~~~~~~~~-{}-{}-{}-{}D~-{}PC~-{}-{}-{}-{}++-{}+~ok}\par
\texttt{F4.1464~013.23.02.0~~~~~~~~~-{}-{}-{}-{}D~-{}PC~-{}-{}-{}-{}++-{}+~ok}\par
\texttt{F4.1465~023.23.02.0~~~~~~~~~-{}-{}-{}-{}D~-{}P-{}~-{}-{}-{}-{}-{}+-{}-{}~1:1011}\par
\texttt{F4.1466~123.23.02.0~~~~~~~~~-{}-{}-{}-{}D~-{}PC~-{}-{}-{}-{}++++~ok}\par
\texttt{F4.1467~01.023.02.0~~~~~~~~~-{}-{}-{}-{}D~-{}PC~-{}-{}-{}-{}++-{}+~ok}\par
\texttt{F4.1468~02.023.02.0~~~~~~~~~-{}-{}-{}-{}D~-{}P-{}~-{}-{}-{}-{}-{}+-{}-{}~1:1010}\par
\texttt{F4.1469~12.023.02.0~~~~~~~~~-{}-{}-{}-{}D~-{}PC~-{}-{}-{}-{}++-{}+~ok}\par
\texttt{F4.1470~03.023.02.0~~~~~~~~~-{}-{}-{}-{}D~-{}-{}-{}~-{}-{}-{}-{}-{}-{}-{}-{}~1:1001}\par
\texttt{F4.1471~13.023.02.0~~~~~~~~~-{}-{}-{}-{}D~-{}PC~-{}-{}-{}-{}++-{}+~ok}\par
\texttt{F4.1472~23.023.02.0~~~~~~~~~-{}-{}-{}-{}D~-{}P-{}~-{}-{}-{}-{}-{}+-{}-{}~1:1011}\par
\texttt{F4.1473~01.123.02.0~~~~~~~~~-{}-{}-{}-{}D~-{}PC~-{}-{}-{}-{}++-{}+~ok}\par
\texttt{F4.1474~02.123.02.0~~~~~~~~~-{}-{}-{}-{}D~-{}-{}-{}~-{}-{}-{}-{}-{}-{}-{}-{}~1:1010}\par
\texttt{F4.1475~12.123.02.0~~~~~~~~~-{}-{}-{}-{}D~-{}PC~-{}-{}-{}-{}++-{}+~ok}\par
\texttt{F4.1476~03.123.02.0~~~~~~~~~-{}-{}-{}-{}D~-{}-{}-{}~-{}-{}-{}-{}-{}-{}-{}-{}~1:1001}\par
\texttt{F4.1477~13.123.02.0~~~~~~~~~-{}-{}-{}-{}D~-{}PC~-{}-{}-{}-{}++-{}+~ok}\par
\texttt{F4.1478~23.123.02.0~~~~~~~~~-{}-{}-{}-{}D~-{}-{}-{}~-{}-{}-{}-{}-{}-{}-{}-{}~1:1011}\par
\texttt{F4.1479~0.0123.02.0~~~~~~~~~-{}-{}4-{}D~-{}-{}-{}~-{}-{}-{}+-{}-{}-{}-{}~1:1000}\par
\texttt{F4.1480~1.0123.02.0~~~~~~~~~-{}-{}-{}-{}D~-{}PC~-{}-{}-{}-{}++-{}+~ok}\par
\texttt{F4.1481~2.0123.02.0~~~~~~~~~-{}-{}-{}-{}D~-{}-{}-{}~-{}-{}-{}-{}-{}-{}-{}-{}~1:1010}\par
\texttt{F4.1482~3.0123.02.0~~~~~~~~~-{}-{}-{}-{}D~-{}-{}-{}~-{}-{}-{}-{}-{}-{}-{}-{}~1:1001}\par
\texttt{F4.1483~012.12.12.0~~~~~~~~~-{}-{}-{}-{}D~-{}-{}-{}~-{}-{}-{}-{}-{}-{}-{}-{}~0:0110}\par
\texttt{F4.1484~013.12.12.0~~~~~~~~~-{}-{}-{}-{}D~-{}-{}-{}~-{}-{}-{}-{}-{}-{}-{}-{}~0:0110}\par
\texttt{F4.1485~123.12.12.0~~~~~~~~~-{}-{}-{}-{}D~-{}-{}-{}~-{}-{}-{}-{}-{}-{}-{}-{}~0:0110}\par
\texttt{F4.1486~01.012.12.0~~~~~~~~~-{}-{}-{}-{}D~-{}P-{}~-{}-{}-{}-{}-{}+-{}-{}~3:1110}\par
\texttt{F4.1487~02.012.12.0~~~~~~~~~-{}-{}-{}-{}D~-{}P-{}~-{}-{}-{}-{}-{}+-{}-{}~3:1110}\par
\texttt{F4.1488~12.012.12.0~~~~~~~~~-{}-{}-{}-{}D~-{}P-{}~-{}-{}-{}-{}-{}+-{}-{}~3:1110}\par
\texttt{F4.1489~03.012.12.0~~~~~~~~~-{}-{}-{}-{}D~-{}-{}-{}~-{}-{}-{}-{}-{}-{}-{}-{}~1:1001}\par
\texttt{F4.1490~13.012.12.0~~~~~~~~~-{}B-{}-{}D~-{}PC~+-{}-{}-{}++-{}+~ok}\par
\texttt{F4.1491~23.012.12.0~~~~~~~~~-{}-{}-{}-{}D~-{}PC~-{}-{}-{}-{}++-{}+~ok}\par
\texttt{F4.1492~0123.3.12.0~~~~~~~~~-{}-{}-{}-{}D~-{}PC~-{}-{}-{}-{}++-{}+~ok}\par
\texttt{F4.1493~012.03.12.0~~~~~~~~~-{}-{}-{}-{}D~-{}PC~-{}-{}-{}-{}++++~ok}\par
\texttt{F4.1494~013.03.12.0~~~~~~~~~-{}-{}-{}-{}D~-{}-{}-{}~-{}-{}-{}-{}-{}-{}-{}-{}~2:1101}\par
\texttt{F4.1495~023.03.12.0~~~~~~~~~-{}-{}-{}-{}D~-{}PC~-{}-{}-{}-{}++-{}+~ok}\par
\texttt{F4.1496~123.03.12.0~~~~~~~~~-{}-{}-{}-{}D~-{}PC~-{}-{}-{}-{}++++~ok}\par
\texttt{F4.1497~012.13.12.0~~~~~~~~~-{}-{}-{}-{}D~-{}PC~-{}-{}-{}-{}++-{}+~ok}\par
\texttt{F4.1498~013.13.12.0~~~~~~~~~-{}-{}-{}-{}D~-{}-{}-{}~-{}-{}-{}-{}-{}-{}-{}-{}~2:1101}\par
\texttt{F4.1499~023.13.12.0~~~~~~~~~-{}-{}-{}-{}D~-{}PC~-{}-{}-{}-{}++-{}+~ok}\par
\texttt{F4.1500~123.13.12.0~~~~~~~~~-{}-{}-{}-{}D~-{}PC~-{}-{}-{}-{}++-{}+~ok}\par
\texttt{F4.1501~01.013.12.0~~~~~~~~~-{}-{}-{}-{}D~-{}-{}-{}~-{}-{}-{}-{}-{}-{}-{}-{}~2:1101}\par
\texttt{F4.1502~02.013.12.0~~~~~~~~~-{}-{}-{}-{}D~-{}PC~-{}-{}-{}-{}++-{}-{}~ok}\par
\texttt{F4.1503~12.013.12.0~~~~~~~~~-{}-{}-{}-{}D~-{}PC~-{}-{}-{}-{}++-{}+~ok}\par
\texttt{F4.1504~03.013.12.0~~~~~~~~~-{}-{}-{}-{}D~-{}-{}-{}~-{}-{}-{}-{}-{}-{}-{}-{}~1:1001}\par
\texttt{F4.1505~13.013.12.0~~~~~~~~~-{}-{}-{}-{}D~-{}-{}-{}~-{}-{}-{}-{}-{}-{}-{}-{}~2:1101}\par
\texttt{F4.1506~23.013.12.0~~~~~~~~~-{}-{}-{}-{}D~-{}PC~-{}-{}-{}-{}++-{}+~ok}\par
\texttt{F4.1507~012.23.12.0~~~~~~~~~-{}-{}-{}-{}D~-{}PC~-{}-{}-{}-{}++-{}+~ok}\par
\texttt{F4.1508~013.23.12.0~~~~~~~~~-{}-{}-{}-{}D~-{}PC~-{}-{}-{}-{}++-{}+~ok}\par
\texttt{F4.1509~023.23.12.0~~~~~~~~~-{}-{}-{}-{}D~-{}PC~-{}-{}-{}-{}++-{}+~ok}\par
\texttt{F4.1510~123.23.12.0~~~~~~~~~-{}-{}-{}-{}D~-{}PC~-{}-{}-{}-{}++-{}+~ok}\par
\texttt{F4.1511~01.023.12.0~~~~~~~~~-{}-{}-{}-{}D~-{}PC~-{}-{}-{}-{}++-{}+~ok}\par
\texttt{F4.1512~02.023.12.0~~~~~~~~~-{}-{}-{}-{}D~-{}PC~-{}-{}-{}-{}++-{}+~ok}\par
\texttt{F4.1513~12.023.12.0~~~~~~~~~-{}-{}-{}-{}D~-{}PC~-{}-{}-{}-{}++++~ok}\par
\texttt{F4.1514~03.023.12.0~~~~~~~~~-{}-{}-{}-{}D~-{}-{}-{}~-{}-{}-{}-{}-{}-{}-{}-{}~1:1001}\par
\texttt{F4.1515~13.023.12.0~~~~~~~~~-{}-{}-{}-{}D~-{}PC~-{}-{}-{}-{}++-{}+~ok}\par
\texttt{F4.1516~23.023.12.0~~~~~~~~~-{}-{}-{}-{}D~-{}PC~-{}-{}-{}-{}++-{}+~ok}\par
\texttt{F4.1517~01.123.12.0~~~~~~~~~-{}-{}-{}-{}D~-{}PC~-{}-{}-{}-{}++-{}+~ok}\par
\texttt{F4.1518~02.123.12.0~~~~~~~~~-{}-{}-{}-{}D~-{}PC~-{}-{}-{}-{}++-{}-{}~ok}\par
\texttt{F4.1519~12.123.12.0~~~~~~~~~-{}-{}-{}-{}D~-{}PC~-{}-{}-{}-{}++-{}+~ok}\par
\texttt{F4.1520~03.123.12.0~~~~~~~~~-{}-{}-{}-{}D~-{}-{}-{}~-{}-{}-{}-{}-{}-{}-{}-{}~1:1001}\par
\texttt{F4.1521~13.123.12.0~~~~~~~~~-{}-{}-{}-{}D~-{}PC~-{}-{}-{}-{}++-{}+~ok}\par
\texttt{F4.1522~23.123.12.0~~~~~~~~~-{}-{}-{}-{}D~-{}PC~-{}-{}-{}-{}++-{}+~ok}\par
\texttt{F4.1523~0.0123.12.0~~~~~~~~~-{}-{}-{}-{}D~-{}-{}-{}~-{}-{}-{}-{}-{}-{}-{}-{}~1:1000}\par
\texttt{F4.1524~2.0123.12.0~~~~~~~~~-{}-{}-{}-{}D~-{}PC~-{}-{}-{}-{}++-{}-{}~ok}\par
\texttt{F4.1525~3.0123.12.0~~~~~~~~~-{}-{}-{}-{}D~-{}-{}-{}~-{}-{}-{}-{}-{}-{}-{}-{}~1:1001}\par
\texttt{F4.1526~0.012.012.0~~~~~~~~~-{}-{}4-{}D~-{}-{}-{}~-{}-{}-{}+-{}-{}-{}-{}~1:1000}\par
\texttt{F4.1527~1.012.012.0~~~~~~~~~-{}-{}-{}-{}D~-{}P-{}~-{}-{}-{}-{}-{}+-{}-{}~3:1110}\par
\texttt{F4.1528~3.012.012.0~~~~~~~~~-{}-{}-{}-{}D~-{}-{}-{}~-{}-{}-{}-{}-{}-{}-{}-{}~1:1001}\par
\texttt{F4.1529~012.3.012.0~~~~~~~~~-{}-{}-{}-{}D~-{}PC~-{}-{}-{}-{}++-{}+~ok}\par
\texttt{F4.1530~013.3.012.0~~~~~~~~~-{}-{}-{}-{}D~-{}-{}-{}~-{}-{}-{}-{}-{}-{}-{}-{}~2:1101}\par
\texttt{F4.1531~023.3.012.0~~~~~~~~~-{}-{}-{}-{}D~-{}PC~-{}-{}-{}-{}++-{}-{}~ok}\par
\texttt{F4.1532~123.3.012.0~~~~~~~~~-{}-{}-{}-{}D~-{}PC~-{}-{}-{}-{}++-{}+~ok}\par
\texttt{F4.1533~01.03.012.0~~~~~~~~~-{}-{}-{}-{}D~-{}-{}-{}~-{}-{}-{}-{}-{}-{}-{}-{}~2:1101}\par
\texttt{F4.1534~02.03.012.0~~~~~~~~~-{}-{}-{}-{}D~-{}PC~-{}-{}-{}-{}++-{}+~ok}\par
\texttt{F4.1535~12.03.012.0~~~~~~~~~-{}-{}-{}-{}D~-{}PC~-{}-{}-{}-{}++++~ok}\par
\texttt{F4.1536~03.03.012.0~~~~~~~~~-{}-{}-{}-{}D~-{}-{}-{}~-{}-{}-{}-{}-{}-{}-{}-{}~1:1001}\par
\texttt{F4.1537~13.03.012.0~~~~~~~~~-{}-{}-{}-{}D~-{}-{}-{}~-{}-{}-{}-{}-{}-{}-{}-{}~2:1101}\par
\texttt{F4.1538~23.03.012.0~~~~~~~~~-{}-{}-{}-{}D~-{}PC~-{}-{}-{}-{}++-{}+~ok}\par
\texttt{F4.1539~01.13.012.0~~~~~~~~~-{}-{}-{}-{}D~-{}-{}-{}~-{}-{}-{}-{}-{}-{}-{}-{}~2:1101}\par
\texttt{F4.1540~02.13.012.0~~~~~~~~~-{}-{}-{}-{}D~-{}PC~-{}-{}-{}-{}++-{}-{}~ok}\par
\texttt{F4.1541~12.13.012.0~~~~~~~~~-{}-{}-{}-{}D~-{}PC~-{}-{}-{}-{}++-{}+~ok}\par
\texttt{F4.1542~03.13.012.0~~~~~~~~~-{}-{}-{}-{}D~-{}-{}-{}~-{}-{}-{}-{}-{}-{}-{}-{}~1:1001}\par
\texttt{F4.1543~13.13.012.0~~~~~~~~~-{}-{}-{}-{}D~-{}-{}-{}~-{}-{}-{}-{}-{}-{}-{}-{}~2:1101}\par
\texttt{F4.1544~23.13.012.0~~~~~~~~~-{}-{}-{}-{}D~-{}PC~-{}-{}-{}-{}++-{}+~ok}\par
\texttt{F4.1545~0.013.012.0~~~~~~~~~-{}-{}-{}-{}D~-{}-{}-{}~-{}-{}-{}-{}-{}-{}-{}-{}~1:1000}\par
\texttt{F4.1546~1.013.012.0~~~~~~~~~-{}-{}-{}-{}D~-{}-{}-{}~-{}-{}-{}-{}-{}-{}-{}-{}~2:1101}\par
\texttt{F4.1547~2.013.012.0~~~~~~~~~-{}-{}-{}-{}D~-{}PC~-{}-{}-{}-{}++-{}-{}~ok}\par
\texttt{F4.1548~3.013.012.0~~~~~~~~~-{}-{}-{}-{}D~-{}-{}-{}~-{}-{}-{}-{}-{}-{}-{}-{}~1:1001}\par
\texttt{F4.1549~01.23.012.0~~~~~~~~~-{}-{}-{}-{}D~-{}PC~-{}-{}-{}-{}++-{}+~ok}\par
\texttt{F4.1550~02.23.012.0~~~~~~~~~-{}-{}-{}-{}D~-{}PC~-{}-{}-{}-{}++-{}+~ok}\par
\texttt{F4.1551~12.23.012.0~~~~~~~~~-{}-{}-{}-{}D~-{}PC~-{}-{}-{}-{}++++~ok}\par
\texttt{F4.1552~03.23.012.0~~~~~~~~~-{}-{}-{}-{}D~-{}-{}-{}~-{}-{}-{}-{}-{}-{}-{}-{}~1:1001}\par
\texttt{F4.1553~13.23.012.0~~~~~~~~~-{}-{}-{}-{}D~-{}PC~-{}-{}-{}-{}++-{}+~ok}\par
\texttt{F4.1554~23.23.012.0~~~~~~~~~-{}-{}-{}-{}D~-{}PC~-{}-{}-{}-{}++-{}+~ok}\par
\texttt{F4.1555~0.023.012.0~~~~~~~~~-{}-{}-{}-{}D~-{}-{}-{}~-{}-{}-{}-{}-{}-{}-{}-{}~1:1000}\par
\texttt{F4.1556~1.023.012.0~~~~~~~~~-{}-{}-{}-{}D~-{}PC~-{}-{}-{}-{}++-{}+~ok}\par
\texttt{F4.1557~2.023.012.0~~~~~~~~~-{}-{}-{}-{}D~-{}PC~-{}-{}-{}-{}++-{}+~ok}\par
\texttt{F4.1558~3.023.012.0~~~~~~~~~-{}-{}-{}-{}D~-{}-{}-{}~-{}-{}-{}-{}-{}-{}-{}-{}~1:1001}\par
\texttt{F4.1559~0.123.012.0~~~~~~~~~-{}-{}-{}-{}D~-{}-{}-{}~-{}-{}-{}-{}-{}-{}-{}-{}~1:1000}\par
\texttt{F4.1560~1.123.012.0~~~~~~~~~-{}-{}-{}-{}D~-{}PC~-{}-{}-{}-{}++-{}+~ok}\par
\texttt{F4.1561~3.123.012.0~~~~~~~~~-{}-{}-{}-{}D~-{}-{}-{}~-{}-{}-{}-{}-{}-{}-{}-{}~1:1001}\par
\texttt{F4.1562~0123.23.3.0~~~~~~~~~-{}-{}-{}-{}D~-{}PC~-{}-{}-{}-{}++-{}+~ok}\par
\texttt{F4.1563~012.023.3.0~~~~~~~~~-{}-{}-{}-{}D~-{}PC~-{}-{}-{}-{}++-{}+~ok}\par
\texttt{F4.1564~023.023.3.0~~~~~~~~~-{}-{}-{}-{}D~-{}P-{}~-{}-{}-{}-{}-{}+-{}-{}~1:1011}\par
\texttt{F4.1565~123.023.3.0~~~~~~~~~-{}-{}-{}-{}D~-{}PC~-{}-{}-{}-{}++-{}+~ok}\par
\texttt{F4.1566~012.123.3.0~~~~~~~~~-{}-{}-{}-{}D~-{}PC~-{}-{}-{}-{}++-{}+~ok}\par
\texttt{F4.1567~023.123.3.0~~~~~~~~~-{}-{}-{}-{}D~-{}-{}-{}~-{}-{}-{}-{}-{}-{}-{}-{}~1:1011}\par
\texttt{F4.1568~123.123.3.0~~~~~~~~~-{}-{}-{}-{}D~-{}PC~-{}-{}-{}-{}++-{}+~ok}\par
\texttt{F4.1569~02.0123.3.0~~~~~~~~~-{}-{}-{}-{}D~-{}-{}-{}~-{}-{}-{}-{}-{}-{}-{}-{}~1:1011}\par
\texttt{F4.1570~12.0123.3.0~~~~~~~~~-{}-{}-{}-{}D~-{}PC~-{}-{}-{}-{}++-{}+~ok}\par
\texttt{F4.1571~23.0123.3.0~~~~~~~~~-{}-{}-{}-{}D~-{}-{}-{}~-{}-{}-{}-{}-{}-{}-{}-{}~1:1011}\par
\texttt{F4.1572~012.03.03.0~~~~~~~~~-{}-{}-{}-{}D~-{}PC~-{}-{}-{}-{}++++~ok}\par
\texttt{F4.1573~013.03.03.0~~~~~~~~~-{}-{}-{}-{}D~-{}P-{}~-{}-{}-{}-{}-{}+-{}-{}~2:1101}\par
\texttt{F4.1574~123.03.03.0~~~~~~~~~-{}-{}-{}-{}D~-{}PC~-{}-{}-{}-{}++++~ok}\par
\texttt{F4.1575~012.13.03.0~~~~~~~~~-{}-{}-{}-{}D~-{}PC~-{}-{}-{}-{}++-{}+~ok}\par
\texttt{F4.1576~013.13.03.0~~~~~~~~~-{}-{}-{}-{}D~-{}P-{}~-{}-{}-{}-{}-{}+-{}-{}~2:1101}\par
\texttt{F4.1577~023.13.03.0~~~~~~~~~-{}-{}-{}-{}D~-{}-{}-{}~-{}-{}-{}-{}-{}-{}-{}-{}~1:1011}\par
\texttt{F4.1578~123.13.03.0~~~~~~~~~-{}-{}-{}-{}D~-{}PC~-{}-{}-{}-{}++++~ok}\par
\texttt{F4.1579~01.013.03.0~~~~~~~~~-{}-{}-{}-{}D~-{}P-{}~-{}-{}-{}-{}-{}+-{}-{}~2:1101}\par
\texttt{F4.1580~02.013.03.0~~~~~~~~~-{}-{}-{}-{}D~-{}-{}-{}~-{}-{}-{}-{}-{}-{}-{}-{}~1:1011}\par
\texttt{F4.1581~12.013.03.0~~~~~~~~~-{}-{}-{}-{}D~-{}PC~-{}-{}-{}-{}++-{}+~ok}\par
\texttt{F4.1582~03.013.03.0~~~~~~~~~-{}-{}-{}-{}D~-{}-{}-{}~-{}-{}-{}-{}-{}-{}-{}-{}~1:1001}\par
\texttt{F4.1583~13.013.03.0~~~~~~~~~-{}-{}-{}-{}D~-{}P-{}~-{}-{}-{}-{}-{}+-{}-{}~2:1101}\par
\texttt{F4.1584~23.013.03.0~~~~~~~~~-{}-{}-{}-{}D~-{}-{}-{}~-{}-{}-{}-{}-{}-{}-{}-{}~1:1011}\par
\texttt{F4.1585~012.23.03.0~~~~~~~~~-{}-{}-{}-{}D~-{}PC~-{}-{}-{}-{}++-{}+~ok}\par
\texttt{F4.1586~013.23.03.0~~~~~~~~~-{}-{}-{}-{}D~-{}PC~-{}-{}-{}-{}++-{}+~ok}\par
\texttt{F4.1587~023.23.03.0~~~~~~~~~-{}-{}-{}-{}D~-{}P-{}~-{}-{}-{}-{}-{}+-{}-{}~1:1011}\par
\texttt{F4.1588~123.23.03.0~~~~~~~~~-{}-{}-{}-{}D~-{}PC~-{}-{}-{}-{}++++~ok}\par
\texttt{F4.1589~01.023.03.0~~~~~~~~~-{}-{}-{}-{}D~-{}PC~-{}-{}-{}-{}++-{}+~ok}\par
\texttt{F4.1590~02.023.03.0~~~~~~~~~-{}-{}-{}-{}D~-{}P-{}~-{}-{}-{}-{}-{}+-{}-{}~1:1011}\par
\texttt{F4.1591~12.023.03.0~~~~~~~~~-{}-{}-{}-{}D~-{}PC~-{}-{}-{}-{}++-{}+~ok}\par
\texttt{F4.1592~03.023.03.0~~~~~~~~~-{}-{}-{}-{}D~-{}P-{}~-{}-{}-{}-{}-{}+-{}-{}~1:1001}\par
\texttt{F4.1593~13.023.03.0~~~~~~~~~-{}-{}-{}-{}D~-{}PC~-{}-{}-{}-{}++-{}+~ok}\par
\texttt{F4.1594~23.023.03.0~~~~~~~~~-{}-{}-{}-{}D~-{}P-{}~-{}-{}-{}-{}-{}+-{}-{}~1:1011}\par
\texttt{F4.1595~01.123.03.0~~~~~~~~~-{}-{}-{}-{}D~-{}PC~-{}-{}-{}-{}++-{}+~ok}\par
\texttt{F4.1596~02.123.03.0~~~~~~~~~-{}-{}-{}-{}D~-{}-{}-{}~-{}-{}-{}-{}-{}-{}-{}-{}~1:1011}\par
\texttt{F4.1597~12.123.03.0~~~~~~~~~-{}-{}-{}-{}D~-{}PC~-{}-{}-{}-{}++-{}+~ok}\par
\texttt{F4.1598~03.123.03.0~~~~~~~~~-{}-{}-{}-{}D~-{}-{}-{}~-{}-{}-{}-{}-{}-{}-{}-{}~1:1001}\par
\texttt{F4.1599~13.123.03.0~~~~~~~~~-{}-{}-{}-{}D~-{}PC~-{}-{}-{}-{}++-{}+~ok}\par
\texttt{F4.1600~23.123.03.0~~~~~~~~~-{}-{}-{}-{}D~-{}-{}-{}~-{}-{}-{}-{}-{}-{}-{}-{}~1:1011}\par
\texttt{F4.1601~0.0123.03.0~~~~~~~~~-{}-{}4-{}D~-{}-{}-{}~-{}-{}-{}+-{}-{}-{}-{}~1:1000}\par
\texttt{F4.1602~3.0123.03.0~~~~~~~~~-{}-{}-{}-{}D~-{}-{}-{}~-{}-{}-{}-{}-{}-{}-{}-{}~1:1001}\par
\texttt{F4.1603~012.13.13.0~~~~~~~~~-{}-{}-{}-{}D~-{}PC~-{}-{}-{}-{}++-{}+~ok}\par
\texttt{F4.1604~013.13.13.0~~~~~~~~~-{}-{}-{}-{}D~-{}P-{}~-{}-{}-{}-{}-{}+-{}-{}~2:1101}\par
\texttt{F4.1605~023.13.13.0~~~~~~~~~-{}-{}-{}-{}D~-{}PC~-{}-{}-{}-{}++-{}+~ok}\par
\texttt{F4.1606~123.13.13.0~~~~~~~~~-{}-{}-{}-{}D~-{}PC~-{}-{}-{}-{}++-{}+~ok}\par
\texttt{F4.1607~01.013.13.0~~~~~~~~~-{}-{}-{}-{}D~-{}P-{}~-{}-{}-{}-{}-{}+-{}-{}~2:1101}\par
\texttt{F4.1608~02.013.13.0~~~~~~~~~-{}-{}-{}-{}D~-{}PC~-{}-{}-{}-{}++-{}+~ok}\par
\texttt{F4.1609~12.013.13.0~~~~~~~~~-{}-{}-{}-{}D~-{}PC~-{}-{}-{}-{}++-{}+~ok}\par
\texttt{F4.1610~03.013.13.0~~~~~~~~~-{}-{}-{}-{}D~-{}-{}-{}~-{}-{}-{}-{}-{}-{}-{}-{}~1:1001}\par
\texttt{F4.1611~13.013.13.0~~~~~~~~~-{}-{}-{}-{}D~-{}P-{}~-{}-{}-{}-{}-{}+-{}-{}~2:1101}\par
\texttt{F4.1612~23.013.13.0~~~~~~~~~-{}-{}-{}-{}D~-{}PC~-{}-{}-{}-{}++-{}+~ok}\par
\texttt{F4.1613~012.23.13.0~~~~~~~~~-{}-{}-{}-{}D~-{}PC~-{}-{}-{}-{}++-{}+~ok}\par
\texttt{F4.1614~013.23.13.0~~~~~~~~~-{}-{}-{}-{}D~-{}PC~-{}-{}-{}-{}++-{}+~ok}\par
\texttt{F4.1615~123.23.13.0~~~~~~~~~-{}-{}-{}-{}D~-{}PC~-{}-{}-{}-{}++++~ok}\par
\texttt{F4.1616~01.023.13.0~~~~~~~~~-{}-{}-{}-{}D~-{}PC~-{}-{}-{}-{}++-{}+~ok}\par
\texttt{F4.1617~02.023.13.0~~~~~~~~~-{}-{}-{}-{}D~-{}PC~-{}-{}-{}-{}++-{}+~ok}\par
\texttt{F4.1618~12.023.13.0~~~~~~~~~-{}-{}-{}-{}D~-{}PC~-{}-{}-{}-{}++-{}+~ok}\par
\texttt{F4.1619~03.023.13.0~~~~~~~~~-{}-{}-{}-{}D~-{}-{}-{}~-{}-{}-{}-{}-{}-{}-{}-{}~1:1001}\par
\texttt{F4.1620~13.023.13.0~~~~~~~~~-{}-{}-{}-{}D~-{}PC~-{}-{}-{}-{}++-{}+~ok}\par
\texttt{F4.1621~23.023.13.0~~~~~~~~~-{}-{}-{}-{}D~-{}PC~-{}-{}-{}-{}++-{}+~ok}\par
\texttt{F4.1622~01.123.13.0~~~~~~~~~-{}-{}-{}-{}D~-{}PC~-{}-{}-{}-{}++-{}+~ok}\par
\texttt{F4.1623~02.123.13.0~~~~~~~~~-{}-{}-{}-{}D~-{}PC~-{}-{}-{}-{}++-{}+~ok}\par
\texttt{F4.1624~12.123.13.0~~~~~~~~~-{}-{}-{}-{}D~-{}PC~-{}-{}-{}-{}++-{}+~ok}\par
\texttt{F4.1625~03.123.13.0~~~~~~~~~-{}-{}-{}-{}D~-{}-{}-{}~-{}-{}-{}-{}-{}-{}-{}-{}~1:1001}\par
\texttt{F4.1626~13.123.13.0~~~~~~~~~-{}-{}-{}-{}D~-{}PC~-{}-{}-{}-{}++-{}+~ok}\par
\texttt{F4.1627~23.123.13.0~~~~~~~~~-{}-{}-{}-{}D~-{}PC~-{}-{}-{}-{}++-{}+~ok}\par
\texttt{F4.1628~0.0123.13.0~~~~~~~~~-{}-{}-{}-{}D~-{}-{}-{}~-{}-{}-{}-{}-{}-{}-{}-{}~1:1000}\par
\texttt{F4.1629~01.13.013.0~~~~~~~~~-{}-{}-{}-{}D~-{}P-{}~-{}-{}-{}-{}-{}+-{}-{}~2:1101}\par
\texttt{F4.1630~02.13.013.0~~~~~~~~~-{}-{}-{}-{}D~-{}PC~-{}-{}-{}-{}++-{}+~ok}\par
\texttt{F4.1631~12.13.013.0~~~~~~~~~-{}-{}-{}-{}D~-{}PC~-{}-{}-{}-{}++-{}+~ok}\par
\texttt{F4.1632~03.13.013.0~~~~~~~~~-{}-{}-{}-{}D~-{}-{}-{}~-{}-{}-{}-{}-{}-{}-{}-{}~1:1001}\par
\texttt{F4.1633~13.13.013.0~~~~~~~~~-{}-{}-{}-{}D~-{}P-{}~-{}-{}-{}-{}-{}+-{}-{}~2:1101}\par
\texttt{F4.1634~23.13.013.0~~~~~~~~~-{}-{}-{}-{}D~-{}PC~-{}-{}-{}-{}++-{}+~ok}\par
\texttt{F4.1635~0.013.013.0~~~~~~~~~-{}-{}4-{}D~-{}-{}-{}~-{}-{}-{}+-{}-{}-{}-{}~1:1000}\par
\texttt{F4.1636~3.013.013.0~~~~~~~~~-{}-{}-{}-{}D~-{}-{}-{}~-{}-{}-{}-{}-{}-{}-{}-{}~1:1001}\par
\texttt{F4.1637~0.023.013.0~~~~~~~~~-{}-{}-{}-{}D~-{}-{}-{}~-{}-{}-{}-{}-{}-{}-{}-{}~1:1000}\par
\texttt{F4.1638~3.023.013.0~~~~~~~~~-{}-{}-{}-{}D~-{}-{}-{}~-{}-{}-{}-{}-{}-{}-{}-{}~1:1001}\par
\texttt{F4.1639~0.123.013.0~~~~~~~~~-{}-{}-{}-{}D~-{}-{}-{}~-{}-{}-{}-{}-{}-{}-{}-{}~1:1000}\par
\texttt{F4.1640~012.13.23.0~~~~~~~~~-{}-{}-{}-{}D~-{}PC~-{}-{}-{}-{}++-{}+~ok}\par
\texttt{F4.1641~013.13.23.0~~~~~~~~~-{}-{}-{}-{}D~-{}-{}-{}~-{}-{}-{}-{}-{}-{}-{}-{}~2:1101}\par
\texttt{F4.1642~123.13.23.0~~~~~~~~~-{}-{}-{}-{}D~-{}PC~-{}-{}-{}-{}++-{}+~ok}\par
\texttt{F4.1643~01.013.23.0~~~~~~~~~-{}-{}-{}-{}D~-{}-{}-{}~-{}-{}-{}-{}-{}-{}-{}-{}~2:1101}\par
\texttt{F4.1644~02.013.23.0~~~~~~~~~-{}-{}-{}-{}D~-{}-{}-{}~-{}-{}-{}-{}-{}-{}-{}-{}~1:1011}\par
\texttt{F4.1645~12.013.23.0~~~~~~~~~-{}-{}-{}-{}D~-{}PC~-{}-{}-{}-{}++-{}+~ok}\par
\texttt{F4.1646~03.013.23.0~~~~~~~~~-{}-{}-{}-{}D~-{}-{}-{}~-{}-{}-{}-{}-{}-{}-{}-{}~1:1001}\par
\texttt{F4.1647~13.013.23.0~~~~~~~~~-{}-{}-{}-{}D~-{}-{}-{}~-{}-{}-{}-{}-{}-{}-{}-{}~2:1101}\par
\texttt{F4.1648~23.013.23.0~~~~~~~~~-{}-{}-{}-{}D~-{}-{}-{}~-{}-{}-{}-{}-{}-{}-{}-{}~1:1011}\par
\texttt{F4.1649~01.123.23.0~~~~~~~~~-{}-{}-{}-{}D~-{}PC~-{}-{}-{}-{}++-{}+~ok}\par
\texttt{F4.1650~02.123.23.0~~~~~~~~~-{}-{}-{}-{}D~-{}-{}-{}~-{}-{}-{}-{}-{}-{}-{}-{}~1:1011}\par
\texttt{F4.1651~12.123.23.0~~~~~~~~~-{}-{}-{}-{}D~-{}PC~-{}-{}-{}-{}++-{}+~ok}\par
\texttt{F4.1652~03.123.23.0~~~~~~~~~-{}-{}-{}-{}D~-{}-{}-{}~-{}-{}-{}-{}-{}-{}-{}-{}~1:1001}\par
\texttt{F4.1653~13.123.23.0~~~~~~~~~-{}-{}-{}-{}D~-{}PC~-{}-{}-{}-{}++-{}+~ok}\par
\texttt{F4.1654~23.123.23.0~~~~~~~~~-{}-{}-{}-{}D~-{}-{}-{}~-{}-{}-{}-{}-{}-{}-{}-{}~1:1011}\par
\texttt{F4.1655~0.0123.23.0~~~~~~~~~-{}-{}-{}-{}D~-{}-{}-{}~-{}-{}-{}-{}-{}-{}-{}-{}~1:1000}\par
\texttt{F4.1656~0.013.023.0~~~~~~~~~-{}-{}4-{}D~-{}-{}-{}~-{}-{}-{}+-{}-{}-{}-{}~1:1000}\par
\texttt{F4.1657~3.013.023.0~~~~~~~~~-{}-{}-{}-{}D~-{}-{}-{}~-{}-{}-{}-{}-{}-{}-{}-{}~1:1001}\par
\texttt{F4.1658~0.123.023.0~~~~~~~~~-{}-{}-{}-{}D~-{}-{}-{}~-{}-{}-{}-{}-{}-{}-{}-{}~1:1000}\par
\texttt{F4.1659~0.123.123.0~~~~~~~~~-{}-{}-{}-{}D~-{}-{}-{}~-{}-{}-{}-{}-{}-{}-{}-{}~1:1000}\par
\texttt{F4.1660~012.1.01.01~~~~~~~~~-{}-{}-{}-{}D~-{}-{}-{}~-{}-{}-{}-{}-{}-{}-{}-{}~0:0100}\par
\texttt{F4.1661~023.1.01.01~~~~~~~~~-{}-{}-{}-{}D~-{}-{}-{}~-{}-{}-{}-{}-{}-{}-{}-{}~2:0100}\par
\texttt{F4.1662~123.1.01.01~~~~~~~~~-{}-{}-{}-{}D~-{}-{}-{}~-{}-{}-{}-{}-{}-{}-{}-{}~0:0100}\par
\texttt{F4.1663~01.01.01.01~~~~~~~~~-{}-{}45D~UP-{}~-{}+++-{}+-{}-{}~2:1100}\par
\texttt{F4.1664~02.01.01.01~~~~~~~~~-{}-{}-{}-{}D~-{}P-{}~-{}-{}-{}-{}-{}+-{}-{}~3:1110}\par
\texttt{F4.1665~12.01.01.01~~~~~~~~~-{}-{}-{}-{}D~-{}P-{}~-{}-{}-{}-{}-{}+-{}-{}~3:1110}\par
\texttt{F4.1666~23.01.01.01~~~~~~~~~-{}-{}-{}-{}D~-{}PC~-{}-{}-{}-{}++++~ok}\par
\texttt{F4.1667~02.02.01.01~~~~~~~~~-{}-{}-{}-{}D~-{}P-{}~-{}-{}-{}-{}-{}+-{}-{}~3:1110}\par
\texttt{F4.1668~12.02.01.01~~~~~~~~~-{}-{}-{}-{}D~-{}P-{}~-{}-{}-{}-{}-{}+-{}-{}~3:1110}\par
\texttt{F4.1669~03.02.01.01~~~~~~~~~-{}-{}-{}-{}D~-{}PC~-{}-{}-{}-{}++-{}+~ok}\par
\texttt{F4.1670~13.02.01.01~~~~~~~~~-{}-{}-{}-{}D~-{}PC~-{}-{}-{}-{}++++~ok}\par
\texttt{F4.1671~23.02.01.01~~~~~~~~~-{}-{}-{}-{}D~-{}PC~-{}-{}-{}-{}++++~ok}\par
\texttt{F4.1672~02.12.01.01~~~~~~~~~-{}-{}-{}-{}D~-{}P-{}~-{}-{}-{}-{}-{}+-{}-{}~3:1110}\par
\texttt{F4.1673~03.12.01.01~~~~~~~~~-{}-{}-{}-{}D~-{}PC~-{}-{}-{}-{}++-{}+~ok}\par
\texttt{F4.1674~23.12.01.01~~~~~~~~~-{}-{}-{}-{}D~-{}PC~-{}-{}-{}-{}++-{}+~ok}\par
\texttt{F4.1675~23.23.01.01~~~~~~~~~-{}B-{}-{}D~UPC~+-{}+-{}++++~ok}\par
\texttt{F4.1676~0123.1.2.01~~~~~~~~~-{}-{}-{}-{}D~-{}-{}-{}~-{}-{}-{}-{}-{}-{}-{}-{}~0:0100}\par
\texttt{F4.1677~012.01.2.01~~~~~~~~~-{}-{}-{}-{}D~-{}-{}-{}~-{}-{}-{}-{}-{}-{}-{}-{}~0:0010}\par
\texttt{F4.1678~013.01.2.01~~~~~~~~~-{}-{}-{}-{}D~-{}PC~-{}-{}-{}-{}++++~ok}\par
\texttt{F4.1679~023.01.2.01~~~~~~~~~-{}-{}-{}-{}D~-{}-{}-{}~-{}-{}-{}-{}-{}-{}-{}-{}~0:0010}\par
\texttt{F4.1680~123.01.2.01~~~~~~~~~-{}-{}-{}-{}D~-{}-{}-{}~-{}-{}-{}-{}-{}-{}-{}-{}~0:0010}\par
\texttt{F4.1681~012.02.2.01~~~~~~~~~-{}-{}-{}-{}D~-{}-{}-{}~-{}-{}-{}-{}-{}-{}-{}-{}~0:0010}\par
\texttt{F4.1682~013.02.2.01~~~~~~~~~-{}-{}-{}-{}D~-{}-{}-{}~-{}-{}-{}-{}-{}-{}-{}-{}~1:0010}\par
\texttt{F4.1683~023.02.2.01~~~~~~~~~-{}-{}-{}-{}D~-{}-{}-{}~-{}-{}-{}-{}-{}-{}-{}-{}~0:0010}\par
\texttt{F4.1684~123.02.2.01~~~~~~~~~-{}-{}-{}-{}D~-{}-{}-{}~-{}-{}-{}-{}-{}-{}-{}-{}~0:0010}\par
\texttt{F4.1685~012.12.2.01~~~~~~~~~-{}-{}-{}-{}D~-{}-{}-{}~-{}-{}-{}-{}-{}-{}-{}-{}~0:0010}\par
\texttt{F4.1686~013.12.2.01~~~~~~~~~-{}-{}-{}-{}D~-{}-{}-{}~-{}-{}-{}-{}-{}-{}-{}-{}~1:0010}\par
\texttt{F4.1687~023.12.2.01~~~~~~~~~-{}-{}-{}-{}D~-{}-{}-{}~-{}-{}-{}-{}-{}-{}-{}-{}~0:0010}\par
\texttt{F4.1688~123.12.2.01~~~~~~~~~-{}-{}-{}-{}D~-{}-{}-{}~-{}-{}-{}-{}-{}-{}-{}-{}~0:0010}\par
\texttt{F4.1689~03.012.2.01~~~~~~~~~-{}-{}-{}-{}D~-{}-{}-{}~-{}-{}-{}-{}-{}-{}-{}-{}~1:0010}\par
\texttt{F4.1690~13.012.2.01~~~~~~~~~-{}-{}-{}-{}D~-{}-{}-{}~-{}-{}-{}-{}-{}-{}-{}-{}~1:0010}\par
\texttt{F4.1691~23.012.2.01~~~~~~~~~-{}-{}-{}-{}D~-{}-{}-{}~-{}-{}-{}-{}-{}-{}-{}-{}~0:0010}\par
\texttt{F4.1692~013.03.2.01~~~~~~~~~-{}B-{}-{}D~-{}PC~+-{}-{}-{}++++~ok}\par
\texttt{F4.1693~123.03.2.01~~~~~~~~~-{}-{}-{}-{}D~-{}-{}-{}~-{}-{}-{}-{}-{}-{}-{}-{}~0:0010}\par
\texttt{F4.1694~013.13.2.01~~~~~~~~~-{}-{}-{}-{}D~-{}PC~-{}-{}-{}-{}++++~ok}\par
\texttt{F4.1695~023.13.2.01~~~~~~~~~-{}-{}-{}-{}D~-{}-{}-{}~-{}-{}-{}-{}-{}-{}-{}-{}~0:0010}\par
\texttt{F4.1696~123.13.2.01~~~~~~~~~-{}-{}-{}-{}D~-{}-{}-{}~-{}-{}-{}-{}-{}-{}-{}-{}~0:0010}\par
\texttt{F4.1697~23.013.2.01~~~~~~~~~-{}-{}-{}-{}D~-{}-{}-{}~-{}-{}-{}-{}-{}-{}-{}-{}~0:0010}\par
\texttt{F4.1698~023.23.2.01~~~~~~~~~-{}-{}-{}-{}D~-{}-{}-{}~-{}-{}-{}-{}-{}-{}-{}-{}~0:0010}\par
\texttt{F4.1699~123.23.2.01~~~~~~~~~-{}-{}-{}-{}D~-{}-{}-{}~-{}-{}-{}-{}-{}-{}-{}-{}~0:0010}\par
\texttt{F4.1700~012.1.02.01~~~~~~~~~-{}-{}-{}-{}D~-{}-{}-{}~-{}-{}-{}-{}-{}-{}-{}-{}~0:0100}\par
\texttt{F4.1701~013.1.02.01~~~~~~~~~-{}-{}-{}-{}D~-{}-{}-{}~-{}-{}-{}-{}-{}-{}-{}-{}~0:0100}\par
\texttt{F4.1702~023.1.02.01~~~~~~~~~-{}-{}-{}-{}D~-{}-{}-{}~-{}-{}-{}-{}-{}-{}-{}-{}~3:0100}\par
\texttt{F4.1703~123.1.02.01~~~~~~~~~-{}-{}-{}-{}D~-{}-{}-{}~-{}-{}-{}-{}-{}-{}-{}-{}~0:0100}\par
\texttt{F4.1704~01.01.02.01~~~~~~~~~-{}-{}-{}-{}D~-{}-{}-{}~-{}-{}-{}-{}-{}-{}-{}-{}~2:1100}\par
\texttt{F4.1705~02.01.02.01~~~~~~~~~-{}-{}-{}-{}D~-{}-{}-{}~-{}-{}-{}-{}-{}-{}-{}-{}~1:1010}\par
\texttt{F4.1706~12.01.02.01~~~~~~~~~-{}-{}-{}-{}D~-{}P-{}~-{}-{}-{}-{}-{}+-{}-{}~3:1110}\par
\texttt{F4.1707~03.01.02.01~~~~~~~~~-{}-{}-{}-{}D~-{}-{}-{}~-{}-{}-{}-{}-{}-{}-{}-{}~2:1101}\par
\texttt{F4.1708~13.01.02.01~~~~~~~~~-{}-{}-{}-{}D~-{}-{}-{}~-{}-{}-{}-{}-{}-{}-{}-{}~2:1101}\par
\texttt{F4.1709~23.01.02.01~~~~~~~~~-{}-{}-{}-{}D~-{}PC~-{}-{}-{}-{}++-{}+~ok}\par
\texttt{F4.1710~01.02.02.01~~~~~~~~~-{}-{}-{}-{}D~-{}P-{}~-{}-{}-{}-{}-{}+-{}-{}~3:1110}\par
\texttt{F4.1711~02.02.02.01~~~~~~~~~-{}-{}-{}-{}D~UP-{}~-{}-{}+-{}-{}+-{}-{}~1:1010}\par
\texttt{F4.1712~12.02.02.01~~~~~~~~~-{}-{}-{}-{}D~-{}P-{}~-{}-{}-{}-{}-{}+-{}-{}~3:1110}\par
\texttt{F4.1713~03.02.02.01~~~~~~~~~-{}-{}-{}-{}D~-{}PC~-{}-{}-{}-{}++-{}+~ok}\par
\texttt{F4.1714~13.02.02.01~~~~~~~~~-{}-{}-{}-{}D~-{}PC~-{}-{}-{}-{}++++~ok}\par
\texttt{F4.1715~23.02.02.01~~~~~~~~~-{}-{}-{}-{}D~-{}PC~-{}-{}-{}-{}++-{}+~ok}\par
\texttt{F4.1716~01.12.02.01~~~~~~~~~-{}-{}-{}-{}D~-{}P-{}~-{}-{}-{}-{}-{}+-{}-{}~3:1110}\par
\texttt{F4.1717~02.12.02.01~~~~~~~~~-{}-{}-{}-{}D~-{}-{}-{}~-{}-{}-{}-{}-{}-{}-{}-{}~1:1010}\par
\texttt{F4.1718~12.12.02.01~~~~~~~~~-{}-{}-{}-{}D~-{}P-{}~-{}-{}-{}-{}-{}+-{}-{}~3:1110}\par
\texttt{F4.1719~03.12.02.01~~~~~~~~~-{}-{}-{}-{}D~-{}PC~-{}-{}-{}-{}++-{}+~ok}\par
\texttt{F4.1720~13.12.02.01~~~~~~~~~-{}-{}-{}-{}D~-{}PC~-{}-{}-{}-{}++-{}+~ok}\par
\texttt{F4.1721~23.12.02.01~~~~~~~~~-{}-{}-{}-{}D~-{}PC~-{}-{}-{}-{}++-{}+~ok}\par
\texttt{F4.1722~0.012.02.01~~~~~~~~~-{}-{}-{}-{}D~-{}-{}-{}~-{}-{}-{}-{}-{}-{}-{}-{}~1:1000}\par
\texttt{F4.1723~01.03.02.01~~~~~~~~~-{}-{}-{}-{}D~-{}-{}-{}~-{}-{}-{}-{}-{}-{}-{}-{}~2:1101}\par
\texttt{F4.1724~02.03.02.01~~~~~~~~~-{}-{}-{}-{}D~-{}-{}-{}~-{}-{}-{}-{}-{}-{}-{}-{}~1:1010}\par
\texttt{F4.1725~12.03.02.01~~~~~~~~~-{}-{}-{}-{}D~-{}PC~-{}-{}-{}-{}++-{}+~ok}\par
\texttt{F4.1726~13.03.02.01~~~~~~~~~-{}-{}-{}-{}D~-{}-{}-{}~-{}-{}-{}-{}-{}-{}-{}-{}~2:1101}\par
\texttt{F4.1727~01.13.02.01~~~~~~~~~-{}-{}-{}-{}D~-{}-{}-{}~-{}-{}-{}-{}-{}-{}-{}-{}~2:1101}\par
\texttt{F4.1728~02.13.02.01~~~~~~~~~-{}-{}-{}-{}D~-{}-{}-{}~-{}-{}-{}-{}-{}-{}-{}-{}~3:1010}\par
\texttt{F4.1729~12.13.02.01~~~~~~~~~-{}-{}-{}-{}D~-{}PC~-{}-{}-{}-{}++-{}+~ok}\par
\texttt{F4.1730~03.13.02.01~~~~~~~~~-{}-{}-{}-{}D~-{}-{}-{}~-{}-{}-{}-{}-{}-{}-{}-{}~2:1101}\par
\texttt{F4.1731~13.13.02.01~~~~~~~~~-{}-{}-{}-{}D~-{}-{}-{}~-{}-{}-{}-{}-{}-{}-{}-{}~2:1101}\par
\texttt{F4.1732~23.13.02.01~~~~~~~~~-{}B-{}-{}D~-{}PC~+-{}-{}-{}++-{}+~ok}\par
\texttt{F4.1733~0.013.02.01~~~~~~~~~-{}-{}-{}-{}D~-{}-{}-{}~-{}-{}-{}-{}-{}-{}-{}-{}~1:1000}\par
\texttt{F4.1734~01.23.02.01~~~~~~~~~-{}-{}-{}-{}D~-{}PC~-{}-{}-{}-{}++++~ok}\par
\texttt{F4.1735~02.23.02.01~~~~~~~~~-{}-{}-{}-{}D~-{}-{}-{}~-{}-{}-{}-{}-{}-{}-{}-{}~1:1010}\par
\texttt{F4.1736~12.23.02.01~~~~~~~~~-{}-{}-{}-{}D~-{}PC~-{}-{}-{}-{}++-{}+~ok}\par
\texttt{F4.1737~13.23.02.01~~~~~~~~~-{}-{}-{}-{}D~-{}PC~-{}-{}-{}-{}++++~ok}\par
\texttt{F4.1738~0.023.02.01~~~~~~~~~-{}-{}-{}-{}D~-{}-{}-{}~-{}-{}-{}-{}-{}-{}-{}-{}~1:1000}\par
\texttt{F4.1739~0.123.02.01~~~~~~~~~-{}-{}-{}-{}D~-{}-{}-{}~-{}-{}-{}-{}-{}-{}-{}-{}~2:1000}\par
\texttt{F4.1740~01.1.012.01~~~~~~~~~-{}-{}4-{}D~-{}-{}-{}~-{}-{}-{}+-{}-{}-{}-{}~0:0100}\par
\texttt{F4.1741~02.1.012.01~~~~~~~~~-{}-{}-{}-{}D~-{}-{}-{}~-{}-{}-{}-{}-{}-{}-{}-{}~2:0100}\par
\texttt{F4.1742~12.1.012.01~~~~~~~~~-{}-{}-{}-{}D~-{}-{}-{}~-{}-{}-{}-{}-{}-{}-{}-{}~0:0100}\par
\texttt{F4.1743~03.1.012.01~~~~~~~~~-{}-{}-{}-{}D~-{}-{}-{}~-{}-{}-{}-{}-{}-{}-{}-{}~2:0100}\par
\texttt{F4.1744~13.1.012.01~~~~~~~~~-{}-{}-{}-{}D~-{}-{}-{}~-{}-{}-{}-{}-{}-{}-{}-{}~0:0100}\par
\texttt{F4.1745~01.02.03.01~~~~~~~~~-{}-{}-{}-{}D~-{}PC~-{}-{}-{}-{}++-{}+~ok}\par
\texttt{F4.1746~12.02.03.01~~~~~~~~~-{}-{}-{}-{}D~-{}PC~-{}-{}-{}-{}++++~ok}\par
\texttt{F4.1747~01.12.03.01~~~~~~~~~-{}-{}-{}-{}D~-{}PC~-{}-{}-{}-{}++-{}+~ok}\par
\texttt{F4.1748~02.12.03.01~~~~~~~~~-{}-{}-{}-{}D~-{}PC~-{}-{}-{}-{}++-{}+~ok}\par
\texttt{F4.1749~12.12.03.01~~~~~~~~~-{}-{}-{}-{}D~-{}PC~-{}-{}-{}-{}++++~ok}\par
\texttt{F4.1750~03.12.03.01~~~~~~~~~-{}-{}-{}-{}D~-{}PC~-{}-{}-{}-{}++-{}+~ok}\par
\texttt{F4.1751~13.12.03.01~~~~~~~~~-{}-{}-{}-{}D~-{}PC~-{}-{}-{}-{}++-{}+~ok}\par
\texttt{F4.1752~0.012.03.01~~~~~~~~~-{}-{}-{}-{}D~-{}-{}-{}~-{}-{}-{}-{}-{}-{}-{}-{}~1:1000}\par
\texttt{F4.1753~12.23.03.01~~~~~~~~~-{}-{}-{}-{}D~-{}PC~-{}-{}-{}-{}++++~ok}\par
\texttt{F4.1754~0.023.03.01~~~~~~~~~-{}-{}-{}-{}D~-{}-{}-{}~-{}-{}-{}-{}-{}-{}-{}-{}~1:1000}\par
\texttt{F4.1755~0.123.03.01~~~~~~~~~-{}-{}-{}-{}D~-{}-{}-{}~-{}-{}-{}-{}-{}-{}-{}-{}~2:1000}\par
\texttt{F4.1756~01.1.013.01~~~~~~~~~-{}-{}4-{}D~-{}-{}-{}~-{}-{}-{}+-{}-{}-{}-{}~0:0100}\par
\texttt{F4.1757~02.1.013.01~~~~~~~~~-{}-{}-{}-{}D~-{}-{}-{}~-{}-{}-{}-{}-{}-{}-{}-{}~2:0100}\par
\texttt{F4.1758~03.1.013.01~~~~~~~~~-{}-{}-{}-{}D~-{}-{}-{}~-{}-{}-{}-{}-{}-{}-{}-{}~2:0100}\par
\texttt{F4.1759~13.1.013.01~~~~~~~~~-{}-{}-{}-{}D~-{}-{}-{}~-{}-{}-{}-{}-{}-{}-{}-{}~0:0100}\par
\texttt{F4.1760~012.1.23.01~~~~~~~~~-{}-{}-{}-{}D~-{}-{}-{}~-{}-{}-{}-{}-{}-{}-{}-{}~0:0100}\par
\texttt{F4.1761~013.1.23.01~~~~~~~~~-{}-{}-{}-{}D~-{}-{}-{}~-{}-{}-{}-{}-{}-{}-{}-{}~0:0100}\par
\texttt{F4.1762~023.1.23.01~~~~~~~~~-{}-{}-{}-{}D~-{}-{}-{}~-{}-{}-{}-{}-{}-{}-{}-{}~3:0100}\par
\texttt{F4.1763~123.1.23.01~~~~~~~~~-{}-{}-{}-{}D~-{}-{}-{}~-{}-{}-{}-{}-{}-{}-{}-{}~0:0100}\par
\texttt{F4.1764~01.01.23.01~~~~~~~~~-{}-{}-{}-{}D~-{}-{}-{}~-{}-{}-{}-{}-{}-{}-{}-{}~2:1100}\par
\texttt{F4.1765~02.01.23.01~~~~~~~~~-{}-{}-{}-{}D~-{}PC~-{}-{}-{}-{}++-{}+~ok}\par
\texttt{F4.1766~03.01.23.01~~~~~~~~~-{}-{}-{}-{}D~-{}-{}-{}~-{}-{}-{}-{}-{}-{}-{}-{}~2:1101}\par
\texttt{F4.1767~23.01.23.01~~~~~~~~~-{}-{}-{}-{}D~-{}PC~-{}-{}-{}-{}++++~ok}\par
\texttt{F4.1768~02.12.23.01~~~~~~~~~-{}-{}-{}-{}D~-{}PC~-{}-{}-{}-{}++-{}+~ok}\par
\texttt{F4.1769~03.12.23.01~~~~~~~~~-{}-{}-{}-{}D~-{}PC~-{}-{}-{}-{}++-{}+~ok}\par
\texttt{F4.1770~03.13.23.01~~~~~~~~~-{}-{}-{}-{}D~-{}-{}-{}~-{}-{}-{}-{}-{}-{}-{}-{}~2:1101}\par
\texttt{F4.1771~01.1.023.01~~~~~~~~~-{}-{}-{}-{}D~-{}-{}-{}~-{}-{}-{}-{}-{}-{}-{}-{}~0:0100}\par
\texttt{F4.1772~02.1.023.01~~~~~~~~~-{}-{}-{}-{}D~-{}-{}-{}~-{}-{}-{}-{}-{}-{}-{}-{}~3:0100}\par
\texttt{F4.1773~03.1.023.01~~~~~~~~~-{}-{}-{}-{}D~-{}-{}-{}~-{}-{}-{}-{}-{}-{}-{}-{}~3:0100}\par
\texttt{F4.1774~13.1.023.01~~~~~~~~~-{}-{}-{}-{}D~-{}-{}-{}~-{}-{}-{}-{}-{}-{}-{}-{}~0:0100}\par
\texttt{F4.1775~0.01.023.01~~~~~~~~~-{}-{}-{}-{}D~-{}-{}-{}~-{}-{}-{}-{}-{}-{}-{}-{}~1:1000}\par
\texttt{F4.1776~0.12.023.01~~~~~~~~~-{}-{}-{}-{}D~-{}-{}-{}~-{}-{}-{}-{}-{}-{}-{}-{}~2:1000}\par
\texttt{F4.1777~0.13.023.01~~~~~~~~~-{}-{}-{}-{}D~-{}-{}-{}~-{}-{}-{}-{}-{}-{}-{}-{}~2:1000}\par
\texttt{F4.1778~0.1.0123.01~~~~~~~~~-{}-{}4-{}D~-{}-{}-{}~-{}-{}-{}+-{}-{}-{}-{}~2:1000}\par
\texttt{F4.1779~012.1.2.012~~~~~~~~~-{}-{}4-{}D~-{}-{}-{}~-{}-{}-{}+-{}-{}-{}-{}~0:0100}\par
\texttt{F4.1780~013.1.2.012~~~~~~~~~-{}-{}-{}-{}D~-{}-{}-{}~-{}-{}-{}-{}-{}-{}-{}-{}~0:0100}\par
\texttt{F4.1781~123.1.2.012~~~~~~~~~-{}-{}-{}-{}D~-{}-{}-{}~-{}-{}-{}-{}-{}-{}-{}-{}~0:0100}\par
\texttt{F4.1782~01.01.2.012~~~~~~~~~-{}-{}4-{}D~-{}-{}-{}~-{}-{}-{}+-{}-{}-{}-{}~3:1100}\par
\texttt{F4.1783~02.01.2.012~~~~~~~~~-{}-{}-{}-{}D~-{}-{}-{}~-{}-{}-{}-{}-{}-{}-{}-{}~0:0010}\par
\texttt{F4.1784~03.01.2.012~~~~~~~~~-{}-{}-{}-{}D~-{}-{}-{}~-{}-{}-{}-{}-{}-{}-{}-{}~3:0010}\par
\texttt{F4.1785~02.12.2.012~~~~~~~~~-{}-{}4-{}D~-{}-{}-{}~-{}-{}-{}+-{}-{}-{}-{}~0:0010}\par
\texttt{F4.1786~03.12.2.012~~~~~~~~~-{}-{}-{}-{}D~-{}-{}-{}~-{}-{}-{}-{}-{}-{}-{}-{}~1:0010}\par
\texttt{F4.1787~03.13.2.012~~~~~~~~~-{}-{}-{}-{}D~-{}-{}-{}~-{}-{}-{}-{}-{}-{}-{}-{}~3:0010}\par
\texttt{F4.1788~0123.01.2.3~~~~~~~~~T-{}-{}-{}D~-{}-{}-{}~-{}-{}-{}-{}-{}-{}-{}-{}~0:0010}\par
\texttt{F4.1789~0123.12.2.3~~~~~~~~~T-{}4-{}D~-{}-{}-{}~-{}-{}-{}+-{}-{}-{}-{}~0:0010}\par
\texttt{F4.1790~012.012.2.3~~~~~~~~~T-{}4-{}D~-{}-{}-{}~-{}-{}-{}+-{}-{}-{}-{}~0:0010}\par
\texttt{F4.1791~013.012.2.3~~~~~~~~~T-{}-{}-{}D~-{}-{}-{}~-{}-{}-{}-{}-{}-{}-{}-{}~1:0010}\par
\texttt{F4.1792~023.012.2.3~~~~~~~~~T-{}-{}-{}D~-{}-{}-{}~-{}-{}-{}-{}-{}-{}-{}-{}~0:0010}\par
\texttt{F4.1793~023.123.2.3~~~~~~~~~T-{}4-{}D~-{}-{}-{}~-{}-{}-{}+-{}-{}-{}-{}~0:0010}\par
\texttt{F4.1794~012.01.02.3~~~~~~~~~TB-{}-{}D~-{}PC~+-{}-{}-{}++++~ok}\par
\texttt{F4.1795~013.01.02.3~~~~~~~~~T-{}-{}-{}D~-{}-{}-{}~-{}-{}-{}-{}-{}-{}-{}-{}~0:0001}\par
\texttt{F4.1796~012.12.02.3~~~~~~~~~T-{}-{}-{}D~-{}PC~-{}-{}-{}-{}++++~ok}\par
\texttt{F4.1797~013.12.02.3~~~~~~~~~T-{}-{}-{}D~-{}-{}-{}~-{}-{}-{}-{}-{}-{}-{}-{}~0:0001}\par
\texttt{F4.1798~023.12.02.3~~~~~~~~~T-{}-{}-{}D~-{}-{}-{}~-{}-{}-{}-{}-{}-{}-{}-{}~0:0001}\par
\texttt{F4.1799~02.012.02.3~~~~~~~~~T-{}4-{}D~-{}-{}-{}~-{}-{}-{}+-{}-{}-{}-{}~1:1010}\par
\texttt{F4.1800~03.012.02.3~~~~~~~~~T-{}-{}-{}D~-{}-{}-{}~-{}-{}-{}-{}-{}-{}-{}-{}~0:0001}\par
\texttt{F4.1801~012.13.02.3~~~~~~~~~T-{}-{}-{}D~-{}-{}-{}~-{}-{}-{}-{}-{}-{}-{}-{}~1:0001}\par
\texttt{F4.1802~013.13.02.3~~~~~~~~~T-{}-{}-{}D~-{}-{}-{}~-{}-{}-{}-{}-{}-{}-{}-{}~0:0001}\par
\texttt{F4.1803~023.13.02.3~~~~~~~~~T-{}-{}-{}D~-{}-{}-{}~-{}-{}-{}-{}-{}-{}-{}-{}~0:0001}\par
\texttt{F4.1804~02.013.02.3~~~~~~~~~T-{}-{}-{}D~-{}-{}-{}~-{}-{}-{}-{}-{}-{}-{}-{}~1:1010}\par
\texttt{F4.1805~03.013.02.3~~~~~~~~~T-{}-{}-{}D~-{}-{}-{}~-{}-{}-{}-{}-{}-{}-{}-{}~0:0001}\par
\texttt{F4.1806~03.123.02.3~~~~~~~~~T-{}-{}-{}D~-{}-{}-{}~-{}-{}-{}-{}-{}-{}-{}-{}~0:0001}\par
\texttt{F4.1807~03.13.012.3~~~~~~~~~T-{}-{}-{}D~-{}-{}-{}~-{}-{}-{}-{}-{}-{}-{}-{}~0:0001}\par
\texttt{F4.1808~013.13.23.3~~~~~~~~~T-{}4-{}D~-{}-{}-{}~-{}-{}-{}+-{}-{}-{}-{}~0:0001}\par
\texttt{F4.1809~01.01.02.03~~~~~~~~~T-{}-{}-{}D~-{}-{}-{}~-{}-{}-{}-{}-{}-{}-{}-{}~2:1100}\par
\texttt{F4.1810~01.12.02.03~~~~~~~~~T-{}-{}-{}D~-{}-{}-{}~-{}-{}-{}-{}-{}-{}-{}-{}~3:1110}\par
\texttt{F4.1811~02.12.02.03~~~~~~~~~T-{}-{}-{}D~-{}-{}-{}~-{}-{}-{}-{}-{}-{}-{}-{}~1:1010}\par
\texttt{F4.1812~03.12.02.03~~~~~~~~~T-{}-{}-{}D~-{}-{}-{}~-{}-{}-{}-{}-{}-{}-{}-{}~2:1001}\par
\texttt{F4.1813~03.12.12.03~~~~~~~~~TB45D~UPC~++++++++~ok}\par
\texttt{F4.1814~02.13.12.03~~~~~~~~~T-{}-{}-{}D~-{}PC~-{}-{}-{}-{}++-{}+~ok}\par
\texttt{F4.1815~0123.0123.0.-{}~~~~~~~-{}-{}-{}-{}-{}~-{}-{}-{}~-{}-{}-{}-{}-{}-{}-{}-{}~0:0000}\par
\texttt{F4.1816~0123.012.01.-{}~~~~~~~-{}-{}-{}-{}-{}~-{}-{}-{}~-{}-{}-{}-{}-{}-{}-{}-{}~0:0000}\par
\texttt{F4.1817~0123.013.01.-{}~~~~~~~-{}-{}-{}-{}-{}~-{}-{}-{}~-{}-{}-{}-{}-{}-{}-{}-{}~0:0000}\par
\texttt{F4.1818~0123.023.01.-{}~~~~~~~-{}-{}-{}-{}-{}~-{}-{}-{}~-{}-{}-{}-{}-{}-{}-{}-{}~0:0000}\par
\texttt{F4.1819~0123.123.01.-{}~~~~~~~-{}-{}-{}-{}-{}~-{}-{}-{}~-{}-{}-{}-{}-{}-{}-{}-{}~0:0000}\par
\texttt{F4.1820~0123.0123.2.-{}~~~~~~~-{}-{}4-{}-{}~-{}-{}-{}~-{}-{}-{}+-{}-{}-{}-{}~0:0000}\par
\texttt{F4.1821~0123.012.02.-{}~~~~~~~-{}-{}-{}-{}-{}~-{}-{}-{}~-{}-{}-{}-{}-{}-{}-{}-{}~0:0000}\par
\texttt{F4.1822~0123.013.02.-{}~~~~~~~-{}-{}-{}-{}-{}~-{}-{}-{}~-{}-{}-{}-{}-{}-{}-{}-{}~0:0000}\par
\texttt{F4.1823~0123.023.02.-{}~~~~~~~-{}-{}-{}-{}-{}~-{}-{}-{}~-{}-{}-{}-{}-{}-{}-{}-{}~0:0000}\par
\texttt{F4.1824~0123.123.02.-{}~~~~~~~-{}-{}-{}-{}-{}~-{}-{}-{}~-{}-{}-{}-{}-{}-{}-{}-{}~0:0000}\par
\texttt{F4.1825~012.0123.02.-{}~~~~~~~-{}-{}-{}-{}-{}~-{}-{}-{}~-{}-{}-{}-{}-{}-{}-{}-{}~1:0000}\par
\texttt{F4.1826~013.0123.02.-{}~~~~~~~-{}-{}-{}-{}-{}~-{}-{}-{}~-{}-{}-{}-{}-{}-{}-{}-{}~0:0000}\par
\texttt{F4.1827~023.0123.02.-{}~~~~~~~-{}-{}-{}-{}-{}~-{}-{}-{}~-{}-{}-{}-{}-{}-{}-{}-{}~0:0000}\par
\texttt{F4.1828~123.0123.02.-{}~~~~~~~-{}-{}-{}-{}-{}~-{}-{}-{}~-{}-{}-{}-{}-{}-{}-{}-{}~0:0000}\par
\texttt{F4.1829~012.012.012.-{}~~~~~~~-{}B45-{}~UPC~++++++++~ok}\par
\texttt{F4.1830~013.012.012.-{}~~~~~~~-{}-{}-{}-{}-{}~-{}-{}-{}~-{}-{}-{}-{}-{}-{}-{}-{}~0:0000}\par
\texttt{F4.1831~123.012.012.-{}~~~~~~~-{}-{}-{}-{}-{}~-{}-{}-{}~-{}-{}-{}-{}-{}-{}-{}-{}~0:0000}\par
\texttt{F4.1832~0123.03.012.-{}~~~~~~~-{}-{}-{}-{}-{}~-{}-{}-{}~-{}-{}-{}-{}-{}-{}-{}-{}~0:0000}\par
\texttt{F4.1833~0123.13.012.-{}~~~~~~~-{}-{}-{}-{}-{}~-{}-{}-{}~-{}-{}-{}-{}-{}-{}-{}-{}~0:0000}\par
\texttt{F4.1834~013.013.012.-{}~~~~~~~-{}-{}-{}-{}-{}~-{}-{}-{}~-{}-{}-{}-{}-{}-{}-{}-{}~0:0000}\par
\texttt{F4.1835~023.013.012.-{}~~~~~~~-{}-{}-{}-{}-{}~-{}-{}-{}~-{}-{}-{}-{}-{}-{}-{}-{}~0:0000}\par
\texttt{F4.1836~123.013.012.-{}~~~~~~~-{}-{}-{}-{}-{}~-{}-{}-{}~-{}-{}-{}-{}-{}-{}-{}-{}~0:0000}\par
\texttt{F4.1837~0123.23.012.-{}~~~~~~~-{}-{}-{}-{}-{}~-{}-{}-{}~-{}-{}-{}-{}-{}-{}-{}-{}~0:0000}\par
\texttt{F4.1838~023.023.012.-{}~~~~~~~-{}-{}-{}-{}-{}~-{}-{}-{}~-{}-{}-{}-{}-{}-{}-{}-{}~0:0000}\par
\texttt{F4.1839~123.023.012.-{}~~~~~~~-{}-{}-{}-{}-{}~-{}-{}-{}~-{}-{}-{}-{}-{}-{}-{}-{}~0:0000}\par
\texttt{F4.1840~023.123.012.-{}~~~~~~~-{}-{}-{}-{}-{}~-{}-{}-{}~-{}-{}-{}-{}-{}-{}-{}-{}~0:0000}\par
\texttt{F4.1841~0123.0123.3.-{}~~~~~~~-{}-{}4-{}-{}~-{}-{}-{}~-{}-{}-{}+-{}-{}-{}-{}~0:0000}\par
\texttt{F4.1842~0123.013.03.-{}~~~~~~~-{}-{}-{}-{}-{}~-{}-{}-{}~-{}-{}-{}-{}-{}-{}-{}-{}~0:0000}\par
\texttt{F4.1843~0123.023.03.-{}~~~~~~~-{}-{}-{}-{}-{}~-{}-{}-{}~-{}-{}-{}-{}-{}-{}-{}-{}~0:0000}\par
\texttt{F4.1844~0123.123.03.-{}~~~~~~~-{}-{}-{}-{}-{}~-{}-{}-{}~-{}-{}-{}-{}-{}-{}-{}-{}~0:0000}\par
\texttt{F4.1845~013.0123.03.-{}~~~~~~~-{}-{}-{}-{}-{}~-{}-{}-{}~-{}-{}-{}-{}-{}-{}-{}-{}~0:0000}\par
\texttt{F4.1846~023.0123.03.-{}~~~~~~~-{}-{}-{}-{}-{}~-{}-{}-{}~-{}-{}-{}-{}-{}-{}-{}-{}~0:0000}\par
\texttt{F4.1847~123.0123.03.-{}~~~~~~~-{}-{}-{}-{}-{}~-{}-{}-{}~-{}-{}-{}-{}-{}-{}-{}-{}~0:0000}\par
\texttt{F4.1848~0123.13.013.-{}~~~~~~~-{}-{}-{}-{}-{}~-{}-{}-{}~-{}-{}-{}-{}-{}-{}-{}-{}~0:0000}\par
\texttt{F4.1849~013.013.013.-{}~~~~~~~-{}-{}4-{}-{}~-{}-{}-{}~-{}-{}-{}+-{}-{}-{}-{}~0:0000}\par
\texttt{F4.1850~023.013.013.-{}~~~~~~~-{}-{}-{}-{}-{}~-{}-{}-{}~-{}-{}-{}-{}-{}-{}-{}-{}~0:0000}\par
\texttt{F4.1851~123.013.013.-{}~~~~~~~-{}-{}-{}-{}-{}~-{}-{}-{}~-{}-{}-{}-{}-{}-{}-{}-{}~0:0000}\par
\texttt{F4.1852~123.023.013.-{}~~~~~~~-{}-{}-{}-{}-{}~-{}-{}-{}~-{}-{}-{}-{}-{}-{}-{}-{}~0:0000}\par
\texttt{F4.1853~023.123.013.-{}~~~~~~~-{}-{}-{}-{}-{}~-{}-{}-{}~-{}-{}-{}-{}-{}-{}-{}-{}~0:0000}\par
\texttt{F4.1854~0123.013.23.-{}~~~~~~~-{}-{}-{}-{}-{}~-{}-{}-{}~-{}-{}-{}-{}-{}-{}-{}-{}~0:0000}\par
\texttt{F4.1855~0123.123.23.-{}~~~~~~~-{}-{}4-{}-{}~-{}-{}-{}~-{}-{}-{}+-{}-{}-{}-{}~0:0000}\par
\texttt{F4.1856~013.013.023.-{}~~~~~~~-{}-{}-{}-{}-{}~-{}-{}-{}~-{}-{}-{}-{}-{}-{}-{}-{}~0:0000}\par
\texttt{F4.1857~013.123.023.-{}~~~~~~~-{}-{}-{}-{}-{}~-{}-{}-{}~-{}-{}-{}-{}-{}-{}-{}-{}~0:0000}\par
\texttt{F4.1858~0123.012.0.0~~~~~~~~-{}-{}-{}-{}D~-{}PC~-{}-{}-{}-{}++++~ok}\par
\texttt{F4.1859~0123.023.0.0~~~~~~~~-{}-{}-{}-{}D~-{}PC~-{}-{}-{}-{}++++~ok}\par
\texttt{F4.1860~0123.123.0.0~~~~~~~~-{}-{}-{}-{}D~-{}PC~-{}-{}-{}-{}++++~ok}\par
\texttt{F4.1861~012.0123.0.0~~~~~~~~-{}-{}-{}-{}D~-{}PC~-{}-{}-{}-{}++-{}+~ok}\par
\texttt{F4.1862~023.0123.0.0~~~~~~~~-{}-{}-{}-{}D~-{}-{}-{}~-{}-{}-{}-{}-{}-{}-{}-{}~1:1011}\par
\texttt{F4.1863~123.0123.0.0~~~~~~~~-{}-{}-{}-{}D~-{}PC~-{}-{}-{}-{}++++~ok}\par
\texttt{F4.1864~0123.012.1.0~~~~~~~~-{}-{}-{}-{}D~-{}PC~-{}-{}-{}-{}++-{}+~ok}\par
\texttt{F4.1865~0123.013.1.0~~~~~~~~-{}-{}-{}-{}D~-{}PC~-{}-{}-{}-{}++-{}+~ok}\par
\texttt{F4.1866~0123.023.1.0~~~~~~~~-{}-{}-{}-{}D~-{}PC~-{}-{}-{}-{}++++~ok}\par
\texttt{F4.1867~0123.123.1.0~~~~~~~~-{}-{}-{}-{}D~-{}PC~-{}-{}-{}-{}++-{}+~ok}\par
\texttt{F4.1868~0123.01.01.0~~~~~~~~-{}-{}-{}-{}D~-{}PC~-{}-{}-{}-{}++++~ok}\par
\texttt{F4.1869~0123.02.01.0~~~~~~~~-{}B-{}-{}D~-{}PC~+-{}-{}-{}++++~ok}\par
\texttt{F4.1870~0123.12.01.0~~~~~~~~-{}-{}-{}-{}D~-{}PC~-{}-{}-{}-{}++-{}+~ok}\par
\texttt{F4.1871~012.012.01.0~~~~~~~~-{}-{}-{}-{}D~-{}P-{}~-{}-{}-{}-{}-{}+-{}-{}~3:1110}\par
\texttt{F4.1872~013.012.01.0~~~~~~~~-{}-{}-{}-{}D~-{}PC~-{}-{}-{}-{}++-{}+~ok}\par
\texttt{F4.1873~023.012.01.0~~~~~~~~-{}-{}-{}-{}D~-{}PC~-{}-{}-{}-{}++-{}+~ok}\par
\texttt{F4.1874~123.012.01.0~~~~~~~~-{}B-{}-{}D~-{}PC~+-{}-{}-{}++++~ok}\par
\texttt{F4.1875~0123.03.01.0~~~~~~~~-{}-{}-{}-{}D~-{}PC~-{}-{}-{}-{}++++~ok}\par
\texttt{F4.1876~0123.13.01.0~~~~~~~~-{}-{}-{}-{}D~-{}PC~-{}-{}-{}-{}++-{}+~ok}\par
\texttt{F4.1877~012.013.01.0~~~~~~~~-{}-{}-{}-{}D~-{}PC~-{}-{}-{}-{}++-{}+~ok}\par
\texttt{F4.1878~013.013.01.0~~~~~~~~-{}-{}-{}-{}D~-{}P-{}~-{}-{}-{}-{}-{}+-{}-{}~2:1101}\par
\texttt{F4.1879~023.013.01.0~~~~~~~~-{}-{}-{}-{}D~-{}PC~-{}-{}-{}-{}++-{}+~ok}\par
\texttt{F4.1880~123.013.01.0~~~~~~~~-{}-{}-{}-{}D~-{}PC~-{}-{}-{}-{}++++~ok}\par
\texttt{F4.1881~0123.23.01.0~~~~~~~~-{}-{}-{}-{}D~-{}PC~-{}-{}-{}-{}++++~ok}\par
\texttt{F4.1882~012.023.01.0~~~~~~~~-{}-{}-{}-{}D~-{}PC~-{}-{}-{}-{}++++~ok}\par
\texttt{F4.1883~013.023.01.0~~~~~~~~-{}-{}-{}-{}D~-{}PC~-{}-{}-{}-{}++-{}+~ok}\par
\texttt{F4.1884~023.023.01.0~~~~~~~~-{}-{}-{}-{}D~-{}PC~-{}-{}-{}-{}++-{}+~ok}\par
\texttt{F4.1885~123.023.01.0~~~~~~~~-{}-{}-{}-{}D~-{}PC~-{}-{}-{}-{}++++~ok}\par
\texttt{F4.1886~012.123.01.0~~~~~~~~-{}-{}-{}-{}D~-{}PC~-{}-{}-{}-{}++-{}+~ok}\par
\texttt{F4.1887~013.123.01.0~~~~~~~~-{}-{}-{}-{}D~-{}PC~-{}-{}-{}-{}++-{}+~ok}\par
\texttt{F4.1888~023.123.01.0~~~~~~~~-{}-{}-{}-{}D~-{}PC~-{}-{}-{}-{}++-{}+~ok}\par
\texttt{F4.1889~123.123.01.0~~~~~~~~-{}-{}-{}-{}D~-{}PC~-{}-{}-{}-{}++++~ok}\par
\texttt{F4.1890~01.0123.01.0~~~~~~~~-{}-{}-{}-{}D~-{}PC~-{}-{}-{}-{}++-{}+~ok}\par
\texttt{F4.1891~02.0123.01.0~~~~~~~~-{}-{}-{}-{}D~-{}PC~-{}-{}-{}-{}++-{}-{}~ok}\par
\texttt{F4.1892~12.0123.01.0~~~~~~~~-{}-{}-{}-{}D~-{}PC~-{}-{}-{}-{}++-{}+~ok}\par
\texttt{F4.1893~03.0123.01.0~~~~~~~~-{}-{}-{}-{}D~-{}-{}-{}~-{}-{}-{}-{}-{}-{}-{}-{}~1:1001}\par
\texttt{F4.1894~13.0123.01.0~~~~~~~~-{}-{}-{}-{}D~-{}PC~-{}-{}-{}-{}++-{}+~ok}\par
\texttt{F4.1895~23.0123.01.0~~~~~~~~-{}-{}-{}-{}D~-{}PC~-{}-{}-{}-{}++-{}+~ok}\par
\texttt{F4.1896~0123.012.2.0~~~~~~~~-{}-{}-{}-{}D~-{}-{}-{}~-{}-{}-{}-{}-{}-{}-{}-{}~0:0010}\par
\texttt{F4.1897~0123.013.2.0~~~~~~~~-{}-{}-{}-{}D~-{}-{}-{}~-{}-{}-{}-{}-{}-{}-{}-{}~0:0010}\par
\texttt{F4.1898~0123.023.2.0~~~~~~~~-{}-{}-{}-{}D~-{}-{}-{}~-{}-{}-{}-{}-{}-{}-{}-{}~0:0010}\par
\texttt{F4.1899~0123.123.2.0~~~~~~~~-{}-{}-{}-{}D~-{}-{}-{}~-{}-{}-{}-{}-{}-{}-{}-{}~0:0010}\par
\texttt{F4.1900~012.0123.2.0~~~~~~~~-{}-{}-{}-{}D~-{}-{}-{}~-{}-{}-{}-{}-{}-{}-{}-{}~0:0010}\par
\texttt{F4.1901~013.0123.2.0~~~~~~~~-{}-{}-{}-{}D~-{}-{}-{}~-{}-{}-{}-{}-{}-{}-{}-{}~1:0010}\par
\texttt{F4.1902~023.0123.2.0~~~~~~~~-{}-{}-{}-{}D~-{}-{}-{}~-{}-{}-{}-{}-{}-{}-{}-{}~0:0010}\par
\texttt{F4.1903~123.0123.2.0~~~~~~~~-{}-{}-{}-{}D~-{}-{}-{}~-{}-{}-{}-{}-{}-{}-{}-{}~0:0010}\par
\texttt{F4.1904~0123.01.02.0~~~~~~~~-{}B-{}-{}D~-{}PC~+-{}-{}-{}++++~ok}\par
\texttt{F4.1905~0123.12.02.0~~~~~~~~-{}-{}-{}-{}D~-{}PC~-{}-{}-{}-{}++-{}+~ok}\par
\texttt{F4.1906~012.012.02.0~~~~~~~~-{}-{}-{}-{}D~-{}P-{}~-{}-{}-{}-{}-{}+-{}-{}~3:1110}\par
\texttt{F4.1907~013.012.02.0~~~~~~~~-{}-{}-{}-{}D~-{}PC~-{}-{}-{}-{}++-{}+~ok}\par
\texttt{F4.1908~023.012.02.0~~~~~~~~-{}-{}-{}-{}D~-{}-{}-{}~-{}-{}-{}-{}-{}-{}-{}-{}~1:1011}\par
\texttt{F4.1909~123.012.02.0~~~~~~~~-{}-{}-{}-{}D~-{}PC~-{}-{}-{}-{}++++~ok}\par
\texttt{F4.1910~0123.03.02.0~~~~~~~~-{}-{}-{}-{}D~-{}PC~-{}-{}-{}-{}++++~ok}\par
\texttt{F4.1911~0123.13.02.0~~~~~~~~-{}-{}-{}-{}D~-{}PC~-{}-{}-{}-{}++-{}+~ok}\par
\texttt{F4.1912~012.013.02.0~~~~~~~~-{}-{}-{}-{}D~-{}PC~-{}-{}-{}-{}++-{}+~ok}\par
\texttt{F4.1913~013.013.02.0~~~~~~~~-{}-{}-{}-{}D~-{}-{}-{}~-{}-{}-{}-{}-{}-{}-{}-{}~2:1101}\par
\texttt{F4.1914~023.013.02.0~~~~~~~~-{}-{}-{}-{}D~-{}-{}-{}~-{}-{}-{}-{}-{}-{}-{}-{}~1:1011}\par
\texttt{F4.1915~123.013.02.0~~~~~~~~-{}-{}-{}-{}D~-{}PC~-{}-{}-{}-{}++++~ok}\par
\texttt{F4.1916~0123.23.02.0~~~~~~~~-{}-{}-{}-{}D~-{}PC~-{}-{}-{}-{}++-{}+~ok}\par
\texttt{F4.1917~012.023.02.0~~~~~~~~-{}-{}-{}-{}D~-{}PC~-{}-{}-{}-{}++-{}+~ok}\par
\texttt{F4.1918~013.023.02.0~~~~~~~~-{}-{}-{}-{}D~-{}PC~-{}-{}-{}-{}++-{}+~ok}\par
\texttt{F4.1919~023.023.02.0~~~~~~~~-{}-{}-{}-{}D~-{}P-{}~-{}-{}-{}-{}-{}+-{}-{}~1:1011}\par
\texttt{F4.1920~123.023.02.0~~~~~~~~-{}-{}-{}-{}D~-{}PC~-{}-{}-{}-{}++++~ok}\par
\texttt{F4.1921~012.123.02.0~~~~~~~~-{}-{}-{}-{}D~-{}PC~-{}-{}-{}-{}++-{}+~ok}\par
\texttt{F4.1922~013.123.02.0~~~~~~~~-{}-{}-{}-{}D~-{}PC~-{}-{}-{}-{}++-{}+~ok}\par
\texttt{F4.1923~023.123.02.0~~~~~~~~-{}-{}-{}-{}D~-{}-{}-{}~-{}-{}-{}-{}-{}-{}-{}-{}~1:1011}\par
\texttt{F4.1924~123.123.02.0~~~~~~~~-{}-{}-{}-{}D~-{}PC~-{}-{}-{}-{}++++~ok}\par
\texttt{F4.1925~01.0123.02.0~~~~~~~~-{}-{}-{}-{}D~-{}PC~-{}-{}-{}-{}++-{}+~ok}\par
\texttt{F4.1926~02.0123.02.0~~~~~~~~-{}-{}-{}-{}D~-{}-{}-{}~-{}-{}-{}-{}-{}-{}-{}-{}~1:1010}\par
\texttt{F4.1927~12.0123.02.0~~~~~~~~-{}-{}-{}-{}D~-{}PC~-{}-{}-{}-{}++-{}+~ok}\par
\texttt{F4.1928~03.0123.02.0~~~~~~~~-{}-{}-{}-{}D~-{}-{}-{}~-{}-{}-{}-{}-{}-{}-{}-{}~1:1001}\par
\texttt{F4.1929~13.0123.02.0~~~~~~~~-{}-{}-{}-{}D~-{}PC~-{}-{}-{}-{}++-{}+~ok}\par
\texttt{F4.1930~23.0123.02.0~~~~~~~~-{}-{}-{}-{}D~-{}-{}-{}~-{}-{}-{}-{}-{}-{}-{}-{}~1:1011}\par
\texttt{F4.1931~0123.12.12.0~~~~~~~~-{}-{}-{}-{}D~-{}-{}-{}~-{}-{}-{}-{}-{}-{}-{}-{}~0:0110}\par
\texttt{F4.1932~012.012.12.0~~~~~~~~-{}-{}-{}-{}D~-{}P-{}~-{}-{}-{}-{}-{}+-{}-{}~3:1110}\par
\texttt{F4.1933~013.012.12.0~~~~~~~~-{}B-{}-{}D~-{}PC~+-{}-{}-{}++-{}+~ok}\par
\texttt{F4.1934~023.012.12.0~~~~~~~~-{}-{}-{}-{}D~-{}PC~-{}-{}-{}-{}++-{}+~ok}\par
\texttt{F4.1935~123.012.12.0~~~~~~~~-{}-{}-{}-{}D~-{}PC~-{}-{}-{}-{}++-{}+~ok}\par
\texttt{F4.1936~0123.03.12.0~~~~~~~~-{}-{}-{}-{}D~-{}PC~-{}-{}-{}-{}++++~ok}\par
\texttt{F4.1937~0123.13.12.0~~~~~~~~-{}-{}-{}-{}D~-{}PC~-{}-{}-{}-{}++-{}+~ok}\par
\texttt{F4.1938~012.013.12.0~~~~~~~~-{}-{}-{}-{}D~-{}PC~-{}-{}-{}-{}++-{}+~ok}\par
\texttt{F4.1939~013.013.12.0~~~~~~~~-{}-{}-{}-{}D~-{}-{}-{}~-{}-{}-{}-{}-{}-{}-{}-{}~2:1101}\par
\texttt{F4.1940~023.013.12.0~~~~~~~~-{}-{}-{}-{}D~-{}PC~-{}-{}-{}-{}++-{}+~ok}\par
\texttt{F4.1941~123.013.12.0~~~~~~~~-{}-{}-{}-{}D~-{}PC~-{}-{}-{}-{}++++~ok}\par
\texttt{F4.1942~0123.23.12.0~~~~~~~~-{}-{}-{}-{}D~-{}PC~-{}-{}-{}-{}++-{}+~ok}\par
\texttt{F4.1943~012.023.12.0~~~~~~~~-{}-{}-{}-{}D~-{}PC~-{}-{}-{}-{}++++~ok}\par
\texttt{F4.1944~013.023.12.0~~~~~~~~-{}-{}-{}-{}D~-{}PC~-{}-{}-{}-{}++-{}+~ok}\par
\texttt{F4.1945~023.023.12.0~~~~~~~~-{}-{}-{}-{}D~-{}PC~-{}-{}-{}-{}++-{}+~ok}\par
\texttt{F4.1946~123.023.12.0~~~~~~~~-{}-{}-{}-{}D~-{}PC~-{}-{}-{}-{}++++~ok}\par
\texttt{F4.1947~012.123.12.0~~~~~~~~-{}-{}-{}-{}D~-{}PC~-{}-{}-{}-{}++-{}+~ok}\par
\texttt{F4.1948~013.123.12.0~~~~~~~~-{}-{}-{}-{}D~-{}PC~-{}-{}-{}-{}++-{}+~ok}\par
\texttt{F4.1949~023.123.12.0~~~~~~~~-{}-{}-{}-{}D~-{}PC~-{}-{}-{}-{}++-{}+~ok}\par
\texttt{F4.1950~123.123.12.0~~~~~~~~-{}-{}-{}-{}D~-{}PC~-{}-{}-{}-{}++-{}+~ok}\par
\texttt{F4.1951~01.0123.12.0~~~~~~~~-{}-{}-{}-{}D~-{}PC~-{}-{}-{}-{}++-{}+~ok}\par
\texttt{F4.1952~02.0123.12.0~~~~~~~~-{}-{}-{}-{}D~-{}PC~-{}-{}-{}-{}++-{}-{}~ok}\par
\texttt{F4.1953~12.0123.12.0~~~~~~~~-{}-{}-{}-{}D~-{}PC~-{}-{}-{}-{}++-{}+~ok}\par
\texttt{F4.1954~03.0123.12.0~~~~~~~~-{}-{}-{}-{}D~-{}-{}-{}~-{}-{}-{}-{}-{}-{}-{}-{}~1:1001}\par
\texttt{F4.1955~13.0123.12.0~~~~~~~~-{}-{}-{}-{}D~-{}PC~-{}-{}-{}-{}++-{}+~ok}\par
\texttt{F4.1956~23.0123.12.0~~~~~~~~-{}-{}-{}-{}D~-{}PC~-{}-{}-{}-{}++-{}+~ok}\par
\texttt{F4.1957~01.012.012.0~~~~~~~~-{}-{}-{}-{}D~-{}P-{}~-{}-{}-{}-{}-{}+-{}-{}~3:1110}\par
\texttt{F4.1958~12.012.012.0~~~~~~~~-{}-{}-{}-{}D~-{}P-{}~-{}-{}-{}-{}-{}+-{}-{}~3:1110}\par
\texttt{F4.1959~03.012.012.0~~~~~~~~-{}-{}-{}-{}D~-{}-{}-{}~-{}-{}-{}-{}-{}-{}-{}-{}~1:1001}\par
\texttt{F4.1960~13.012.012.0~~~~~~~~-{}-{}-{}-{}D~-{}PC~-{}-{}-{}-{}++-{}+~ok}\par
\texttt{F4.1961~0123.3.012.0~~~~~~~~-{}-{}-{}-{}D~-{}PC~-{}-{}-{}-{}++-{}+~ok}\par
\texttt{F4.1962~012.03.012.0~~~~~~~~-{}-{}-{}-{}D~-{}PC~-{}-{}-{}-{}++++~ok}\par
\texttt{F4.1963~013.03.012.0~~~~~~~~-{}-{}-{}-{}D~-{}-{}-{}~-{}-{}-{}-{}-{}-{}-{}-{}~2:1101}\par
\texttt{F4.1964~023.03.012.0~~~~~~~~-{}-{}-{}-{}D~-{}PC~-{}-{}-{}-{}++-{}+~ok}\par
\texttt{F4.1965~123.03.012.0~~~~~~~~-{}-{}-{}-{}D~-{}PC~-{}-{}-{}-{}++++~ok}\par
\texttt{F4.1966~012.13.012.0~~~~~~~~-{}-{}-{}-{}D~-{}PC~-{}-{}-{}-{}++-{}+~ok}\par
\texttt{F4.1967~013.13.012.0~~~~~~~~-{}-{}-{}-{}D~-{}-{}-{}~-{}-{}-{}-{}-{}-{}-{}-{}~2:1101}\par
\texttt{F4.1968~023.13.012.0~~~~~~~~-{}-{}-{}-{}D~-{}PC~-{}-{}-{}-{}++-{}+~ok}\par
\texttt{F4.1969~123.13.012.0~~~~~~~~-{}-{}-{}-{}D~-{}PC~-{}-{}-{}-{}++-{}+~ok}\par
\texttt{F4.1970~01.013.012.0~~~~~~~~-{}-{}-{}-{}D~-{}-{}-{}~-{}-{}-{}-{}-{}-{}-{}-{}~2:1101}\par
\texttt{F4.1971~02.013.012.0~~~~~~~~-{}-{}-{}-{}D~-{}PC~-{}-{}-{}-{}++-{}-{}~ok}\par
\texttt{F4.1972~12.013.012.0~~~~~~~~-{}-{}-{}-{}D~-{}PC~-{}-{}-{}-{}++-{}+~ok}\par
\texttt{F4.1973~03.013.012.0~~~~~~~~-{}-{}-{}-{}D~-{}-{}-{}~-{}-{}-{}-{}-{}-{}-{}-{}~1:1001}\par
\texttt{F4.1974~13.013.012.0~~~~~~~~-{}-{}-{}-{}D~-{}-{}-{}~-{}-{}-{}-{}-{}-{}-{}-{}~2:1101}\par
\texttt{F4.1975~23.013.012.0~~~~~~~~-{}-{}-{}-{}D~-{}PC~-{}-{}-{}-{}++-{}+~ok}\par
\texttt{F4.1976~012.23.012.0~~~~~~~~-{}-{}-{}-{}D~-{}PC~-{}-{}-{}-{}++++~ok}\par
\texttt{F4.1977~013.23.012.0~~~~~~~~-{}-{}-{}-{}D~-{}PC~-{}-{}-{}-{}++-{}+~ok}\par
\texttt{F4.1978~023.23.012.0~~~~~~~~-{}-{}-{}-{}D~-{}PC~-{}-{}-{}-{}++-{}+~ok}\par
\texttt{F4.1979~123.23.012.0~~~~~~~~-{}-{}-{}-{}D~-{}PC~-{}-{}-{}-{}++++~ok}\par
\texttt{F4.1980~01.023.012.0~~~~~~~~-{}-{}-{}-{}D~-{}PC~-{}-{}-{}-{}++-{}+~ok}\par
\texttt{F4.1981~02.023.012.0~~~~~~~~-{}-{}-{}-{}D~-{}PC~-{}-{}-{}-{}++-{}+~ok}\par
\texttt{F4.1982~12.023.012.0~~~~~~~~-{}-{}-{}-{}D~-{}PC~-{}-{}-{}-{}++++~ok}\par
\texttt{F4.1983~03.023.012.0~~~~~~~~-{}-{}-{}-{}D~-{}-{}-{}~-{}-{}-{}-{}-{}-{}-{}-{}~1:1001}\par
\texttt{F4.1984~13.023.012.0~~~~~~~~-{}-{}-{}-{}D~-{}PC~-{}-{}-{}-{}++-{}+~ok}\par
\texttt{F4.1985~23.023.012.0~~~~~~~~-{}-{}-{}-{}D~-{}PC~-{}-{}-{}-{}++-{}+~ok}\par
\texttt{F4.1986~01.123.012.0~~~~~~~~-{}-{}-{}-{}D~-{}PC~-{}-{}-{}-{}++-{}+~ok}\par
\texttt{F4.1987~02.123.012.0~~~~~~~~-{}-{}-{}-{}D~-{}PC~-{}-{}-{}-{}++-{}-{}~ok}\par
\texttt{F4.1988~12.123.012.0~~~~~~~~-{}-{}-{}-{}D~-{}PC~-{}-{}-{}-{}++-{}+~ok}\par
\texttt{F4.1989~03.123.012.0~~~~~~~~-{}-{}-{}-{}D~-{}-{}-{}~-{}-{}-{}-{}-{}-{}-{}-{}~1:1001}\par
\texttt{F4.1990~13.123.012.0~~~~~~~~-{}-{}-{}-{}D~-{}PC~-{}-{}-{}-{}++-{}+~ok}\par
\texttt{F4.1991~23.123.012.0~~~~~~~~-{}-{}-{}-{}D~-{}PC~-{}-{}-{}-{}++-{}+~ok}\par
\texttt{F4.1992~0.0123.012.0~~~~~~~~-{}-{}-{}-{}D~-{}-{}-{}~-{}-{}-{}-{}-{}-{}-{}-{}~1:1000}\par
\texttt{F4.1993~1.0123.012.0~~~~~~~~-{}-{}-{}-{}D~-{}PC~-{}-{}-{}-{}++-{}+~ok}\par
\texttt{F4.1994~2.0123.012.0~~~~~~~~-{}-{}-{}-{}D~-{}PC~-{}-{}-{}-{}++-{}-{}~ok}\par
\texttt{F4.1995~3.0123.012.0~~~~~~~~-{}-{}-{}-{}D~-{}-{}-{}~-{}-{}-{}-{}-{}-{}-{}-{}~1:1001}\par
\texttt{F4.1996~0123.023.3.0~~~~~~~~-{}-{}-{}-{}D~-{}PC~-{}-{}-{}-{}++-{}+~ok}\par
\texttt{F4.1997~0123.123.3.0~~~~~~~~-{}-{}-{}-{}D~-{}PC~-{}-{}-{}-{}++-{}+~ok}\par
\texttt{F4.1998~012.0123.3.0~~~~~~~~-{}-{}-{}-{}D~-{}PC~-{}-{}-{}-{}++-{}+~ok}\par
\texttt{F4.1999~023.0123.3.0~~~~~~~~-{}-{}-{}-{}D~-{}-{}-{}~-{}-{}-{}-{}-{}-{}-{}-{}~1:1011}\par
\texttt{F4.2000~123.0123.3.0~~~~~~~~-{}-{}-{}-{}D~-{}PC~-{}-{}-{}-{}++-{}+~ok}\par
\texttt{F4.2001~0123.03.03.0~~~~~~~~-{}-{}-{}-{}D~-{}PC~-{}-{}-{}-{}++++~ok}\par
\texttt{F4.2002~0123.13.03.0~~~~~~~~-{}-{}-{}-{}D~-{}PC~-{}-{}-{}-{}++-{}+~ok}\par
\texttt{F4.2003~012.013.03.0~~~~~~~~-{}-{}-{}-{}D~-{}PC~-{}-{}-{}-{}++-{}+~ok}\par
\texttt{F4.2004~013.013.03.0~~~~~~~~-{}-{}-{}-{}D~-{}P-{}~-{}-{}-{}-{}-{}+-{}-{}~2:1101}\par
\texttt{F4.2005~023.013.03.0~~~~~~~~-{}-{}-{}-{}D~-{}-{}-{}~-{}-{}-{}-{}-{}-{}-{}-{}~1:1011}\par
\texttt{F4.2006~123.013.03.0~~~~~~~~-{}-{}-{}-{}D~-{}PC~-{}-{}-{}-{}++++~ok}\par
\texttt{F4.2007~0123.23.03.0~~~~~~~~-{}-{}-{}-{}D~-{}PC~-{}-{}-{}-{}++-{}+~ok}\par
\texttt{F4.2008~012.023.03.0~~~~~~~~-{}-{}-{}-{}D~-{}PC~-{}-{}-{}-{}++-{}+~ok}\par
\texttt{F4.2009~013.023.03.0~~~~~~~~-{}-{}-{}-{}D~-{}PC~-{}-{}-{}-{}++-{}+~ok}\par
\texttt{F4.2010~023.023.03.0~~~~~~~~-{}-{}-{}-{}D~-{}P-{}~-{}-{}-{}-{}-{}+-{}-{}~1:1011}\par
\texttt{F4.2011~123.023.03.0~~~~~~~~-{}-{}-{}-{}D~-{}PC~-{}-{}-{}-{}++++~ok}\par
\texttt{F4.2012~012.123.03.0~~~~~~~~-{}-{}-{}-{}D~-{}PC~-{}-{}-{}-{}++-{}+~ok}\par
\texttt{F4.2013~013.123.03.0~~~~~~~~-{}-{}-{}-{}D~-{}PC~-{}-{}-{}-{}++-{}+~ok}\par
\texttt{F4.2014~023.123.03.0~~~~~~~~-{}-{}-{}-{}D~-{}-{}-{}~-{}-{}-{}-{}-{}-{}-{}-{}~1:1011}\par
\texttt{F4.2015~123.123.03.0~~~~~~~~-{}-{}-{}-{}D~-{}PC~-{}-{}-{}-{}++++~ok}\par
\texttt{F4.2016~01.0123.03.0~~~~~~~~-{}-{}-{}-{}D~-{}PC~-{}-{}-{}-{}++-{}+~ok}\par
\texttt{F4.2017~02.0123.03.0~~~~~~~~-{}-{}-{}-{}D~-{}-{}-{}~-{}-{}-{}-{}-{}-{}-{}-{}~1:1011}\par
\texttt{F4.2018~12.0123.03.0~~~~~~~~-{}-{}-{}-{}D~-{}PC~-{}-{}-{}-{}++-{}+~ok}\par
\texttt{F4.2019~03.0123.03.0~~~~~~~~-{}-{}-{}-{}D~-{}-{}-{}~-{}-{}-{}-{}-{}-{}-{}-{}~1:1001}\par
\texttt{F4.2020~13.0123.03.0~~~~~~~~-{}-{}-{}-{}D~-{}PC~-{}-{}-{}-{}++-{}+~ok}\par
\texttt{F4.2021~23.0123.03.0~~~~~~~~-{}-{}-{}-{}D~-{}-{}-{}~-{}-{}-{}-{}-{}-{}-{}-{}~1:1011}\par
\texttt{F4.2022~0123.13.13.0~~~~~~~~-{}-{}-{}-{}D~-{}PC~-{}-{}-{}-{}++-{}+~ok}\par
\texttt{F4.2023~012.013.13.0~~~~~~~~-{}-{}-{}-{}D~-{}PC~-{}-{}-{}-{}++-{}+~ok}\par
\texttt{F4.2024~013.013.13.0~~~~~~~~-{}-{}-{}-{}D~-{}P-{}~-{}-{}-{}-{}-{}+-{}-{}~2:1101}\par
\texttt{F4.2025~023.013.13.0~~~~~~~~-{}-{}-{}-{}D~-{}PC~-{}-{}-{}-{}++-{}+~ok}\par
\texttt{F4.2026~123.013.13.0~~~~~~~~-{}-{}-{}-{}D~-{}PC~-{}-{}-{}-{}++++~ok}\par
\texttt{F4.2027~0123.23.13.0~~~~~~~~-{}-{}-{}-{}D~-{}PC~-{}-{}-{}-{}++++~ok}\par
\texttt{F4.2028~012.023.13.0~~~~~~~~-{}-{}-{}-{}D~-{}PC~-{}-{}-{}-{}++-{}+~ok}\par
\texttt{F4.2029~013.023.13.0~~~~~~~~-{}-{}-{}-{}D~-{}PC~-{}-{}-{}-{}++-{}+~ok}\par
\texttt{F4.2030~023.023.13.0~~~~~~~~-{}-{}-{}-{}D~-{}PC~-{}-{}-{}-{}++-{}+~ok}\par
\texttt{F4.2031~123.023.13.0~~~~~~~~-{}-{}-{}-{}D~-{}PC~-{}-{}-{}-{}++++~ok}\par
\texttt{F4.2032~012.123.13.0~~~~~~~~-{}-{}-{}-{}D~-{}PC~-{}-{}-{}-{}++-{}+~ok}\par
\texttt{F4.2033~013.123.13.0~~~~~~~~-{}-{}-{}-{}D~-{}PC~-{}-{}-{}-{}++-{}+~ok}\par
\texttt{F4.2034~023.123.13.0~~~~~~~~-{}-{}-{}-{}D~-{}PC~-{}-{}-{}-{}++-{}+~ok}\par
\texttt{F4.2035~123.123.13.0~~~~~~~~-{}-{}-{}-{}D~-{}PC~-{}-{}-{}-{}++++~ok}\par
\texttt{F4.2036~01.0123.13.0~~~~~~~~-{}-{}-{}-{}D~-{}PC~-{}-{}-{}-{}++-{}+~ok}\par
\texttt{F4.2037~02.0123.13.0~~~~~~~~-{}-{}-{}-{}D~-{}PC~-{}-{}-{}-{}++-{}+~ok}\par
\texttt{F4.2038~12.0123.13.0~~~~~~~~-{}-{}-{}-{}D~-{}PC~-{}-{}-{}-{}++-{}+~ok}\par
\texttt{F4.2039~03.0123.13.0~~~~~~~~-{}-{}-{}-{}D~-{}-{}-{}~-{}-{}-{}-{}-{}-{}-{}-{}~1:1001}\par
\texttt{F4.2040~13.0123.13.0~~~~~~~~-{}-{}-{}-{}D~-{}PC~-{}-{}-{}-{}++-{}+~ok}\par
\texttt{F4.2041~23.0123.13.0~~~~~~~~-{}-{}-{}-{}D~-{}PC~-{}-{}-{}-{}++-{}+~ok}\par
\texttt{F4.2042~012.13.013.0~~~~~~~~-{}-{}-{}-{}D~-{}PC~-{}-{}-{}-{}++-{}+~ok}\par
\texttt{F4.2043~013.13.013.0~~~~~~~~-{}-{}-{}-{}D~-{}P-{}~-{}-{}-{}-{}-{}+-{}-{}~2:1101}\par
\texttt{F4.2044~023.13.013.0~~~~~~~~-{}-{}-{}-{}D~-{}PC~-{}-{}-{}-{}++-{}+~ok}\par
\texttt{F4.2045~123.13.013.0~~~~~~~~-{}-{}-{}-{}D~-{}PC~-{}-{}-{}-{}++-{}+~ok}\par
\texttt{F4.2046~01.013.013.0~~~~~~~~-{}-{}-{}-{}D~-{}P-{}~-{}-{}-{}-{}-{}+-{}-{}~2:1101}\par
\texttt{F4.2047~02.013.013.0~~~~~~~~-{}-{}-{}-{}D~-{}PC~-{}-{}-{}-{}++-{}+~ok}\par
\texttt{F4.2048~12.013.013.0~~~~~~~~-{}-{}-{}-{}D~-{}PC~-{}-{}-{}-{}++-{}+~ok}\par
\texttt{F4.2049~03.013.013.0~~~~~~~~-{}-{}-{}-{}D~-{}-{}-{}~-{}-{}-{}-{}-{}-{}-{}-{}~1:1001}\par
\texttt{F4.2050~13.013.013.0~~~~~~~~-{}-{}-{}-{}D~-{}P-{}~-{}-{}-{}-{}-{}+-{}-{}~2:1101}\par
\texttt{F4.2051~23.013.013.0~~~~~~~~-{}-{}-{}-{}D~-{}PC~-{}-{}-{}-{}++-{}+~ok}\par
\texttt{F4.2052~01.023.013.0~~~~~~~~-{}-{}-{}-{}D~-{}PC~-{}-{}-{}-{}++-{}+~ok}\par
\texttt{F4.2053~12.023.013.0~~~~~~~~-{}-{}-{}-{}D~-{}PC~-{}-{}-{}-{}++-{}+~ok}\par
\texttt{F4.2054~03.023.013.0~~~~~~~~-{}-{}-{}-{}D~-{}-{}-{}~-{}-{}-{}-{}-{}-{}-{}-{}~1:1001}\par
\texttt{F4.2055~13.023.013.0~~~~~~~~-{}-{}-{}-{}D~-{}PC~-{}-{}-{}-{}++-{}+~ok}\par
\texttt{F4.2056~01.123.013.0~~~~~~~~-{}-{}-{}-{}D~-{}PC~-{}-{}-{}-{}++-{}+~ok}\par
\texttt{F4.2057~02.123.013.0~~~~~~~~-{}-{}-{}-{}D~-{}PC~-{}-{}-{}-{}++-{}+~ok}\par
\texttt{F4.2058~12.123.013.0~~~~~~~~-{}-{}-{}-{}D~-{}PC~-{}-{}-{}-{}++-{}+~ok}\par
\texttt{F4.2059~03.123.013.0~~~~~~~~-{}-{}-{}-{}D~-{}-{}-{}~-{}-{}-{}-{}-{}-{}-{}-{}~1:1001}\par
\texttt{F4.2060~13.123.013.0~~~~~~~~-{}-{}-{}-{}D~-{}PC~-{}-{}-{}-{}++-{}+~ok}\par
\texttt{F4.2061~23.123.013.0~~~~~~~~-{}-{}-{}-{}D~-{}PC~-{}-{}-{}-{}++-{}+~ok}\par
\texttt{F4.2062~0.0123.013.0~~~~~~~~-{}-{}-{}-{}D~-{}-{}-{}~-{}-{}-{}-{}-{}-{}-{}-{}~1:1000}\par
\texttt{F4.2063~3.0123.013.0~~~~~~~~-{}-{}-{}-{}D~-{}-{}-{}~-{}-{}-{}-{}-{}-{}-{}-{}~1:1001}\par
\texttt{F4.2064~0123.13.23.0~~~~~~~~-{}-{}-{}-{}D~-{}PC~-{}-{}-{}-{}++-{}+~ok}\par
\texttt{F4.2065~012.013.23.0~~~~~~~~-{}-{}-{}-{}D~-{}PC~-{}-{}-{}-{}++-{}+~ok}\par
\texttt{F4.2066~013.013.23.0~~~~~~~~-{}-{}-{}-{}D~-{}-{}-{}~-{}-{}-{}-{}-{}-{}-{}-{}~2:1101}\par
\texttt{F4.2067~023.013.23.0~~~~~~~~-{}-{}-{}-{}D~-{}-{}-{}~-{}-{}-{}-{}-{}-{}-{}-{}~1:1011}\par
\texttt{F4.2068~123.013.23.0~~~~~~~~-{}-{}-{}-{}D~-{}PC~-{}-{}-{}-{}++-{}+~ok}\par
\texttt{F4.2069~012.123.23.0~~~~~~~~-{}-{}-{}-{}D~-{}PC~-{}-{}-{}-{}++-{}+~ok}\par
\texttt{F4.2070~013.123.23.0~~~~~~~~-{}-{}-{}-{}D~-{}PC~-{}-{}-{}-{}++-{}+~ok}\par
\texttt{F4.2071~023.123.23.0~~~~~~~~-{}-{}-{}-{}D~-{}-{}-{}~-{}-{}-{}-{}-{}-{}-{}-{}~1:1011}\par
\texttt{F4.2072~123.123.23.0~~~~~~~~-{}-{}-{}-{}D~-{}PC~-{}-{}-{}-{}++-{}+~ok}\par
\texttt{F4.2073~01.0123.23.0~~~~~~~~-{}-{}-{}-{}D~-{}PC~-{}-{}-{}-{}++-{}+~ok}\par
\texttt{F4.2074~02.0123.23.0~~~~~~~~-{}-{}-{}-{}D~-{}-{}-{}~-{}-{}-{}-{}-{}-{}-{}-{}~1:1011}\par
\texttt{F4.2075~12.0123.23.0~~~~~~~~-{}-{}-{}-{}D~-{}PC~-{}-{}-{}-{}++-{}+~ok}\par
\texttt{F4.2076~03.0123.23.0~~~~~~~~-{}-{}-{}-{}D~-{}-{}-{}~-{}-{}-{}-{}-{}-{}-{}-{}~1:1001}\par
\texttt{F4.2077~13.0123.23.0~~~~~~~~-{}-{}-{}-{}D~-{}PC~-{}-{}-{}-{}++-{}+~ok}\par
\texttt{F4.2078~23.0123.23.0~~~~~~~~-{}-{}-{}-{}D~-{}-{}-{}~-{}-{}-{}-{}-{}-{}-{}-{}~1:1011}\par
\texttt{F4.2079~01.013.023.0~~~~~~~~-{}-{}-{}-{}D~-{}-{}-{}~-{}-{}-{}-{}-{}-{}-{}-{}~2:1101}\par
\texttt{F4.2080~12.013.023.0~~~~~~~~-{}-{}-{}-{}D~-{}PC~-{}-{}-{}-{}++-{}+~ok}\par
\texttt{F4.2081~03.013.023.0~~~~~~~~-{}-{}-{}-{}D~-{}-{}-{}~-{}-{}-{}-{}-{}-{}-{}-{}~1:1001}\par
\texttt{F4.2082~13.013.023.0~~~~~~~~-{}-{}-{}-{}D~-{}-{}-{}~-{}-{}-{}-{}-{}-{}-{}-{}~2:1101}\par
\texttt{F4.2083~01.123.023.0~~~~~~~~-{}-{}-{}-{}D~-{}PC~-{}-{}-{}-{}++-{}+~ok}\par
\texttt{F4.2084~02.123.023.0~~~~~~~~-{}-{}-{}-{}D~-{}-{}-{}~-{}-{}-{}-{}-{}-{}-{}-{}~1:1011}\par
\texttt{F4.2085~12.123.023.0~~~~~~~~-{}-{}-{}-{}D~-{}PC~-{}-{}-{}-{}++-{}+~ok}\par
\texttt{F4.2086~03.123.023.0~~~~~~~~-{}-{}-{}-{}D~-{}-{}-{}~-{}-{}-{}-{}-{}-{}-{}-{}~1:1001}\par
\texttt{F4.2087~13.123.023.0~~~~~~~~-{}-{}-{}-{}D~-{}PC~-{}-{}-{}-{}++-{}+~ok}\par
\texttt{F4.2088~23.123.023.0~~~~~~~~-{}-{}-{}-{}D~-{}-{}-{}~-{}-{}-{}-{}-{}-{}-{}-{}~1:1011}\par
\texttt{F4.2089~0.0123.023.0~~~~~~~~-{}-{}4-{}D~-{}-{}-{}~-{}-{}-{}+-{}-{}-{}-{}~1:1000}\par
\texttt{F4.2090~3.0123.023.0~~~~~~~~-{}-{}-{}-{}D~-{}-{}-{}~-{}-{}-{}-{}-{}-{}-{}-{}~1:1001}\par
\texttt{F4.2091~01.123.123.0~~~~~~~~-{}-{}-{}-{}D~-{}PC~-{}-{}-{}-{}++-{}+~ok}\par
\texttt{F4.2092~12.123.123.0~~~~~~~~-{}-{}-{}-{}D~-{}PC~-{}-{}-{}-{}++-{}+~ok}\par
\texttt{F4.2093~03.123.123.0~~~~~~~~-{}-{}-{}-{}D~-{}-{}-{}~-{}-{}-{}-{}-{}-{}-{}-{}~1:1001}\par
\texttt{F4.2094~13.123.123.0~~~~~~~~-{}-{}-{}-{}D~-{}PC~-{}-{}-{}-{}++-{}+~ok}\par
\texttt{F4.2095~0.0123.123.0~~~~~~~~-{}-{}-{}-{}D~-{}-{}-{}~-{}-{}-{}-{}-{}-{}-{}-{}~1:1000}\par
\texttt{F4.2096~0123.1.01.01~~~~~~~~-{}-{}-{}-{}D~-{}-{}-{}~-{}-{}-{}-{}-{}-{}-{}-{}~0:0100}\par
\texttt{F4.2097~012.01.01.01~~~~~~~~-{}-{}-{}-{}D~-{}P-{}~-{}-{}-{}-{}-{}+-{}-{}~3:1110}\par
\texttt{F4.2098~023.01.01.01~~~~~~~~-{}-{}-{}-{}D~-{}PC~-{}-{}-{}-{}++++~ok}\par
\texttt{F4.2099~123.01.01.01~~~~~~~~-{}-{}-{}-{}D~-{}PC~-{}-{}-{}-{}++++~ok}\par
\texttt{F4.2100~012.02.01.01~~~~~~~~-{}-{}-{}-{}D~-{}P-{}~-{}-{}-{}-{}-{}+-{}-{}~3:1110}\par
\texttt{F4.2101~013.02.01.01~~~~~~~~-{}-{}-{}-{}D~-{}PC~-{}-{}-{}-{}++++~ok}\par
\texttt{F4.2102~023.02.01.01~~~~~~~~-{}-{}-{}-{}D~-{}PC~-{}-{}-{}-{}++++~ok}\par
\texttt{F4.2103~123.02.01.01~~~~~~~~-{}-{}-{}-{}D~-{}PC~-{}-{}-{}-{}++++~ok}\par
\texttt{F4.2104~012.12.01.01~~~~~~~~-{}-{}-{}-{}D~-{}P-{}~-{}-{}-{}-{}-{}+-{}-{}~3:1110}\par
\texttt{F4.2105~013.12.01.01~~~~~~~~-{}-{}-{}-{}D~-{}PC~-{}-{}-{}-{}++-{}+~ok}\par
\texttt{F4.2106~023.12.01.01~~~~~~~~-{}-{}-{}-{}D~-{}PC~-{}-{}-{}-{}++-{}+~ok}\par
\texttt{F4.2107~123.12.01.01~~~~~~~~-{}-{}-{}-{}D~-{}PC~-{}-{}-{}-{}++-{}+~ok}\par
\texttt{F4.2108~23.012.01.01~~~~~~~~-{}-{}-{}-{}D~-{}PC~-{}-{}-{}-{}++++~ok}\par
\texttt{F4.2109~023.23.01.01~~~~~~~~-{}B-{}-{}D~-{}PC~+-{}-{}-{}++++~ok}\par
\texttt{F4.2110~123.23.01.01~~~~~~~~-{}-{}-{}-{}D~-{}PC~-{}-{}-{}-{}++++~ok}\par
\texttt{F4.2111~0123.01.2.01~~~~~~~~-{}-{}-{}-{}D~-{}-{}-{}~-{}-{}-{}-{}-{}-{}-{}-{}~0:0010}\par
\texttt{F4.2112~0123.02.2.01~~~~~~~~-{}-{}-{}-{}D~-{}-{}-{}~-{}-{}-{}-{}-{}-{}-{}-{}~0:0010}\par
\texttt{F4.2113~0123.12.2.01~~~~~~~~-{}-{}-{}-{}D~-{}-{}-{}~-{}-{}-{}-{}-{}-{}-{}-{}~0:0010}\par
\texttt{F4.2114~012.012.2.01~~~~~~~~-{}-{}-{}-{}D~-{}-{}-{}~-{}-{}-{}-{}-{}-{}-{}-{}~0:0010}\par
\texttt{F4.2115~013.012.2.01~~~~~~~~-{}-{}-{}-{}D~-{}-{}-{}~-{}-{}-{}-{}-{}-{}-{}-{}~1:0010}\par
\texttt{F4.2116~023.012.2.01~~~~~~~~-{}-{}-{}-{}D~-{}-{}-{}~-{}-{}-{}-{}-{}-{}-{}-{}~0:0010}\par
\texttt{F4.2117~123.012.2.01~~~~~~~~-{}-{}-{}-{}D~-{}-{}-{}~-{}-{}-{}-{}-{}-{}-{}-{}~0:0010}\par
\texttt{F4.2118~0123.03.2.01~~~~~~~~-{}-{}-{}-{}D~-{}-{}-{}~-{}-{}-{}-{}-{}-{}-{}-{}~0:0010}\par
\texttt{F4.2119~0123.13.2.01~~~~~~~~-{}-{}-{}-{}D~-{}-{}-{}~-{}-{}-{}-{}-{}-{}-{}-{}~0:0010}\par
\texttt{F4.2120~013.013.2.01~~~~~~~~-{}B-{}-{}D~-{}PC~+-{}-{}-{}++++~ok}\par
\texttt{F4.2121~023.013.2.01~~~~~~~~-{}-{}-{}-{}D~-{}-{}-{}~-{}-{}-{}-{}-{}-{}-{}-{}~0:0010}\par
\texttt{F4.2122~123.013.2.01~~~~~~~~-{}-{}-{}-{}D~-{}-{}-{}~-{}-{}-{}-{}-{}-{}-{}-{}~0:0010}\par
\texttt{F4.2123~0123.23.2.01~~~~~~~~-{}-{}-{}-{}D~-{}-{}-{}~-{}-{}-{}-{}-{}-{}-{}-{}~0:0010}\par
\texttt{F4.2124~023.023.2.01~~~~~~~~-{}-{}-{}-{}D~-{}-{}-{}~-{}-{}-{}-{}-{}-{}-{}-{}~0:0010}\par
\texttt{F4.2125~123.023.2.01~~~~~~~~-{}-{}-{}-{}D~-{}-{}-{}~-{}-{}-{}-{}-{}-{}-{}-{}~0:0010}\par
\texttt{F4.2126~023.123.2.01~~~~~~~~-{}-{}-{}-{}D~-{}-{}-{}~-{}-{}-{}-{}-{}-{}-{}-{}~0:0010}\par
\texttt{F4.2127~0123.1.02.01~~~~~~~~-{}-{}-{}-{}D~-{}-{}-{}~-{}-{}-{}-{}-{}-{}-{}-{}~0:0100}\par
\texttt{F4.2128~012.01.02.01~~~~~~~~-{}-{}-{}-{}D~-{}P-{}~-{}-{}-{}-{}-{}+-{}-{}~3:1110}\par
\texttt{F4.2129~013.01.02.01~~~~~~~~-{}-{}-{}-{}D~-{}-{}-{}~-{}-{}-{}-{}-{}-{}-{}-{}~2:1101}\par
\texttt{F4.2130~023.01.02.01~~~~~~~~-{}-{}-{}-{}D~-{}PC~-{}-{}-{}-{}++-{}+~ok}\par
\texttt{F4.2131~123.01.02.01~~~~~~~~-{}-{}-{}-{}D~-{}PC~-{}-{}-{}-{}++++~ok}\par
\texttt{F4.2132~012.02.02.01~~~~~~~~-{}-{}-{}-{}D~-{}P-{}~-{}-{}-{}-{}-{}+-{}-{}~3:1110}\par
\texttt{F4.2133~013.02.02.01~~~~~~~~-{}-{}-{}-{}D~-{}PC~-{}-{}-{}-{}++++~ok}\par
\texttt{F4.2134~023.02.02.01~~~~~~~~-{}-{}-{}-{}D~-{}PC~-{}-{}-{}-{}++-{}+~ok}\par
\texttt{F4.2135~123.02.02.01~~~~~~~~-{}-{}-{}-{}D~-{}PC~-{}-{}-{}-{}++++~ok}\par
\texttt{F4.2136~012.12.02.01~~~~~~~~-{}-{}-{}-{}D~-{}P-{}~-{}-{}-{}-{}-{}+-{}-{}~3:1110}\par
\texttt{F4.2137~013.12.02.01~~~~~~~~-{}-{}-{}-{}D~-{}PC~-{}-{}-{}-{}++-{}+~ok}\par
\texttt{F4.2138~023.12.02.01~~~~~~~~-{}-{}-{}-{}D~-{}PC~-{}-{}-{}-{}++-{}+~ok}\par
\texttt{F4.2139~123.12.02.01~~~~~~~~-{}-{}-{}-{}D~-{}PC~-{}-{}-{}-{}++-{}+~ok}\par
\texttt{F4.2140~01.012.02.01~~~~~~~~-{}-{}-{}-{}D~-{}P-{}~-{}-{}-{}-{}-{}+-{}-{}~3:1110}\par
\texttt{F4.2141~02.012.02.01~~~~~~~~-{}-{}-{}-{}D~-{}-{}-{}~-{}-{}-{}-{}-{}-{}-{}-{}~1:1010}\par
\texttt{F4.2142~12.012.02.01~~~~~~~~-{}-{}-{}-{}D~-{}P-{}~-{}-{}-{}-{}-{}+-{}-{}~3:1110}\par
\texttt{F4.2143~03.012.02.01~~~~~~~~-{}-{}-{}-{}D~-{}PC~-{}-{}-{}-{}++-{}+~ok}\par
\texttt{F4.2144~13.012.02.01~~~~~~~~-{}-{}-{}-{}D~-{}PC~-{}-{}-{}-{}++++~ok}\par
\texttt{F4.2145~23.012.02.01~~~~~~~~-{}-{}-{}-{}D~-{}PC~-{}-{}-{}-{}++-{}+~ok}\par
\texttt{F4.2146~012.03.02.01~~~~~~~~-{}-{}-{}-{}D~-{}PC~-{}-{}-{}-{}++-{}+~ok}\par
\texttt{F4.2147~013.03.02.01~~~~~~~~-{}-{}-{}-{}D~-{}-{}-{}~-{}-{}-{}-{}-{}-{}-{}-{}~2:1101}\par
\texttt{F4.2148~123.03.02.01~~~~~~~~-{}B-{}-{}D~-{}PC~+-{}-{}-{}++++~ok}\par
\texttt{F4.2149~012.13.02.01~~~~~~~~-{}-{}-{}-{}D~-{}PC~-{}-{}-{}-{}++-{}+~ok}\par
\texttt{F4.2150~013.13.02.01~~~~~~~~-{}-{}-{}-{}D~-{}-{}-{}~-{}-{}-{}-{}-{}-{}-{}-{}~2:1101}\par
\texttt{F4.2151~023.13.02.01~~~~~~~~-{}B-{}-{}D~-{}PC~+-{}-{}-{}++-{}+~ok}\par
\texttt{F4.2152~123.13.02.01~~~~~~~~-{}-{}-{}-{}D~-{}PC~-{}-{}-{}-{}++-{}+~ok}\par
\texttt{F4.2153~01.013.02.01~~~~~~~~-{}-{}-{}-{}D~-{}-{}-{}~-{}-{}-{}-{}-{}-{}-{}-{}~2:1101}\par
\texttt{F4.2154~02.013.02.01~~~~~~~~-{}-{}-{}-{}D~-{}-{}-{}~-{}-{}-{}-{}-{}-{}-{}-{}~1:1010}\par
\texttt{F4.2155~12.013.02.01~~~~~~~~-{}-{}-{}-{}D~-{}PC~-{}-{}-{}-{}++-{}+~ok}\par
\texttt{F4.2156~03.013.02.01~~~~~~~~-{}-{}-{}-{}D~-{}-{}-{}~-{}-{}-{}-{}-{}-{}-{}-{}~2:1101}\par
\texttt{F4.2157~13.013.02.01~~~~~~~~-{}-{}-{}-{}D~-{}-{}-{}~-{}-{}-{}-{}-{}-{}-{}-{}~2:1101}\par
\texttt{F4.2158~23.013.02.01~~~~~~~~-{}-{}-{}-{}D~-{}PC~-{}-{}-{}-{}++-{}+~ok}\par
\texttt{F4.2159~012.23.02.01~~~~~~~~-{}-{}-{}-{}D~-{}PC~-{}-{}-{}-{}++-{}+~ok}\par
\texttt{F4.2160~013.23.02.01~~~~~~~~-{}-{}-{}-{}D~-{}PC~-{}-{}-{}-{}++++~ok}\par
\texttt{F4.2161~023.23.02.01~~~~~~~~-{}-{}-{}-{}D~-{}PC~-{}-{}-{}-{}++-{}+~ok}\par
\texttt{F4.2162~123.23.02.01~~~~~~~~-{}-{}-{}-{}D~-{}PC~-{}-{}-{}-{}++++~ok}\par
\texttt{F4.2163~01.023.02.01~~~~~~~~-{}-{}-{}-{}D~-{}PC~-{}-{}-{}-{}++++~ok}\par
\texttt{F4.2164~02.023.02.01~~~~~~~~-{}-{}-{}-{}D~-{}-{}-{}~-{}-{}-{}-{}-{}-{}-{}-{}~1:1010}\par
\texttt{F4.2165~12.023.02.01~~~~~~~~-{}-{}-{}-{}D~-{}PC~-{}-{}-{}-{}++-{}+~ok}\par
\texttt{F4.2166~03.023.02.01~~~~~~~~-{}-{}-{}-{}D~-{}PC~-{}-{}-{}-{}++-{}+~ok}\par
\texttt{F4.2167~13.023.02.01~~~~~~~~-{}-{}-{}-{}D~-{}PC~-{}-{}-{}-{}++++~ok}\par
\texttt{F4.2168~23.023.02.01~~~~~~~~-{}-{}-{}-{}D~-{}PC~-{}-{}-{}-{}++-{}+~ok}\par
\texttt{F4.2169~01.123.02.01~~~~~~~~-{}-{}-{}-{}D~-{}PC~-{}-{}-{}-{}++-{}+~ok}\par
\texttt{F4.2170~02.123.02.01~~~~~~~~-{}-{}-{}-{}D~-{}-{}-{}~-{}-{}-{}-{}-{}-{}-{}-{}~1:1010}\par
\texttt{F4.2171~12.123.02.01~~~~~~~~-{}-{}-{}-{}D~-{}PC~-{}-{}-{}-{}++-{}+~ok}\par
\texttt{F4.2172~03.123.02.01~~~~~~~~-{}-{}-{}-{}D~-{}PC~-{}-{}-{}-{}++-{}+~ok}\par
\texttt{F4.2173~13.123.02.01~~~~~~~~-{}-{}-{}-{}D~-{}PC~-{}-{}-{}-{}++++~ok}\par
\texttt{F4.2174~23.123.02.01~~~~~~~~-{}-{}-{}-{}D~-{}PC~-{}-{}-{}-{}++-{}+~ok}\par
\texttt{F4.2175~0.0123.02.01~~~~~~~~-{}-{}-{}-{}D~-{}-{}-{}~-{}-{}-{}-{}-{}-{}-{}-{}~1:1000}\par
\texttt{F4.2176~012.1.012.01~~~~~~~~-{}-{}-{}-{}D~-{}-{}-{}~-{}-{}-{}-{}-{}-{}-{}-{}~0:0100}\par
\texttt{F4.2177~013.1.012.01~~~~~~~~-{}-{}-{}-{}D~-{}-{}-{}~-{}-{}-{}-{}-{}-{}-{}-{}~0:0100}\par
\texttt{F4.2178~023.1.012.01~~~~~~~~-{}-{}-{}-{}D~-{}-{}-{}~-{}-{}-{}-{}-{}-{}-{}-{}~2:0100}\par
\texttt{F4.2179~123.1.012.01~~~~~~~~-{}-{}-{}-{}D~-{}-{}-{}~-{}-{}-{}-{}-{}-{}-{}-{}~0:0100}\par
\texttt{F4.2180~01.01.012.01~~~~~~~~-{}-{}4-{}D~-{}-{}-{}~-{}-{}-{}+-{}-{}-{}-{}~2:1100}\par
\texttt{F4.2181~02.01.012.01~~~~~~~~-{}-{}-{}-{}D~-{}P-{}~-{}-{}-{}-{}-{}+-{}-{}~3:1110}\par
\texttt{F4.2182~12.01.012.01~~~~~~~~-{}-{}-{}-{}D~-{}P-{}~-{}-{}-{}-{}-{}+-{}-{}~3:1110}\par
\texttt{F4.2183~03.01.012.01~~~~~~~~-{}-{}-{}-{}D~-{}-{}-{}~-{}-{}-{}-{}-{}-{}-{}-{}~2:1101}\par
\texttt{F4.2184~13.01.012.01~~~~~~~~-{}-{}-{}-{}D~-{}-{}-{}~-{}-{}-{}-{}-{}-{}-{}-{}~2:1101}\par
\texttt{F4.2185~23.01.012.01~~~~~~~~-{}-{}-{}-{}D~-{}PC~-{}-{}-{}-{}++++~ok}\par
\texttt{F4.2186~02.02.012.01~~~~~~~~-{}-{}-{}-{}D~-{}P-{}~-{}-{}-{}-{}-{}+-{}-{}~3:1110}\par
\texttt{F4.2187~12.02.012.01~~~~~~~~-{}-{}-{}-{}D~-{}P-{}~-{}-{}-{}-{}-{}+-{}-{}~3:1110}\par
\texttt{F4.2188~03.02.012.01~~~~~~~~-{}-{}-{}-{}D~-{}PC~-{}-{}-{}-{}++-{}+~ok}\par
\texttt{F4.2189~13.02.012.01~~~~~~~~-{}-{}-{}-{}D~-{}PC~-{}-{}-{}-{}++++~ok}\par
\texttt{F4.2190~23.02.012.01~~~~~~~~-{}-{}-{}-{}D~-{}PC~-{}-{}-{}-{}++++~ok}\par
\texttt{F4.2191~02.12.012.01~~~~~~~~-{}-{}-{}-{}D~-{}P-{}~-{}-{}-{}-{}-{}+-{}-{}~3:1110}\par
\texttt{F4.2192~03.12.012.01~~~~~~~~-{}-{}-{}-{}D~-{}PC~-{}-{}-{}-{}++-{}+~ok}\par
\texttt{F4.2193~13.12.012.01~~~~~~~~-{}-{}-{}-{}D~-{}PC~-{}-{}-{}-{}++-{}+~ok}\par
\texttt{F4.2194~03.03.012.01~~~~~~~~-{}-{}-{}-{}D~-{}-{}-{}~-{}-{}-{}-{}-{}-{}-{}-{}~2:1101}\par
\texttt{F4.2195~13.03.012.01~~~~~~~~-{}-{}-{}-{}D~-{}-{}-{}~-{}-{}-{}-{}-{}-{}-{}-{}~2:1101}\par
\texttt{F4.2196~23.03.012.01~~~~~~~~-{}-{}-{}-{}D~-{}PC~-{}-{}-{}-{}++++~ok}\par
\texttt{F4.2197~03.13.012.01~~~~~~~~-{}-{}-{}-{}D~-{}-{}-{}~-{}-{}-{}-{}-{}-{}-{}-{}~2:1101}\par
\texttt{F4.2198~012.02.03.01~~~~~~~~-{}-{}-{}-{}D~-{}PC~-{}-{}-{}-{}++++~ok}\par
\texttt{F4.2199~123.02.03.01~~~~~~~~-{}-{}-{}-{}D~-{}PC~-{}-{}-{}-{}++++~ok}\par
\texttt{F4.2200~012.12.03.01~~~~~~~~-{}-{}-{}-{}D~-{}PC~-{}-{}-{}-{}++++~ok}\par
\texttt{F4.2201~013.12.03.01~~~~~~~~-{}-{}-{}-{}D~-{}PC~-{}-{}-{}-{}++-{}+~ok}\par
\texttt{F4.2202~023.12.03.01~~~~~~~~-{}-{}-{}-{}D~-{}PC~-{}-{}-{}-{}++++~ok}\par
\texttt{F4.2203~123.12.03.01~~~~~~~~-{}-{}-{}-{}D~-{}PC~-{}-{}-{}-{}++++~ok}\par
\texttt{F4.2204~01.012.03.01~~~~~~~~-{}-{}-{}-{}D~-{}PC~-{}-{}-{}-{}++-{}+~ok}\par
\texttt{F4.2205~02.012.03.01~~~~~~~~-{}-{}-{}-{}D~-{}PC~-{}-{}-{}-{}++-{}+~ok}\par
\texttt{F4.2206~12.012.03.01~~~~~~~~-{}-{}-{}-{}D~-{}PC~-{}-{}-{}-{}++++~ok}\par
\texttt{F4.2207~03.012.03.01~~~~~~~~-{}-{}-{}-{}D~-{}PC~-{}-{}-{}-{}++-{}+~ok}\par
\texttt{F4.2208~13.012.03.01~~~~~~~~-{}-{}-{}-{}D~-{}PC~-{}-{}-{}-{}++-{}+~ok}\par
\texttt{F4.2209~23.012.03.01~~~~~~~~-{}-{}-{}-{}D~-{}PC~-{}-{}-{}-{}++++~ok}\par
\texttt{F4.2210~012.23.03.01~~~~~~~~-{}-{}-{}-{}D~-{}PC~-{}-{}-{}-{}++++~ok}\par
\texttt{F4.2211~123.23.03.01~~~~~~~~-{}-{}-{}-{}D~-{}PC~-{}-{}-{}-{}++++~ok}\par
\texttt{F4.2212~01.023.03.01~~~~~~~~-{}-{}-{}-{}D~-{}PC~-{}-{}-{}-{}++++~ok}\par
\texttt{F4.2213~02.023.03.01~~~~~~~~-{}-{}-{}-{}D~-{}PC~-{}-{}-{}-{}++-{}+~ok}\par
\texttt{F4.2214~03.023.03.01~~~~~~~~-{}-{}-{}-{}D~-{}PC~-{}-{}-{}-{}++-{}+~ok}\par
\texttt{F4.2215~13.023.03.01~~~~~~~~-{}-{}-{}-{}D~-{}PC~-{}-{}-{}-{}++++~ok}\par
\texttt{F4.2216~23.023.03.01~~~~~~~~-{}-{}-{}-{}D~-{}PC~-{}-{}-{}-{}++++~ok}\par
\texttt{F4.2217~01.123.03.01~~~~~~~~-{}-{}-{}-{}D~-{}PC~-{}-{}-{}-{}++-{}+~ok}\par
\texttt{F4.2218~02.123.03.01~~~~~~~~-{}-{}-{}-{}D~-{}PC~-{}-{}-{}-{}++-{}+~ok}\par
\texttt{F4.2219~03.123.03.01~~~~~~~~-{}-{}-{}-{}D~-{}PC~-{}-{}-{}-{}++-{}+~ok}\par
\texttt{F4.2220~0.0123.03.01~~~~~~~~-{}-{}-{}-{}D~-{}-{}-{}~-{}-{}-{}-{}-{}-{}-{}-{}~1:1000}\par
\texttt{F4.2221~012.1.013.01~~~~~~~~-{}-{}-{}-{}D~-{}-{}-{}~-{}-{}-{}-{}-{}-{}-{}-{}~0:0100}\par
\texttt{F4.2222~013.1.013.01~~~~~~~~-{}-{}-{}-{}D~-{}-{}-{}~-{}-{}-{}-{}-{}-{}-{}-{}~0:0100}\par
\texttt{F4.2223~023.1.013.01~~~~~~~~-{}-{}-{}-{}D~-{}-{}-{}~-{}-{}-{}-{}-{}-{}-{}-{}~2:0100}\par
\texttt{F4.2224~123.1.013.01~~~~~~~~-{}-{}-{}-{}D~-{}-{}-{}~-{}-{}-{}-{}-{}-{}-{}-{}~0:0100}\par
\texttt{F4.2225~01.01.013.01~~~~~~~~-{}-{}4-{}D~UP-{}~-{}-{}++-{}+-{}-{}~2:1100}\par
\texttt{F4.2226~02.01.013.01~~~~~~~~-{}-{}-{}-{}D~-{}PC~-{}-{}-{}-{}++-{}+~ok}\par
\texttt{F4.2227~03.01.013.01~~~~~~~~-{}-{}-{}-{}D~-{}P-{}~-{}-{}-{}-{}-{}+-{}-{}~2:1101}\par
\texttt{F4.2228~13.01.013.01~~~~~~~~-{}-{}-{}-{}D~-{}P-{}~-{}-{}-{}-{}-{}+-{}-{}~2:1101}\par
\texttt{F4.2229~23.01.013.01~~~~~~~~-{}-{}-{}-{}D~-{}PC~-{}-{}-{}-{}++++~ok}\par
\texttt{F4.2230~02.12.013.01~~~~~~~~-{}-{}-{}-{}D~-{}PC~-{}-{}-{}-{}++++~ok}\par
\texttt{F4.2231~03.12.013.01~~~~~~~~-{}-{}-{}-{}D~-{}PC~-{}-{}-{}-{}++-{}+~ok}\par
\texttt{F4.2232~13.12.013.01~~~~~~~~-{}-{}-{}-{}D~-{}PC~-{}-{}-{}-{}++++~ok}\par
\texttt{F4.2233~13.03.013.01~~~~~~~~-{}-{}-{}-{}D~-{}P-{}~-{}-{}-{}-{}-{}+-{}-{}~2:1101}\par
\texttt{F4.2234~03.13.013.01~~~~~~~~-{}-{}-{}-{}D~-{}P-{}~-{}-{}-{}-{}-{}+-{}-{}~2:1101}\par
\texttt{F4.2235~0123.1.23.01~~~~~~~~-{}-{}-{}-{}D~-{}-{}-{}~-{}-{}-{}-{}-{}-{}-{}-{}~0:0100}\par
\texttt{F4.2236~012.01.23.01~~~~~~~~-{}-{}-{}-{}D~-{}PC~-{}-{}-{}-{}++-{}+~ok}\par
\texttt{F4.2237~013.01.23.01~~~~~~~~-{}-{}-{}-{}D~-{}-{}-{}~-{}-{}-{}-{}-{}-{}-{}-{}~2:1101}\par
\texttt{F4.2238~023.01.23.01~~~~~~~~-{}-{}-{}-{}D~-{}PC~-{}-{}-{}-{}++++~ok}\par
\texttt{F4.2239~123.01.23.01~~~~~~~~-{}-{}-{}-{}D~-{}PC~-{}-{}-{}-{}++++~ok}\par
\texttt{F4.2240~012.12.23.01~~~~~~~~-{}-{}-{}-{}D~-{}PC~-{}-{}-{}-{}++-{}+~ok}\par
\texttt{F4.2241~013.12.23.01~~~~~~~~-{}-{}-{}-{}D~-{}PC~-{}-{}-{}-{}++-{}+~ok}\par
\texttt{F4.2242~023.12.23.01~~~~~~~~-{}-{}-{}-{}D~-{}PC~-{}-{}-{}-{}++-{}+~ok}\par
\texttt{F4.2243~123.12.23.01~~~~~~~~-{}-{}-{}-{}D~-{}PC~-{}-{}-{}-{}++-{}+~ok}\par
\texttt{F4.2244~03.012.23.01~~~~~~~~-{}-{}-{}-{}D~-{}PC~-{}-{}-{}-{}++-{}+~ok}\par
\texttt{F4.2245~23.012.23.01~~~~~~~~-{}-{}-{}-{}D~-{}PC~-{}-{}-{}-{}++++~ok}\par
\texttt{F4.2246~013.13.23.01~~~~~~~~-{}-{}-{}-{}D~-{}-{}-{}~-{}-{}-{}-{}-{}-{}-{}-{}~2:1101}\par
\texttt{F4.2247~023.13.23.01~~~~~~~~-{}-{}-{}-{}D~-{}PC~-{}-{}-{}-{}++-{}+~ok}\par
\texttt{F4.2248~123.13.23.01~~~~~~~~-{}-{}-{}-{}D~-{}PC~-{}-{}-{}-{}++-{}+~ok}\par
\texttt{F4.2249~23.013.23.01~~~~~~~~-{}-{}-{}-{}D~-{}PC~-{}-{}-{}-{}++++~ok}\par
\texttt{F4.2250~012.1.023.01~~~~~~~~-{}-{}-{}-{}D~-{}-{}-{}~-{}-{}-{}-{}-{}-{}-{}-{}~0:0100}\par
\texttt{F4.2251~013.1.023.01~~~~~~~~-{}-{}-{}-{}D~-{}-{}-{}~-{}-{}-{}-{}-{}-{}-{}-{}~0:0100}\par
\texttt{F4.2252~023.1.023.01~~~~~~~~-{}-{}-{}-{}D~-{}-{}-{}~-{}-{}-{}-{}-{}-{}-{}-{}~3:0100}\par
\texttt{F4.2253~123.1.023.01~~~~~~~~-{}-{}-{}-{}D~-{}-{}-{}~-{}-{}-{}-{}-{}-{}-{}-{}~0:0100}\par
\texttt{F4.2254~01.01.023.01~~~~~~~~-{}-{}-{}-{}D~-{}-{}-{}~-{}-{}-{}-{}-{}-{}-{}-{}~2:1100}\par
\texttt{F4.2255~02.01.023.01~~~~~~~~-{}-{}-{}-{}D~-{}PC~-{}-{}-{}-{}++-{}+~ok}\par
\texttt{F4.2256~03.01.023.01~~~~~~~~-{}-{}-{}-{}D~-{}-{}-{}~-{}-{}-{}-{}-{}-{}-{}-{}~2:1101}\par
\texttt{F4.2257~13.01.023.01~~~~~~~~-{}-{}-{}-{}D~-{}-{}-{}~-{}-{}-{}-{}-{}-{}-{}-{}~2:1101}\par
\texttt{F4.2258~01.12.023.01~~~~~~~~-{}-{}-{}-{}D~-{}PC~-{}-{}-{}-{}++-{}+~ok}\par
\texttt{F4.2259~02.12.023.01~~~~~~~~-{}-{}-{}-{}D~-{}PC~-{}-{}-{}-{}++-{}+~ok}\par
\texttt{F4.2260~03.12.023.01~~~~~~~~-{}-{}-{}-{}D~-{}PC~-{}-{}-{}-{}++-{}+~ok}\par
\texttt{F4.2261~13.12.023.01~~~~~~~~-{}-{}-{}-{}D~-{}PC~-{}-{}-{}-{}++++~ok}\par
\texttt{F4.2262~0.012.023.01~~~~~~~~-{}-{}-{}-{}D~-{}-{}-{}~-{}-{}-{}-{}-{}-{}-{}-{}~1:1000}\par
\texttt{F4.2263~01.13.023.01~~~~~~~~-{}-{}-{}-{}D~-{}-{}-{}~-{}-{}-{}-{}-{}-{}-{}-{}~2:1101}\par
\texttt{F4.2264~02.13.023.01~~~~~~~~-{}-{}-{}-{}D~-{}PC~-{}-{}-{}-{}++-{}+~ok}\par
\texttt{F4.2265~03.13.023.01~~~~~~~~-{}-{}-{}-{}D~-{}-{}-{}~-{}-{}-{}-{}-{}-{}-{}-{}~2:1101}\par
\texttt{F4.2266~0.013.023.01~~~~~~~~-{}-{}-{}-{}D~-{}-{}-{}~-{}-{}-{}-{}-{}-{}-{}-{}~1:1000}\par
\texttt{F4.2267~0.023.023.01~~~~~~~~-{}-{}-{}-{}D~-{}-{}-{}~-{}-{}-{}-{}-{}-{}-{}-{}~1:1000}\par
\texttt{F4.2268~0.123.023.01~~~~~~~~-{}-{}-{}-{}D~-{}-{}-{}~-{}-{}-{}-{}-{}-{}-{}-{}~2:1000}\par
\texttt{F4.2269~01.1.0123.01~~~~~~~~-{}-{}4-{}D~-{}-{}-{}~-{}-{}-{}+-{}-{}-{}-{}~0:0100}\par
\texttt{F4.2270~02.1.0123.01~~~~~~~~-{}-{}-{}-{}D~-{}-{}-{}~-{}-{}-{}-{}-{}-{}-{}-{}~2:0100}\par
\texttt{F4.2271~03.1.0123.01~~~~~~~~-{}-{}-{}-{}D~-{}-{}-{}~-{}-{}-{}-{}-{}-{}-{}-{}~2:0100}\par
\texttt{F4.2272~13.1.0123.01~~~~~~~~-{}-{}-{}-{}D~-{}-{}-{}~-{}-{}-{}-{}-{}-{}-{}-{}~0:0100}\par
\texttt{F4.2273~0123.1.2.012~~~~~~~~-{}-{}-{}-{}D~-{}-{}-{}~-{}-{}-{}-{}-{}-{}-{}-{}~0:0100}\par
\texttt{F4.2274~012.01.2.012~~~~~~~~-{}-{}-{}-{}D~-{}-{}-{}~-{}-{}-{}-{}-{}-{}-{}-{}~0:0010}\par
\texttt{F4.2275~013.01.2.012~~~~~~~~-{}-{}-{}-{}D~-{}-{}-{}~-{}-{}-{}-{}-{}-{}-{}-{}~3:0010}\par
\texttt{F4.2276~023.01.2.012~~~~~~~~-{}-{}-{}-{}D~-{}-{}-{}~-{}-{}-{}-{}-{}-{}-{}-{}~0:0010}\par
\texttt{F4.2277~123.01.2.012~~~~~~~~-{}-{}-{}-{}D~-{}-{}-{}~-{}-{}-{}-{}-{}-{}-{}-{}~0:0010}\par
\texttt{F4.2278~012.12.2.012~~~~~~~~-{}-{}4-{}D~-{}-{}-{}~-{}-{}-{}+-{}-{}-{}-{}~0:0010}\par
\texttt{F4.2279~013.12.2.012~~~~~~~~-{}-{}-{}-{}D~-{}-{}-{}~-{}-{}-{}-{}-{}-{}-{}-{}~1:0010}\par
\texttt{F4.2280~023.12.2.012~~~~~~~~-{}-{}-{}-{}D~-{}-{}-{}~-{}-{}-{}-{}-{}-{}-{}-{}~0:0010}\par
\texttt{F4.2281~123.12.2.012~~~~~~~~-{}-{}-{}-{}D~-{}-{}-{}~-{}-{}-{}-{}-{}-{}-{}-{}~0:0010}\par
\texttt{F4.2282~03.012.2.012~~~~~~~~-{}-{}-{}-{}D~-{}-{}-{}~-{}-{}-{}-{}-{}-{}-{}-{}~1:0010}\par
\texttt{F4.2283~013.13.2.012~~~~~~~~-{}-{}-{}-{}D~-{}-{}-{}~-{}-{}-{}-{}-{}-{}-{}-{}~3:0010}\par
\texttt{F4.2284~023.13.2.012~~~~~~~~-{}-{}-{}-{}D~-{}-{}-{}~-{}-{}-{}-{}-{}-{}-{}-{}~0:0010}\par
\texttt{F4.2285~01.01.02.012~~~~~~~~-{}-{}-{}-{}D~-{}-{}-{}~-{}-{}-{}-{}-{}-{}-{}-{}~2:1100}\par
\texttt{F4.2286~03.01.02.012~~~~~~~~-{}-{}-{}-{}D~-{}PC~-{}-{}-{}-{}++-{}+~ok}\par
\texttt{F4.2287~01.12.02.012~~~~~~~~-{}-{}-{}-{}D~-{}P-{}~-{}-{}-{}-{}-{}+-{}-{}~3:1110}\par
\texttt{F4.2288~03.12.02.012~~~~~~~~-{}-{}-{}-{}D~-{}PC~-{}-{}-{}-{}++-{}+~ok}\par
\texttt{F4.2289~02.13.02.012~~~~~~~~-{}-{}-{}-{}D~-{}-{}-{}~-{}-{}-{}-{}-{}-{}-{}-{}~3:1010}\par
\texttt{F4.2290~03.13.02.012~~~~~~~~-{}-{}-{}-{}D~-{}PC~-{}-{}-{}-{}++-{}+~ok}\par
\texttt{F4.2291~03.13.23.012~~~~~~~~-{}B-{}-{}D~-{}PC~+-{}-{}-{}++++~ok}\par
\texttt{F4.2292~0123.012.2.3~~~~~~~~T-{}-{}-{}D~-{}-{}-{}~-{}-{}-{}-{}-{}-{}-{}-{}~0:0010}\par
\texttt{F4.2293~0123.123.2.3~~~~~~~~T-{}4-{}D~-{}-{}-{}~-{}-{}-{}+-{}-{}-{}-{}~0:0010}\par
\texttt{F4.2294~0123.01.02.3~~~~~~~~T-{}-{}-{}D~-{}-{}-{}~-{}-{}-{}-{}-{}-{}-{}-{}~0:0001}\par
\texttt{F4.2295~0123.12.02.3~~~~~~~~T-{}-{}-{}D~-{}-{}-{}~-{}-{}-{}-{}-{}-{}-{}-{}~0:0001}\par
\texttt{F4.2296~012.012.02.3~~~~~~~~T-{}-{}-{}D~-{}PC~-{}-{}-{}-{}++++~ok}\par
\texttt{F4.2297~013.012.02.3~~~~~~~~T-{}-{}-{}D~-{}-{}-{}~-{}-{}-{}-{}-{}-{}-{}-{}~0:0001}\par
\texttt{F4.2298~023.012.02.3~~~~~~~~T-{}-{}-{}D~-{}-{}-{}~-{}-{}-{}-{}-{}-{}-{}-{}~0:0001}\par
\texttt{F4.2299~0123.13.02.3~~~~~~~~T-{}-{}-{}D~-{}-{}-{}~-{}-{}-{}-{}-{}-{}-{}-{}~0:0001}\par
\texttt{F4.2300~012.013.02.3~~~~~~~~T-{}-{}-{}D~-{}-{}-{}~-{}-{}-{}-{}-{}-{}-{}-{}~1:0001}\par
\texttt{F4.2301~013.013.02.3~~~~~~~~T-{}-{}-{}D~-{}-{}-{}~-{}-{}-{}-{}-{}-{}-{}-{}~0:0001}\par
\texttt{F4.2302~023.013.02.3~~~~~~~~T-{}-{}-{}D~-{}-{}-{}~-{}-{}-{}-{}-{}-{}-{}-{}~0:0001}\par
\texttt{F4.2303~012.123.02.3~~~~~~~~T-{}-{}-{}D~-{}-{}-{}~-{}-{}-{}-{}-{}-{}-{}-{}~1:0001}\par
\texttt{F4.2304~013.123.02.3~~~~~~~~T-{}-{}-{}D~-{}-{}-{}~-{}-{}-{}-{}-{}-{}-{}-{}~0:0001}\par
\texttt{F4.2305~023.123.02.3~~~~~~~~T-{}-{}-{}D~-{}-{}-{}~-{}-{}-{}-{}-{}-{}-{}-{}~0:0001}\par
\texttt{F4.2306~02.0123.02.3~~~~~~~~T-{}4-{}D~-{}-{}-{}~-{}-{}-{}+-{}-{}-{}-{}~1:1010}\par
\texttt{F4.2307~03.0123.02.3~~~~~~~~T-{}-{}-{}D~-{}-{}-{}~-{}-{}-{}-{}-{}-{}-{}-{}~0:0001}\par
\texttt{F4.2308~03.012.012.3~~~~~~~~T-{}-{}-{}D~-{}-{}-{}~-{}-{}-{}-{}-{}-{}-{}-{}~0:0001}\par
\texttt{F4.2309~013.13.012.3~~~~~~~~T-{}-{}-{}D~-{}-{}-{}~-{}-{}-{}-{}-{}-{}-{}-{}~0:0001}\par
\texttt{F4.2310~023.13.012.3~~~~~~~~T-{}-{}-{}D~-{}-{}-{}~-{}-{}-{}-{}-{}-{}-{}-{}~0:0001}\par
\texttt{F4.2311~0123.13.23.3~~~~~~~~T-{}4-{}D~-{}-{}-{}~-{}-{}-{}+-{}-{}-{}-{}~0:0001}\par
\texttt{F4.2312~013.013.23.3~~~~~~~~T-{}4-{}D~-{}-{}-{}~-{}-{}-{}+-{}-{}-{}-{}~0:0001}\par
\texttt{F4.2313~023.013.23.3~~~~~~~~T-{}-{}-{}D~-{}-{}-{}~-{}-{}-{}-{}-{}-{}-{}-{}~0:0001}\par
\texttt{F4.2314~023.123.23.3~~~~~~~~T-{}4-{}D~-{}-{}-{}~-{}-{}-{}+-{}-{}-{}-{}~0:0001}\par
\texttt{F4.2315~012.01.02.03~~~~~~~~T-{}-{}-{}D~-{}-{}-{}~-{}-{}-{}-{}-{}-{}-{}-{}~3:1110}\par
\texttt{F4.2316~012.12.02.03~~~~~~~~T-{}-{}-{}D~-{}-{}-{}~-{}-{}-{}-{}-{}-{}-{}-{}~3:1110}\par
\texttt{F4.2317~013.12.02.03~~~~~~~~T-{}-{}-{}D~-{}PC~-{}-{}-{}-{}++-{}+~ok}\par
\texttt{F4.2318~023.12.02.03~~~~~~~~T-{}-{}-{}D~-{}-{}-{}~-{}-{}-{}-{}-{}-{}-{}-{}~1:1011}\par
\texttt{F4.2319~01.012.02.03~~~~~~~~T-{}-{}-{}D~-{}-{}-{}~-{}-{}-{}-{}-{}-{}-{}-{}~3:1110}\par
\texttt{F4.2320~02.012.02.03~~~~~~~~T-{}-{}-{}D~-{}-{}-{}~-{}-{}-{}-{}-{}-{}-{}-{}~1:1010}\par
\texttt{F4.2321~03.012.02.03~~~~~~~~T-{}-{}-{}D~-{}-{}-{}~-{}-{}-{}-{}-{}-{}-{}-{}~1:1001}\par
\texttt{F4.2322~01.123.02.03~~~~~~~~T-{}-{}-{}D~-{}PC~-{}-{}-{}-{}++-{}+~ok}\par
\texttt{F4.2323~02.123.02.03~~~~~~~~T-{}-{}-{}D~-{}-{}-{}~-{}-{}-{}-{}-{}-{}-{}-{}~1:1010}\par
\texttt{F4.2324~012.12.12.03~~~~~~~~T-{}-{}-{}D~-{}-{}-{}~-{}-{}-{}-{}-{}-{}-{}-{}~0:0110}\par
\texttt{F4.2325~013.12.12.03~~~~~~~~T-{}-{}-{}D~-{}-{}-{}~-{}-{}-{}-{}-{}-{}-{}-{}~0:0110}\par
\texttt{F4.2326~02.012.12.03~~~~~~~~T-{}-{}-{}D~-{}-{}-{}~-{}-{}-{}-{}-{}-{}-{}-{}~3:1110}\par
\texttt{F4.2327~012.13.12.03~~~~~~~~T-{}-{}-{}D~-{}PC~-{}-{}-{}-{}++-{}+~ok}\par
\texttt{F4.2328~023.13.12.03~~~~~~~~T-{}-{}-{}D~-{}PC~-{}-{}-{}-{}++-{}+~ok}\par
\texttt{F4.2329~01.13.012.03~~~~~~~~T-{}-{}-{}D~-{}-{}-{}~-{}-{}-{}-{}-{}-{}-{}-{}~2:1101}\par
\texttt{F4.2330~0123.0123.01.-{}~~~~~~-{}-{}-{}-{}-{}~-{}-{}-{}~-{}-{}-{}-{}-{}-{}-{}-{}~0:0000}\par
\texttt{F4.2331~0123.0123.02.-{}~~~~~~-{}-{}-{}-{}-{}~-{}-{}-{}~-{}-{}-{}-{}-{}-{}-{}-{}~0:0000}\par
\texttt{F4.2332~0123.012.012.-{}~~~~~~-{}-{}-{}-{}-{}~-{}-{}-{}~-{}-{}-{}-{}-{}-{}-{}-{}~0:0000}\par
\texttt{F4.2333~0123.013.012.-{}~~~~~~-{}-{}-{}-{}-{}~-{}-{}-{}~-{}-{}-{}-{}-{}-{}-{}-{}~0:0000}\par
\texttt{F4.2334~0123.023.012.-{}~~~~~~-{}-{}-{}-{}-{}~-{}-{}-{}~-{}-{}-{}-{}-{}-{}-{}-{}~0:0000}\par
\texttt{F4.2335~0123.123.012.-{}~~~~~~-{}-{}-{}-{}-{}~-{}-{}-{}~-{}-{}-{}-{}-{}-{}-{}-{}~0:0000}\par
\texttt{F4.2336~0123.0123.03.-{}~~~~~~-{}-{}-{}-{}-{}~-{}-{}-{}~-{}-{}-{}-{}-{}-{}-{}-{}~0:0000}\par
\texttt{F4.2337~0123.013.013.-{}~~~~~~-{}-{}-{}-{}-{}~-{}-{}-{}~-{}-{}-{}-{}-{}-{}-{}-{}~0:0000}\par
\texttt{F4.2338~0123.023.013.-{}~~~~~~-{}-{}-{}-{}-{}~-{}-{}-{}~-{}-{}-{}-{}-{}-{}-{}-{}~0:0000}\par
\texttt{F4.2339~0123.123.013.-{}~~~~~~-{}-{}-{}-{}-{}~-{}-{}-{}~-{}-{}-{}-{}-{}-{}-{}-{}~0:0000}\par
\texttt{F4.2340~0123.0123.23.-{}~~~~~~-{}-{}4-{}-{}~-{}-{}-{}~-{}-{}-{}+-{}-{}-{}-{}~0:0000}\par
\texttt{F4.2341~0123.013.023.-{}~~~~~~-{}-{}-{}-{}-{}~-{}-{}-{}~-{}-{}-{}-{}-{}-{}-{}-{}~0:0000}\par
\texttt{F4.2342~0123.123.023.-{}~~~~~~-{}-{}-{}-{}-{}~-{}-{}-{}~-{}-{}-{}-{}-{}-{}-{}-{}~0:0000}\par
\texttt{F4.2343~023.0123.023.-{}~~~~~~-{}-{}4-{}-{}~-{}-{}-{}~-{}-{}-{}+-{}-{}-{}-{}~0:0000}\par
\texttt{F4.2344~0123.0123.0.0~~~~~~~-{}-{}-{}-{}D~-{}PC~-{}-{}-{}-{}++++~ok}\par
\texttt{F4.2345~0123.0123.1.0~~~~~~~-{}-{}-{}-{}D~-{}PC~-{}-{}-{}-{}++++~ok}\par
\texttt{F4.2346~0123.012.01.0~~~~~~~-{}B-{}-{}D~-{}PC~+-{}-{}-{}++++~ok}\par
\texttt{F4.2347~0123.013.01.0~~~~~~~-{}-{}-{}-{}D~-{}PC~-{}-{}-{}-{}++++~ok}\par
\texttt{F4.2348~0123.023.01.0~~~~~~~-{}-{}-{}-{}D~-{}PC~-{}-{}-{}-{}++++~ok}\par
\texttt{F4.2349~0123.123.01.0~~~~~~~-{}-{}-{}-{}D~-{}PC~-{}-{}-{}-{}++++~ok}\par
\texttt{F4.2350~012.0123.01.0~~~~~~~-{}-{}-{}-{}D~-{}PC~-{}-{}-{}-{}++-{}+~ok}\par
\texttt{F4.2351~013.0123.01.0~~~~~~~-{}-{}-{}-{}D~-{}PC~-{}-{}-{}-{}++-{}+~ok}\par
\texttt{F4.2352~023.0123.01.0~~~~~~~-{}-{}-{}-{}D~-{}PC~-{}-{}-{}-{}++-{}+~ok}\par
\texttt{F4.2353~123.0123.01.0~~~~~~~-{}-{}-{}-{}D~-{}PC~-{}-{}-{}-{}++++~ok}\par
\texttt{F4.2354~0123.0123.2.0~~~~~~~-{}-{}-{}-{}D~-{}-{}-{}~-{}-{}-{}-{}-{}-{}-{}-{}~0:0010}\par
\texttt{F4.2355~0123.012.02.0~~~~~~~-{}-{}-{}-{}D~-{}PC~-{}-{}-{}-{}++++~ok}\par
\texttt{F4.2356~0123.013.02.0~~~~~~~-{}-{}-{}-{}D~-{}PC~-{}-{}-{}-{}++++~ok}\par
\texttt{F4.2357~0123.023.02.0~~~~~~~-{}-{}-{}-{}D~-{}PC~-{}-{}-{}-{}++++~ok}\par
\texttt{F4.2358~0123.123.02.0~~~~~~~-{}-{}-{}-{}D~-{}PC~-{}-{}-{}-{}++++~ok}\par
\texttt{F4.2359~012.0123.02.0~~~~~~~-{}-{}-{}-{}D~-{}PC~-{}-{}-{}-{}++-{}+~ok}\par
\texttt{F4.2360~013.0123.02.0~~~~~~~-{}-{}-{}-{}D~-{}PC~-{}-{}-{}-{}++-{}+~ok}\par
\texttt{F4.2361~023.0123.02.0~~~~~~~-{}-{}-{}-{}D~-{}-{}-{}~-{}-{}-{}-{}-{}-{}-{}-{}~1:1011}\par
\texttt{F4.2362~123.0123.02.0~~~~~~~-{}-{}-{}-{}D~-{}PC~-{}-{}-{}-{}++++~ok}\par
\texttt{F4.2363~0123.012.12.0~~~~~~~-{}-{}-{}-{}D~-{}PC~-{}-{}-{}-{}++-{}+~ok}\par
\texttt{F4.2364~0123.013.12.0~~~~~~~-{}-{}-{}-{}D~-{}PC~-{}-{}-{}-{}++++~ok}\par
\texttt{F4.2365~0123.023.12.0~~~~~~~-{}-{}-{}-{}D~-{}PC~-{}-{}-{}-{}++++~ok}\par
\texttt{F4.2366~0123.123.12.0~~~~~~~-{}-{}-{}-{}D~-{}PC~-{}-{}-{}-{}++-{}+~ok}\par
\texttt{F4.2367~012.0123.12.0~~~~~~~-{}-{}-{}-{}D~-{}PC~-{}-{}-{}-{}++-{}+~ok}\par
\texttt{F4.2368~013.0123.12.0~~~~~~~-{}-{}-{}-{}D~-{}PC~-{}-{}-{}-{}++-{}+~ok}\par
\texttt{F4.2369~023.0123.12.0~~~~~~~-{}-{}-{}-{}D~-{}PC~-{}-{}-{}-{}++-{}+~ok}\par
\texttt{F4.2370~123.0123.12.0~~~~~~~-{}-{}-{}-{}D~-{}PC~-{}-{}-{}-{}++++~ok}\par
\texttt{F4.2371~012.012.012.0~~~~~~~-{}-{}-{}5D~UP-{}~-{}++-{}-{}+-{}-{}~3:1110}\par
\texttt{F4.2372~013.012.012.0~~~~~~~-{}-{}-{}-{}D~-{}PC~-{}-{}-{}-{}++-{}+~ok}\par
\texttt{F4.2373~123.012.012.0~~~~~~~-{}B-{}-{}D~-{}PC~+-{}-{}-{}++++~ok}\par
\texttt{F4.2374~0123.03.012.0~~~~~~~-{}-{}-{}-{}D~-{}PC~-{}-{}-{}-{}++++~ok}\par
\texttt{F4.2375~0123.13.012.0~~~~~~~-{}-{}-{}-{}D~-{}PC~-{}-{}-{}-{}++-{}+~ok}\par
\texttt{F4.2376~012.013.012.0~~~~~~~-{}-{}-{}-{}D~-{}PC~-{}-{}-{}-{}++-{}+~ok}\par
\texttt{F4.2377~013.013.012.0~~~~~~~-{}-{}-{}-{}D~-{}-{}-{}~-{}-{}-{}-{}-{}-{}-{}-{}~2:1101}\par
\texttt{F4.2378~023.013.012.0~~~~~~~-{}-{}-{}-{}D~-{}PC~-{}-{}-{}-{}++-{}+~ok}\par
\texttt{F4.2379~123.013.012.0~~~~~~~-{}-{}-{}-{}D~-{}PC~-{}-{}-{}-{}++++~ok}\par
\texttt{F4.2380~0123.23.012.0~~~~~~~-{}-{}-{}-{}D~-{}PC~-{}-{}-{}-{}++++~ok}\par
\texttt{F4.2381~012.023.012.0~~~~~~~-{}-{}-{}-{}D~-{}PC~-{}-{}-{}-{}++++~ok}\par
\texttt{F4.2382~013.023.012.0~~~~~~~-{}-{}-{}-{}D~-{}PC~-{}-{}-{}-{}++-{}+~ok}\par
\texttt{F4.2383~023.023.012.0~~~~~~~-{}-{}-{}-{}D~-{}PC~-{}-{}-{}-{}++-{}+~ok}\par
\texttt{F4.2384~123.023.012.0~~~~~~~-{}-{}-{}-{}D~-{}PC~-{}-{}-{}-{}++++~ok}\par
\texttt{F4.2385~012.123.012.0~~~~~~~-{}-{}-{}-{}D~-{}PC~-{}-{}-{}-{}++-{}+~ok}\par
\texttt{F4.2386~013.123.012.0~~~~~~~-{}-{}-{}-{}D~-{}PC~-{}-{}-{}-{}++-{}+~ok}\par
\texttt{F4.2387~023.123.012.0~~~~~~~-{}-{}-{}-{}D~-{}PC~-{}-{}-{}-{}++-{}+~ok}\par
\texttt{F4.2388~123.123.012.0~~~~~~~-{}-{}-{}-{}D~-{}PC~-{}-{}-{}-{}++++~ok}\par
\texttt{F4.2389~01.0123.012.0~~~~~~~-{}-{}-{}-{}D~-{}PC~-{}-{}-{}-{}++-{}+~ok}\par
\texttt{F4.2390~02.0123.012.0~~~~~~~-{}-{}-{}-{}D~-{}PC~-{}-{}-{}-{}++-{}-{}~ok}\par
\texttt{F4.2391~12.0123.012.0~~~~~~~-{}-{}-{}-{}D~-{}PC~-{}-{}-{}-{}++-{}+~ok}\par
\texttt{F4.2392~03.0123.012.0~~~~~~~-{}-{}-{}-{}D~-{}-{}-{}~-{}-{}-{}-{}-{}-{}-{}-{}~1:1001}\par
\texttt{F4.2393~13.0123.012.0~~~~~~~-{}-{}-{}-{}D~-{}PC~-{}-{}-{}-{}++-{}+~ok}\par
\texttt{F4.2394~23.0123.012.0~~~~~~~-{}-{}-{}-{}D~-{}PC~-{}-{}-{}-{}++-{}+~ok}\par
\texttt{F4.2395~0123.0123.3.0~~~~~~~-{}-{}-{}-{}D~-{}PC~-{}-{}-{}-{}++-{}+~ok}\par
\texttt{F4.2396~0123.013.03.0~~~~~~~-{}-{}-{}-{}D~-{}PC~-{}-{}-{}-{}++++~ok}\par
\texttt{F4.2397~0123.023.03.0~~~~~~~-{}-{}-{}-{}D~-{}PC~-{}-{}-{}-{}++++~ok}\par
\texttt{F4.2398~0123.123.03.0~~~~~~~-{}-{}-{}-{}D~-{}PC~-{}-{}-{}-{}++++~ok}\par
\texttt{F4.2399~012.0123.03.0~~~~~~~-{}-{}-{}-{}D~-{}PC~-{}-{}-{}-{}++-{}+~ok}\par
\texttt{F4.2400~013.0123.03.0~~~~~~~-{}-{}-{}-{}D~-{}PC~-{}-{}-{}-{}++-{}+~ok}\par
\texttt{F4.2401~023.0123.03.0~~~~~~~-{}-{}-{}-{}D~-{}-{}-{}~-{}-{}-{}-{}-{}-{}-{}-{}~1:1011}\par
\texttt{F4.2402~123.0123.03.0~~~~~~~-{}-{}-{}-{}D~-{}PC~-{}-{}-{}-{}++++~ok}\par
\texttt{F4.2403~0123.013.13.0~~~~~~~-{}-{}-{}-{}D~-{}PC~-{}-{}-{}-{}++-{}+~ok}\par
\texttt{F4.2404~0123.023.13.0~~~~~~~-{}-{}-{}-{}D~-{}PC~-{}-{}-{}-{}++++~ok}\par
\texttt{F4.2405~0123.123.13.0~~~~~~~-{}-{}-{}-{}D~-{}PC~-{}-{}-{}-{}++++~ok}\par
\texttt{F4.2406~012.0123.13.0~~~~~~~-{}-{}-{}-{}D~-{}PC~-{}-{}-{}-{}++-{}+~ok}\par
\texttt{F4.2407~013.0123.13.0~~~~~~~-{}-{}-{}-{}D~-{}PC~-{}-{}-{}-{}++-{}+~ok}\par
\texttt{F4.2408~023.0123.13.0~~~~~~~-{}-{}-{}-{}D~-{}PC~-{}-{}-{}-{}++-{}+~ok}\par
\texttt{F4.2409~123.0123.13.0~~~~~~~-{}-{}-{}-{}D~-{}PC~-{}-{}-{}-{}++++~ok}\par
\texttt{F4.2410~0123.13.013.0~~~~~~~-{}-{}-{}-{}D~-{}PC~-{}-{}-{}-{}++-{}+~ok}\par
\texttt{F4.2411~012.013.013.0~~~~~~~-{}-{}-{}-{}D~-{}PC~-{}-{}-{}-{}++-{}+~ok}\par
\texttt{F4.2412~013.013.013.0~~~~~~~-{}-{}-{}-{}D~-{}P-{}~-{}-{}-{}-{}-{}+-{}-{}~2:1101}\par
\texttt{F4.2413~023.013.013.0~~~~~~~-{}-{}-{}-{}D~-{}PC~-{}-{}-{}-{}++-{}+~ok}\par
\texttt{F4.2414~123.013.013.0~~~~~~~-{}-{}-{}-{}D~-{}PC~-{}-{}-{}-{}++++~ok}\par
\texttt{F4.2415~012.023.013.0~~~~~~~-{}-{}-{}-{}D~-{}PC~-{}-{}-{}-{}++-{}+~ok}\par
\texttt{F4.2416~013.023.013.0~~~~~~~-{}-{}-{}-{}D~-{}PC~-{}-{}-{}-{}++-{}+~ok}\par
\texttt{F4.2417~123.023.013.0~~~~~~~-{}-{}-{}-{}D~-{}PC~-{}-{}-{}-{}++++~ok}\par
\texttt{F4.2418~012.123.013.0~~~~~~~-{}-{}-{}-{}D~-{}PC~-{}-{}-{}-{}++-{}+~ok}\par
\texttt{F4.2419~013.123.013.0~~~~~~~-{}-{}-{}-{}D~-{}PC~-{}-{}-{}-{}++-{}+~ok}\par
\texttt{F4.2420~023.123.013.0~~~~~~~-{}-{}-{}-{}D~-{}PC~-{}-{}-{}-{}++-{}+~ok}\par
\texttt{F4.2421~123.123.013.0~~~~~~~-{}-{}-{}-{}D~-{}PC~-{}-{}-{}-{}++++~ok}\par
\texttt{F4.2422~01.0123.013.0~~~~~~~-{}-{}-{}-{}D~-{}PC~-{}-{}-{}-{}++-{}+~ok}\par
\texttt{F4.2423~02.0123.013.0~~~~~~~-{}-{}-{}-{}D~-{}PC~-{}-{}-{}-{}++-{}+~ok}\par
\texttt{F4.2424~12.0123.013.0~~~~~~~-{}-{}-{}-{}D~-{}PC~-{}-{}-{}-{}++-{}+~ok}\par
\texttt{F4.2425~03.0123.013.0~~~~~~~-{}-{}-{}-{}D~-{}-{}-{}~-{}-{}-{}-{}-{}-{}-{}-{}~1:1001}\par
\texttt{F4.2426~13.0123.013.0~~~~~~~-{}-{}-{}-{}D~-{}PC~-{}-{}-{}-{}++-{}+~ok}\par
\texttt{F4.2427~23.0123.013.0~~~~~~~-{}-{}-{}-{}D~-{}PC~-{}-{}-{}-{}++-{}+~ok}\par
\texttt{F4.2428~0123.013.23.0~~~~~~~-{}-{}-{}-{}D~-{}PC~-{}-{}-{}-{}++-{}+~ok}\par
\texttt{F4.2429~0123.123.23.0~~~~~~~-{}-{}-{}-{}D~-{}PC~-{}-{}-{}-{}++-{}+~ok}\par
\texttt{F4.2430~012.0123.23.0~~~~~~~-{}-{}-{}-{}D~-{}PC~-{}-{}-{}-{}++-{}+~ok}\par
\texttt{F4.2431~013.0123.23.0~~~~~~~-{}-{}-{}-{}D~-{}PC~-{}-{}-{}-{}++-{}+~ok}\par
\texttt{F4.2432~023.0123.23.0~~~~~~~-{}-{}-{}-{}D~-{}-{}-{}~-{}-{}-{}-{}-{}-{}-{}-{}~1:1011}\par
\texttt{F4.2433~123.0123.23.0~~~~~~~-{}-{}-{}-{}D~-{}PC~-{}-{}-{}-{}++-{}+~ok}\par
\texttt{F4.2434~012.013.023.0~~~~~~~-{}-{}-{}-{}D~-{}PC~-{}-{}-{}-{}++-{}+~ok}\par
\texttt{F4.2435~013.013.023.0~~~~~~~-{}-{}-{}-{}D~-{}-{}-{}~-{}-{}-{}-{}-{}-{}-{}-{}~2:1101}\par
\texttt{F4.2436~123.013.023.0~~~~~~~-{}-{}-{}-{}D~-{}PC~-{}-{}-{}-{}++++~ok}\par
\texttt{F4.2437~012.123.023.0~~~~~~~-{}-{}-{}-{}D~-{}PC~-{}-{}-{}-{}++-{}+~ok}\par
\texttt{F4.2438~013.123.023.0~~~~~~~-{}-{}-{}-{}D~-{}PC~-{}-{}-{}-{}++-{}+~ok}\par
\texttt{F4.2439~023.123.023.0~~~~~~~-{}-{}-{}-{}D~-{}-{}-{}~-{}-{}-{}-{}-{}-{}-{}-{}~1:1011}\par
\texttt{F4.2440~123.123.023.0~~~~~~~-{}-{}-{}-{}D~-{}PC~-{}-{}-{}-{}++++~ok}\par
\texttt{F4.2441~01.0123.023.0~~~~~~~-{}-{}-{}-{}D~-{}PC~-{}-{}-{}-{}++-{}+~ok}\par
\texttt{F4.2442~02.0123.023.0~~~~~~~-{}-{}-{}-{}D~-{}-{}-{}~-{}-{}-{}-{}-{}-{}-{}-{}~1:1011}\par
\texttt{F4.2443~12.0123.023.0~~~~~~~-{}-{}-{}-{}D~-{}PC~-{}-{}-{}-{}++-{}+~ok}\par
\texttt{F4.2444~03.0123.023.0~~~~~~~-{}-{}-{}-{}D~-{}-{}-{}~-{}-{}-{}-{}-{}-{}-{}-{}~1:1001}\par
\texttt{F4.2445~13.0123.023.0~~~~~~~-{}-{}-{}-{}D~-{}PC~-{}-{}-{}-{}++-{}+~ok}\par
\texttt{F4.2446~23.0123.023.0~~~~~~~-{}-{}-{}-{}D~-{}-{}-{}~-{}-{}-{}-{}-{}-{}-{}-{}~1:1011}\par
\texttt{F4.2447~012.123.123.0~~~~~~~-{}-{}-{}-{}D~-{}PC~-{}-{}-{}-{}++-{}+~ok}\par
\texttt{F4.2448~013.123.123.0~~~~~~~-{}-{}-{}-{}D~-{}PC~-{}-{}-{}-{}++-{}+~ok}\par
\texttt{F4.2449~123.123.123.0~~~~~~~-{}-{}-{}-{}D~-{}PC~-{}-{}-{}-{}++++~ok}\par
\texttt{F4.2450~01.0123.123.0~~~~~~~-{}-{}-{}-{}D~-{}PC~-{}-{}-{}-{}++-{}+~ok}\par
\texttt{F4.2451~02.0123.123.0~~~~~~~-{}-{}-{}-{}D~-{}PC~-{}-{}-{}-{}++-{}+~ok}\par
\texttt{F4.2452~12.0123.123.0~~~~~~~-{}-{}-{}-{}D~-{}PC~-{}-{}-{}-{}++-{}+~ok}\par
\texttt{F4.2453~03.0123.123.0~~~~~~~-{}-{}-{}-{}D~-{}-{}-{}~-{}-{}-{}-{}-{}-{}-{}-{}~1:1001}\par
\texttt{F4.2454~13.0123.123.0~~~~~~~-{}-{}-{}-{}D~-{}PC~-{}-{}-{}-{}++-{}+~ok}\par
\texttt{F4.2455~23.0123.123.0~~~~~~~-{}-{}-{}-{}D~-{}PC~-{}-{}-{}-{}++-{}+~ok}\par
\texttt{F4.2456~0.0123.0123.0~~~~~~~-{}-{}4-{}D~-{}-{}-{}~-{}-{}-{}+-{}-{}-{}-{}~1:1000}\par
\texttt{F4.2457~3.0123.0123.0~~~~~~~-{}-{}-{}-{}D~-{}-{}-{}~-{}-{}-{}-{}-{}-{}-{}-{}~1:1001}\par
\texttt{F4.2458~0123.01.01.01~~~~~~~-{}-{}-{}-{}D~-{}PC~-{}-{}-{}-{}++++~ok}\par
\texttt{F4.2459~0123.02.01.01~~~~~~~-{}-{}-{}-{}D~-{}PC~-{}-{}-{}-{}++++~ok}\par
\texttt{F4.2460~0123.12.01.01~~~~~~~-{}-{}-{}-{}D~-{}PC~-{}-{}-{}-{}++-{}+~ok}\par
\texttt{F4.2461~012.012.01.01~~~~~~~-{}-{}-{}-{}D~-{}P-{}~-{}-{}-{}-{}-{}+-{}-{}~3:1110}\par
\texttt{F4.2462~013.012.01.01~~~~~~~-{}-{}-{}-{}D~-{}PC~-{}-{}-{}-{}++++~ok}\par
\texttt{F4.2463~023.012.01.01~~~~~~~-{}-{}-{}-{}D~-{}PC~-{}-{}-{}-{}++++~ok}\par
\texttt{F4.2464~123.012.01.01~~~~~~~-{}-{}-{}-{}D~-{}PC~-{}-{}-{}-{}++++~ok}\par
\texttt{F4.2465~0123.23.01.01~~~~~~~-{}-{}-{}-{}D~-{}PC~-{}-{}-{}-{}++++~ok}\par
\texttt{F4.2466~023.023.01.01~~~~~~~-{}-{}-{}-{}D~-{}PC~-{}-{}-{}-{}++++~ok}\par
\texttt{F4.2467~123.023.01.01~~~~~~~-{}B-{}-{}D~-{}PC~+-{}-{}-{}++++~ok}\par
\texttt{F4.2468~023.123.01.01~~~~~~~-{}B-{}-{}D~-{}PC~+-{}-{}-{}++++~ok}\par
\texttt{F4.2469~0123.012.2.01~~~~~~~-{}-{}-{}-{}D~-{}-{}-{}~-{}-{}-{}-{}-{}-{}-{}-{}~0:0010}\par
\texttt{F4.2470~0123.013.2.01~~~~~~~-{}-{}-{}-{}D~-{}-{}-{}~-{}-{}-{}-{}-{}-{}-{}-{}~0:0010}\par
\texttt{F4.2471~0123.023.2.01~~~~~~~-{}-{}-{}-{}D~-{}-{}-{}~-{}-{}-{}-{}-{}-{}-{}-{}~0:0010}\par
\texttt{F4.2472~0123.123.2.01~~~~~~~-{}-{}-{}-{}D~-{}-{}-{}~-{}-{}-{}-{}-{}-{}-{}-{}~0:0010}\par
\texttt{F4.2473~0123.01.02.01~~~~~~~-{}-{}-{}-{}D~-{}PC~-{}-{}-{}-{}++++~ok}\par
\texttt{F4.2474~0123.02.02.01~~~~~~~-{}-{}-{}-{}D~-{}PC~-{}-{}-{}-{}++++~ok}\par
\texttt{F4.2475~0123.12.02.01~~~~~~~-{}-{}-{}-{}D~-{}PC~-{}-{}-{}-{}++-{}+~ok}\par
\texttt{F4.2476~012.012.02.01~~~~~~~-{}-{}-{}-{}D~-{}P-{}~-{}-{}-{}-{}-{}+-{}-{}~3:1110}\par
\texttt{F4.2477~013.012.02.01~~~~~~~-{}-{}-{}-{}D~-{}PC~-{}-{}-{}-{}++++~ok}\par
\texttt{F4.2478~023.012.02.01~~~~~~~-{}-{}-{}-{}D~-{}PC~-{}-{}-{}-{}++-{}+~ok}\par
\texttt{F4.2479~123.012.02.01~~~~~~~-{}-{}-{}-{}D~-{}PC~-{}-{}-{}-{}++++~ok}\par
\texttt{F4.2480~0123.03.02.01~~~~~~~-{}B-{}-{}D~-{}PC~+-{}-{}-{}++++~ok}\par
\texttt{F4.2481~0123.13.02.01~~~~~~~-{}-{}-{}-{}D~-{}PC~-{}-{}-{}-{}++-{}+~ok}\par
\texttt{F4.2482~012.013.02.01~~~~~~~-{}-{}-{}-{}D~-{}PC~-{}-{}-{}-{}++-{}+~ok}\par
\texttt{F4.2483~013.013.02.01~~~~~~~-{}-{}-{}-{}D~-{}-{}-{}~-{}-{}-{}-{}-{}-{}-{}-{}~2:1101}\par
\texttt{F4.2484~023.013.02.01~~~~~~~-{}-{}-{}-{}D~-{}PC~-{}-{}-{}-{}++-{}+~ok}\par
\texttt{F4.2485~123.013.02.01~~~~~~~-{}B-{}-{}D~-{}PC~+-{}-{}-{}++++~ok}\par
\texttt{F4.2486~0123.23.02.01~~~~~~~-{}-{}-{}-{}D~-{}PC~-{}-{}-{}-{}++++~ok}\par
\texttt{F4.2487~012.023.02.01~~~~~~~-{}-{}-{}-{}D~-{}PC~-{}-{}-{}-{}++-{}+~ok}\par
\texttt{F4.2488~013.023.02.01~~~~~~~-{}-{}-{}-{}D~-{}PC~-{}-{}-{}-{}++++~ok}\par
\texttt{F4.2489~023.023.02.01~~~~~~~-{}-{}-{}-{}D~-{}PC~-{}-{}-{}-{}++-{}+~ok}\par
\texttt{F4.2490~123.023.02.01~~~~~~~-{}-{}-{}-{}D~-{}PC~-{}-{}-{}-{}++++~ok}\par
\texttt{F4.2491~012.123.02.01~~~~~~~-{}-{}-{}-{}D~-{}PC~-{}-{}-{}-{}++-{}+~ok}\par
\texttt{F4.2492~013.123.02.01~~~~~~~-{}-{}-{}-{}D~-{}PC~-{}-{}-{}-{}++++~ok}\par
\texttt{F4.2493~023.123.02.01~~~~~~~-{}-{}-{}-{}D~-{}PC~-{}-{}-{}-{}++-{}+~ok}\par
\texttt{F4.2494~123.123.02.01~~~~~~~-{}-{}-{}-{}D~-{}PC~-{}-{}-{}-{}++++~ok}\par
\texttt{F4.2495~01.0123.02.01~~~~~~~-{}-{}-{}-{}D~-{}PC~-{}-{}-{}-{}++-{}+~ok}\par
\texttt{F4.2496~02.0123.02.01~~~~~~~-{}-{}-{}-{}D~-{}-{}-{}~-{}-{}-{}-{}-{}-{}-{}-{}~1:1010}\par
\texttt{F4.2497~12.0123.02.01~~~~~~~-{}-{}-{}-{}D~-{}PC~-{}-{}-{}-{}++-{}+~ok}\par
\texttt{F4.2498~03.0123.02.01~~~~~~~-{}-{}-{}-{}D~-{}PC~-{}-{}-{}-{}++-{}+~ok}\par
\texttt{F4.2499~13.0123.02.01~~~~~~~-{}-{}-{}-{}D~-{}PC~-{}-{}-{}-{}++++~ok}\par
\texttt{F4.2500~23.0123.02.01~~~~~~~-{}-{}-{}-{}D~-{}PC~-{}-{}-{}-{}++-{}+~ok}\par
\texttt{F4.2501~0123.1.012.01~~~~~~~-{}-{}-{}-{}D~-{}-{}-{}~-{}-{}-{}-{}-{}-{}-{}-{}~0:0100}\par
\texttt{F4.2502~012.01.012.01~~~~~~~-{}-{}-{}-{}D~-{}P-{}~-{}-{}-{}-{}-{}+-{}-{}~3:1110}\par
\texttt{F4.2503~013.01.012.01~~~~~~~-{}-{}-{}-{}D~-{}-{}-{}~-{}-{}-{}-{}-{}-{}-{}-{}~2:1101}\par
\texttt{F4.2504~023.01.012.01~~~~~~~-{}-{}-{}-{}D~-{}PC~-{}-{}-{}-{}++++~ok}\par
\texttt{F4.2505~123.01.012.01~~~~~~~-{}-{}-{}-{}D~-{}PC~-{}-{}-{}-{}++++~ok}\par
\texttt{F4.2506~012.02.012.01~~~~~~~-{}-{}-{}-{}D~-{}P-{}~-{}-{}-{}-{}-{}+-{}-{}~3:1110}\par
\texttt{F4.2507~013.02.012.01~~~~~~~-{}-{}-{}-{}D~-{}PC~-{}-{}-{}-{}++++~ok}\par
\texttt{F4.2508~023.02.012.01~~~~~~~-{}-{}-{}-{}D~-{}PC~-{}-{}-{}-{}++++~ok}\par
\texttt{F4.2509~123.02.012.01~~~~~~~-{}-{}-{}-{}D~-{}PC~-{}-{}-{}-{}++++~ok}\par
\texttt{F4.2510~012.12.012.01~~~~~~~-{}-{}-{}-{}D~-{}P-{}~-{}-{}-{}-{}-{}+-{}-{}~3:1110}\par
\texttt{F4.2511~013.12.012.01~~~~~~~-{}-{}-{}-{}D~-{}PC~-{}-{}-{}-{}++-{}+~ok}\par
\texttt{F4.2512~023.12.012.01~~~~~~~-{}-{}-{}-{}D~-{}PC~-{}-{}-{}-{}++-{}+~ok}\par
\texttt{F4.2513~123.12.012.01~~~~~~~-{}-{}-{}-{}D~-{}PC~-{}-{}-{}-{}++-{}+~ok}\par
\texttt{F4.2514~03.012.012.01~~~~~~~-{}-{}-{}-{}D~-{}PC~-{}-{}-{}-{}++-{}+~ok}\par
\texttt{F4.2515~13.012.012.01~~~~~~~-{}-{}-{}-{}D~-{}PC~-{}-{}-{}-{}++++~ok}\par
\texttt{F4.2516~23.012.012.01~~~~~~~-{}-{}-{}-{}D~-{}PC~-{}-{}-{}-{}++++~ok}\par
\texttt{F4.2517~013.03.012.01~~~~~~~-{}-{}-{}-{}D~-{}-{}-{}~-{}-{}-{}-{}-{}-{}-{}-{}~2:1101}\par
\texttt{F4.2518~023.03.012.01~~~~~~~-{}-{}-{}-{}D~-{}PC~-{}-{}-{}-{}++++~ok}\par
\texttt{F4.2519~123.03.012.01~~~~~~~-{}-{}-{}-{}D~-{}PC~-{}-{}-{}-{}++++~ok}\par
\texttt{F4.2520~013.13.012.01~~~~~~~-{}-{}-{}-{}D~-{}-{}-{}~-{}-{}-{}-{}-{}-{}-{}-{}~2:1101}\par
\texttt{F4.2521~023.13.012.01~~~~~~~-{}-{}-{}-{}D~-{}PC~-{}-{}-{}-{}++-{}+~ok}\par
\texttt{F4.2522~123.13.012.01~~~~~~~-{}-{}-{}-{}D~-{}PC~-{}-{}-{}-{}++-{}+~ok}\par
\texttt{F4.2523~23.013.012.01~~~~~~~-{}-{}-{}-{}D~-{}PC~-{}-{}-{}-{}++++~ok}\par
\texttt{F4.2524~023.23.012.01~~~~~~~-{}B-{}-{}D~-{}PC~+-{}-{}-{}++++~ok}\par
\texttt{F4.2525~123.23.012.01~~~~~~~-{}-{}-{}-{}D~-{}PC~-{}-{}-{}-{}++++~ok}\par
\texttt{F4.2526~0123.02.03.01~~~~~~~-{}-{}-{}-{}D~-{}PC~-{}-{}-{}-{}++++~ok}\par
\texttt{F4.2527~0123.12.03.01~~~~~~~-{}-{}-{}-{}D~-{}PC~-{}-{}-{}-{}++++~ok}\par
\texttt{F4.2528~012.012.03.01~~~~~~~-{}-{}-{}-{}D~-{}PC~-{}-{}-{}-{}++++~ok}\par
\texttt{F4.2529~013.012.03.01~~~~~~~-{}-{}-{}-{}D~-{}PC~-{}-{}-{}-{}++-{}+~ok}\par
\texttt{F4.2530~023.012.03.01~~~~~~~-{}-{}-{}-{}D~-{}PC~-{}-{}-{}-{}++++~ok}\par
\texttt{F4.2531~123.012.03.01~~~~~~~-{}-{}-{}-{}D~-{}PC~-{}-{}-{}-{}++++~ok}\par
\texttt{F4.2532~0123.23.03.01~~~~~~~-{}-{}-{}-{}D~-{}PC~-{}-{}-{}-{}++++~ok}\par
\texttt{F4.2533~012.023.03.01~~~~~~~-{}-{}-{}-{}D~-{}PC~-{}-{}-{}-{}++++~ok}\par
\texttt{F4.2534~013.023.03.01~~~~~~~-{}-{}-{}-{}D~-{}PC~-{}-{}-{}-{}++++~ok}\par
\texttt{F4.2535~023.023.03.01~~~~~~~-{}-{}-{}-{}D~-{}PC~-{}-{}-{}-{}++++~ok}\par
\texttt{F4.2536~123.023.03.01~~~~~~~-{}-{}-{}-{}D~-{}PC~-{}-{}-{}-{}++++~ok}\par
\texttt{F4.2537~012.123.03.01~~~~~~~-{}-{}-{}-{}D~-{}PC~-{}-{}-{}-{}++++~ok}\par
\texttt{F4.2538~013.123.03.01~~~~~~~-{}-{}-{}-{}D~-{}PC~-{}-{}-{}-{}++-{}+~ok}\par
\texttt{F4.2539~023.123.03.01~~~~~~~-{}-{}-{}-{}D~-{}PC~-{}-{}-{}-{}++++~ok}\par
\texttt{F4.2540~123.123.03.01~~~~~~~-{}-{}-{}-{}D~-{}PC~-{}-{}-{}-{}++++~ok}\par
\texttt{F4.2541~01.0123.03.01~~~~~~~-{}-{}-{}-{}D~-{}PC~-{}-{}-{}-{}++-{}+~ok}\par
\texttt{F4.2542~02.0123.03.01~~~~~~~-{}-{}-{}-{}D~-{}PC~-{}-{}-{}-{}++-{}+~ok}\par
\texttt{F4.2543~03.0123.03.01~~~~~~~-{}-{}-{}-{}D~-{}PC~-{}-{}-{}-{}++-{}+~ok}\par
\texttt{F4.2544~13.0123.03.01~~~~~~~-{}-{}-{}-{}D~-{}PC~-{}-{}-{}-{}++-{}+~ok}\par
\texttt{F4.2545~23.0123.03.01~~~~~~~-{}-{}-{}-{}D~-{}PC~-{}-{}-{}-{}++++~ok}\par
\texttt{F4.2546~0123.1.013.01~~~~~~~-{}-{}-{}-{}D~-{}-{}-{}~-{}-{}-{}-{}-{}-{}-{}-{}~0:0100}\par
\texttt{F4.2547~012.01.013.01~~~~~~~-{}-{}-{}-{}D~-{}PC~-{}-{}-{}-{}++-{}+~ok}\par
\texttt{F4.2548~013.01.013.01~~~~~~~-{}-{}-{}-{}D~-{}P-{}~-{}-{}-{}-{}-{}+-{}-{}~2:1101}\par
\texttt{F4.2549~023.01.013.01~~~~~~~-{}-{}-{}-{}D~-{}PC~-{}-{}-{}-{}++++~ok}\par
\texttt{F4.2550~123.01.013.01~~~~~~~-{}-{}-{}-{}D~-{}PC~-{}-{}-{}-{}++++~ok}\par
\texttt{F4.2551~012.12.013.01~~~~~~~-{}-{}-{}-{}D~-{}PC~-{}-{}-{}-{}++++~ok}\par
\texttt{F4.2552~013.12.013.01~~~~~~~-{}-{}-{}-{}D~-{}PC~-{}-{}-{}-{}++++~ok}\par
\texttt{F4.2553~023.12.013.01~~~~~~~-{}-{}-{}-{}D~-{}PC~-{}-{}-{}-{}++++~ok}\par
\texttt{F4.2554~123.12.013.01~~~~~~~-{}-{}-{}-{}D~-{}PC~-{}-{}-{}-{}++++~ok}\par
\texttt{F4.2555~03.012.013.01~~~~~~~-{}-{}-{}-{}D~-{}PC~-{}-{}-{}-{}++-{}+~ok}\par
\texttt{F4.2556~13.012.013.01~~~~~~~-{}-{}-{}-{}D~-{}PC~-{}-{}-{}-{}++++~ok}\par
\texttt{F4.2557~23.012.013.01~~~~~~~-{}-{}-{}-{}D~-{}PC~-{}-{}-{}-{}++++~ok}\par
\texttt{F4.2558~013.03.013.01~~~~~~~-{}-{}-{}-{}D~-{}P-{}~-{}-{}-{}-{}-{}+-{}-{}~2:1101}\par
\texttt{F4.2559~123.03.013.01~~~~~~~-{}-{}-{}-{}D~-{}PC~-{}-{}-{}-{}++++~ok}\par
\texttt{F4.2560~013.13.013.01~~~~~~~-{}-{}-{}-{}D~-{}P-{}~-{}-{}-{}-{}-{}+-{}-{}~2:1101}\par
\texttt{F4.2561~023.13.013.01~~~~~~~-{}-{}-{}-{}D~-{}PC~-{}-{}-{}-{}++-{}+~ok}\par
\texttt{F4.2562~123.13.013.01~~~~~~~-{}-{}-{}-{}D~-{}PC~-{}-{}-{}-{}++-{}+~ok}\par
\texttt{F4.2563~23.013.013.01~~~~~~~-{}-{}-{}-{}D~-{}PC~-{}-{}-{}-{}++++~ok}\par
\texttt{F4.2564~123.23.013.01~~~~~~~-{}-{}-{}-{}D~-{}PC~-{}-{}-{}-{}++++~ok}\par
\texttt{F4.2565~0123.01.23.01~~~~~~~-{}-{}-{}-{}D~-{}PC~-{}-{}-{}-{}++++~ok}\par
\texttt{F4.2566~0123.12.23.01~~~~~~~-{}-{}-{}-{}D~-{}PC~-{}-{}-{}-{}++-{}+~ok}\par
\texttt{F4.2567~012.012.23.01~~~~~~~-{}-{}-{}-{}D~-{}PC~-{}-{}-{}-{}++-{}+~ok}\par
\texttt{F4.2568~013.012.23.01~~~~~~~-{}-{}-{}-{}D~-{}PC~-{}-{}-{}-{}++++~ok}\par
\texttt{F4.2569~023.012.23.01~~~~~~~-{}-{}-{}-{}D~-{}PC~-{}-{}-{}-{}++++~ok}\par
\texttt{F4.2570~123.012.23.01~~~~~~~-{}-{}-{}-{}D~-{}PC~-{}-{}-{}-{}++++~ok}\par
\texttt{F4.2571~0123.13.23.01~~~~~~~-{}-{}-{}-{}D~-{}PC~-{}-{}-{}-{}++-{}+~ok}\par
\texttt{F4.2572~013.013.23.01~~~~~~~-{}-{}-{}-{}D~-{}-{}-{}~-{}-{}-{}-{}-{}-{}-{}-{}~2:1101}\par
\texttt{F4.2573~023.013.23.01~~~~~~~-{}-{}-{}-{}D~-{}PC~-{}-{}-{}-{}++++~ok}\par
\texttt{F4.2574~123.013.23.01~~~~~~~-{}-{}-{}-{}D~-{}PC~-{}-{}-{}-{}++++~ok}\par
\texttt{F4.2575~0123.23.23.01~~~~~~~-{}-{}-{}-{}D~-{}PC~-{}-{}-{}-{}++++~ok}\par
\texttt{F4.2576~023.023.23.01~~~~~~~-{}-{}-{}-{}D~-{}PC~-{}-{}-{}-{}++++~ok}\par
\texttt{F4.2577~123.023.23.01~~~~~~~-{}-{}-{}-{}D~-{}PC~-{}-{}-{}-{}++++~ok}\par
\texttt{F4.2578~023.123.23.01~~~~~~~-{}-{}-{}-{}D~-{}PC~-{}-{}-{}-{}++++~ok}\par
\texttt{F4.2579~0123.1.023.01~~~~~~~-{}-{}-{}-{}D~-{}-{}-{}~-{}-{}-{}-{}-{}-{}-{}-{}~0:0100}\par
\texttt{F4.2580~012.01.023.01~~~~~~~-{}-{}-{}-{}D~-{}PC~-{}-{}-{}-{}++-{}+~ok}\par
\texttt{F4.2581~013.01.023.01~~~~~~~-{}-{}-{}-{}D~-{}-{}-{}~-{}-{}-{}-{}-{}-{}-{}-{}~2:1101}\par
\texttt{F4.2582~023.01.023.01~~~~~~~-{}-{}-{}-{}D~-{}PC~-{}-{}-{}-{}++++~ok}\par
\texttt{F4.2583~123.01.023.01~~~~~~~-{}-{}-{}-{}D~-{}PC~-{}-{}-{}-{}++++~ok}\par
\texttt{F4.2584~012.12.023.01~~~~~~~-{}-{}-{}-{}D~-{}PC~-{}-{}-{}-{}++-{}+~ok}\par
\texttt{F4.2585~013.12.023.01~~~~~~~-{}-{}-{}-{}D~-{}PC~-{}-{}-{}-{}++++~ok}\par
\texttt{F4.2586~023.12.023.01~~~~~~~-{}-{}-{}-{}D~-{}PC~-{}-{}-{}-{}++++~ok}\par
\texttt{F4.2587~123.12.023.01~~~~~~~-{}-{}-{}-{}D~-{}PC~-{}-{}-{}-{}++++~ok}\par
\texttt{F4.2588~01.012.023.01~~~~~~~-{}-{}-{}-{}D~-{}PC~-{}-{}-{}-{}++-{}+~ok}\par
\texttt{F4.2589~02.012.023.01~~~~~~~-{}-{}-{}-{}D~-{}PC~-{}-{}-{}-{}++-{}+~ok}\par
\texttt{F4.2590~03.012.023.01~~~~~~~-{}-{}-{}-{}D~-{}PC~-{}-{}-{}-{}++-{}+~ok}\par
\texttt{F4.2591~13.012.023.01~~~~~~~-{}-{}-{}-{}D~-{}PC~-{}-{}-{}-{}++++~ok}\par
\texttt{F4.2592~23.012.023.01~~~~~~~-{}-{}-{}-{}D~-{}PC~-{}-{}-{}-{}++++~ok}\par
\texttt{F4.2593~012.13.023.01~~~~~~~-{}-{}-{}-{}D~-{}PC~-{}-{}-{}-{}++-{}+~ok}\par
\texttt{F4.2594~013.13.023.01~~~~~~~-{}-{}-{}-{}D~-{}-{}-{}~-{}-{}-{}-{}-{}-{}-{}-{}~2:1101}\par
\texttt{F4.2595~023.13.023.01~~~~~~~-{}-{}-{}-{}D~-{}PC~-{}-{}-{}-{}++-{}+~ok}\par
\texttt{F4.2596~123.13.023.01~~~~~~~-{}-{}-{}-{}D~-{}PC~-{}-{}-{}-{}++-{}+~ok}\par
\texttt{F4.2597~01.013.023.01~~~~~~~-{}-{}-{}-{}D~-{}-{}-{}~-{}-{}-{}-{}-{}-{}-{}-{}~2:1101}\par
\texttt{F4.2598~02.013.023.01~~~~~~~-{}-{}-{}-{}D~-{}PC~-{}-{}-{}-{}++-{}+~ok}\par
\texttt{F4.2599~03.013.023.01~~~~~~~-{}-{}-{}-{}D~-{}-{}-{}~-{}-{}-{}-{}-{}-{}-{}-{}~2:1101}\par
\texttt{F4.2600~13.013.023.01~~~~~~~-{}-{}-{}-{}D~-{}-{}-{}~-{}-{}-{}-{}-{}-{}-{}-{}~2:1101}\par
\texttt{F4.2601~23.013.023.01~~~~~~~-{}-{}-{}-{}D~-{}PC~-{}-{}-{}-{}++++~ok}\par
\texttt{F4.2602~012.23.023.01~~~~~~~-{}-{}-{}-{}D~-{}PC~-{}-{}-{}-{}++++~ok}\par
\texttt{F4.2603~01.023.023.01~~~~~~~-{}-{}-{}-{}D~-{}PC~-{}-{}-{}-{}++++~ok}\par
\texttt{F4.2604~02.023.023.01~~~~~~~-{}-{}-{}-{}D~-{}PC~-{}-{}-{}-{}++-{}+~ok}\par
\texttt{F4.2605~03.023.023.01~~~~~~~-{}-{}-{}-{}D~-{}PC~-{}-{}-{}-{}++-{}+~ok}\par
\texttt{F4.2606~13.023.023.01~~~~~~~-{}-{}-{}-{}D~-{}PC~-{}-{}-{}-{}++++~ok}\par
\texttt{F4.2607~01.123.023.01~~~~~~~-{}-{}-{}-{}D~-{}PC~-{}-{}-{}-{}++++~ok}\par
\texttt{F4.2608~02.123.023.01~~~~~~~-{}-{}-{}-{}D~-{}PC~-{}-{}-{}-{}++-{}+~ok}\par
\texttt{F4.2609~03.123.023.01~~~~~~~-{}-{}-{}-{}D~-{}PC~-{}-{}-{}-{}++-{}+~ok}\par
\texttt{F4.2610~0.0123.023.01~~~~~~~-{}-{}-{}-{}D~-{}-{}-{}~-{}-{}-{}-{}-{}-{}-{}-{}~1:1000}\par
\texttt{F4.2611~012.1.0123.01~~~~~~~-{}-{}-{}-{}D~-{}-{}-{}~-{}-{}-{}-{}-{}-{}-{}-{}~0:0100}\par
\texttt{F4.2612~013.1.0123.01~~~~~~~-{}-{}-{}-{}D~-{}-{}-{}~-{}-{}-{}-{}-{}-{}-{}-{}~0:0100}\par
\texttt{F4.2613~023.1.0123.01~~~~~~~-{}-{}-{}-{}D~-{}-{}-{}~-{}-{}-{}-{}-{}-{}-{}-{}~2:0100}\par
\texttt{F4.2614~123.1.0123.01~~~~~~~-{}-{}-{}-{}D~-{}-{}-{}~-{}-{}-{}-{}-{}-{}-{}-{}~0:0100}\par
\texttt{F4.2615~01.01.0123.01~~~~~~~-{}-{}4-{}D~-{}-{}-{}~-{}-{}-{}+-{}-{}-{}-{}~2:1100}\par
\texttt{F4.2616~02.01.0123.01~~~~~~~-{}-{}-{}-{}D~-{}PC~-{}-{}-{}-{}++-{}+~ok}\par
\texttt{F4.2617~03.01.0123.01~~~~~~~-{}-{}-{}-{}D~-{}-{}-{}~-{}-{}-{}-{}-{}-{}-{}-{}~2:1101}\par
\texttt{F4.2618~13.01.0123.01~~~~~~~-{}-{}-{}-{}D~-{}-{}-{}~-{}-{}-{}-{}-{}-{}-{}-{}~2:1101}\par
\texttt{F4.2619~02.12.0123.01~~~~~~~-{}-{}-{}-{}D~-{}PC~-{}-{}-{}-{}++-{}+~ok}\par
\texttt{F4.2620~03.12.0123.01~~~~~~~-{}-{}-{}-{}D~-{}PC~-{}-{}-{}-{}++-{}+~ok}\par
\texttt{F4.2621~13.12.0123.01~~~~~~~-{}-{}-{}-{}D~-{}PC~-{}-{}-{}-{}++++~ok}\par
\texttt{F4.2622~13.03.0123.01~~~~~~~-{}-{}-{}-{}D~-{}-{}-{}~-{}-{}-{}-{}-{}-{}-{}-{}~2:1101}\par
\texttt{F4.2623~03.13.0123.01~~~~~~~-{}-{}-{}-{}D~-{}-{}-{}~-{}-{}-{}-{}-{}-{}-{}-{}~2:1101}\par
\texttt{F4.2624~0123.01.2.012~~~~~~~-{}-{}-{}-{}D~-{}-{}-{}~-{}-{}-{}-{}-{}-{}-{}-{}~0:0010}\par
\texttt{F4.2625~0123.12.2.012~~~~~~~-{}-{}-{}-{}D~-{}-{}-{}~-{}-{}-{}-{}-{}-{}-{}-{}~0:0010}\par
\texttt{F4.2626~012.012.2.012~~~~~~~-{}-{}4-{}D~-{}-{}-{}~-{}-{}-{}+-{}-{}-{}-{}~0:0010}\par
\texttt{F4.2627~013.012.2.012~~~~~~~-{}-{}-{}-{}D~-{}-{}-{}~-{}-{}-{}-{}-{}-{}-{}-{}~1:0010}\par
\texttt{F4.2628~023.012.2.012~~~~~~~-{}-{}-{}-{}D~-{}-{}-{}~-{}-{}-{}-{}-{}-{}-{}-{}~0:0010}\par
\texttt{F4.2629~123.012.2.012~~~~~~~-{}-{}-{}-{}D~-{}-{}-{}~-{}-{}-{}-{}-{}-{}-{}-{}~0:0010}\par
\texttt{F4.2630~0123.13.2.012~~~~~~~-{}-{}-{}-{}D~-{}-{}-{}~-{}-{}-{}-{}-{}-{}-{}-{}~0:0010}\par
\texttt{F4.2631~013.013.2.012~~~~~~~-{}-{}-{}-{}D~-{}-{}-{}~-{}-{}-{}-{}-{}-{}-{}-{}~3:0010}\par
\texttt{F4.2632~023.013.2.012~~~~~~~-{}-{}-{}-{}D~-{}-{}-{}~-{}-{}-{}-{}-{}-{}-{}-{}~0:0010}\par
\texttt{F4.2633~123.013.2.012~~~~~~~-{}-{}-{}-{}D~-{}-{}-{}~-{}-{}-{}-{}-{}-{}-{}-{}~0:0010}\par
\texttt{F4.2634~123.023.2.012~~~~~~~-{}-{}-{}-{}D~-{}-{}-{}~-{}-{}-{}-{}-{}-{}-{}-{}~0:0010}\par
\texttt{F4.2635~023.123.2.012~~~~~~~-{}-{}-{}-{}D~-{}-{}-{}~-{}-{}-{}-{}-{}-{}-{}-{}~0:0010}\par
\texttt{F4.2636~012.01.02.012~~~~~~~-{}-{}-{}-{}D~-{}P-{}~-{}-{}-{}-{}-{}+-{}-{}~3:1110}\par
\texttt{F4.2637~013.01.02.012~~~~~~~-{}-{}-{}-{}D~-{}PC~-{}-{}-{}-{}++-{}+~ok}\par
\texttt{F4.2638~123.01.02.012~~~~~~~-{}-{}-{}-{}D~-{}PC~-{}-{}-{}-{}++++~ok}\par
\texttt{F4.2639~012.12.02.012~~~~~~~-{}-{}-{}-{}D~-{}P-{}~-{}-{}-{}-{}-{}+-{}-{}~3:1110}\par
\texttt{F4.2640~013.12.02.012~~~~~~~-{}-{}-{}-{}D~-{}PC~-{}-{}-{}-{}++-{}+~ok}\par
\texttt{F4.2641~023.12.02.012~~~~~~~-{}-{}-{}-{}D~-{}PC~-{}-{}-{}-{}++-{}+~ok}\par
\texttt{F4.2642~123.12.02.012~~~~~~~-{}-{}-{}-{}D~-{}PC~-{}-{}-{}-{}++-{}+~ok}\par
\texttt{F4.2643~02.012.02.012~~~~~~~-{}-{}4-{}D~-{}-{}-{}~-{}-{}-{}+-{}-{}-{}-{}~1:1010}\par
\texttt{F4.2644~03.012.02.012~~~~~~~-{}-{}-{}-{}D~-{}PC~-{}-{}-{}-{}++-{}+~ok}\par
\texttt{F4.2645~012.13.02.012~~~~~~~-{}-{}-{}-{}D~-{}PC~-{}-{}-{}-{}++-{}+~ok}\par
\texttt{F4.2646~013.13.02.012~~~~~~~-{}-{}-{}-{}D~-{}PC~-{}-{}-{}-{}++++~ok}\par
\texttt{F4.2647~023.13.02.012~~~~~~~-{}-{}-{}-{}D~-{}PC~-{}-{}-{}-{}++-{}+~ok}\par
\texttt{F4.2648~123.13.02.012~~~~~~~-{}-{}-{}-{}D~-{}PC~-{}-{}-{}-{}++++~ok}\par
\texttt{F4.2649~02.013.02.012~~~~~~~-{}-{}-{}-{}D~-{}-{}-{}~-{}-{}-{}-{}-{}-{}-{}-{}~1:1010}\par
\texttt{F4.2650~03.013.02.012~~~~~~~-{}-{}-{}-{}D~-{}PC~-{}-{}-{}-{}++-{}+~ok}\par
\texttt{F4.2651~02.023.02.012~~~~~~~-{}-{}-{}-{}D~-{}-{}-{}~-{}-{}-{}-{}-{}-{}-{}-{}~1:1010}\par
\texttt{F4.2652~03.123.02.012~~~~~~~-{}-{}-{}-{}D~-{}PC~-{}-{}-{}-{}++-{}+~ok}\par
\texttt{F4.2653~03.13.012.012~~~~~~~-{}-{}-{}-{}D~-{}PC~-{}-{}-{}-{}++++~ok}\par
\texttt{F4.2654~013.13.23.012~~~~~~~-{}-{}-{}-{}D~-{}PC~-{}-{}-{}-{}++++~ok}\par
\texttt{F4.2655~0123.0123.2.3~~~~~~~T-{}4-{}D~-{}-{}-{}~-{}-{}-{}+-{}-{}-{}-{}~0:0010}\par
\texttt{F4.2656~0123.012.02.3~~~~~~~T-{}-{}-{}D~-{}-{}-{}~-{}-{}-{}-{}-{}-{}-{}-{}~0:0001}\par
\texttt{F4.2657~0123.013.02.3~~~~~~~T-{}-{}-{}D~-{}-{}-{}~-{}-{}-{}-{}-{}-{}-{}-{}~0:0001}\par
\texttt{F4.2658~0123.123.02.3~~~~~~~T-{}-{}-{}D~-{}-{}-{}~-{}-{}-{}-{}-{}-{}-{}-{}~0:0001}\par
\texttt{F4.2659~012.0123.02.3~~~~~~~T-{}-{}-{}D~-{}-{}-{}~-{}-{}-{}-{}-{}-{}-{}-{}~1:0001}\par
\texttt{F4.2660~013.0123.02.3~~~~~~~T-{}-{}-{}D~-{}-{}-{}~-{}-{}-{}-{}-{}-{}-{}-{}~0:0001}\par
\texttt{F4.2661~023.0123.02.3~~~~~~~T-{}-{}-{}D~-{}-{}-{}~-{}-{}-{}-{}-{}-{}-{}-{}~0:0001}\par
\texttt{F4.2662~012.012.012.3~~~~~~~TB45D~UPC~++++++++~ok}\par
\texttt{F4.2663~013.012.012.3~~~~~~~T-{}-{}-{}D~-{}-{}-{}~-{}-{}-{}-{}-{}-{}-{}-{}~0:0001}\par
\texttt{F4.2664~0123.13.012.3~~~~~~~T-{}-{}-{}D~-{}-{}-{}~-{}-{}-{}-{}-{}-{}-{}-{}~0:0001}\par
\texttt{F4.2665~013.013.012.3~~~~~~~T-{}-{}-{}D~-{}-{}-{}~-{}-{}-{}-{}-{}-{}-{}-{}~0:0001}\par
\texttt{F4.2666~023.013.012.3~~~~~~~T-{}-{}-{}D~-{}-{}-{}~-{}-{}-{}-{}-{}-{}-{}-{}~0:0001}\par
\texttt{F4.2667~023.123.012.3~~~~~~~T-{}-{}-{}D~-{}-{}-{}~-{}-{}-{}-{}-{}-{}-{}-{}~0:0001}\par
\texttt{F4.2668~0123.013.23.3~~~~~~~T-{}-{}-{}D~-{}-{}-{}~-{}-{}-{}-{}-{}-{}-{}-{}~0:0001}\par
\texttt{F4.2669~0123.123.23.3~~~~~~~T-{}4-{}D~-{}-{}-{}~-{}-{}-{}+-{}-{}-{}-{}~0:0001}\par
\texttt{F4.2670~013.013.023.3~~~~~~~T-{}-{}-{}D~-{}-{}-{}~-{}-{}-{}-{}-{}-{}-{}-{}~0:0001}\par
\texttt{F4.2671~013.123.023.3~~~~~~~T-{}-{}-{}D~-{}-{}-{}~-{}-{}-{}-{}-{}-{}-{}-{}~0:0001}\par
\texttt{F4.2672~0123.01.02.03~~~~~~~TB-{}-{}D~-{}PC~+-{}-{}-{}++++~ok}\par
\texttt{F4.2673~0123.12.02.03~~~~~~~T-{}-{}-{}D~-{}PC~-{}-{}-{}-{}++-{}+~ok}\par
\texttt{F4.2674~012.012.02.03~~~~~~~T-{}-{}-{}D~-{}-{}-{}~-{}-{}-{}-{}-{}-{}-{}-{}~3:1110}\par
\texttt{F4.2675~013.012.02.03~~~~~~~T-{}-{}-{}D~-{}PC~-{}-{}-{}-{}++-{}+~ok}\par
\texttt{F4.2676~023.012.02.03~~~~~~~T-{}-{}-{}D~-{}-{}-{}~-{}-{}-{}-{}-{}-{}-{}-{}~1:1011}\par
\texttt{F4.2677~012.123.02.03~~~~~~~T-{}-{}-{}D~-{}PC~-{}-{}-{}-{}++-{}+~ok}\par
\texttt{F4.2678~023.123.02.03~~~~~~~T-{}-{}-{}D~-{}-{}-{}~-{}-{}-{}-{}-{}-{}-{}-{}~1:1011}\par
\texttt{F4.2679~01.0123.02.03~~~~~~~T-{}-{}-{}D~-{}PC~-{}-{}-{}-{}++-{}+~ok}\par
\texttt{F4.2680~02.0123.02.03~~~~~~~T-{}-{}-{}D~-{}-{}-{}~-{}-{}-{}-{}-{}-{}-{}-{}~1:1010}\par
\texttt{F4.2681~0123.12.12.03~~~~~~~T-{}-{}-{}D~-{}-{}-{}~-{}-{}-{}-{}-{}-{}-{}-{}~0:0110}\par
\texttt{F4.2682~012.012.12.03~~~~~~~T-{}-{}-{}D~-{}-{}-{}~-{}-{}-{}-{}-{}-{}-{}-{}~3:1110}\par
\texttt{F4.2683~013.012.12.03~~~~~~~TB-{}-{}D~-{}PC~+-{}-{}-{}++-{}+~ok}\par
\texttt{F4.2684~023.012.12.03~~~~~~~T-{}-{}-{}D~-{}PC~-{}-{}-{}-{}++-{}+~ok}\par
\texttt{F4.2685~0123.13.12.03~~~~~~~T-{}-{}-{}D~-{}PC~-{}-{}-{}-{}++-{}+~ok}\par
\texttt{F4.2686~012.013.12.03~~~~~~~T-{}-{}-{}D~-{}PC~-{}-{}-{}-{}++++~ok}\par
\texttt{F4.2687~023.013.12.03~~~~~~~T-{}-{}-{}D~-{}PC~-{}-{}-{}-{}++-{}+~ok}\par
\texttt{F4.2688~023.123.12.03~~~~~~~T-{}-{}-{}D~-{}PC~-{}-{}-{}-{}++-{}+~ok}\par
\texttt{F4.2689~01.012.012.03~~~~~~~T-{}-{}-{}D~-{}-{}-{}~-{}-{}-{}-{}-{}-{}-{}-{}~3:1110}\par
\texttt{F4.2690~03.012.012.03~~~~~~~T-{}-{}-{}D~-{}-{}-{}~-{}-{}-{}-{}-{}-{}-{}-{}~1:1001}\par
\texttt{F4.2691~012.13.012.03~~~~~~~T-{}-{}-{}D~-{}PC~-{}-{}-{}-{}++-{}+~ok}\par
\texttt{F4.2692~013.13.012.03~~~~~~~T-{}-{}-{}D~-{}-{}-{}~-{}-{}-{}-{}-{}-{}-{}-{}~2:1101}\par
\texttt{F4.2693~023.13.012.03~~~~~~~T-{}-{}-{}D~-{}PC~-{}-{}-{}-{}++-{}+~ok}\par
\texttt{F4.2694~01.013.012.03~~~~~~~T-{}-{}-{}D~-{}-{}-{}~-{}-{}-{}-{}-{}-{}-{}-{}~2:1101}\par
\texttt{F4.2695~02.013.012.03~~~~~~~T-{}-{}-{}D~-{}PC~-{}-{}-{}-{}++-{}+~ok}\par
\texttt{F4.2696~03.013.012.03~~~~~~~T-{}-{}-{}D~-{}-{}-{}~-{}-{}-{}-{}-{}-{}-{}-{}~1:1001}\par
\texttt{F4.2697~01.123.012.03~~~~~~~T-{}-{}-{}D~-{}PC~-{}-{}-{}-{}++-{}+~ok}\par
\texttt{F4.2698~03.123.012.03~~~~~~~T-{}-{}-{}D~-{}-{}-{}~-{}-{}-{}-{}-{}-{}-{}-{}~1:1001}\par
\texttt{F4.2699~012.123.23.03~~~~~~~T-{}-{}-{}D~-{}PC~-{}-{}-{}-{}++-{}+~ok}\par
\texttt{F4.2700~023.123.23.03~~~~~~~T-{}-{}-{}D~-{}-{}-{}~-{}-{}-{}-{}-{}-{}-{}-{}~1:1011}\par
\texttt{F4.2701~02.0123.23.03~~~~~~~T-{}-{}-{}D~-{}-{}-{}~-{}-{}-{}-{}-{}-{}-{}-{}~1:1011}\par
\texttt{F4.2702~03.013.023.03~~~~~~~T-{}4-{}D~-{}-{}-{}~-{}-{}-{}+-{}-{}-{}-{}~1:1001}\par
\texttt{F4.2703~0123.0123.012.-{}~~~~~-{}-{}-{}-{}-{}~-{}-{}-{}~-{}-{}-{}-{}-{}-{}-{}-{}~0:0000}\par
\texttt{F4.2704~0123.0123.013.-{}~~~~~-{}-{}-{}-{}-{}~-{}-{}-{}~-{}-{}-{}-{}-{}-{}-{}-{}~0:0000}\par
\texttt{F4.2705~0123.0123.023.-{}~~~~~-{}-{}-{}-{}-{}~-{}-{}-{}~-{}-{}-{}-{}-{}-{}-{}-{}~0:0000}\par
\texttt{F4.2706~0123.0123.01.0~~~~~~-{}-{}-{}-{}D~-{}PC~-{}-{}-{}-{}++++~ok}\par
\texttt{F4.2707~0123.0123.02.0~~~~~~-{}-{}-{}-{}D~-{}PC~-{}-{}-{}-{}++++~ok}\par
\texttt{F4.2708~0123.0123.12.0~~~~~~-{}-{}-{}-{}D~-{}PC~-{}-{}-{}-{}++++~ok}\par
\texttt{F4.2709~0123.012.012.0~~~~~~-{}B-{}-{}D~-{}PC~+-{}-{}-{}++++~ok}\par
\texttt{F4.2710~0123.013.012.0~~~~~~-{}-{}-{}-{}D~-{}PC~-{}-{}-{}-{}++++~ok}\par
\texttt{F4.2711~0123.023.012.0~~~~~~-{}-{}-{}-{}D~-{}PC~-{}-{}-{}-{}++++~ok}\par
\texttt{F4.2712~0123.123.012.0~~~~~~-{}-{}-{}-{}D~-{}PC~-{}-{}-{}-{}++++~ok}\par
\texttt{F4.2713~012.0123.012.0~~~~~~-{}-{}-{}-{}D~-{}PC~-{}-{}-{}-{}++-{}+~ok}\par
\texttt{F4.2714~013.0123.012.0~~~~~~-{}-{}-{}-{}D~-{}PC~-{}-{}-{}-{}++-{}+~ok}\par
\texttt{F4.2715~023.0123.012.0~~~~~~-{}-{}-{}-{}D~-{}PC~-{}-{}-{}-{}++-{}+~ok}\par
\texttt{F4.2716~123.0123.012.0~~~~~~-{}-{}-{}-{}D~-{}PC~-{}-{}-{}-{}++++~ok}\par
\texttt{F4.2717~0123.0123.03.0~~~~~~-{}-{}-{}-{}D~-{}PC~-{}-{}-{}-{}++++~ok}\par
\texttt{F4.2718~0123.0123.13.0~~~~~~-{}-{}-{}-{}D~-{}PC~-{}-{}-{}-{}++++~ok}\par
\texttt{F4.2719~0123.013.013.0~~~~~~-{}-{}-{}-{}D~-{}PC~-{}-{}-{}-{}++++~ok}\par
\texttt{F4.2720~0123.023.013.0~~~~~~-{}-{}-{}-{}D~-{}PC~-{}-{}-{}-{}++++~ok}\par
\texttt{F4.2721~0123.123.013.0~~~~~~-{}-{}-{}-{}D~-{}PC~-{}-{}-{}-{}++++~ok}\par
\texttt{F4.2722~012.0123.013.0~~~~~~-{}-{}-{}-{}D~-{}PC~-{}-{}-{}-{}++-{}+~ok}\par
\texttt{F4.2723~013.0123.013.0~~~~~~-{}-{}-{}-{}D~-{}PC~-{}-{}-{}-{}++-{}+~ok}\par
\texttt{F4.2724~023.0123.013.0~~~~~~-{}-{}-{}-{}D~-{}PC~-{}-{}-{}-{}++-{}+~ok}\par
\texttt{F4.2725~123.0123.013.0~~~~~~-{}-{}-{}-{}D~-{}PC~-{}-{}-{}-{}++++~ok}\par
\texttt{F4.2726~0123.0123.23.0~~~~~~-{}-{}-{}-{}D~-{}PC~-{}-{}-{}-{}++-{}+~ok}\par
\texttt{F4.2727~0123.013.023.0~~~~~~-{}-{}-{}-{}D~-{}PC~-{}-{}-{}-{}++++~ok}\par
\texttt{F4.2728~0123.123.023.0~~~~~~-{}-{}-{}-{}D~-{}PC~-{}-{}-{}-{}++++~ok}\par
\texttt{F4.2729~012.0123.023.0~~~~~~-{}-{}-{}-{}D~-{}PC~-{}-{}-{}-{}++-{}+~ok}\par
\texttt{F4.2730~013.0123.023.0~~~~~~-{}-{}-{}-{}D~-{}PC~-{}-{}-{}-{}++-{}+~ok}\par
\texttt{F4.2731~023.0123.023.0~~~~~~-{}-{}-{}-{}D~-{}-{}-{}~-{}-{}-{}-{}-{}-{}-{}-{}~1:1011}\par
\texttt{F4.2732~123.0123.023.0~~~~~~-{}-{}-{}-{}D~-{}PC~-{}-{}-{}-{}++++~ok}\par
\texttt{F4.2733~0123.123.123.0~~~~~~-{}-{}-{}-{}D~-{}PC~-{}-{}-{}-{}++++~ok}\par
\texttt{F4.2734~012.0123.123.0~~~~~~-{}-{}-{}-{}D~-{}PC~-{}-{}-{}-{}++-{}+~ok}\par
\texttt{F4.2735~013.0123.123.0~~~~~~-{}-{}-{}-{}D~-{}PC~-{}-{}-{}-{}++-{}+~ok}\par
\texttt{F4.2736~023.0123.123.0~~~~~~-{}-{}-{}-{}D~-{}PC~-{}-{}-{}-{}++-{}+~ok}\par
\texttt{F4.2737~123.0123.123.0~~~~~~-{}-{}-{}-{}D~-{}PC~-{}-{}-{}-{}++++~ok}\par
\texttt{F4.2738~01.0123.0123.0~~~~~~-{}-{}-{}-{}D~-{}PC~-{}-{}-{}-{}++-{}+~ok}\par
\texttt{F4.2739~12.0123.0123.0~~~~~~-{}-{}-{}-{}D~-{}PC~-{}-{}-{}-{}++-{}+~ok}\par
\texttt{F4.2740~03.0123.0123.0~~~~~~-{}-{}-{}-{}D~-{}-{}-{}~-{}-{}-{}-{}-{}-{}-{}-{}~1:1001}\par
\texttt{F4.2741~13.0123.0123.0~~~~~~-{}-{}-{}-{}D~-{}PC~-{}-{}-{}-{}++-{}+~ok}\par
\texttt{F4.2742~0123.012.01.01~~~~~~-{}-{}-{}-{}D~-{}PC~-{}-{}-{}-{}++++~ok}\par
\texttt{F4.2743~0123.023.01.01~~~~~~-{}B-{}-{}D~-{}PC~+-{}-{}-{}++++~ok}\par
\texttt{F4.2744~0123.123.01.01~~~~~~-{}-{}-{}-{}D~-{}PC~-{}-{}-{}-{}++++~ok}\par
\texttt{F4.2745~0123.0123.2.01~~~~~~-{}-{}-{}-{}D~-{}-{}-{}~-{}-{}-{}-{}-{}-{}-{}-{}~0:0010}\par
\texttt{F4.2746~0123.012.02.01~~~~~~-{}-{}-{}-{}D~-{}PC~-{}-{}-{}-{}++++~ok}\par
\texttt{F4.2747~0123.013.02.01~~~~~~-{}B-{}-{}D~-{}PC~+-{}-{}-{}++++~ok}\par
\texttt{F4.2748~0123.023.02.01~~~~~~-{}-{}-{}-{}D~-{}PC~-{}-{}-{}-{}++++~ok}\par
\texttt{F4.2749~0123.123.02.01~~~~~~-{}-{}-{}-{}D~-{}PC~-{}-{}-{}-{}++++~ok}\par
\texttt{F4.2750~012.0123.02.01~~~~~~-{}-{}-{}-{}D~-{}PC~-{}-{}-{}-{}++-{}+~ok}\par
\texttt{F4.2751~013.0123.02.01~~~~~~-{}-{}-{}-{}D~-{}PC~-{}-{}-{}-{}++++~ok}\par
\texttt{F4.2752~023.0123.02.01~~~~~~-{}-{}-{}-{}D~-{}PC~-{}-{}-{}-{}++-{}+~ok}\par
\texttt{F4.2753~123.0123.02.01~~~~~~-{}-{}-{}-{}D~-{}PC~-{}-{}-{}-{}++++~ok}\par
\texttt{F4.2754~0123.01.012.01~~~~~~-{}-{}-{}-{}D~-{}PC~-{}-{}-{}-{}++++~ok}\par
\texttt{F4.2755~0123.02.012.01~~~~~~-{}-{}-{}-{}D~-{}PC~-{}-{}-{}-{}++++~ok}\par
\texttt{F4.2756~0123.12.012.01~~~~~~-{}-{}-{}-{}D~-{}PC~-{}-{}-{}-{}++-{}+~ok}\par
\texttt{F4.2757~012.012.012.01~~~~~~-{}-{}-{}5D~UP-{}~-{}++-{}-{}+-{}-{}~3:1110}\par
\texttt{F4.2758~013.012.012.01~~~~~~-{}-{}-{}-{}D~-{}PC~-{}-{}-{}-{}++++~ok}\par
\texttt{F4.2759~023.012.012.01~~~~~~-{}-{}-{}-{}D~-{}PC~-{}-{}-{}-{}++++~ok}\par
\texttt{F4.2760~123.012.012.01~~~~~~-{}-{}-{}-{}D~-{}PC~-{}-{}-{}-{}++++~ok}\par
\texttt{F4.2761~0123.03.012.01~~~~~~-{}-{}-{}-{}D~-{}PC~-{}-{}-{}-{}++++~ok}\par
\texttt{F4.2762~0123.13.012.01~~~~~~-{}-{}-{}-{}D~-{}PC~-{}-{}-{}-{}++-{}+~ok}\par
\texttt{F4.2763~013.013.012.01~~~~~~-{}-{}-{}-{}D~-{}-{}-{}~-{}-{}-{}-{}-{}-{}-{}-{}~2:1101}\par
\texttt{F4.2764~023.013.012.01~~~~~~-{}-{}-{}-{}D~-{}PC~-{}-{}-{}-{}++++~ok}\par
\texttt{F4.2765~123.013.012.01~~~~~~-{}-{}-{}-{}D~-{}PC~-{}-{}-{}-{}++++~ok}\par
\texttt{F4.2766~0123.23.012.01~~~~~~-{}-{}-{}-{}D~-{}PC~-{}-{}-{}-{}++++~ok}\par
\texttt{F4.2767~023.023.012.01~~~~~~-{}-{}-{}-{}D~-{}PC~-{}-{}-{}-{}++++~ok}\par
\texttt{F4.2768~123.023.012.01~~~~~~-{}B-{}-{}D~-{}PC~+-{}-{}-{}++++~ok}\par
\texttt{F4.2769~023.123.012.01~~~~~~-{}B-{}-{}D~-{}PC~+-{}-{}-{}++++~ok}\par
\texttt{F4.2770~0123.012.03.01~~~~~~-{}-{}-{}-{}D~-{}PC~-{}-{}-{}-{}++++~ok}\par
\texttt{F4.2771~0123.023.03.01~~~~~~-{}-{}-{}-{}D~-{}PC~-{}-{}-{}-{}++++~ok}\par
\texttt{F4.2772~0123.123.03.01~~~~~~-{}-{}-{}-{}D~-{}PC~-{}-{}-{}-{}++++~ok}\par
\texttt{F4.2773~012.0123.03.01~~~~~~-{}-{}-{}-{}D~-{}PC~-{}-{}-{}-{}++++~ok}\par
\texttt{F4.2774~013.0123.03.01~~~~~~-{}-{}-{}-{}D~-{}PC~-{}-{}-{}-{}++-{}+~ok}\par
\texttt{F4.2775~023.0123.03.01~~~~~~-{}-{}-{}-{}D~-{}PC~-{}-{}-{}-{}++++~ok}\par
\texttt{F4.2776~123.0123.03.01~~~~~~-{}-{}-{}-{}D~-{}PC~-{}-{}-{}-{}++++~ok}\par
\texttt{F4.2777~0123.01.013.01~~~~~~-{}-{}-{}-{}D~-{}PC~-{}-{}-{}-{}++++~ok}\par
\texttt{F4.2778~0123.12.013.01~~~~~~-{}-{}-{}-{}D~-{}PC~-{}-{}-{}-{}++++~ok}\par
\texttt{F4.2779~012.012.013.01~~~~~~-{}-{}-{}-{}D~-{}PC~-{}-{}-{}-{}++++~ok}\par
\texttt{F4.2780~013.012.013.01~~~~~~-{}-{}-{}-{}D~-{}PC~-{}-{}-{}-{}++++~ok}\par
\texttt{F4.2781~023.012.013.01~~~~~~-{}-{}-{}-{}D~-{}PC~-{}-{}-{}-{}++++~ok}\par
\texttt{F4.2782~123.012.013.01~~~~~~-{}-{}-{}-{}D~-{}PC~-{}-{}-{}-{}++++~ok}\par
\texttt{F4.2783~0123.03.013.01~~~~~~-{}-{}-{}-{}D~-{}PC~-{}-{}-{}-{}++++~ok}\par
\texttt{F4.2784~0123.13.013.01~~~~~~-{}-{}-{}-{}D~-{}PC~-{}-{}-{}-{}++-{}+~ok}\par
\texttt{F4.2785~013.013.013.01~~~~~~-{}-{}-{}-{}D~-{}P-{}~-{}-{}-{}-{}-{}+-{}-{}~2:1101}\par
\texttt{F4.2786~023.013.013.01~~~~~~-{}-{}-{}-{}D~-{}PC~-{}-{}-{}-{}++++~ok}\par
\texttt{F4.2787~123.013.013.01~~~~~~-{}-{}-{}-{}D~-{}PC~-{}-{}-{}-{}++++~ok}\par
\texttt{F4.2788~0123.23.013.01~~~~~~-{}-{}-{}-{}D~-{}PC~-{}-{}-{}-{}++++~ok}\par
\texttt{F4.2789~023.023.013.01~~~~~~-{}-{}-{}-{}D~-{}PC~-{}-{}-{}-{}++++~ok}\par
\texttt{F4.2790~123.023.013.01~~~~~~-{}-{}-{}-{}D~-{}PC~-{}-{}-{}-{}++++~ok}\par
\texttt{F4.2791~023.123.013.01~~~~~~-{}-{}-{}-{}D~-{}PC~-{}-{}-{}-{}++++~ok}\par
\texttt{F4.2792~0123.012.23.01~~~~~~-{}-{}-{}-{}D~-{}PC~-{}-{}-{}-{}++++~ok}\par
\texttt{F4.2793~0123.013.23.01~~~~~~-{}-{}-{}-{}D~-{}PC~-{}-{}-{}-{}++++~ok}\par
\texttt{F4.2794~0123.023.23.01~~~~~~-{}-{}-{}-{}D~-{}PC~-{}-{}-{}-{}++++~ok}\par
\texttt{F4.2795~0123.123.23.01~~~~~~-{}-{}-{}-{}D~-{}PC~-{}-{}-{}-{}++++~ok}\par
\texttt{F4.2796~0123.01.023.01~~~~~~-{}-{}-{}-{}D~-{}PC~-{}-{}-{}-{}++++~ok}\par
\texttt{F4.2797~0123.12.023.01~~~~~~-{}-{}-{}-{}D~-{}PC~-{}-{}-{}-{}++++~ok}\par
\texttt{F4.2798~012.012.023.01~~~~~~-{}-{}-{}-{}D~-{}PC~-{}-{}-{}-{}++-{}+~ok}\par
\texttt{F4.2799~013.012.023.01~~~~~~-{}-{}-{}-{}D~-{}PC~-{}-{}-{}-{}++++~ok}\par
\texttt{F4.2800~023.012.023.01~~~~~~-{}-{}-{}-{}D~-{}PC~-{}-{}-{}-{}++++~ok}\par
\texttt{F4.2801~123.012.023.01~~~~~~-{}-{}-{}-{}D~-{}PC~-{}-{}-{}-{}++++~ok}\par
\texttt{F4.2802~0123.13.023.01~~~~~~-{}-{}-{}-{}D~-{}PC~-{}-{}-{}-{}++-{}+~ok}\par
\texttt{F4.2803~012.013.023.01~~~~~~-{}-{}-{}-{}D~-{}PC~-{}-{}-{}-{}++++~ok}\par
\texttt{F4.2804~013.013.023.01~~~~~~-{}-{}-{}-{}D~-{}-{}-{}~-{}-{}-{}-{}-{}-{}-{}-{}~2:1101}\par
\texttt{F4.2805~023.013.023.01~~~~~~-{}-{}-{}-{}D~-{}PC~-{}-{}-{}-{}++++~ok}\par
\texttt{F4.2806~123.013.023.01~~~~~~-{}-{}-{}-{}D~-{}PC~-{}-{}-{}-{}++++~ok}\par
\texttt{F4.2807~0123.23.023.01~~~~~~-{}-{}-{}-{}D~-{}PC~-{}-{}-{}-{}++++~ok}\par
\texttt{F4.2808~012.023.023.01~~~~~~-{}-{}-{}-{}D~-{}PC~-{}-{}-{}-{}++++~ok}\par
\texttt{F4.2809~013.023.023.01~~~~~~-{}-{}-{}-{}D~-{}PC~-{}-{}-{}-{}++++~ok}\par
\texttt{F4.2810~023.023.023.01~~~~~~-{}-{}-{}-{}D~-{}PC~-{}-{}-{}-{}++++~ok}\par
\texttt{F4.2811~123.023.023.01~~~~~~-{}-{}-{}-{}D~-{}PC~-{}-{}-{}-{}++++~ok}\par
\texttt{F4.2812~012.123.023.01~~~~~~-{}-{}-{}-{}D~-{}PC~-{}-{}-{}-{}++++~ok}\par
\texttt{F4.2813~013.123.023.01~~~~~~-{}-{}-{}-{}D~-{}PC~-{}-{}-{}-{}++++~ok}\par
\texttt{F4.2814~023.123.023.01~~~~~~-{}-{}-{}-{}D~-{}PC~-{}-{}-{}-{}++++~ok}\par
\texttt{F4.2815~123.123.023.01~~~~~~-{}-{}-{}-{}D~-{}PC~-{}-{}-{}-{}++++~ok}\par
\texttt{F4.2816~01.0123.023.01~~~~~~-{}-{}-{}-{}D~-{}PC~-{}-{}-{}-{}++++~ok}\par
\texttt{F4.2817~02.0123.023.01~~~~~~-{}-{}-{}-{}D~-{}PC~-{}-{}-{}-{}++-{}+~ok}\par
\texttt{F4.2818~03.0123.023.01~~~~~~-{}-{}-{}-{}D~-{}PC~-{}-{}-{}-{}++-{}+~ok}\par
\texttt{F4.2819~13.0123.023.01~~~~~~-{}-{}-{}-{}D~-{}PC~-{}-{}-{}-{}++++~ok}\par
\texttt{F4.2820~23.0123.023.01~~~~~~-{}-{}-{}-{}D~-{}PC~-{}-{}-{}-{}++++~ok}\par
\texttt{F4.2821~0123.1.0123.01~~~~~~-{}-{}-{}-{}D~-{}-{}-{}~-{}-{}-{}-{}-{}-{}-{}-{}~0:0100}\par
\texttt{F4.2822~012.01.0123.01~~~~~~-{}-{}-{}-{}D~-{}PC~-{}-{}-{}-{}++-{}+~ok}\par
\texttt{F4.2823~013.01.0123.01~~~~~~-{}-{}-{}-{}D~-{}-{}-{}~-{}-{}-{}-{}-{}-{}-{}-{}~2:1101}\par
\texttt{F4.2824~023.01.0123.01~~~~~~-{}-{}-{}-{}D~-{}PC~-{}-{}-{}-{}++++~ok}\par
\texttt{F4.2825~123.01.0123.01~~~~~~-{}-{}-{}-{}D~-{}PC~-{}-{}-{}-{}++++~ok}\par
\texttt{F4.2826~012.12.0123.01~~~~~~-{}-{}-{}-{}D~-{}PC~-{}-{}-{}-{}++-{}+~ok}\par
\texttt{F4.2827~013.12.0123.01~~~~~~-{}-{}-{}-{}D~-{}PC~-{}-{}-{}-{}++++~ok}\par
\texttt{F4.2828~023.12.0123.01~~~~~~-{}-{}-{}-{}D~-{}PC~-{}-{}-{}-{}++++~ok}\par
\texttt{F4.2829~123.12.0123.01~~~~~~-{}-{}-{}-{}D~-{}PC~-{}-{}-{}-{}++++~ok}\par
\texttt{F4.2830~03.012.0123.01~~~~~~-{}-{}-{}-{}D~-{}PC~-{}-{}-{}-{}++-{}+~ok}\par
\texttt{F4.2831~13.012.0123.01~~~~~~-{}-{}-{}-{}D~-{}PC~-{}-{}-{}-{}++++~ok}\par
\texttt{F4.2832~013.03.0123.01~~~~~~-{}-{}-{}-{}D~-{}-{}-{}~-{}-{}-{}-{}-{}-{}-{}-{}~2:1101}\par
\texttt{F4.2833~123.03.0123.01~~~~~~-{}-{}-{}-{}D~-{}PC~-{}-{}-{}-{}++++~ok}\par
\texttt{F4.2834~013.13.0123.01~~~~~~-{}-{}-{}-{}D~-{}-{}-{}~-{}-{}-{}-{}-{}-{}-{}-{}~2:1101}\par
\texttt{F4.2835~023.13.0123.01~~~~~~-{}-{}-{}-{}D~-{}PC~-{}-{}-{}-{}++-{}+~ok}\par
\texttt{F4.2836~123.13.0123.01~~~~~~-{}-{}-{}-{}D~-{}PC~-{}-{}-{}-{}++-{}+~ok}\par
\texttt{F4.2837~0123.012.2.012~~~~~~-{}-{}-{}-{}D~-{}-{}-{}~-{}-{}-{}-{}-{}-{}-{}-{}~0:0010}\par
\texttt{F4.2838~0123.013.2.012~~~~~~-{}-{}-{}-{}D~-{}-{}-{}~-{}-{}-{}-{}-{}-{}-{}-{}~0:0010}\par
\texttt{F4.2839~0123.023.2.012~~~~~~-{}-{}-{}-{}D~-{}-{}-{}~-{}-{}-{}-{}-{}-{}-{}-{}~0:0010}\par
\texttt{F4.2840~0123.123.2.012~~~~~~-{}-{}-{}-{}D~-{}-{}-{}~-{}-{}-{}-{}-{}-{}-{}-{}~0:0010}\par
\texttt{F4.2841~0123.01.02.012~~~~~~-{}-{}-{}-{}D~-{}PC~-{}-{}-{}-{}++++~ok}\par
\texttt{F4.2842~0123.12.02.012~~~~~~-{}-{}-{}-{}D~-{}PC~-{}-{}-{}-{}++-{}+~ok}\par
\texttt{F4.2843~012.012.02.012~~~~~~-{}-{}-{}-{}D~-{}P-{}~-{}-{}-{}-{}-{}+-{}-{}~3:1110}\par
\texttt{F4.2844~013.012.02.012~~~~~~-{}-{}-{}-{}D~-{}PC~-{}-{}-{}-{}++++~ok}\par
\texttt{F4.2845~023.012.02.012~~~~~~-{}-{}-{}-{}D~-{}PC~-{}-{}-{}-{}++-{}+~ok}\par
\texttt{F4.2846~123.012.02.012~~~~~~-{}-{}-{}-{}D~-{}PC~-{}-{}-{}-{}++++~ok}\par
\texttt{F4.2847~0123.13.02.012~~~~~~-{}-{}-{}-{}D~-{}PC~-{}-{}-{}-{}++++~ok}\par
\texttt{F4.2848~012.013.02.012~~~~~~-{}-{}-{}-{}D~-{}PC~-{}-{}-{}-{}++-{}+~ok}\par
\texttt{F4.2849~013.013.02.012~~~~~~-{}-{}-{}-{}D~-{}PC~-{}-{}-{}-{}++++~ok}\par
\texttt{F4.2850~023.013.02.012~~~~~~-{}-{}-{}-{}D~-{}PC~-{}-{}-{}-{}++-{}+~ok}\par
\texttt{F4.2851~123.013.02.012~~~~~~-{}-{}-{}-{}D~-{}PC~-{}-{}-{}-{}++++~ok}\par
\texttt{F4.2852~012.023.02.012~~~~~~-{}-{}-{}-{}D~-{}PC~-{}-{}-{}-{}++-{}+~ok}\par
\texttt{F4.2853~013.023.02.012~~~~~~-{}-{}-{}-{}D~-{}PC~-{}-{}-{}-{}++++~ok}\par
\texttt{F4.2854~123.023.02.012~~~~~~-{}-{}-{}-{}D~-{}PC~-{}-{}-{}-{}++++~ok}\par
\texttt{F4.2855~012.123.02.012~~~~~~-{}-{}-{}-{}D~-{}PC~-{}-{}-{}-{}++-{}+~ok}\par
\texttt{F4.2856~013.123.02.012~~~~~~-{}-{}-{}-{}D~-{}PC~-{}-{}-{}-{}++++~ok}\par
\texttt{F4.2857~023.123.02.012~~~~~~-{}-{}-{}-{}D~-{}PC~-{}-{}-{}-{}++-{}+~ok}\par
\texttt{F4.2858~123.123.02.012~~~~~~-{}-{}-{}-{}D~-{}PC~-{}-{}-{}-{}++++~ok}\par
\texttt{F4.2859~02.0123.02.012~~~~~~-{}-{}-{}-{}D~-{}-{}-{}~-{}-{}-{}-{}-{}-{}-{}-{}~1:1010}\par
\texttt{F4.2860~03.0123.02.012~~~~~~-{}-{}-{}-{}D~-{}PC~-{}-{}-{}-{}++-{}+~ok}\par
\texttt{F4.2861~03.012.012.012~~~~~~-{}-{}-{}-{}D~-{}PC~-{}-{}-{}-{}++++~ok}\par
\texttt{F4.2862~013.13.012.012~~~~~~-{}-{}-{}-{}D~-{}PC~-{}-{}-{}-{}++++~ok}\par
\texttt{F4.2863~023.13.012.012~~~~~~-{}-{}-{}-{}D~-{}PC~-{}-{}-{}-{}++++~ok}\par
\texttt{F4.2864~0123.13.23.012~~~~~~-{}-{}-{}-{}D~-{}PC~-{}-{}-{}-{}++++~ok}\par
\texttt{F4.2865~013.013.23.012~~~~~~-{}B-{}-{}D~-{}PC~+-{}-{}-{}++++~ok}\par
\texttt{F4.2866~023.013.23.012~~~~~~-{}-{}-{}-{}D~-{}PC~-{}-{}-{}-{}++++~ok}\par
\texttt{F4.2867~023.123.23.012~~~~~~-{}-{}-{}-{}D~-{}PC~-{}-{}-{}-{}++++~ok}\par
\texttt{F4.2868~0123.0123.02.3~~~~~~T-{}-{}-{}D~-{}-{}-{}~-{}-{}-{}-{}-{}-{}-{}-{}~0:0001}\par
\texttt{F4.2869~0123.012.012.3~~~~~~T-{}-{}-{}D~-{}-{}-{}~-{}-{}-{}-{}-{}-{}-{}-{}~0:0001}\par
\texttt{F4.2870~0123.013.012.3~~~~~~T-{}-{}-{}D~-{}-{}-{}~-{}-{}-{}-{}-{}-{}-{}-{}~0:0001}\par
\texttt{F4.2871~0123.123.012.3~~~~~~T-{}-{}-{}D~-{}-{}-{}~-{}-{}-{}-{}-{}-{}-{}-{}~0:0001}\par
\texttt{F4.2872~0123.0123.23.3~~~~~~T-{}4-{}D~-{}-{}-{}~-{}-{}-{}+-{}-{}-{}-{}~0:0001}\par
\texttt{F4.2873~0123.013.023.3~~~~~~T-{}-{}-{}D~-{}-{}-{}~-{}-{}-{}-{}-{}-{}-{}-{}~0:0001}\par
\texttt{F4.2874~0123.123.023.3~~~~~~T-{}-{}-{}D~-{}-{}-{}~-{}-{}-{}-{}-{}-{}-{}-{}~0:0001}\par
\texttt{F4.2875~023.0123.023.3~~~~~~T-{}4-{}D~-{}-{}-{}~-{}-{}-{}+-{}-{}-{}-{}~0:0001}\par
\texttt{F4.2876~0123.012.02.03~~~~~~T-{}-{}-{}D~-{}PC~-{}-{}-{}-{}++++~ok}\par
\texttt{F4.2877~0123.123.02.03~~~~~~T-{}-{}-{}D~-{}PC~-{}-{}-{}-{}++++~ok}\par
\texttt{F4.2878~012.0123.02.03~~~~~~T-{}-{}-{}D~-{}PC~-{}-{}-{}-{}++-{}+~ok}\par
\texttt{F4.2879~023.0123.02.03~~~~~~T-{}-{}-{}D~-{}-{}-{}~-{}-{}-{}-{}-{}-{}-{}-{}~1:1011}\par
\texttt{F4.2880~0123.012.12.03~~~~~~T-{}-{}-{}D~-{}PC~-{}-{}-{}-{}++-{}+~ok}\par
\texttt{F4.2881~0123.013.12.03~~~~~~T-{}-{}-{}D~-{}PC~-{}-{}-{}-{}++++~ok}\par
\texttt{F4.2882~0123.123.12.03~~~~~~T-{}-{}-{}D~-{}PC~-{}-{}-{}-{}++-{}+~ok}\par
\texttt{F4.2883~012.012.012.03~~~~~~T-{}-{}-{}D~-{}-{}-{}~-{}-{}-{}-{}-{}-{}-{}-{}~3:1110}\par
\texttt{F4.2884~013.012.012.03~~~~~~T-{}-{}-{}D~-{}PC~-{}-{}-{}-{}++-{}+~ok}\par
\texttt{F4.2885~0123.13.012.03~~~~~~T-{}-{}-{}D~-{}PC~-{}-{}-{}-{}++-{}+~ok}\par
\texttt{F4.2886~012.013.012.03~~~~~~T-{}-{}-{}D~-{}PC~-{}-{}-{}-{}++++~ok}\par
\texttt{F4.2887~013.013.012.03~~~~~~T-{}-{}-{}D~-{}-{}-{}~-{}-{}-{}-{}-{}-{}-{}-{}~2:1101}\par
\texttt{F4.2888~023.013.012.03~~~~~~T-{}-{}-{}D~-{}PC~-{}-{}-{}-{}++-{}+~ok}\par
\texttt{F4.2889~012.123.012.03~~~~~~T-{}-{}-{}D~-{}PC~-{}-{}-{}-{}++++~ok}\par
\texttt{F4.2890~013.123.012.03~~~~~~T-{}-{}-{}D~-{}PC~-{}-{}-{}-{}++-{}+~ok}\par
\texttt{F4.2891~023.123.012.03~~~~~~T-{}-{}-{}D~-{}PC~-{}-{}-{}-{}++-{}+~ok}\par
\texttt{F4.2892~01.0123.012.03~~~~~~T-{}-{}-{}D~-{}PC~-{}-{}-{}-{}++-{}+~ok}\par
\texttt{F4.2893~02.0123.012.03~~~~~~T-{}-{}-{}D~-{}PC~-{}-{}-{}-{}++-{}+~ok}\par
\texttt{F4.2894~03.0123.012.03~~~~~~T-{}-{}-{}D~-{}-{}-{}~-{}-{}-{}-{}-{}-{}-{}-{}~1:1001}\par
\texttt{F4.2895~0123.123.23.03~~~~~~T-{}-{}-{}D~-{}PC~-{}-{}-{}-{}++-{}+~ok}\par
\texttt{F4.2896~012.0123.23.03~~~~~~T-{}-{}-{}D~-{}PC~-{}-{}-{}-{}++-{}+~ok}\par
\texttt{F4.2897~023.0123.23.03~~~~~~T-{}-{}-{}D~-{}-{}-{}~-{}-{}-{}-{}-{}-{}-{}-{}~1:1011}\par
\texttt{F4.2898~012.013.023.03~~~~~~T-{}-{}-{}D~-{}PC~-{}-{}-{}-{}++++~ok}\par
\texttt{F4.2899~013.013.023.03~~~~~~T-{}-{}-{}D~-{}-{}-{}~-{}-{}-{}-{}-{}-{}-{}-{}~2:1101}\par
\texttt{F4.2900~012.123.023.03~~~~~~T-{}-{}-{}D~-{}PC~-{}-{}-{}-{}++++~ok}\par
\texttt{F4.2901~013.123.023.03~~~~~~T-{}-{}-{}D~-{}PC~-{}-{}-{}-{}++-{}+~ok}\par
\texttt{F4.2902~023.123.023.03~~~~~~T-{}-{}-{}D~-{}-{}-{}~-{}-{}-{}-{}-{}-{}-{}-{}~1:1011}\par
\texttt{F4.2903~03.0123.023.03~~~~~~T-{}4-{}D~-{}-{}-{}~-{}-{}-{}+-{}-{}-{}-{}~1:1001}\par
\texttt{F4.2904~012.123.123.03~~~~~~T-{}-{}-{}D~-{}PC~-{}-{}-{}-{}++++~ok}\par
\texttt{F4.2905~013.123.123.03~~~~~~T-{}-{}-{}D~-{}PC~-{}-{}-{}-{}++-{}+~ok}\par
\texttt{F4.2906~0123.0123.0123.-{}~~~~-{}-{}4-{}-{}~-{}-{}-{}~-{}-{}-{}+-{}-{}-{}-{}~0:0000}\par
\texttt{F4.2907~0123.0123.012.0~~~~~-{}-{}-{}-{}D~-{}PC~-{}-{}-{}-{}++++~ok}\par
\texttt{F4.2908~0123.0123.013.0~~~~~-{}-{}-{}-{}D~-{}PC~-{}-{}-{}-{}++++~ok}\par
\texttt{F4.2909~0123.0123.023.0~~~~~-{}-{}-{}-{}D~-{}PC~-{}-{}-{}-{}++++~ok}\par
\texttt{F4.2910~0123.0123.123.0~~~~~-{}-{}-{}-{}D~-{}PC~-{}-{}-{}-{}++++~ok}\par
\texttt{F4.2911~012.0123.0123.0~~~~~-{}-{}-{}-{}D~-{}PC~-{}-{}-{}-{}++-{}+~ok}\par
\texttt{F4.2912~013.0123.0123.0~~~~~-{}-{}-{}-{}D~-{}PC~-{}-{}-{}-{}++-{}+~ok}\par
\texttt{F4.2913~123.0123.0123.0~~~~~-{}-{}-{}-{}D~-{}PC~-{}-{}-{}-{}++++~ok}\par
\texttt{F4.2914~0123.0123.01.01~~~~~-{}B-{}-{}D~-{}PC~+-{}-{}-{}++++~ok}\par
\texttt{F4.2915~0123.0123.02.01~~~~~-{}-{}-{}-{}D~-{}PC~-{}-{}-{}-{}++++~ok}\par
\texttt{F4.2916~0123.012.012.01~~~~~-{}-{}-{}-{}D~-{}PC~-{}-{}-{}-{}++++~ok}\par
\texttt{F4.2917~0123.013.012.01~~~~~-{}-{}-{}-{}D~-{}PC~-{}-{}-{}-{}++++~ok}\par
\texttt{F4.2918~0123.023.012.01~~~~~-{}B-{}-{}D~-{}PC~+-{}-{}-{}++++~ok}\par
\texttt{F4.2919~0123.123.012.01~~~~~-{}-{}-{}-{}D~-{}PC~-{}-{}-{}-{}++++~ok}\par
\texttt{F4.2920~0123.0123.03.01~~~~~-{}-{}-{}-{}D~-{}PC~-{}-{}-{}-{}++++~ok}\par
\texttt{F4.2921~0123.012.013.01~~~~~-{}-{}-{}-{}D~-{}PC~-{}-{}-{}-{}++++~ok}\par
\texttt{F4.2922~0123.013.013.01~~~~~-{}-{}-{}-{}D~-{}PC~-{}-{}-{}-{}++++~ok}\par
\texttt{F4.2923~0123.023.013.01~~~~~-{}-{}-{}-{}D~-{}PC~-{}-{}-{}-{}++++~ok}\par
\texttt{F4.2924~0123.123.013.01~~~~~-{}-{}-{}-{}D~-{}PC~-{}-{}-{}-{}++++~ok}\par
\texttt{F4.2925~0123.0123.23.01~~~~~-{}-{}-{}-{}D~-{}PC~-{}-{}-{}-{}++++~ok}\par
\texttt{F4.2926~0123.012.023.01~~~~~-{}-{}-{}-{}D~-{}PC~-{}-{}-{}-{}++++~ok}\par
\texttt{F4.2927~0123.013.023.01~~~~~-{}-{}-{}-{}D~-{}PC~-{}-{}-{}-{}++++~ok}\par
\texttt{F4.2928~0123.023.023.01~~~~~-{}-{}-{}-{}D~-{}PC~-{}-{}-{}-{}++++~ok}\par
\texttt{F4.2929~0123.123.023.01~~~~~-{}-{}-{}-{}D~-{}PC~-{}-{}-{}-{}++++~ok}\par
\texttt{F4.2930~012.0123.023.01~~~~~-{}-{}-{}-{}D~-{}PC~-{}-{}-{}-{}++++~ok}\par
\texttt{F4.2931~013.0123.023.01~~~~~-{}-{}-{}-{}D~-{}PC~-{}-{}-{}-{}++++~ok}\par
\texttt{F4.2932~023.0123.023.01~~~~~-{}-{}-{}-{}D~-{}PC~-{}-{}-{}-{}++++~ok}\par
\texttt{F4.2933~123.0123.023.01~~~~~-{}-{}-{}-{}D~-{}PC~-{}-{}-{}-{}++++~ok}\par
\texttt{F4.2934~0123.01.0123.01~~~~~-{}-{}-{}-{}D~-{}PC~-{}-{}-{}-{}++++~ok}\par
\texttt{F4.2935~0123.12.0123.01~~~~~-{}-{}-{}-{}D~-{}PC~-{}-{}-{}-{}++++~ok}\par
\texttt{F4.2936~012.012.0123.01~~~~~-{}-{}-{}-{}D~-{}PC~-{}-{}-{}-{}++-{}+~ok}\par
\texttt{F4.2937~013.012.0123.01~~~~~-{}-{}-{}-{}D~-{}PC~-{}-{}-{}-{}++++~ok}\par
\texttt{F4.2938~023.012.0123.01~~~~~-{}-{}-{}-{}D~-{}PC~-{}-{}-{}-{}++++~ok}\par
\texttt{F4.2939~123.012.0123.01~~~~~-{}-{}-{}-{}D~-{}PC~-{}-{}-{}-{}++++~ok}\par
\texttt{F4.2940~0123.03.0123.01~~~~~-{}-{}-{}-{}D~-{}PC~-{}-{}-{}-{}++++~ok}\par
\texttt{F4.2941~0123.13.0123.01~~~~~-{}-{}-{}-{}D~-{}PC~-{}-{}-{}-{}++-{}+~ok}\par
\texttt{F4.2942~013.013.0123.01~~~~~-{}-{}-{}-{}D~-{}-{}-{}~-{}-{}-{}-{}-{}-{}-{}-{}~2:1101}\par
\texttt{F4.2943~023.013.0123.01~~~~~-{}-{}-{}-{}D~-{}PC~-{}-{}-{}-{}++++~ok}\par
\texttt{F4.2944~123.013.0123.01~~~~~-{}-{}-{}-{}D~-{}PC~-{}-{}-{}-{}++++~ok}\par
\texttt{F4.2945~0123.23.0123.01~~~~~-{}-{}-{}-{}D~-{}PC~-{}-{}-{}-{}++++~ok}\par
\texttt{F4.2946~023.023.0123.01~~~~~-{}-{}-{}-{}D~-{}PC~-{}-{}-{}-{}++++~ok}\par
\texttt{F4.2947~123.023.0123.01~~~~~-{}-{}-{}-{}D~-{}PC~-{}-{}-{}-{}++++~ok}\par
\texttt{F4.2948~023.123.0123.01~~~~~-{}-{}-{}-{}D~-{}PC~-{}-{}-{}-{}++++~ok}\par
\texttt{F4.2949~0123.0123.2.012~~~~~-{}-{}-{}-{}D~-{}-{}-{}~-{}-{}-{}-{}-{}-{}-{}-{}~0:0010}\par
\texttt{F4.2950~0123.012.02.012~~~~~-{}-{}-{}-{}D~-{}PC~-{}-{}-{}-{}++++~ok}\par
\texttt{F4.2951~0123.013.02.012~~~~~-{}-{}-{}-{}D~-{}PC~-{}-{}-{}-{}++++~ok}\par
\texttt{F4.2952~0123.023.02.012~~~~~-{}-{}-{}-{}D~-{}PC~-{}-{}-{}-{}++++~ok}\par
\texttt{F4.2953~0123.123.02.012~~~~~-{}-{}-{}-{}D~-{}PC~-{}-{}-{}-{}++++~ok}\par
\texttt{F4.2954~012.0123.02.012~~~~~-{}-{}-{}-{}D~-{}PC~-{}-{}-{}-{}++-{}+~ok}\par
\texttt{F4.2955~013.0123.02.012~~~~~-{}-{}-{}-{}D~-{}PC~-{}-{}-{}-{}++++~ok}\par
\texttt{F4.2956~023.0123.02.012~~~~~-{}-{}-{}-{}D~-{}PC~-{}-{}-{}-{}++-{}+~ok}\par
\texttt{F4.2957~123.0123.02.012~~~~~-{}-{}-{}-{}D~-{}PC~-{}-{}-{}-{}++++~ok}\par
\texttt{F4.2958~012.012.012.012~~~~~-{}-{}45D~UP-{}~-{}+++-{}+-{}-{}~3:1110}\par
\texttt{F4.2959~013.012.012.012~~~~~-{}-{}-{}-{}D~-{}PC~-{}-{}-{}-{}++++~ok}\par
\texttt{F4.2960~123.012.012.012~~~~~-{}-{}-{}-{}D~-{}PC~-{}-{}-{}-{}++++~ok}\par
\texttt{F4.2961~0123.13.012.012~~~~~-{}-{}-{}-{}D~-{}PC~-{}-{}-{}-{}++++~ok}\par
\texttt{F4.2962~013.013.012.012~~~~~-{}-{}-{}-{}D~-{}PC~-{}-{}-{}-{}++++~ok}\par
\texttt{F4.2963~023.013.012.012~~~~~-{}-{}-{}-{}D~-{}PC~-{}-{}-{}-{}++++~ok}\par
\texttt{F4.2964~123.013.012.012~~~~~-{}-{}-{}-{}D~-{}PC~-{}-{}-{}-{}++++~ok}\par
\texttt{F4.2965~023.023.012.012~~~~~-{}-{}-{}-{}D~-{}PC~-{}-{}-{}-{}++++~ok}\par
\texttt{F4.2966~123.023.012.012~~~~~-{}-{}-{}-{}D~-{}PC~-{}-{}-{}-{}++++~ok}\par
\texttt{F4.2967~023.123.012.012~~~~~-{}-{}-{}-{}D~-{}PC~-{}-{}-{}-{}++++~ok}\par
\texttt{F4.2968~123.023.013.012~~~~~-{}B-{}-{}D~-{}PC~+-{}-{}-{}++++~ok}\par
\texttt{F4.2969~023.123.013.012~~~~~-{}B-{}-{}D~-{}PC~+-{}-{}-{}++++~ok}\par
\texttt{F4.2970~0123.013.23.012~~~~~-{}-{}-{}-{}D~-{}PC~-{}-{}-{}-{}++++~ok}\par
\texttt{F4.2971~0123.123.23.012~~~~~-{}-{}-{}-{}D~-{}PC~-{}-{}-{}-{}++++~ok}\par
\texttt{F4.2972~013.013.023.012~~~~~-{}-{}-{}-{}D~-{}PC~-{}-{}-{}-{}++++~ok}\par
\texttt{F4.2973~013.123.023.012~~~~~-{}-{}-{}-{}D~-{}PC~-{}-{}-{}-{}++++~ok}\par
\texttt{F4.2974~0123.0123.012.3~~~~~T-{}-{}-{}D~-{}-{}-{}~-{}-{}-{}-{}-{}-{}-{}-{}~0:0001}\par
\texttt{F4.2975~0123.0123.023.3~~~~~T-{}-{}-{}D~-{}-{}-{}~-{}-{}-{}-{}-{}-{}-{}-{}~0:0001}\par
\texttt{F4.2976~0123.0123.02.03~~~~~T-{}-{}-{}D~-{}PC~-{}-{}-{}-{}++++~ok}\par
\texttt{F4.2977~0123.0123.12.03~~~~~T-{}-{}-{}D~-{}PC~-{}-{}-{}-{}++++~ok}\par
\texttt{F4.2978~0123.012.012.03~~~~~TB-{}-{}D~-{}PC~+-{}-{}-{}++++~ok}\par
\texttt{F4.2979~0123.013.012.03~~~~~T-{}-{}-{}D~-{}PC~-{}-{}-{}-{}++++~ok}\par
\texttt{F4.2980~0123.123.012.03~~~~~T-{}-{}-{}D~-{}PC~-{}-{}-{}-{}++++~ok}\par
\texttt{F4.2981~012.0123.012.03~~~~~T-{}-{}-{}D~-{}PC~-{}-{}-{}-{}++++~ok}\par
\texttt{F4.2982~013.0123.012.03~~~~~T-{}-{}-{}D~-{}PC~-{}-{}-{}-{}++-{}+~ok}\par
\texttt{F4.2983~023.0123.012.03~~~~~T-{}-{}-{}D~-{}PC~-{}-{}-{}-{}++-{}+~ok}\par
\texttt{F4.2984~0123.0123.23.03~~~~~T-{}-{}-{}D~-{}PC~-{}-{}-{}-{}++-{}+~ok}\par
\texttt{F4.2985~0123.013.023.03~~~~~T-{}-{}-{}D~-{}PC~-{}-{}-{}-{}++++~ok}\par
\texttt{F4.2986~0123.123.023.03~~~~~T-{}-{}-{}D~-{}PC~-{}-{}-{}-{}++++~ok}\par
\texttt{F4.2987~012.0123.023.03~~~~~T-{}-{}-{}D~-{}PC~-{}-{}-{}-{}++++~ok}\par
\texttt{F4.2988~013.0123.023.03~~~~~T-{}-{}-{}D~-{}PC~-{}-{}-{}-{}++-{}+~ok}\par
\texttt{F4.2989~023.0123.023.03~~~~~T-{}-{}-{}D~-{}-{}-{}~-{}-{}-{}-{}-{}-{}-{}-{}~1:1011}\par
\texttt{F4.2990~0123.123.123.03~~~~~T-{}-{}-{}D~-{}PC~-{}-{}-{}-{}++++~ok}\par
\texttt{F4.2991~012.0123.123.03~~~~~T-{}-{}-{}D~-{}PC~-{}-{}-{}-{}++++~ok}\par
\texttt{F4.2992~013.0123.123.03~~~~~T-{}-{}-{}D~-{}PC~-{}-{}-{}-{}++-{}+~ok}\par
\texttt{F4.2993~023.0123.123.03~~~~~T-{}-{}-{}D~-{}PC~-{}-{}-{}-{}++-{}+~ok}\par
\texttt{F4.2994~03.0123.0123.03~~~~~T-{}4-{}D~-{}-{}-{}~-{}-{}-{}+-{}-{}-{}-{}~1:1001}\par
\texttt{F4.2995~012.012.012.013~~~~~T-{}-{}-{}D~-{}-{}-{}~-{}-{}-{}-{}-{}-{}-{}-{}~3:1110}\par
\texttt{F4.2996~013.012.012.013~~~~~T-{}-{}-{}D~-{}PC~-{}-{}-{}-{}++++~ok}\par
\texttt{F4.2997~023.012.012.013~~~~~T-{}-{}-{}D~-{}PC~-{}-{}-{}-{}++++~ok}\par
\texttt{F4.2998~023.123.012.013~~~~~TB-{}-{}D~-{}PC~+-{}-{}-{}++++~ok}\par
\texttt{F4.2999~012.012.023.013~~~~~T-{}-{}-{}D~-{}PC~-{}-{}-{}-{}++++~ok}\par
\texttt{F4.3000~012.123.023.013~~~~~T-{}-{}-{}D~-{}PC~-{}-{}-{}-{}++++~ok}\par
\texttt{F4.3001~0123.0123.0123.0~~~~-{}-{}-{}-{}D~-{}PC~-{}-{}-{}-{}++++~ok}\par
\texttt{F4.3002~0123.0123.012.01~~~~-{}B-{}-{}D~-{}PC~+-{}-{}-{}++++~ok}\par
\texttt{F4.3003~0123.0123.013.01~~~~-{}-{}-{}-{}D~-{}PC~-{}-{}-{}-{}++++~ok}\par
\texttt{F4.3004~0123.0123.023.01~~~~-{}-{}-{}-{}D~-{}PC~-{}-{}-{}-{}++++~ok}\par
\texttt{F4.3005~0123.012.0123.01~~~~-{}-{}-{}-{}D~-{}PC~-{}-{}-{}-{}++++~ok}\par
\texttt{F4.3006~0123.013.0123.01~~~~-{}-{}-{}-{}D~-{}PC~-{}-{}-{}-{}++++~ok}\par
\texttt{F4.3007~0123.023.0123.01~~~~-{}-{}-{}-{}D~-{}PC~-{}-{}-{}-{}++++~ok}\par
\texttt{F4.3008~0123.123.0123.01~~~~-{}-{}-{}-{}D~-{}PC~-{}-{}-{}-{}++++~ok}\par
\texttt{F4.3009~0123.0123.02.012~~~~-{}-{}-{}-{}D~-{}PC~-{}-{}-{}-{}++++~ok}\par
\texttt{F4.3010~0123.012.012.012~~~~-{}-{}-{}-{}D~-{}PC~-{}-{}-{}-{}++++~ok}\par
\texttt{F4.3011~0123.013.012.012~~~~-{}-{}-{}-{}D~-{}PC~-{}-{}-{}-{}++++~ok}\par
\texttt{F4.3012~0123.023.012.012~~~~-{}-{}-{}-{}D~-{}PC~-{}-{}-{}-{}++++~ok}\par
\texttt{F4.3013~0123.123.012.012~~~~-{}-{}-{}-{}D~-{}PC~-{}-{}-{}-{}++++~ok}\par
\texttt{F4.3014~0123.023.013.012~~~~-{}B-{}-{}D~-{}PC~+-{}-{}-{}++++~ok}\par
\texttt{F4.3015~0123.123.013.012~~~~-{}-{}-{}-{}D~-{}PC~-{}-{}-{}-{}++++~ok}\par
\texttt{F4.3016~0123.0123.23.012~~~~-{}-{}-{}-{}D~-{}PC~-{}-{}-{}-{}++++~ok}\par
\texttt{F4.3017~0123.013.023.012~~~~-{}B-{}-{}D~-{}PC~+-{}-{}-{}++++~ok}\par
\texttt{F4.3018~0123.123.023.012~~~~-{}-{}-{}-{}D~-{}PC~-{}-{}-{}-{}++++~ok}\par
\texttt{F4.3019~023.0123.023.012~~~~-{}-{}-{}-{}D~-{}PC~-{}-{}-{}-{}++++~ok}\par
\texttt{F4.3020~0123.0123.0123.3~~~~T-{}4-{}D~-{}-{}-{}~-{}-{}-{}+-{}-{}-{}-{}~0:0001}\par
\texttt{F4.3021~0123.0123.012.03~~~~T-{}-{}-{}D~-{}PC~-{}-{}-{}-{}++++~ok}\par
\texttt{F4.3022~0123.0123.023.03~~~~T-{}-{}-{}D~-{}PC~-{}-{}-{}-{}++++~ok}\par
\texttt{F4.3023~0123.0123.123.03~~~~T-{}-{}-{}D~-{}PC~-{}-{}-{}-{}++++~ok}\par
\texttt{F4.3024~012.0123.0123.03~~~~T-{}-{}-{}D~-{}PC~-{}-{}-{}-{}++++~ok}\par
\texttt{F4.3025~013.0123.0123.03~~~~T-{}-{}-{}D~-{}PC~-{}-{}-{}-{}++-{}+~ok}\par
\texttt{F4.3026~0123.012.012.013~~~~T-{}-{}-{}D~-{}PC~-{}-{}-{}-{}++++~ok}\par
\texttt{F4.3027~0123.123.012.013~~~~T-{}-{}-{}D~-{}PC~-{}-{}-{}-{}++++~ok}\par
\texttt{F4.3028~0123.012.023.013~~~~T-{}-{}-{}D~-{}PC~-{}-{}-{}-{}++++~ok}\par
\texttt{F4.3029~0123.123.023.013~~~~T-{}-{}-{}D~-{}PC~-{}-{}-{}-{}++++~ok}\par
\texttt{F4.3030~013.0123.023.013~~~~T-{}-{}-{}D~-{}PC~-{}-{}-{}-{}++++~ok}\par
\texttt{F4.3031~023.0123.023.013~~~~T-{}-{}-{}D~-{}PC~-{}-{}-{}-{}++++~ok}\par
\texttt{F4.3032~013.013.0123.013~~~~T-{}4-{}D~-{}-{}-{}~-{}-{}-{}+-{}-{}-{}-{}~2:1101}\par
\texttt{F4.3033~0123.0123.0123.01~~~-{}-{}-{}-{}D~-{}PC~-{}-{}-{}-{}++++~ok}\par
\texttt{F4.3034~0123.0123.012.012~~~-{}-{}-{}-{}D~-{}PC~-{}-{}-{}-{}++++~ok}\par
\texttt{F4.3035~0123.0123.013.012~~~-{}B-{}-{}D~-{}PC~+-{}-{}-{}++++~ok}\par
\texttt{F4.3036~0123.0123.023.012~~~-{}-{}-{}-{}D~-{}PC~-{}-{}-{}-{}++++~ok}\par
\texttt{F4.3037~0123.0123.0123.03~~~T-{}-{}-{}D~-{}PC~-{}-{}-{}-{}++++~ok}\par
\texttt{F4.3038~0123.0123.012.013~~~TB-{}-{}D~-{}PC~+-{}-{}-{}++++~ok}\par
\texttt{F4.3039~0123.0123.023.013~~~T-{}-{}-{}D~-{}PC~-{}-{}-{}-{}++++~ok}\par
\texttt{F4.3040~0123.013.0123.013~~~T-{}-{}-{}D~-{}PC~-{}-{}-{}-{}++++~ok}\par
\texttt{F4.3041~0123.123.0123.013~~~T-{}-{}-{}D~-{}PC~-{}-{}-{}-{}++++~ok}\par
\texttt{F4.3042~0123.0123.0123.012~~-{}B-{}-{}D~-{}PC~+-{}-{}-{}++++~ok}\par
\texttt{F4.3043~0123.0123.0123.013~~T-{}-{}-{}D~-{}PC~-{}-{}-{}-{}++++~ok}\par
\texttt{F4.3044~0123.0123.0123.0123~TB45D~UPC~++++++++~ok}\par
\end{multicols}
\endgroup


\end{document}
